\section{二刷复盘总览}

这一组材料用于第二遍和第三遍复习。结构上不再把 265 道题直接堆成题号列表，而是按题型分成若干大节；每个题型大节内部只保留复盘目标、核心题卡和自测标准几个小节。

\section{数组与原地维护：二刷复盘卡}
这一大节只围绕一个题型复盘，不把题号直接铺成目录。阅读时先看复盘目标，再读核心题卡，最后用自测标准检查自己是否能独立重讲。

\subsection{复盘目标}

追问通常会落在原地修改、稳定顺序、输出规模和整数溢出上。请准备说明为什么保存下标比保存指针更稳，为什么某个位置一旦写入就不会再被未来元素破坏。

\subsection{核心题卡}

\noindent\textbf{LeetCode 5 最长回文子串。} 模型：回文围绕中心对称，中心可以是字符也可以是两个字符之间。推导重点：中心扩展避免枚举子串后再逐个检查；每次扩展只比较新增左右边界。

\textbf{二刷读题。} 读题时先问答案落在哪些下标上，是保留元素、重排元素、计算片段，还是从两侧合成贡献。本题可以压成“回文围绕中心对称，中心可以是字符也可以是两个字符之间”，所以后续所有指针都必须有区间语义。先遮住答案，只保留题面，复述候选对象、合法性条件和输出目标。

\textbf{暴力回放。} 慢版本从下标枚举开始：把每个可能位置、片段、中心或边界都按题意检查一遍。它能直接验证回文围绕中心对称，中心可以是字符也可以是两个字符之间，缺点是同一段信息会被重复读取、重复搬移或重复比较。二刷时先写慢版本或口述慢版本，用它确认不会漏答案。

\textbf{特例回放。} 拿 aba 和 abba 手算，回文不是任意子串都要重查，而是围绕中心向两边同步比较。aba 的中心是 b，abba 的中心在两个 b 中间；每扩一层只新增左右两个字符。规律是枚举 2n-1 个中心，每个中心只在新增边界相等时继续扩展。把这个特例继续拆开看：先标出已经确认的前缀或后缀，再标出未知区；能被丢掉的候选，必须是以后再也不会改变已确认区域语义的候选。这样才算从例子推出数组不变量，而不是只记一次操作。复盘时把这个特例压成状态字段：当前下标、已处理前缀边界、未知区边界、答案变量；若有前后缀贡献，还要说明左侧累计值和右侧累计值分别何时更新；再用边界 偶数长度回文、单字符、全相同字符、多个同长答案 反查初始化、更新顺序和退出条件。

\textbf{状态回放。} 手算路线：手算时把下标语义写在纸上：当前候选来自哪个位置、哪段区间、哪个中心或哪条边界；每轮状态变化后，都要能指出哪些候选已经确定、哪些候选还没有看。状态字段：当前下标、已处理前缀边界、未知区边界、答案变量；若有前后缀贡献，还要说明左侧累计值和右侧累计值分别何时更新。状态升级：把重复工作收进边界变量、前后缀值、中心扩展结果或慢指针中；状态只保存已经被证明稳定的信息，不让相同区间或相同位置被反复完整检查。

\textbf{证明和边界。} 证明从区间不变量开始：已处理区域已经满足题意，未知区域还没有承诺，每次推进不会破坏已处理区域。边界：偶数长度回文、单字符、全相同字符、多个同长答案。变体：若要求回文子序列，应转成区间 DP，不能用中心扩展。

\noindent\textbf{LeetCode 26 删除有序数组重复项。} 模型：有序数组让重复元素相邻，慢指针表示下一个唯一值写入位置。推导重点：快指针扫描，只有发现新值才写入慢指针并推进。

\textbf{二刷读题。} 读题时先问答案落在哪些下标上，是保留元素、重排元素、计算片段，还是从两侧合成贡献。本题可以压成“有序数组让重复元素相邻，慢指针表示下一个唯一值写入位置”，所以后续所有指针都必须有区间语义。先遮住答案，只保留题面，复述候选对象、合法性条件和输出目标。

\textbf{暴力回放。} 慢版本从下标枚举开始：把每个可能位置、片段、中心或边界都按题意检查一遍。它能直接验证有序数组让重复元素相邻，慢指针表示下一个唯一值写入位置，缺点是同一段信息会被重复读取、重复搬移或重复比较。二刷时先写慢版本或口述慢版本，用它确认不会漏答案。

\textbf{特例回放。} 拿 [1,1,2] 手算，write 先停在第一个 1 后面；读到第二个 1 时，它和已保留尾值相同，不写入；读到 2 时写到 write 位置，前缀变成 [1,2]。规律是慢指针左侧始终是不重复前缀，快指针只在遇到新值时推进慢指针；空数组和全重复数组只改变初始化，不改变不变量。把这个特例继续拆开看：先标出已经确认的前缀或后缀，再标出未知区；能被丢掉的候选，必须是以后再也不会改变已确认区域语义的候选。这样才算从例子推出数组不变量，而不是只记一次操作。复盘时把这个特例压成状态字段：当前下标、已处理前缀边界、未知区边界、答案变量；若有前后缀贡献，还要说明左侧累计值和右侧累计值分别何时更新；再用边界 空数组、全重复、无重复、返回长度不等于物理容量 反查初始化、更新顺序和退出条件。

\textbf{状态回放。} 手算路线：手算时把下标语义写在纸上：当前候选来自哪个位置、哪段区间、哪个中心或哪条边界；每轮状态变化后，都要能指出哪些候选已经确定、哪些候选还没有看。状态字段：当前下标、已处理前缀边界、未知区边界、答案变量；若有前后缀贡献，还要说明左侧累计值和右侧累计值分别何时更新。状态升级：把重复工作收进边界变量、前后缀值、中心扩展结果或慢指针中；状态只保存已经被证明稳定的信息，不让相同区间或相同位置被反复完整检查。

\textbf{证明和边界。} 证明从区间不变量开始：已处理区域已经满足题意，未知区域还没有承诺，每次推进不会破坏已处理区域。边界：空数组、全重复、无重复、返回长度不等于物理容量。变体：若允许每个元素最多保留两次，只需改变写入条件为和前两个比较。

\noindent\textbf{LeetCode 27 移除元素。} 模型：原地过滤，慢指针左侧始终是不等于目标值的稳定前缀。推导重点：保持顺序用快慢指针；不保持顺序可用左右交换减少写入。

\textbf{二刷读题。} 读题时先问答案落在哪些下标上，是保留元素、重排元素、计算片段，还是从两侧合成贡献。本题可以压成“原地过滤，慢指针左侧始终是不等于目标值的稳定前缀”，所以后续所有指针都必须有区间语义。先遮住答案，只保留题面，复述候选对象、合法性条件和输出目标。

\textbf{暴力回放。} 慢版本从下标枚举开始：把每个可能位置、片段、中心或边界都按题意检查一遍。它能直接验证原地过滤，慢指针左侧始终是不等于目标值的稳定前缀，缺点是同一段信息会被重复读取、重复搬移或重复比较。二刷时先写慢版本或口述慢版本，用它确认不会漏答案。

\textbf{特例回放。} 拿 [3,2,2,3]、val=3 手算，稳定写法读到 3 不写，读到两个 2 依次写入前缀，最后有效区间是 [2,2]。若题目不要求顺序，也可以把左侧 3 和右侧非 3 交换来减少写入。规律是先确认是否稳定，再决定快慢指针还是左右交换；所有元素都被移除时返回长度 0。把这个特例继续拆开看：先标出已经确认的前缀或后缀，再标出未知区；能被丢掉的候选，必须是以后再也不会改变已确认区域语义的候选。这样才算从例子推出数组不变量，而不是只记一次操作。复盘时把这个特例压成状态字段：当前下标、已处理前缀边界、未知区边界、答案变量；若有前后缀贡献，还要说明左侧累计值和右侧累计值分别何时更新；再用边界 所有元素都被移除、没有元素被移除、目标值在尾部密集 反查初始化、更新顺序和退出条件。

\textbf{状态回放。} 手算路线：手算时把下标语义写在纸上：当前候选来自哪个位置、哪段区间、哪个中心或哪条边界；每轮状态变化后，都要能指出哪些候选已经确定、哪些候选还没有看。状态字段：当前下标、已处理前缀边界、未知区边界、答案变量；若有前后缀贡献，还要说明左侧累计值和右侧累计值分别何时更新。状态升级：把重复工作收进边界变量、前后缀值、中心扩展结果或慢指针中；状态只保存已经被证明稳定的信息，不让相同区间或相同位置被反复完整检查。

\textbf{证明和边界。} 证明从区间不变量开始：已处理区域已经满足题意，未知区域还没有承诺，每次推进不会破坏已处理区域。边界：所有元素都被移除、没有元素被移除、目标值在尾部密集。变体：工程里要先确认是否要求稳定顺序，不同要求对应不同写法。

\noindent\textbf{LeetCode 31 下一个排列。} 模型：字典序结构由最长非递增后缀决定，枢轴是后缀前一个位置。推导重点：从右找第一个上升点，再从右找刚好更大的元素交换，最后反转后缀。

\textbf{二刷读题。} 读题时先问答案落在哪些下标上，是保留元素、重排元素、计算片段，还是从两侧合成贡献。本题可以压成“字典序结构由最长非递增后缀决定，枢轴是后缀前一个位置”，所以后续所有指针都必须有区间语义。先遮住答案，只保留题面，复述候选对象、合法性条件和输出目标。

\textbf{暴力回放。} 慢版本从下标枚举开始：把每个可能位置、片段、中心或边界都按题意检查一遍。它能直接验证字典序结构由最长非递增后缀决定，枢轴是后缀前一个位置，缺点是同一段信息会被重复读取、重复搬移或重复比较。二刷时先写慢版本或口述慢版本，用它确认不会漏答案。

\textbf{特例回放。} 拿 [1,3,2] 手算，后缀 [3,2] 已经非递增，枢轴是 1；从右找刚好比 1 大的 2 交换，得到 [2,3,1]，再反转后缀成 [2,1,3]。规律是最长非递增后缀已经是该前缀下的最大排列，必须提升枢轴并把后缀变成最小结构；完全降序时整体反转成最小排列。把这个特例继续拆开看：先标出已经确认的前缀或后缀，再标出未知区；能被丢掉的候选，必须是以后再也不会改变已确认区域语义的候选。这样才算从例子推出数组不变量，而不是只记一次操作。复盘时把这个特例压成状态字段：当前下标、已处理前缀边界、未知区边界、答案变量；若有前后缀贡献，还要说明左侧累计值和右侧累计值分别何时更新；再用边界 完全降序、重复数字、只有一个元素、交换后后缀必须最小化 反查初始化、更新顺序和退出条件。

\textbf{状态回放。} 手算路线：手算时把下标语义写在纸上：当前候选来自哪个位置、哪段区间、哪个中心或哪条边界；每轮状态变化后，都要能指出哪些候选已经确定、哪些候选还没有看。状态字段：当前下标、已处理前缀边界、未知区边界、答案变量；若有前后缀贡献，还要说明左侧累计值和右侧累计值分别何时更新。状态升级：把重复工作收进边界变量、前后缀值、中心扩展结果或慢指针中；状态只保存已经被证明稳定的信息，不让相同区间或相同位置被反复完整检查。

\textbf{证明和边界。} 证明从区间不变量开始：已处理区域已经满足题意，未知区域还没有承诺，每次推进不会破坏已处理区域。边界：完全降序、重复数字、只有一个元素、交换后后缀必须最小化。变体：若要求上一个排列，比较方向和后缀处理对称变化。

\noindent\textbf{LeetCode 41 缺失的第一个正数。} 模型：答案只可能落在一到 n 加一之间，数组本身可以充当值到位置的映射。推导重点：把合法正数 x 尽量交换到下标 x 减一，最后第一个位置不匹配就是答案。

\textbf{二刷读题。} 读题时先问答案落在哪些下标上，是保留元素、重排元素、计算片段，还是从两侧合成贡献。本题可以压成“答案只可能落在一到 n 加一之间，数组本身可以充当值到位置的映射”，所以后续所有指针都必须有区间语义。先遮住答案，只保留题面，复述候选对象、合法性条件和输出目标。

\textbf{暴力回放。} 慢版本从下标枚举开始：把每个可能位置、片段、中心或边界都按题意检查一遍。它能直接验证答案只可能落在一到 n 加一之间，数组本身可以充当值到位置的映射，缺点是同一段信息会被重复读取、重复搬移或重复比较。二刷时先写慢版本或口述慢版本，用它确认不会漏答案。

\textbf{特例回放。} 拿 [3,4,-1,1] 手算。答案只可能在 1 到 n+1；值 3 应放到下标 2，值 4 放到下标 3，-1 不管，1 放到下标 0。放完后第一个位置不等于 i+1 的地方就是缺失值。再试 [1,1]，若不判断目标位置已有相同值会无限交换。规律是数组本身当哈希表，只处理合法正数且避免重复交换。把这个特例继续拆开看：先标出已经确认的前缀或后缀，再标出未知区；能被丢掉的候选，必须是以后再也不会改变已确认区域语义的候选。这样才算从例子推出数组不变量，而不是只记一次操作。复盘时把这个特例压成状态字段：当前下标、已处理前缀边界、未知区边界、答案变量；若有前后缀贡献，还要说明左侧累计值和右侧累计值分别何时更新；再用边界 非正数、大于 n 的数、重复值导致无限交换、数组已经完整包含一到 n 反查初始化、更新顺序和退出条件。

\textbf{状态回放。} 手算路线：手算时把下标语义写在纸上：当前候选来自哪个位置、哪段区间、哪个中心或哪条边界；每轮状态变化后，都要能指出哪些候选已经确定、哪些候选还没有看。状态字段：当前下标、已处理前缀边界、未知区边界、答案变量；若有前后缀贡献，还要说明左侧累计值和右侧累计值分别何时更新。状态升级：把重复工作收进边界变量、前后缀值、中心扩展结果或慢指针中；状态只保存已经被证明稳定的信息，不让相同区间或相同位置被反复完整检查。

\textbf{证明和边界。} 证明从区间不变量开始：已处理区域已经满足题意，未知区域还没有承诺，每次推进不会破坏已处理区域。边界：非正数、大于 n 的数、重复值导致无限交换、数组已经完整包含一到 n。变体：若不能修改输入，只能使用哈希集合或额外布尔数组换取更简单边界。

\noindent\textbf{LeetCode 54 螺旋矩阵。} 模型：四个边界收缩模拟一圈一圈访问。推导重点：每走完一条边就收缩对应边界，并在走反方向前检查是否仍合法。

\textbf{二刷读题。} 读题时先问答案落在哪些下标上，是保留元素、重排元素、计算片段，还是从两侧合成贡献。本题可以压成“四个边界收缩模拟一圈一圈访问”，所以后续所有指针都必须有区间语义。先遮住答案，只保留题面，复述候选对象、合法性条件和输出目标。

\textbf{暴力回放。} 慢版本从下标枚举开始：把每个可能位置、片段、中心或边界都按题意检查一遍。它能直接验证四个边界收缩模拟一圈一圈访问，缺点是同一段信息会被重复读取、重复搬移或重复比较。二刷时先写慢版本或口述慢版本，用它确认不会漏答案。

\textbf{特例回放。} 拿 3x3 矩阵手算，先走上边一整行，再走右边一整列，然后收缩 top 和 right；反向走下边和左边前要检查 top<=bottom、left<=right。规律是四个边界每走完一条边就收缩一次，最后剩一行或一列时不能重复访问。把这个特例继续拆开看：先标出已经确认的前缀或后缀，再标出未知区；能被丢掉的候选，必须是以后再也不会改变已确认区域语义的候选。这样才算从例子推出数组不变量，而不是只记一次操作。复盘时把这个特例压成状态字段：当前下标、已处理前缀边界、未知区边界、答案变量；若有前后缀贡献，还要说明左侧累计值和右侧累计值分别何时更新；再用边界 单行、单列、空矩阵、最后一圈只剩一个元素 反查初始化、更新顺序和退出条件。

\textbf{状态回放。} 手算路线：手算时把下标语义写在纸上：当前候选来自哪个位置、哪段区间、哪个中心或哪条边界；每轮状态变化后，都要能指出哪些候选已经确定、哪些候选还没有看。状态字段：当前下标、已处理前缀边界、未知区边界、答案变量；若有前后缀贡献，还要说明左侧累计值和右侧累计值分别何时更新。状态升级：把重复工作收进边界变量、前后缀值、中心扩展结果或慢指针中；状态只保存已经被证明稳定的信息，不让相同区间或相同位置被反复完整检查。

\textbf{证明和边界。} 证明从区间不变量开始：已处理区域已经满足题意，未知区域还没有承诺，每次推进不会破坏已处理区域。边界：单行、单列、空矩阵、最后一圈只剩一个元素。变体：若改成生成矩阵，边界逻辑相同，只是读变成写。

\noindent\textbf{LeetCode 75 颜色分类。} 模型：三指针维护零区、一号区、未知区和二区。推导重点：遇到二时只移动右边界，不移动扫描指针，因为换回来的元素未知。

\textbf{二刷读题。} 读题时先问答案落在哪些下标上，是保留元素、重排元素、计算片段，还是从两侧合成贡献。本题可以压成“三指针维护零区、一号区、未知区和二区”，所以后续所有指针都必须有区间语义。先遮住答案，只保留题面，复述候选对象、合法性条件和输出目标。

\textbf{暴力回放。} 慢版本从下标枚举开始：把每个可能位置、片段、中心或边界都按题意检查一遍。它能直接验证三指针维护零区、一号区、未知区和二区，缺点是同一段信息会被重复读取、重复搬移或重复比较。二刷时先写慢版本或口述慢版本，用它确认不会漏答案。

\textbf{特例回放。} 拿 [2,0,1] 手算，mid 指向 2 时和右边界交换，换回来的元素未知，所以 mid 不能前进；再看到 1 才前进。规律是 [0,low) 是 0， [low,mid) 是 1， (high,n) 是 2，未知区只在 mid 到 high 之间收缩。把这个特例继续拆开看：先标出已经确认的前缀或后缀，再标出未知区；能被丢掉的候选，必须是以后再也不会改变已确认区域语义的候选。这样才算从例子推出数组不变量，而不是只记一次操作。复盘时把这个特例压成状态字段：当前下标、已处理前缀边界、未知区边界、答案变量；若有前后缀贡献，还要说明左侧累计值和右侧累计值分别何时更新；再用边界 全零、全二、已经有序、逆序、交换后漏检查 反查初始化、更新顺序和退出条件。

\textbf{状态回放。} 手算路线：手算时把下标语义写在纸上：当前候选来自哪个位置、哪段区间、哪个中心或哪条边界；每轮状态变化后，都要能指出哪些候选已经确定、哪些候选还没有看。状态字段：当前下标、已处理前缀边界、未知区边界、答案变量；若有前后缀贡献，还要说明左侧累计值和右侧累计值分别何时更新。状态升级：把重复工作收进边界变量、前后缀值、中心扩展结果或慢指针中；状态只保存已经被证明稳定的信息，不让相同区间或相同位置被反复完整检查。

\textbf{证明和边界。} 证明从区间不变量开始：已处理区域已经满足题意，未知区域还没有承诺，每次推进不会破坏已处理区域。边界：全零、全二、已经有序、逆序、交换后漏检查。变体：这是荷兰国旗问题，能推广到三类分区的原地维护。

\noindent\textbf{LeetCode 88 合并两个有序数组。} 模型：第一个数组尾部有空位，适合从后往前写入最终最大值。推导重点：逆向填充避免覆盖尚未读取的 nums1 元素。

\textbf{二刷读题。} 读题时先问答案落在哪些下标上，是保留元素、重排元素、计算片段，还是从两侧合成贡献。本题可以压成“第一个数组尾部有空位，适合从后往前写入最终最大值”，所以后续所有指针都必须有区间语义。先遮住答案，只保留题面，复述候选对象、合法性条件和输出目标。

\textbf{暴力回放。} 慢版本从下标枚举开始：把每个可能位置、片段、中心或边界都按题意检查一遍。它能直接验证第一个数组尾部有空位，适合从后往前写入最终最大值，缺点是同一段信息会被重复读取、重复搬移或重复比较。二刷时先写慢版本或口述慢版本，用它确认不会漏答案。

\textbf{特例回放。} 拿 nums1=[1,3,0]、nums2=[2] 手算，如果从前往后写，2 会覆盖 nums1 中尚未比较的 3；从尾部写，先比较 3 和 2，把 3 放到最后，再把 2 放到中间。规律是利用尾部空位逆向填充，避免覆盖未读元素。把这个特例继续拆开看：先标出已经确认的前缀或后缀，再标出未知区；能被丢掉的候选，必须是以后再也不会改变已确认区域语义的候选。这样才算从例子推出数组不变量，而不是只记一次操作。复盘时把这个特例压成状态字段：当前下标、已处理前缀边界、未知区边界、答案变量；若有前后缀贡献，还要说明左侧累计值和右侧累计值分别何时更新；再用边界 第二个数组为空、第一个有效元素为空、重复值、尾部剩余 反查初始化、更新顺序和退出条件。

\textbf{状态回放。} 手算路线：手算时把下标语义写在纸上：当前候选来自哪个位置、哪段区间、哪个中心或哪条边界；每轮状态变化后，都要能指出哪些候选已经确定、哪些候选还没有看。状态字段：当前下标、已处理前缀边界、未知区边界、答案变量；若有前后缀贡献，还要说明左侧累计值和右侧累计值分别何时更新。状态升级：把重复工作收进边界变量、前后缀值、中心扩展结果或慢指针中；状态只保存已经被证明稳定的信息，不让相同区间或相同位置被反复完整检查。

\textbf{证明和边界。} 证明从区间不变量开始：已处理区域已经满足题意，未知区域还没有承诺，每次推进不会破坏已处理区域。边界：第二个数组为空、第一个有效元素为空、重复值、尾部剩余。变体：若没有额外空位，只能新开数组或使用更复杂原地归并。

\noindent\textbf{LeetCode 189 轮转数组。} 模型：右移 k 位等价于把数组切成两段后换序。推导重点：三次反转能原地完成：整体、前 k、后 n 减 k。

\textbf{二刷读题。} 读题时先问答案落在哪些下标上，是保留元素、重排元素、计算片段，还是从两侧合成贡献。本题可以压成“右移 k 位等价于把数组切成两段后换序”，所以后续所有指针都必须有区间语义。先遮住答案，只保留题面，复述候选对象、合法性条件和输出目标。

\textbf{暴力回放。} 慢版本从下标枚举开始：把每个可能位置、片段、中心或边界都按题意检查一遍。它能直接验证右移 k 位等价于把数组切成两段后换序，缺点是同一段信息会被重复读取、重复搬移或重复比较。二刷时先写慢版本或口述慢版本，用它确认不会漏答案。

\textbf{特例回放。} 拿 [1,2,3,4]、k=2 手算，目标是 [3,4,1,2]；整体反转得 [4,3,2,1]，再反转前 2 个得 [3,4,2,1]，反转后半得 [3,4,1,2]。规律是右移 k 等价于两段换序，k 要先对 n 取模。把这个特例继续拆开看：先标出已经确认的前缀或后缀，再标出未知区；能被丢掉的候选，必须是以后再也不会改变已确认区域语义的候选。这样才算从例子推出数组不变量，而不是只记一次操作。复盘时把这个特例压成状态字段：当前下标、已处理前缀边界、未知区边界、答案变量；若有前后缀贡献，还要说明左侧累计值和右侧累计值分别何时更新；再用边界 k 为零、k 大于 n、空数组、单元素 反查初始化、更新顺序和退出条件。

\textbf{状态回放。} 手算路线：手算时把下标语义写在纸上：当前候选来自哪个位置、哪段区间、哪个中心或哪条边界；每轮状态变化后，都要能指出哪些候选已经确定、哪些候选还没有看。状态字段：当前下标、已处理前缀边界、未知区边界、答案变量；若有前后缀贡献，还要说明左侧累计值和右侧累计值分别何时更新。状态升级：把重复工作收进边界变量、前后缀值、中心扩展结果或慢指针中；状态只保存已经被证明稳定的信息，不让相同区间或相同位置被反复完整检查。

\textbf{证明和边界。} 证明从区间不变量开始：已处理区域已经满足题意，未知区域还没有承诺，每次推进不会破坏已处理区域。边界：k 为零、k 大于 n、空数组、单元素。变体：环状替换也可做，但要处理最大公约数个循环。

\noindent\textbf{LeetCode 238 除自身以外数组的乘积。} 模型：不能用除法时，答案由左侧乘积和右侧乘积组成。推导重点：先写入左侧前缀积，再从右向左乘右侧后缀积。

\textbf{二刷读题。} 读题时先问答案落在哪些下标上，是保留元素、重排元素、计算片段，还是从两侧合成贡献。本题可以压成“不能用除法时，答案由左侧乘积和右侧乘积组成”，所以后续所有指针都必须有区间语义。先遮住答案，只保留题面，复述候选对象、合法性条件和输出目标。

\textbf{暴力回放。} 慢版本从下标枚举开始：把每个可能位置、片段、中心或边界都按题意检查一遍。它能直接验证不能用除法时，答案由左侧乘积和右侧乘积组成，缺点是同一段信息会被重复读取、重复搬移或重复比较。二刷时先写慢版本或口述慢版本，用它确认不会漏答案。

\textbf{特例回放。} 拿 [1,2,3,4] 手算，第一遍把每个位置左侧乘积写进答案，得到 [1,1,2,6]；第二遍从右侧累乘 4、12、24 并乘回去。规律是答案等于左侧乘积乘右侧乘积，不能使用除法时就用两次扫描合并贡献。把这个特例继续拆开看：先标出已经确认的前缀或后缀，再标出未知区；能被丢掉的候选，必须是以后再也不会改变已确认区域语义的候选。这样才算从例子推出数组不变量，而不是只记一次操作。复盘时把这个特例压成状态字段：当前下标、已处理前缀边界、未知区边界、答案变量；若有前后缀贡献，还要说明左侧累计值和右侧累计值分别何时更新；再用边界 一个零、多个零、负数、乘积溢出 反查初始化、更新顺序和退出条件。

\textbf{状态回放。} 手算路线：手算时把下标语义写在纸上：当前候选来自哪个位置、哪段区间、哪个中心或哪条边界；每轮状态变化后，都要能指出哪些候选已经确定、哪些候选还没有看。状态字段：当前下标、已处理前缀边界、未知区边界、答案变量；若有前后缀贡献，还要说明左侧累计值和右侧累计值分别何时更新。状态升级：把重复工作收进边界变量、前后缀值、中心扩展结果或慢指针中；状态只保存已经被证明稳定的信息，不让相同区间或相同位置被反复完整检查。

\textbf{证明和边界。} 证明从区间不变量开始：已处理区域已经满足题意，未知区域还没有承诺，每次推进不会破坏已处理区域。边界：一个零、多个零、负数、乘积溢出。变体：这个题展示输出数组可作为状态存储的一部分。

\noindent\textbf{LeetCode 283 移动零。} 模型：稳定保留非零元素顺序，把零集中到尾部。推导重点：慢指针写入下一个非零位置，最后补零或用交换。

\textbf{二刷读题。} 读题时先问答案落在哪些下标上，是保留元素、重排元素、计算片段，还是从两侧合成贡献。本题可以压成“稳定保留非零元素顺序，把零集中到尾部”，所以后续所有指针都必须有区间语义。先遮住答案，只保留题面，复述候选对象、合法性条件和输出目标。

\textbf{暴力回放。} 慢版本从下标枚举开始：把每个可能位置、片段、中心或边界都按题意检查一遍。它能直接验证稳定保留非零元素顺序，把零集中到尾部，缺点是同一段信息会被重复读取、重复搬移或重复比较。二刷时先写慢版本或口述慢版本，用它确认不会漏答案。

\textbf{特例回放。} 拿 [0,1,0,3] 手算，write 指向下一个非零写入位置，读到 1 写到 0 号位，读到 3 写到 1 号位，最后把后面补 0。规律是先稳定压缩非零前缀，再填零；不能在扫描时盲目交换导致顺序变化。把这个特例继续拆开看：先标出已经确认的前缀或后缀，再标出未知区；能被丢掉的候选，必须是以后再也不会改变已确认区域语义的候选。这样才算从例子推出数组不变量，而不是只记一次操作。复盘时把这个特例压成状态字段：当前下标、已处理前缀边界、未知区边界、答案变量；若有前后缀贡献，还要说明左侧累计值和右侧累计值分别何时更新；再用边界 全零、无零、零在开头结尾、稳定顺序要求 反查初始化、更新顺序和退出条件。

\textbf{状态回放。} 手算路线：手算时把下标语义写在纸上：当前候选来自哪个位置、哪段区间、哪个中心或哪条边界；每轮状态变化后，都要能指出哪些候选已经确定、哪些候选还没有看。状态字段：当前下标、已处理前缀边界、未知区边界、答案变量；若有前后缀贡献，还要说明左侧累计值和右侧累计值分别何时更新。状态升级：把重复工作收进边界变量、前后缀值、中心扩展结果或慢指针中；状态只保存已经被证明稳定的信息，不让相同区间或相同位置被反复完整检查。

\textbf{证明和边界。} 证明从区间不变量开始：已处理区域已经满足题意，未知区域还没有承诺，每次推进不会破坏已处理区域。边界：全零、无零、零在开头结尾、稳定顺序要求。变体：若目标是移动指定值，模型与移除元素相同。

\noindent\textbf{LeetCode 581 最短无序连续子数组。} 模型：要找破坏全局有序性的最左和最右位置。推导重点：从左维护历史最大值确定右边界，从右维护历史最小值确定左边界。

\textbf{二刷读题。} 读题时先问答案落在哪些下标上，是保留元素、重排元素、计算片段，还是从两侧合成贡献。本题可以压成“要找破坏全局有序性的最左和最右位置”，所以后续所有指针都必须有区间语义。先遮住答案，只保留题面，复述候选对象、合法性条件和输出目标。

\textbf{暴力回放。} 慢版本从下标枚举开始：把每个可能位置、片段、中心或边界都按题意检查一遍。它能直接验证要找破坏全局有序性的最左和最右位置，缺点是同一段信息会被重复读取、重复搬移或重复比较。二刷时先写慢版本或口述慢版本，用它确认不会漏答案。

\textbf{特例回放。} 拿 [2,6,4,8,10,9,15] 手算，从左扫描时 4 小于此前最大 6，说明右边界至少到 4；9 小于此前最大 10，右边界更新到 9。规律是左右扫描分别找破坏有序的位置，最后得到最短需要排序的连续段。把这个特例继续拆开看：先标出已经确认的前缀或后缀，再标出未知区；能被丢掉的候选，必须是以后再也不会改变已确认区域语义的候选。这样才算从例子推出数组不变量，而不是只记一次操作。复盘时把这个特例压成状态字段：当前下标、已处理前缀边界、未知区边界、答案变量；若有前后缀贡献，还要说明左侧累计值和右侧累计值分别何时更新；再用边界 已经有序、逆序、重复值、局部乱序 反查初始化、更新顺序和退出条件。

\textbf{状态回放。} 手算路线：手算时把下标语义写在纸上：当前候选来自哪个位置、哪段区间、哪个中心或哪条边界；每轮状态变化后，都要能指出哪些候选已经确定、哪些候选还没有看。状态字段：当前下标、已处理前缀边界、未知区边界、答案变量；若有前后缀贡献，还要说明左侧累计值和右侧累计值分别何时更新。状态升级：把重复工作收进边界变量、前后缀值、中心扩展结果或慢指针中；状态只保存已经被证明稳定的信息，不让相同区间或相同位置被反复完整检查。

\textbf{证明和边界。} 证明从区间不变量开始：已处理区域已经满足题意，未知区域还没有承诺，每次推进不会破坏已处理区域。边界：已经有序、逆序、重复值、局部乱序。变体：排序比较法更直观但复杂度更高。

\noindent\textbf{LeetCode 647 回文子串。} 模型：每个回文都有唯一中心或中心间隙。推导重点：对每个中心向两边扩展，扩展成功一次就计数一次。

\textbf{二刷读题。} 读题时先问答案落在哪些下标上，是保留元素、重排元素、计算片段，还是从两侧合成贡献。本题可以压成“每个回文都有唯一中心或中心间隙”，所以后续所有指针都必须有区间语义。先遮住答案，只保留题面，复述候选对象、合法性条件和输出目标。

\textbf{暴力回放。} 慢版本从下标枚举开始：把每个可能位置、片段、中心或边界都按题意检查一遍。它能直接验证每个回文都有唯一中心或中心间隙，缺点是同一段信息会被重复读取、重复搬移或重复比较。二刷时先写慢版本或口述慢版本，用它确认不会漏答案。

\textbf{特例回放。} 拿 aaa 手算，以每个字符为中心有 3 个单字符回文，以两个字符中间为中心有 2 个 aa，以中间 a 为中心还有 aaa。规律是枚举字符中心和间隙中心，每次扩展成功就新增一个回文子串。把这个特例继续拆开看：先标出已经确认的前缀或后缀，再标出未知区；能被丢掉的候选，必须是以后再也不会改变已确认区域语义的候选。这样才算从例子推出数组不变量，而不是只记一次操作。复盘时把这个特例压成状态字段：当前下标、已处理前缀边界、未知区边界、答案变量；若有前后缀贡献，还要说明左侧累计值和右侧累计值分别何时更新；再用边界 空串、全相同字符、偶数回文、奇数回文 反查初始化、更新顺序和退出条件。

\textbf{状态回放。} 手算路线：手算时把下标语义写在纸上：当前候选来自哪个位置、哪段区间、哪个中心或哪条边界；每轮状态变化后，都要能指出哪些候选已经确定、哪些候选还没有看。状态字段：当前下标、已处理前缀边界、未知区边界、答案变量；若有前后缀贡献，还要说明左侧累计值和右侧累计值分别何时更新。状态升级：把重复工作收进边界变量、前后缀值、中心扩展结果或慢指针中；状态只保存已经被证明稳定的信息，不让相同区间或相同位置被反复完整检查。

\textbf{证明和边界。} 证明从区间不变量开始：已处理区域已经满足题意，未知区域还没有承诺，每次推进不会破坏已处理区域。边界：空串、全相同字符、偶数回文、奇数回文。变体：Manacher 算法可线性解决，但中心扩展更适合面试基础。

\subsection{自测标准}

二刷时不要只确认自己见过这道题。请遮住答案，先讲题面模型，再讲暴力基线和瓶颈，然后说出优化结构保存了什么、不变量如何排除错误候选、边界实验如何打破错误实现。

\section{滑动窗口与双指针：二刷复盘卡}
这一大节只围绕一个题型复盘，不把题号直接铺成目录。阅读时先看复盘目标，再读核心题卡，最后用自测标准检查自己是否能独立重讲。

\subsection{复盘目标}

追问通常会问窗口为什么不会漏掉左端、为什么有负数会失效、固定窗口和可变窗口有什么区别。请准备一个反例证明错误窗口条件会漏解。

\subsection{核心题卡}

\noindent\textbf{LeetCode 3 无重复字符的最长子串。} 模型：窗口保存当前无重复子串，右端每次加入一个字符。推导重点：用上次出现位置直接跳过无效左端，左边界只能向右不能回退。

\textbf{二刷读题。} 读题时先确认左右端的语义：它们可能表示连续窗口，也可能表示排序后的双指针候选。本题可以压成“窗口保存当前无重复子串，右端每次加入一个字符”，因此左右端不是模板变量，而是候选排除和状态变化的触发点。先遮住答案，只保留题面，复述候选对象、合法性条件和输出目标。

\textbf{暴力回放。} 慢版本枚举所有左端和右端，或在排序后枚举固定点与另一侧配对。它先完整覆盖题目要求的窗口、片段或有序候选；真正浪费的是相邻候选之间有大量状态可以复用，却被重新统计或重新搜索。二刷时先写慢版本或口述慢版本，用它确认不会漏答案。

\textbf{特例回放。} 拿字符串 abca 手算。窗口 abc 合法，右端加入第二个 a 后，真正冲突的是上一次 a 的位置；左端从 0 跳到 1 后，bca 重新合法。再试 abba：处理第二个 b 时左端跳到 2，后面遇到 a 的上次位置在窗口左侧之外，不能让 left 回退。规律是左边界只向右，状态保存每个字符最近出现位置。把这个特例继续拆开看：右端进入只带来一个新元素，左端离开只删掉一个旧元素；只有当旧左端被移走后不可能再产生更优答案时，窗口才可以单向推进。若这个单调淘汰条件不存在，就要换成前缀和、队列或其他状态。复盘时把这个特例压成状态字段：左端、右端、窗口计数或当前双指针和、合法性指标、当前答案；每次移动边界后，状态必须和候选区间或窗口内容一致；再用边界 空串、全相同字符、重复字符出现在窗口左侧之外、扩展字符集 反查初始化、更新顺序和退出条件。

\textbf{状态回放。} 手算路线：手算时画一张表，列出 left、right、窗口内容、窗口状态和答案。每次右端进入只加一个元素，每次左端离开只删一个元素。状态字段：左端、右端、窗口计数或当前双指针和、合法性指标、当前答案；每次移动边界后，状态必须和候选区间或窗口内容一致。状态升级：把完整窗口检查改成增量维护；右端负责引入新信息，左端负责恢复合法性或压缩答案，窗口状态必须始终和边界一致。

\textbf{证明和边界。} 证明从左端淘汰开始：被移走的左端对应窗口已经检查完，或继续保留只会让状态更差。边界：空串、全相同字符、重复字符出现在窗口左侧之外、扩展字符集。变体：若要求恰好包含 k 种字符，窗口条件从无重复变成种类计数。

\noindent\textbf{LeetCode 11 盛最多水的容器。} 模型：左右指针代表当前选取的两条边，面积由短板和宽度共同决定。推导重点：每次移动短板，因为移动长板只会缩小宽度且不会提高短板。

\textbf{二刷读题。} 读题时先确认左右端的语义：它们可能表示连续窗口，也可能表示排序后的双指针候选。本题可以压成“左右指针代表当前选取的两条边，面积由短板和宽度共同决定”，因此左右端不是模板变量，而是候选排除和状态变化的触发点。先遮住答案，只保留题面，复述候选对象、合法性条件和输出目标。

\textbf{暴力回放。} 慢版本枚举所有左端和右端，或在排序后枚举固定点与另一侧配对。它先完整覆盖题目要求的窗口、片段或有序候选；真正浪费的是相邻候选之间有大量状态可以复用，却被重新统计或重新搜索。二刷时先写慢版本或口述慢版本，用它确认不会漏答案。

\textbf{特例回放。} 拿高度 [1,8,6,2] 的两端手算。若固定宽度从最外侧开始，面积受短板限制；移动高板只会让宽度变小，短板仍不变或更差，不可能得到更大面积。只有移动短板，才有机会提高限制高度。规律是每轮淘汰短板那一侧，因为所有保留它且宽度更小的候选都被当前面积上界支配。把这个特例继续拆开看：右端进入只带来一个新元素，左端离开只删掉一个旧元素；只有当旧左端被移走后不可能再产生更优答案时，窗口才可以单向推进。若这个单调淘汰条件不存在，就要换成前缀和、队列或其他状态。复盘时把这个特例压成状态字段：左端、右端、窗口计数或当前双指针和、合法性指标、当前答案；每次移动边界后，状态必须和候选区间或窗口内容一致；再用边界 两端高度相等、数组长度为二、面积乘法可能溢出 反查初始化、更新顺序和退出条件。

\textbf{状态回放。} 手算路线：手算时画一张表，列出 left、right、窗口内容、窗口状态和答案。每次右端进入只加一个元素，每次左端离开只删一个元素。状态字段：左端、右端、窗口计数或当前双指针和、合法性指标、当前答案；每次移动边界后，状态必须和候选区间或窗口内容一致。状态升级：把完整窗口检查改成增量维护；右端负责引入新信息，左端负责恢复合法性或压缩答案，窗口状态必须始终和边界一致。

\textbf{证明和边界。} 证明从左端淘汰开始：被移走的左端对应窗口已经检查完，或继续保留只会让状态更差。边界：两端高度相等、数组长度为二、面积乘法可能溢出。变体：接雨水计算每个位置贡献，不能把本题双指针证明直接搬过去。

\noindent\textbf{LeetCode 15 三数之和。} 模型：排序后固定第一个数，剩余两数转成有序双指针。推导重点：和太小移动左指针，和太大移动右指针；固定值和左右值都要跳过重复。

\textbf{二刷读题。} 读题时先确认左右端的语义：它们可能表示连续窗口，也可能表示排序后的双指针候选。本题可以压成“排序后固定第一个数，剩余两数转成有序双指针”，因此左右端不是模板变量，而是候选排除和状态变化的触发点。先遮住答案，只保留题面，复述候选对象、合法性条件和输出目标。

\textbf{暴力回放。} 慢版本枚举所有左端和右端，或在排序后枚举固定点与另一侧配对。它先完整覆盖题目要求的窗口、片段或有序候选；真正浪费的是相邻候选之间有大量状态可以复用，却被重新统计或重新搜索。二刷时先写慢版本或口述慢版本，用它确认不会漏答案。

\textbf{特例回放。} 拿 [-1,0,1,2,-1,-4] 排序后手算。固定第一个 -1，剩下问题变成有序两数之和；和太小只能左指针右移，和太大只能右指针左移。再遇到第二个 -1 作为固定点时，会生成重复三元组，所以同层固定点要跳过重复。规律是排序把三数问题拆成枚举固定点加双指针，并在每层处理去重。把这个特例继续拆开看：右端进入只带来一个新元素，左端离开只删掉一个旧元素；只有当旧左端被移走后不可能再产生更优答案时，窗口才可以单向推进。若这个单调淘汰条件不存在，就要换成前缀和、队列或其他状态。复盘时把这个特例压成状态字段：左端、右端、窗口计数或当前双指针和、合法性指标、当前答案；每次移动边界后，状态必须和候选区间或窗口内容一致；再用边界 全零、重复元素很多、没有答案、结果顺序不要求但组合不能重复 反查初始化、更新顺序和退出条件。

\textbf{状态回放。} 手算路线：手算时画一张表，列出 left、right、窗口内容、窗口状态和答案。每次右端进入只加一个元素，每次左端离开只删一个元素。状态字段：左端、右端、窗口计数或当前双指针和、合法性指标、当前答案；每次移动边界后，状态必须和候选区间或窗口内容一致。状态升级：把完整窗口检查改成增量维护；右端负责引入新信息，左端负责恢复合法性或压缩答案，窗口状态必须始终和边界一致。

\textbf{证明和边界。} 证明从左端淘汰开始：被移走的左端对应窗口已经检查完，或继续保留只会让状态更差。边界：全零、重复元素很多、没有答案、结果顺序不要求但组合不能重复。变体：四数之和再外层固定一个数，并用 long long 防止目标和溢出。

\noindent\textbf{LeetCode 16 最接近的三数之和。} 模型：排序后固定一数，双指针寻找最接近目标的两数和。推导重点：每次根据当前和与目标的大小移动指针，同时更新最小绝对差。

\textbf{二刷读题。} 读题时先确认左右端的语义：它们可能表示连续窗口，也可能表示排序后的双指针候选。本题可以压成“排序后固定一数，双指针寻找最接近目标的两数和”，因此左右端不是模板变量，而是候选排除和状态变化的触发点。先遮住答案，只保留题面，复述候选对象、合法性条件和输出目标。

\textbf{暴力回放。} 慢版本枚举所有左端和右端，或在排序后枚举固定点与另一侧配对。它先完整覆盖题目要求的窗口、片段或有序候选；真正浪费的是相邻候选之间有大量状态可以复用，却被重新统计或重新搜索。二刷时先写慢版本或口述慢版本，用它确认不会漏答案。

\textbf{特例回放。} 拿 [-1,2,1,-4]、target=1 手算，排序后固定 -1，左右指针在 1 和 2 上得到和 2，距离目标只差 1；若和偏大，只能让右指针左移，因为右侧数再大只会更远。规律是每次先更新最小差，再按当前和与目标的大小排除一侧候选；差值相等和三数和溢出要单独检查。把这个特例继续拆开看：右端进入只带来一个新元素，左端离开只删掉一个旧元素；只有当旧左端被移走后不可能再产生更优答案时，窗口才可以单向推进。若这个单调淘汰条件不存在，就要换成前缀和、队列或其他状态。复盘时把这个特例压成状态字段：左端、右端、窗口计数或当前双指针和、合法性指标、当前答案；每次移动边界后，状态必须和候选区间或窗口内容一致；再用边界 差值相等、负数、目标很大、三数和溢出 反查初始化、更新顺序和退出条件。

\textbf{状态回放。} 手算路线：手算时画一张表，列出 left、right、窗口内容、窗口状态和答案。每次右端进入只加一个元素，每次左端离开只删一个元素。状态字段：左端、右端、窗口计数或当前双指针和、合法性指标、当前答案；每次移动边界后，状态必须和候选区间或窗口内容一致。状态升级：把完整窗口检查改成增量维护；右端负责引入新信息，左端负责恢复合法性或压缩答案，窗口状态必须始终和边界一致。

\textbf{证明和边界。} 证明从左端淘汰开始：被移走的左端对应窗口已经检查完，或继续保留只会让状态更差。边界：差值相等、负数、目标很大、三数和溢出。变体：若要求返回所有最接近组合，需要记录同差值并去重。

\noindent\textbf{LeetCode 18 四数之和。} 模型：两层固定加一层双指针，关键是剪枝和去重。推导重点：排序提供单调性，外层可用最小可能和和最大可能和提前跳过。

\textbf{二刷读题。} 读题时先确认左右端的语义：它们可能表示连续窗口，也可能表示排序后的双指针候选。本题可以压成“两层固定加一层双指针，关键是剪枝和去重”，因此左右端不是模板变量，而是候选排除和状态变化的触发点。先遮住答案，只保留题面，复述候选对象、合法性条件和输出目标。

\textbf{暴力回放。} 慢版本枚举所有左端和右端，或在排序后枚举固定点与另一侧配对。它先完整覆盖题目要求的窗口、片段或有序候选；真正浪费的是相邻候选之间有大量状态可以复用，却被重新统计或重新搜索。二刷时先写慢版本或口述慢版本，用它确认不会漏答案。

\textbf{特例回放。} 拿 [1,0,-1,0,-2,2]、target=0 手算，排序后先固定 -2，再固定 -1，剩余两数用双指针找 3。找到 [-2,-1,1,2] 后，左右指针相同值必须跳过，否则会重复输出。规律是两层固定把四数降成有序两数，外层用最小可能和与最大可能和剪枝，内外层都要去重。把这个特例继续拆开看：右端进入只带来一个新元素，左端离开只删掉一个旧元素；只有当旧左端被移走后不可能再产生更优答案时，窗口才可以单向推进。若这个单调淘汰条件不存在，就要换成前缀和、队列或其他状态。复盘时把这个特例压成状态字段：左端、右端、窗口计数或当前双指针和、合法性指标、当前答案；每次移动边界后，状态必须和候选区间或窗口内容一致；再用边界 重复元素、目标超出 int、答案很多导致输出空间主导复杂度 反查初始化、更新顺序和退出条件。

\textbf{状态回放。} 手算路线：手算时画一张表，列出 left、right、窗口内容、窗口状态和答案。每次右端进入只加一个元素，每次左端离开只删一个元素。状态字段：左端、右端、窗口计数或当前双指针和、合法性指标、当前答案；每次移动边界后，状态必须和候选区间或窗口内容一致。状态升级：把完整窗口检查改成增量维护；右端负责引入新信息，左端负责恢复合法性或压缩答案，窗口状态必须始终和边界一致。

\textbf{证明和边界。} 证明从左端淘汰开始：被移走的左端对应窗口已经检查完，或继续保留只会让状态更差。边界：重复元素、目标超出 int、答案很多导致输出空间主导复杂度。变体：k 数之和可以递归降维，但工程实现要控制重复和边界。

\noindent\textbf{LeetCode 76 最小覆盖子串。} 模型：窗口内字符计数必须覆盖目标计数，满足后尽量收缩左端。推导重点：用 formed 记录满足需求的字符种类，避免每次全量比较。

\textbf{二刷读题。} 读题时先确认左右端的语义：它们可能表示连续窗口，也可能表示排序后的双指针候选。本题可以压成“窗口内字符计数必须覆盖目标计数，满足后尽量收缩左端”，因此左右端不是模板变量，而是候选排除和状态变化的触发点。先遮住答案，只保留题面，复述候选对象、合法性条件和输出目标。

\textbf{暴力回放。} 慢版本枚举所有左端和右端，或在排序后枚举固定点与另一侧配对。它先完整覆盖题目要求的窗口、片段或有序候选；真正浪费的是相邻候选之间有大量状态可以复用，却被重新统计或重新搜索。二刷时先写慢版本或口述慢版本，用它确认不会漏答案。

\textbf{特例回放。} 拿 s=ADOBEC、t=ABC 手算，右端扩展直到第一次覆盖 A/B/C，此时窗口可行；再移动左端，去掉 A 后窗口不再覆盖，就必须停止收缩。规律是右端负责获得可行，左端负责在可行时尽量缩短；计数表要区分已有字符和需求字符。把这个特例继续拆开看：右端进入只带来一个新元素，左端离开只删掉一个旧元素；只有当旧左端被移走后不可能再产生更优答案时，窗口才可以单向推进。若这个单调淘汰条件不存在，就要换成前缀和、队列或其他状态。复盘时把这个特例压成状态字段：左端、右端、窗口计数或当前双指针和、合法性指标、当前答案；每次移动边界后，状态必须和候选区间或窗口内容一致；再用边界 目标有重复字符、源串缺字符、最小窗口在开头或结尾、大小写 反查初始化、更新顺序和退出条件。

\textbf{状态回放。} 手算路线：手算时画一张表，列出 left、right、窗口内容、窗口状态和答案。每次右端进入只加一个元素，每次左端离开只删一个元素。状态字段：左端、右端、窗口计数或当前双指针和、合法性指标、当前答案；每次移动边界后，状态必须和候选区间或窗口内容一致。状态升级：把完整窗口检查改成增量维护；右端负责引入新信息，左端负责恢复合法性或压缩答案，窗口状态必须始终和边界一致。

\textbf{证明和边界。} 证明从左端淘汰开始：被移走的左端对应窗口已经检查完，或继续保留只会让状态更差。边界：目标有重复字符、源串缺字符、最小窗口在开头或结尾、大小写。变体：若字符集固定，数组计数更快；若是 Unicode，哈希表更通用。

\noindent\textbf{LeetCode 167 两数之和 II 输入有序数组。} 模型：有序数组中左右指针的和随指针移动单调变化。推导重点：当前和太小左指针右移，太大右指针左移，等于目标返回。

\textbf{二刷读题。} 读题时先确认左右端的语义：它们可能表示连续窗口，也可能表示排序后的双指针候选。本题可以压成“有序数组中左右指针的和随指针移动单调变化”，因此左右端不是模板变量，而是候选排除和状态变化的触发点。先遮住答案，只保留题面，复述候选对象、合法性条件和输出目标。

\textbf{暴力回放。} 慢版本枚举所有左端和右端，或在排序后枚举固定点与另一侧配对。它先完整覆盖题目要求的窗口、片段或有序候选；真正浪费的是相邻候选之间有大量状态可以复用，却被重新统计或重新搜索。二刷时先写慢版本或口述慢版本，用它确认不会漏答案。

\textbf{特例回放。} 拿 [2,7,11,15]、target=9 手算，左右和为 17 太大，右指针左移；2+7 正好等于 9。规律是有序数组让和随左指针右移变大、随右指针左移变小，因此每次能排除一侧候选。把这个特例继续拆开看：右端进入只带来一个新元素，左端离开只删掉一个旧元素；只有当旧左端被移走后不可能再产生更优答案时，窗口才可以单向推进。若这个单调淘汰条件不存在，就要换成前缀和、队列或其他状态。复盘时把这个特例压成状态字段：左端、右端、窗口计数或当前双指针和、合法性指标、当前答案；每次移动边界后，状态必须和候选区间或窗口内容一致；再用边界 两个数在边界、重复值、题目下标从一开始、保证有解 反查初始化、更新顺序和退出条件。

\textbf{状态回放。} 手算路线：手算时画一张表，列出 left、right、窗口内容、窗口状态和答案。每次右端进入只加一个元素，每次左端离开只删一个元素。状态字段：左端、右端、窗口计数或当前双指针和、合法性指标、当前答案；每次移动边界后，状态必须和候选区间或窗口内容一致。状态升级：把完整窗口检查改成增量维护；右端负责引入新信息，左端负责恢复合法性或压缩答案，窗口状态必须始终和边界一致。

\textbf{证明和边界。} 证明从左端淘汰开始：被移走的左端对应窗口已经检查完，或继续保留只会让状态更差。边界：两个数在边界、重复值、题目下标从一开始、保证有解。变体：无序版本通常用哈希表，不能直接套双指针。

\noindent\textbf{LeetCode 209 长度最小的子数组。} 模型：正数数组让窗口和随右端增加、随左端减少。推导重点：右端扩展到满足目标后，不断左移缩短并更新答案。

\textbf{二刷读题。} 读题时先确认左右端的语义：它们可能表示连续窗口，也可能表示排序后的双指针候选。本题可以压成“正数数组让窗口和随右端增加、随左端减少”，因此左右端不是模板变量，而是候选排除和状态变化的触发点。先遮住答案，只保留题面，复述候选对象、合法性条件和输出目标。

\textbf{暴力回放。} 慢版本枚举所有左端和右端，或在排序后枚举固定点与另一侧配对。它先完整覆盖题目要求的窗口、片段或有序候选；真正浪费的是相邻候选之间有大量状态可以复用，却被重新统计或重新搜索。二刷时先写慢版本或口述慢版本，用它确认不会漏答案。

\textbf{特例回放。} 拿 [2,3,1,2,4,3]、target=7 手算，右端扩到 4 时窗口和达到 12，可以不断左缩；最后 [4,3] 长度 2 仍满足。规律是全正数保证右扩和变大、左缩和变小，所以每个左端只需被移动一次。把这个特例继续拆开看：右端进入只带来一个新元素，左端离开只删掉一个旧元素；只有当旧左端被移走后不可能再产生更优答案时，窗口才可以单向推进。若这个单调淘汰条件不存在，就要换成前缀和、队列或其他状态。复盘时把这个特例压成状态字段：左端、右端、窗口计数或当前双指针和、合法性指标、当前答案；每次移动边界后，状态必须和候选区间或窗口内容一致；再用边界 不存在答案、单元素达到目标、全正是必要条件 反查初始化、更新顺序和退出条件。

\textbf{状态回放。} 手算路线：手算时画一张表，列出 left、right、窗口内容、窗口状态和答案。每次右端进入只加一个元素，每次左端离开只删一个元素。状态字段：左端、右端、窗口计数或当前双指针和、合法性指标、当前答案；每次移动边界后，状态必须和候选区间或窗口内容一致。状态升级：把完整窗口检查改成增量维护；右端负责引入新信息，左端负责恢复合法性或压缩答案，窗口状态必须始终和边界一致。

\textbf{证明和边界。} 证明从左端淘汰开始：被移走的左端对应窗口已经检查完，或继续保留只会让状态更差。边界：不存在答案、单元素达到目标、全正是必要条件。变体：若含负数，需要前缀和加单调队列。

\noindent\textbf{LeetCode 438 找到字符串中所有字母异位词。} 模型：固定长度窗口内字符频次等于目标频次即为答案。推导重点：右进左出维护计数，窗口长度达到目标长度后检查差异。

\textbf{二刷读题。} 读题时先确认左右端的语义：它们可能表示连续窗口，也可能表示排序后的双指针候选。本题可以压成“固定长度窗口内字符频次等于目标频次即为答案”，因此左右端不是模板变量，而是候选排除和状态变化的触发点。先遮住答案，只保留题面，复述候选对象、合法性条件和输出目标。

\textbf{暴力回放。} 慢版本枚举所有左端和右端，或在排序后枚举固定点与另一侧配对。它先完整覆盖题目要求的窗口、片段或有序候选；真正浪费的是相邻候选之间有大量状态可以复用，却被重新统计或重新搜索。二刷时先写慢版本或口述慢版本，用它确认不会漏答案。

\textbf{特例回放。} 拿 s=cbaebabacd、p=abc 手算，长度为 3 的窗口 cba 频次和 abc 相同，起点 0 是答案；窗口右移时只删除 c、加入 e。规律是固定长度窗口维护字符频次差，重叠答案不能漏。把这个特例继续拆开看：右端进入只带来一个新元素，左端离开只删掉一个旧元素；只有当旧左端被移走后不可能再产生更优答案时，窗口才可以单向推进。若这个单调淘汰条件不存在，就要换成前缀和、队列或其他状态。复盘时把这个特例压成状态字段：左端、右端、窗口计数或当前双指针和、合法性指标、当前答案；每次移动边界后，状态必须和候选区间或窗口内容一致；再用边界 目标比源串长、重复字符、字符集固定、答案重叠 反查初始化、更新顺序和退出条件。

\textbf{状态回放。} 手算路线：手算时画一张表，列出 left、right、窗口内容、窗口状态和答案。每次右端进入只加一个元素，每次左端离开只删一个元素。状态字段：左端、右端、窗口计数或当前双指针和、合法性指标、当前答案；每次移动边界后，状态必须和候选区间或窗口内容一致。状态升级：把完整窗口检查改成增量维护；右端负责引入新信息，左端负责恢复合法性或压缩答案，窗口状态必须始终和边界一致。

\textbf{证明和边界。} 证明从左端淘汰开始：被移走的左端对应窗口已经检查完，或继续保留只会让状态更差。边界：目标比源串长、重复字符、字符集固定、答案重叠。变体：维护差异计数可以避免每次比较整个数组。

\noindent\textbf{LeetCode 567 字符串的排列。} 模型：目标串任意排列出现在源串中，等价于固定长度窗口频次相同。推导重点：窗口长度固定为目标长度，右进左出维护字符计数差异。

\textbf{二刷读题。} 读题时先确认左右端的语义：它们可能表示连续窗口，也可能表示排序后的双指针候选。本题可以压成“目标串任意排列出现在源串中，等价于固定长度窗口频次相同”，因此左右端不是模板变量，而是候选排除和状态变化的触发点。先遮住答案，只保留题面，复述候选对象、合法性条件和输出目标。

\textbf{暴力回放。} 慢版本枚举所有左端和右端，或在排序后枚举固定点与另一侧配对。它先完整覆盖题目要求的窗口、片段或有序候选；真正浪费的是相邻候选之间有大量状态可以复用，却被重新统计或重新搜索。二刷时先写慢版本或口述慢版本，用它确认不会漏答案。

\textbf{特例回放。} 拿 s=eidbaooo、p=ab 手算，窗口 ba 的频次和 ab 相同，返回真；窗口长度固定为 2，每步只进一个字符出一个字符。规律是排列出现等价于固定长度频次相同，不需要枚举所有排列。把这个特例继续拆开看：右端进入只带来一个新元素，左端离开只删掉一个旧元素；只有当旧左端被移走后不可能再产生更优答案时，窗口才可以单向推进。若这个单调淘汰条件不存在，就要换成前缀和、队列或其他状态。复盘时把这个特例压成状态字段：左端、右端、窗口计数或当前双指针和、合法性指标、当前答案；每次移动边界后，状态必须和候选区间或窗口内容一致；再用边界 目标比源串长、重复字符、字符集固定、窗口重叠 反查初始化、更新顺序和退出条件。

\textbf{状态回放。} 手算路线：手算时画一张表，列出 left、right、窗口内容、窗口状态和答案。每次右端进入只加一个元素，每次左端离开只删一个元素。状态字段：左端、右端、窗口计数或当前双指针和、合法性指标、当前答案；每次移动边界后，状态必须和候选区间或窗口内容一致。状态升级：把完整窗口检查改成增量维护；右端负责引入新信息，左端负责恢复合法性或压缩答案，窗口状态必须始终和边界一致。

\textbf{证明和边界。} 证明从左端淘汰开始：被移走的左端对应窗口已经检查完，或继续保留只会让状态更差。边界：目标比源串长、重复字符、字符集固定、窗口重叠。变体：与找到所有异位词同型，只是返回布尔值。

\noindent\textbf{LeetCode 862 和至少为 K 的最短子数组。} 模型：数组允许负数，普通滑动窗口的和不再随左端移动单调变化。推导重点：在前缀和数组上用单调队列维护可能成为最优左端的下标。

\textbf{二刷读题。} 读题时先确认左右端的语义：它们可能表示连续窗口，也可能表示排序后的双指针候选。本题可以压成“数组允许负数，普通滑动窗口的和不再随左端移动单调变化”，因此左右端不是模板变量，而是候选排除和状态变化的触发点。先遮住答案，只保留题面，复述候选对象、合法性条件和输出目标。

\textbf{暴力回放。} 慢版本枚举所有左端和右端，或在排序后枚举固定点与另一侧配对。它先完整覆盖题目要求的窗口、片段或有序候选；真正浪费的是相邻候选之间有大量状态可以复用，却被重新统计或重新搜索。二刷时先写慢版本或口述慢版本，用它确认不会漏答案。

\textbf{特例回放。} 拿 [2,-1,2]、K=3 手算，普通正数窗口会被 -1 破坏单调性；前缀和为 0,2,1,3，只有 3-0 达到 3。规律是在前缀和上用单调队列维护可能的最优左端，队首能满足时结算长度。把这个特例继续拆开看：右端进入只带来一个新元素，左端离开只删掉一个旧元素；只有当旧左端被移走后不可能再产生更优答案时，窗口才可以单向推进。若这个单调淘汰条件不存在，就要换成前缀和、队列或其他状态。复盘时把这个特例压成状态字段：左端、右端、窗口计数或当前双指针和、合法性指标、当前答案；每次移动边界后，状态必须和候选区间或窗口内容一致；再用边界 负数、答案不存在、K 为负或零、队尾支配关系 反查初始化、更新顺序和退出条件。

\textbf{状态回放。} 手算路线：手算时画一张表，列出 left、right、窗口内容、窗口状态和答案。每次右端进入只加一个元素，每次左端离开只删一个元素。状态字段：左端、右端、窗口计数或当前双指针和、合法性指标、当前答案；每次移动边界后，状态必须和候选区间或窗口内容一致。状态升级：把完整窗口检查改成增量维护；右端负责引入新信息，左端负责恢复合法性或压缩答案，窗口状态必须始终和边界一致。

\textbf{证明和边界。} 证明从左端淘汰开始：被移走的左端对应窗口已经检查完，或继续保留只会让状态更差。边界：负数、答案不存在、K 为负或零、队尾支配关系。变体：这是从窗口题升级到前缀和加单调队列的代表。

\noindent\textbf{LeetCode 1438 绝对差不超过限制的最长连续子数组。} 模型：窗口合法性取决于窗口内最大值和最小值之差。推导重点：两个单调队列分别维护最大和最小，超过限制时移动左端并清理过期下标。

\textbf{二刷读题。} 读题时先确认左右端的语义：它们可能表示连续窗口，也可能表示排序后的双指针候选。本题可以压成“窗口合法性取决于窗口内最大值和最小值之差”，因此左右端不是模板变量，而是候选排除和状态变化的触发点。先遮住答案，只保留题面，复述候选对象、合法性条件和输出目标。

\textbf{暴力回放。} 慢版本枚举所有左端和右端，或在排序后枚举固定点与另一侧配对。它先完整覆盖题目要求的窗口、片段或有序候选；真正浪费的是相邻候选之间有大量状态可以复用，却被重新统计或重新搜索。二刷时先写慢版本或口述慢版本，用它确认不会漏答案。

\textbf{特例回放。} 拿 [8,2,4,7]、limit=4 手算，窗口 [8,2] 最大最小差 6，不合法，左端必须右移；窗口 [2,4] 合法。规律是两个单调队列分别维护最大值和最小值，差值超限时移动左端并清理过期下标。把这个特例继续拆开看：右端进入只带来一个新元素，左端离开只删掉一个旧元素；只有当旧左端被移走后不可能再产生更优答案时，窗口才可以单向推进。若这个单调淘汰条件不存在，就要换成前缀和、队列或其他状态。复盘时把这个特例压成状态字段：左端、右端、窗口计数或当前双指针和、合法性指标、当前答案；每次移动边界后，状态必须和候选区间或窗口内容一致；再用边界 limit 为零、重复元素、最大最小同时过期、窗口收缩多次 反查初始化、更新顺序和退出条件。

\textbf{状态回放。} 手算路线：手算时画一张表，列出 left、right、窗口内容、窗口状态和答案。每次右端进入只加一个元素，每次左端离开只删一个元素。状态字段：左端、右端、窗口计数或当前双指针和、合法性指标、当前答案；每次移动边界后，状态必须和候选区间或窗口内容一致。状态升级：把完整窗口检查改成增量维护；右端负责引入新信息，左端负责恢复合法性或压缩答案，窗口状态必须始终和边界一致。

\textbf{证明和边界。} 证明从左端淘汰开始：被移走的左端对应窗口已经检查完，或继续保留只会让状态更差。边界：limit 为零、重复元素、最大最小同时过期、窗口收缩多次。变体：它说明窗口状态可以是动态极值，不只是频次或总和。

\subsection{自测标准}

二刷时不要只确认自己见过这道题。请遮住答案，先讲题面模型，再讲暴力基线和瓶颈，然后说出优化结构保存了什么、不变量如何排除错误候选、边界实验如何打破错误实现。

\section{前缀和与哈希表：二刷复盘卡}
这一大节只围绕一个题型复盘，不把题号直接铺成目录。阅读时先看复盘目标，再读核心题卡，最后用自测标准检查自己是否能独立重讲。

\subsection{复盘目标}

追问通常会问先查询还是先更新、哈希表保存次数还是最早下标、为什么负数不影响前缀和。请准备说明从任意合法区间如何反推到两个前缀状态。

\subsection{核心题卡}

\noindent\textbf{LeetCode 1 两数之和。} 模型：当前元素只需要寻找唯一补数，历史值到下标的映射就是最小状态。推导重点：先查再插入，避免同一个下标被用两次；若数组有序才考虑双指针。

\textbf{二刷读题。} 读题时先把条件改写成当前状态与历史状态的关系。本题可以压成“当前元素只需要寻找唯一补数，历史值到下标的映射就是最小状态”，所以重点是历史表的 key 和 value 分别代表什么。先遮住答案，只保留题面，复述候选对象、合法性条件和输出目标。

\textbf{暴力回放。} 慢版本枚举所有区间或候选对，再逐个检查条件。把检查式写出来后，可以看到当前状态只缺一个历史状态；这正是从平方枚举走向哈希表的入口。二刷时先写慢版本或口述慢版本，用它确认不会漏答案。

\textbf{特例回放。} 拿数组 [2, 7]、target=9 手算：站在 2 时历史为空，不能配对；站在 7 时只需要问历史里有没有 2。再试 [3, 3]、target=6：如果先把当前 3 插入再查，就会把同一个下标用两次。于是规律不是“用哈希表”四个字，而是每个当前位置只和它左侧历史配对，代码必须先查补数再插入当前值。把这个特例继续拆开看：每个合法答案都要还原成当前状态和某个历史状态的配对；哈希表的 key 负责定位历史状态，value 负责表达次数、最早位置或存在性。先查还是先写，必须由是否允许当前前缀和自己配对来决定。复盘时把这个特例压成状态字段：当前前缀状态、需要查询的历史 key、哈希表 value 语义、答案累计值；初始化的空前缀必须单独说明；再用边界 重复数字、目标为偶数且补数等于当前值、没有答案时返回空结果 反查初始化、更新顺序和退出条件。

\textbf{状态回放。} 手算路线：手算时按下标列出当前状态、需要的历史 key、查到的 value、更新后的哈希表。这样能看出先查后插还是先插后查。状态字段：当前前缀状态、需要查询的历史 key、哈希表 value 语义、答案累计值；初始化的空前缀必须单独说明。状态升级：把回头扫描历史改成查表；key 是当前状态需要的历史状态，value 根据目标选择次数、最早位置、最近位置或存在性。

\textbf{证明和边界。} 证明从代数等价开始：任意合法答案都对应当前状态和某个历史 key，查到的历史 key 也能反推出合法答案。边界：重复数字、目标为偶数且补数等于当前值、没有答案时返回空结果。变体：若改成返回所有不重复数对，要先排序或用频次表处理重复。

\noindent\textbf{LeetCode 49 字母异位词分组。} 模型：异位词共享同一个规范表示，可以是排序字符串或字符计数。推导重点：哈希表按规范键分桶，避免两两比较字符串。

\textbf{二刷读题。} 读题时先把条件改写成当前状态与历史状态的关系。本题可以压成“异位词共享同一个规范表示，可以是排序字符串或字符计数”，所以重点是历史表的 key 和 value 分别代表什么。先遮住答案，只保留题面，复述候选对象、合法性条件和输出目标。

\textbf{暴力回放。} 慢版本枚举所有区间或候选对，再逐个检查条件。把检查式写出来后，可以看到当前状态只缺一个历史状态；这正是从平方枚举走向哈希表的入口。二刷时先写慢版本或口述慢版本，用它确认不会漏答案。

\textbf{特例回放。} 拿 eat、tea、tan 手算，eat 和 tea 排序后都是 aet，或计数向量都是 a1e1t1；tan 的键不同。两两比较会重复比较很多字符串，哈希分桶只比较规范键。规律是异位词共享规范表示，字符集固定时计数键更直接；空字符串和重复单词也必须落到同一个键语义里。把这个特例继续拆开看：每个合法答案都要还原成当前状态和某个历史状态的配对；哈希表的 key 负责定位历史状态，value 负责表达次数、最早位置或存在性。先查还是先写，必须由是否允许当前前缀和自己配对来决定。复盘时把这个特例压成状态字段：当前前缀状态、需要查询的历史 key、哈希表 value 语义、答案累计值；初始化的空前缀必须单独说明；再用边界 空字符串、重复单词、字符集不止小写字母、计数键必须无歧义 反查初始化、更新顺序和退出条件。

\textbf{状态回放。} 手算路线：手算时按下标列出当前状态、需要的历史 key、查到的 value、更新后的哈希表。这样能看出先查后插还是先插后查。状态字段：当前前缀状态、需要查询的历史 key、哈希表 value 语义、答案累计值；初始化的空前缀必须单独说明。状态升级：把回头扫描历史改成查表；key 是当前状态需要的历史状态，value 根据目标选择次数、最早位置、最近位置或存在性。

\textbf{证明和边界。} 证明从代数等价开始：任意合法答案都对应当前状态和某个历史 key，查到的历史 key 也能反推出合法答案。边界：空字符串、重复单词、字符集不止小写字母、计数键必须无歧义。变体：若字符串很长且字符集固定，计数键通常比排序键更稳定。

\noindent\textbf{LeetCode 128 最长连续序列。} 模型：连续段只需要从没有前驱的数开始扩展。推导重点：哈希集合判断 x 减一 是否存在，跳过非起点避免重复扫描。

\textbf{二刷读题。} 读题时先把条件改写成当前状态与历史状态的关系。本题可以压成“连续段只需要从没有前驱的数开始扩展”，所以重点是历史表的 key 和 value 分别代表什么。先遮住答案，只保留题面，复述候选对象、合法性条件和输出目标。

\textbf{暴力回放。} 慢版本枚举所有区间或候选对，再逐个检查条件。把检查式写出来后，可以看到当前状态只缺一个历史状态；这正是从平方枚举走向哈希表的入口。二刷时先写慢版本或口述慢版本，用它确认不会漏答案。

\textbf{特例回放。} 拿 [100,4,200,1,3,2] 手算，4 不是起点，因为 3 存在；1 没有前驱，从 1 扩到 4 得到长度 4。规律是只从没有 x-1 的数开始扩展，避免每个连续段被重复扫描；重复元素先由集合去掉。把这个特例继续拆开看：每个合法答案都要还原成当前状态和某个历史状态的配对；哈希表的 key 负责定位历史状态，value 负责表达次数、最早位置或存在性。先查还是先写，必须由是否允许当前前缀和自己配对来决定。复盘时把这个特例压成状态字段：当前前缀状态、需要查询的历史 key、哈希表 value 语义、答案累计值；初始化的空前缀必须单独说明；再用边界 重复元素、负数、空数组、整数边界 反查初始化、更新顺序和退出条件。

\textbf{状态回放。} 手算路线：手算时按下标列出当前状态、需要的历史 key、查到的 value、更新后的哈希表。这样能看出先查后插还是先插后查。状态字段：当前前缀状态、需要查询的历史 key、哈希表 value 语义、答案累计值；初始化的空前缀必须单独说明。状态升级：把回头扫描历史改成查表；key 是当前状态需要的历史状态，value 根据目标选择次数、最早位置、最近位置或存在性。

\textbf{证明和边界。} 证明从代数等价开始：任意合法答案都对应当前状态和某个历史 key，查到的历史 key 也能反推出合法答案。边界：重复元素、负数、空数组、整数边界。变体：若数据流动态插入，可以用并查集或区间合并维护长度。

\noindent\textbf{LeetCode 202 快乐数。} 模型：反复变换数字会进入一或循环，状态空间最终有限。推导重点：用集合检测重复状态，或用快慢指针看成函数图。

\textbf{二刷读题。} 读题时先把条件改写成当前状态与历史状态的关系。本题可以压成“反复变换数字会进入一或循环，状态空间最终有限”，所以重点是历史表的 key 和 value 分别代表什么。先遮住答案，只保留题面，复述候选对象、合法性条件和输出目标。

\textbf{暴力回放。} 慢版本枚举所有区间或候选对，再逐个检查条件。把检查式写出来后，可以看到当前状态只缺一个历史状态；这正是从平方枚举走向哈希表的入口。二刷时先写慢版本或口述慢版本，用它确认不会漏答案。

\textbf{特例回放。} 拿 19 手算，1 的平方加 9 的平方得到 82，82 再变成 68、100、1，所以快乐；拿 2 会进入循环而不到 1。规律是用集合记录出现过的中间值，或用快慢指针检测平方和序列是否成环。把这个特例继续拆开看：每个合法答案都要还原成当前状态和某个历史状态的配对；哈希表的 key 负责定位历史状态，value 负责表达次数、最早位置或存在性。先查还是先写，必须由是否允许当前前缀和自己配对来决定。复盘时把这个特例压成状态字段：当前前缀状态、需要查询的历史 key、哈希表 value 语义、答案累计值；初始化的空前缀必须单独说明；再用边界 输入为一、进入循环、位平方和函数、负数不在题意内 反查初始化、更新顺序和退出条件。

\textbf{状态回放。} 手算路线：手算时按下标列出当前状态、需要的历史 key、查到的 value、更新后的哈希表。这样能看出先查后插还是先插后查。状态字段：当前前缀状态、需要查询的历史 key、哈希表 value 语义、答案累计值；初始化的空前缀必须单独说明。状态升级：把回头扫描历史改成查表；key 是当前状态需要的历史状态，value 根据目标选择次数、最早位置、最近位置或存在性。

\textbf{证明和边界。} 证明从代数等价开始：任意合法答案都对应当前状态和某个历史 key，查到的历史 key 也能反推出合法答案。边界：输入为一、进入循环、位平方和函数、负数不在题意内。变体：快慢指针版本能训练把数值迭代看成链表。

\noindent\textbf{LeetCode 217 存在重复元素。} 模型：判断是否出现过是哈希集合最直接的职责。推导重点：扫描时若插入前已经存在则立即返回真。

\textbf{二刷读题。} 读题时先把条件改写成当前状态与历史状态的关系。本题可以压成“判断是否出现过是哈希集合最直接的职责”，所以重点是历史表的 key 和 value 分别代表什么。先遮住答案，只保留题面，复述候选对象、合法性条件和输出目标。

\textbf{暴力回放。} 慢版本枚举所有区间或候选对，再逐个检查条件。把检查式写出来后，可以看到当前状态只缺一个历史状态；这正是从平方枚举走向哈希表的入口。二刷时先写慢版本或口述慢版本，用它确认不会漏答案。

\textbf{特例回放。} 拿 [1,2,3,1] 手算，扫描到最后一个 1 时，哈希集合里已经有 1，立即返回重复。规律是集合保存已经出现过的历史值，插入前查询可以在第一次重复处停止；空数组和无重复数组返回假。把这个特例继续拆开看：每个合法答案都要还原成当前状态和某个历史状态的配对；哈希表的 key 负责定位历史状态，value 负责表达次数、最早位置或存在性。先查还是先写，必须由是否允许当前前缀和自己配对来决定。复盘时把这个特例压成状态字段：当前前缀状态、需要查询的历史 key、哈希表 value 语义、答案累计值；初始化的空前缀必须单独说明；再用边界 空数组、全唯一、重复在末尾、值域很小 反查初始化、更新顺序和退出条件。

\textbf{状态回放。} 手算路线：手算时按下标列出当前状态、需要的历史 key、查到的 value、更新后的哈希表。这样能看出先查后插还是先插后查。状态字段：当前前缀状态、需要查询的历史 key、哈希表 value 语义、答案累计值；初始化的空前缀必须单独说明。状态升级：把回头扫描历史改成查表；key 是当前状态需要的历史状态，value 根据目标选择次数、最早位置、最近位置或存在性。

\textbf{证明和边界。} 证明从代数等价开始：任意合法答案都对应当前状态和某个历史 key，查到的历史 key 也能反推出合法答案。边界：空数组、全唯一、重复在末尾、值域很小。变体：若允许修改输入，排序后看相邻可以降低额外空间。

\noindent\textbf{LeetCode 242 有效的字母异位词。} 模型：两个字符串是异位词等价于每个字符出现次数相同。推导重点：固定小写字母集时用定长计数数组，字符集未知时用哈希表。

\textbf{二刷读题。} 读题时先把条件改写成当前状态与历史状态的关系。本题可以压成“两个字符串是异位词等价于每个字符出现次数相同”，所以重点是历史表的 key 和 value 分别代表什么。先遮住答案，只保留题面，复述候选对象、合法性条件和输出目标。

\textbf{暴力回放。} 慢版本枚举所有区间或候选对，再逐个检查条件。把检查式写出来后，可以看到当前状态只缺一个历史状态；这正是从平方枚举走向哈希表的入口。二刷时先写慢版本或口述慢版本，用它确认不会漏答案。

\textbf{特例回放。} 拿 anagram 和 nagaram 手算，两个字符串每个字符计数都相同；拿 rat 和 car，r 和 c 的计数不同。规律是异位词比较的是字符频次，不是字符顺序；字符集固定可用数组，泛化字符集用哈希表。把这个特例继续拆开看：每个合法答案都要还原成当前状态和某个历史状态的配对；哈希表的 key 负责定位历史状态，value 负责表达次数、最早位置或存在性。先查还是先写，必须由是否允许当前前缀和自己配对来决定。复盘时把这个特例压成状态字段：当前前缀状态、需要查询的历史 key、哈希表 value 语义、答案累计值；初始化的空前缀必须单独说明；再用边界 长度不同、空串、Unicode 字符、计数减到负数 反查初始化、更新顺序和退出条件。

\textbf{状态回放。} 手算路线：手算时按下标列出当前状态、需要的历史 key、查到的 value、更新后的哈希表。这样能看出先查后插还是先插后查。状态字段：当前前缀状态、需要查询的历史 key、哈希表 value 语义、答案累计值；初始化的空前缀必须单独说明。状态升级：把回头扫描历史改成查表；key 是当前状态需要的历史状态，value 根据目标选择次数、最早位置、最近位置或存在性。

\textbf{证明和边界。} 证明从代数等价开始：任意合法答案都对应当前状态和某个历史 key，查到的历史 key 也能反推出合法答案。边界：长度不同、空串、Unicode 字符、计数减到负数。变体：异位词分组把这个二元比较扩展成规范 key 分桶。

\noindent\textbf{LeetCode 325 和等于 k 的最长子数组长度。} 模型：最长区间要求在同一前缀差关系下尽量拉大左右端距离。推导重点：哈希表保存每个前缀和最早出现下标，当前前缀查询 prefix 减 k。

\textbf{二刷读题。} 读题时先把条件改写成当前状态与历史状态的关系。本题可以压成“最长区间要求在同一前缀差关系下尽量拉大左右端距离”，所以重点是历史表的 key 和 value 分别代表什么。先遮住答案，只保留题面，复述候选对象、合法性条件和输出目标。

\textbf{暴力回放。} 慢版本枚举所有区间或候选对，再逐个检查条件。把检查式写出来后，可以看到当前状态只缺一个历史状态；这正是从平方枚举走向哈希表的入口。二刷时先写慢版本或口述慢版本，用它确认不会漏答案。

\textbf{特例回放。} 拿 [1,-1,5,-2,3]、k=3 手算，到前缀和 3 时需要历史前缀 0；到前缀和 4 时需要历史前缀 1，最早的 1 能给出更长区间。规律是最长区间要保存每个前缀和的最早下标，而不是最新下标。把这个特例继续拆开看：每个合法答案都要还原成当前状态和某个历史状态的配对；哈希表的 key 负责定位历史状态，value 负责表达次数、最早位置或存在性。先查还是先写，必须由是否允许当前前缀和自己配对来决定。复盘时把这个特例压成状态字段：当前前缀状态、需要查询的历史 key、哈希表 value 语义、答案累计值；初始化的空前缀必须单独说明；再用边界 从下标零开始、负数、前缀重复时不能覆盖最早位置、答案不存在 反查初始化、更新顺序和退出条件。

\textbf{状态回放。} 手算路线：手算时按下标列出当前状态、需要的历史 key、查到的 value、更新后的哈希表。这样能看出先查后插还是先插后查。状态字段：当前前缀状态、需要查询的历史 key、哈希表 value 语义、答案累计值；初始化的空前缀必须单独说明。状态升级：把回头扫描历史改成查表；key 是当前状态需要的历史状态，value 根据目标选择次数、最早位置、最近位置或存在性。

\textbf{证明和边界。} 证明从代数等价开始：任意合法答案都对应当前状态和某个历史 key，查到的历史 key 也能反推出合法答案。边界：从下标零开始、负数、前缀重复时不能覆盖最早位置、答案不存在。变体：计数版本保存频次，最长版本保存最早位置，这是核心差别。

\noindent\textbf{LeetCode 454 四数相加 II。} 模型：四个数组的和为零可以拆成两个二数组和互为相反数。推导重点：先统计 A 和 B 的所有两数和频次，再枚举 C 和 D 查询相反数。

\textbf{二刷读题。} 读题时先把条件改写成当前状态与历史状态的关系。本题可以压成“四个数组的和为零可以拆成两个二数组和互为相反数”，所以重点是历史表的 key 和 value 分别代表什么。先遮住答案，只保留题面，复述候选对象、合法性条件和输出目标。

\textbf{暴力回放。} 慢版本枚举所有区间或候选对，再逐个检查条件。把检查式写出来后，可以看到当前状态只缺一个历史状态；这正是从平方枚举走向哈希表的入口。二刷时先写慢版本或口述慢版本，用它确认不会漏答案。

\textbf{特例回放。} 拿 A=[1,2]、B=[-2,-1]、C=[-1,2]、D=[0,2] 手算，先统计 A+B 的和，后面枚举 C+D 时只查相反数。规律是四重枚举拆成两个二重枚举，用哈希表保存一半的和频次。把这个特例继续拆开看：每个合法答案都要还原成当前状态和某个历史状态的配对；哈希表的 key 负责定位历史状态，value 负责表达次数、最早位置或存在性。先查还是先写，必须由是否允许当前前缀和自己配对来决定。复盘时把这个特例压成状态字段：当前前缀状态、需要查询的历史 key、哈希表 value 语义、答案累计值；初始化的空前缀必须单独说明；再用边界 重复值、负数、答案数量很大、两数和范围 反查初始化、更新顺序和退出条件。

\textbf{状态回放。} 手算路线：手算时按下标列出当前状态、需要的历史 key、查到的 value、更新后的哈希表。这样能看出先查后插还是先插后查。状态字段：当前前缀状态、需要查询的历史 key、哈希表 value 语义、答案累计值；初始化的空前缀必须单独说明。状态升级：把回头扫描历史改成查表；key 是当前状态需要的历史状态，value 根据目标选择次数、最早位置、最近位置或存在性。

\textbf{证明和边界。} 证明从代数等价开始：任意合法答案都对应当前状态和某个历史 key，查到的历史 key 也能反推出合法答案。边界：重复值、负数、答案数量很大、两数和范围。变体：这是 meet-in-the-middle 思想，比四层枚举少两个维度。

\noindent\textbf{LeetCode 523 连续的子数组和。} 模型：长度至少二的子数组和为 k 的倍数，等价于两个前缀和同余。推导重点：哈希表保存某个余数最早出现下标，当前余数若重复且距离至少二则成立。

\textbf{二刷读题。} 读题时先把条件改写成当前状态与历史状态的关系。本题可以压成“长度至少二的子数组和为 k 的倍数，等价于两个前缀和同余”，所以重点是历史表的 key 和 value 分别代表什么。先遮住答案，只保留题面，复述候选对象、合法性条件和输出目标。

\textbf{暴力回放。} 慢版本枚举所有区间或候选对，再逐个检查条件。把检查式写出来后，可以看到当前状态只缺一个历史状态；这正是从平方枚举走向哈希表的入口。二刷时先写慢版本或口述慢版本，用它确认不会漏答案。

\textbf{特例回放。} 拿 [23,2,4,6,7]、k=6 手算，前缀余数 5 在下标 0 和 2 重复，中间 [2,4] 的和能被 6 整除。规律是两个前缀余数相同表示区间和为 k 的倍数，且区间长度至少为 2。把这个特例继续拆开看：每个合法答案都要还原成当前状态和某个历史状态的配对；哈希表的 key 负责定位历史状态，value 负责表达次数、最早位置或存在性。先查还是先写，必须由是否允许当前前缀和自己配对来决定。复盘时把这个特例压成状态字段：当前前缀状态、需要查询的历史 key、哈希表 value 语义、答案累计值；初始化的空前缀必须单独说明；再用边界 k 为零、负数余数、长度限制、从开头开始的区间 反查初始化、更新顺序和退出条件。

\textbf{状态回放。} 手算路线：手算时按下标列出当前状态、需要的历史 key、查到的 value、更新后的哈希表。这样能看出先查后插还是先插后查。状态字段：当前前缀状态、需要查询的历史 key、哈希表 value 语义、答案累计值；初始化的空前缀必须单独说明。状态升级：把回头扫描历史改成查表；key 是当前状态需要的历史状态，value 根据目标选择次数、最早位置、最近位置或存在性。

\textbf{证明和边界。} 证明从代数等价开始：任意合法答案都对应当前状态和某个历史 key，查到的历史 key 也能反推出合法答案。边界：k 为零、负数余数、长度限制、从开头开始的区间。变体：保存最早下标是为了满足长度条件并保留最长距离。

\noindent\textbf{LeetCode 525 连续数组。} 模型：把零看成负一后，零和一数量相等等价于前缀和相同。推导重点：哈希表保存每个前缀和最早下标，重复出现时更新最长长度。

\textbf{二刷读题。} 读题时先把条件改写成当前状态与历史状态的关系。本题可以压成“把零看成负一后，零和一数量相等等价于前缀和相同”，所以重点是历史表的 key 和 value 分别代表什么。先遮住答案，只保留题面，复述候选对象、合法性条件和输出目标。

\textbf{暴力回放。} 慢版本枚举所有区间或候选对，再逐个检查条件。把检查式写出来后，可以看到当前状态只缺一个历史状态；这正是从平方枚举走向哈希表的入口。二刷时先写慢版本或口述慢版本，用它确认不会漏答案。

\textbf{特例回放。} 拿 [0,1,0] 手算，把 0 当 -1，前缀差值序列 -1、0、-1；差值再次出现说明中间 0 和 1 数量相同。规律是保存差值第一次出现的位置，最长区间用最早位置。把这个特例继续拆开看：每个合法答案都要还原成当前状态和某个历史状态的配对；哈希表的 key 负责定位历史状态，value 负责表达次数、最早位置或存在性。先查还是先写，必须由是否允许当前前缀和自己配对来决定。复盘时把这个特例压成状态字段：当前前缀状态、需要查询的历史 key、哈希表 value 语义、答案累计值；初始化的空前缀必须单独说明；再用边界 全零、全一、从下标零开始、不要覆盖最早位置 反查初始化、更新顺序和退出条件。

\textbf{状态回放。} 手算路线：手算时按下标列出当前状态、需要的历史 key、查到的 value、更新后的哈希表。这样能看出先查后插还是先插后查。状态字段：当前前缀状态、需要查询的历史 key、哈希表 value 语义、答案累计值；初始化的空前缀必须单独说明。状态升级：把回头扫描历史改成查表；key 是当前状态需要的历史状态，value 根据目标选择次数、最早位置、最近位置或存在性。

\textbf{证明和边界。} 证明从代数等价开始：任意合法答案都对应当前状态和某个历史 key，查到的历史 key 也能反推出合法答案。边界：全零、全一、从下标零开始、不要覆盖最早位置。变体：这是前缀和最早位置语义，不是频次语义。

\noindent\textbf{LeetCode 560 和为 K 的子数组。} 模型：当前前缀和减去目标值，就是需要匹配的历史前缀。推导重点：哈希表保存前缀和出现次数，先查询再更新。

\textbf{二刷读题。} 读题时先把条件改写成当前状态与历史状态的关系。本题可以压成“当前前缀和减去目标值，就是需要匹配的历史前缀”，所以重点是历史表的 key 和 value 分别代表什么。先遮住答案，只保留题面，复述候选对象、合法性条件和输出目标。

\textbf{暴力回放。} 慢版本枚举所有区间或候选对，再逐个检查条件。把检查式写出来后，可以看到当前状态只缺一个历史状态；这正是从平方枚举走向哈希表的入口。二刷时先写慢版本或口述慢版本，用它确认不会漏答案。

\textbf{特例回放。} 拿 [1,2,3]、k=3 手算，到前缀和 3 时，需要历史前缀 0；到前缀和 6 时，需要历史前缀 3。每个历史前缀出现几次，就新增几个区间。规律是哈希表保存前缀和频次，且初始要放 0:1。把这个特例继续拆开看：每个合法答案都要还原成当前状态和某个历史状态的配对；哈希表的 key 负责定位历史状态，value 负责表达次数、最早位置或存在性。先查还是先写，必须由是否允许当前前缀和自己配对来决定。复盘时把这个特例压成状态字段：当前前缀状态、需要查询的历史 key、哈希表 value 语义、答案累计值；初始化的空前缀必须单独说明；再用边界 负数、目标为零、从下标零开始的区间、重复前缀 反查初始化、更新顺序和退出条件。

\textbf{状态回放。} 手算路线：手算时按下标列出当前状态、需要的历史 key、查到的 value、更新后的哈希表。这样能看出先查后插还是先插后查。状态字段：当前前缀状态、需要查询的历史 key、哈希表 value 语义、答案累计值；初始化的空前缀必须单独说明。状态升级：把回头扫描历史改成查表；key 是当前状态需要的历史状态，value 根据目标选择次数、最早位置、最近位置或存在性。

\textbf{证明和边界。} 证明从代数等价开始：任意合法答案都对应当前状态和某个历史 key，查到的历史 key 也能反推出合法答案。边界：负数、目标为零、从下标零开始的区间、重复前缀。变体：若要求最长区间，哈希表应保存最早位置而不是次数。

\noindent\textbf{LeetCode 974 和可被 K 整除的子数组。} 模型：区间和能被 K 整除等价于两个前缀和余数相同。推导重点：统计每个余数出现次数，当前余数能和所有历史同余前缀组成答案。

\textbf{二刷读题。} 读题时先把条件改写成当前状态与历史状态的关系。本题可以压成“区间和能被 K 整除等价于两个前缀和余数相同”，所以重点是历史表的 key 和 value 分别代表什么。先遮住答案，只保留题面，复述候选对象、合法性条件和输出目标。

\textbf{暴力回放。} 慢版本枚举所有区间或候选对，再逐个检查条件。把检查式写出来后，可以看到当前状态只缺一个历史状态；这正是从平方枚举走向哈希表的入口。二刷时先写慢版本或口述慢版本，用它确认不会漏答案。

\textbf{特例回放。} 拿 [4,5,0,-2,-3,1]、k=5 手算，前缀余数重复时，中间区间和能被 5 整除；负数前缀要先把余数归一化。规律是哈希表保存余数频次，每到一个余数就增加已有次数。把这个特例继续拆开看：每个合法答案都要还原成当前状态和某个历史状态的配对；哈希表的 key 负责定位历史状态，value 负责表达次数、最早位置或存在性。先查还是先写，必须由是否允许当前前缀和自己配对来决定。复盘时把这个特例压成状态字段：当前前缀状态、需要查询的历史 key、哈希表 value 语义、答案累计值；初始化的空前缀必须单独说明；再用边界 负数余数归一化、K 为一、从开头开始、重复余数很多 反查初始化、更新顺序和退出条件。

\textbf{状态回放。} 手算路线：手算时按下标列出当前状态、需要的历史 key、查到的 value、更新后的哈希表。这样能看出先查后插还是先插后查。状态字段：当前前缀状态、需要查询的历史 key、哈希表 value 语义、答案累计值；初始化的空前缀必须单独说明。状态升级：把回头扫描历史改成查表；key 是当前状态需要的历史状态，value 根据目标选择次数、最早位置、最近位置或存在性。

\textbf{证明和边界。} 证明从代数等价开始：任意合法答案都对应当前状态和某个历史 key，查到的历史 key 也能反推出合法答案。边界：负数余数归一化、K 为一、从开头开始、重复余数很多。变体：如果问最长长度，保存最早位置而不是次数。

\noindent\textbf{LeetCode 1074 元素和为目标值的子矩阵数量。} 模型：二维矩阵子区域和可以固定上下边界后转成一维子数组和。推导重点：枚举行边界，压缩每列区间和，再用前缀和哈希统计目标和子数组。

\textbf{二刷读题。} 读题时先把条件改写成当前状态与历史状态的关系。本题可以压成“二维矩阵子区域和可以固定上下边界后转成一维子数组和”，所以重点是历史表的 key 和 value 分别代表什么。先遮住答案，只保留题面，复述候选对象、合法性条件和输出目标。

\textbf{暴力回放。} 慢版本枚举所有区间或候选对，再逐个检查条件。把检查式写出来后，可以看到当前状态只缺一个历史状态；这正是从平方枚举走向哈希表的入口。二刷时先写慢版本或口述慢版本，用它确认不会漏答案。

\textbf{特例回放。} 拿 3x3 矩阵和 target=0 手算，先固定上下两行，把每列压成列和，问题变成一维子数组和为 0 的计数。规律是枚举行边界，列方向用前缀和加哈希计数，二维问题降成多次一维问题。把这个特例继续拆开看：每个合法答案都要还原成当前状态和某个历史状态的配对；哈希表的 key 负责定位历史状态，value 负责表达次数、最早位置或存在性。先查还是先写，必须由是否允许当前前缀和自己配对来决定。复盘时把这个特例压成状态字段：当前前缀状态、需要查询的历史 key、哈希表 value 语义、答案累计值；初始化的空前缀必须单独说明；再用边界 单行、单列、负数、目标为零、矩阵较大 反查初始化、更新顺序和退出条件。

\textbf{状态回放。} 手算路线：手算时按下标列出当前状态、需要的历史 key、查到的 value、更新后的哈希表。这样能看出先查后插还是先插后查。状态字段：当前前缀状态、需要查询的历史 key、哈希表 value 语义、答案累计值；初始化的空前缀必须单独说明。状态升级：把回头扫描历史改成查表；key 是当前状态需要的历史状态，value 根据目标选择次数、最早位置、最近位置或存在性。

\textbf{证明和边界。} 证明从代数等价开始：任意合法答案都对应当前状态和某个历史 key，查到的历史 key 也能反推出合法答案。边界：单行、单列、负数、目标为零、矩阵较大。变体：这是二维前缀思想和一维哈希计数的组合。

\noindent\textbf{LeetCode 1248 统计优美子数组。} 模型：恰好 K 个奇数可转成前缀奇数计数差为 K。推导重点：哈希表保存某个奇数计数出现次数，当前计数查询 count - K。

\textbf{二刷读题。} 读题时先把条件改写成当前状态与历史状态的关系。本题可以压成“恰好 K 个奇数可转成前缀奇数计数差为 K”，所以重点是历史表的 key 和 value 分别代表什么。先遮住答案，只保留题面，复述候选对象、合法性条件和输出目标。

\textbf{暴力回放。} 慢版本枚举所有区间或候选对，再逐个检查条件。把检查式写出来后，可以看到当前状态只缺一个历史状态；这正是从平方枚举走向哈希表的入口。二刷时先写慢版本或口述慢版本，用它确认不会漏答案。

\textbf{特例回放。} 拿 [1,1,2,1,1]、k=3 手算，奇数个数前缀从 0 到 3 时，需要历史奇数个数 0；每个匹配历史前缀对应一个优美子数组。规律是把奇数看成 1、偶数看成 0，转成前缀和计数。把这个特例继续拆开看：每个合法答案都要还原成当前状态和某个历史状态的配对；哈希表的 key 负责定位历史状态，value 负责表达次数、最早位置或存在性。先查还是先写，必须由是否允许当前前缀和自己配对来决定。复盘时把这个特例压成状态字段：当前前缀状态、需要查询的历史 key、哈希表 value 语义、答案累计值；初始化的空前缀必须单独说明；再用边界 K 为零、全偶数、全奇数、从开头开始的区间 反查初始化、更新顺序和退出条件。

\textbf{状态回放。} 手算路线：手算时按下标列出当前状态、需要的历史 key、查到的 value、更新后的哈希表。这样能看出先查后插还是先插后查。状态字段：当前前缀状态、需要查询的历史 key、哈希表 value 语义、答案累计值；初始化的空前缀必须单独说明。状态升级：把回头扫描历史改成查表；key 是当前状态需要的历史状态，value 根据目标选择次数、最早位置、最近位置或存在性。

\textbf{证明和边界。} 证明从代数等价开始：任意合法答案都对应当前状态和某个历史 key，查到的历史 key 也能反推出合法答案。边界：K 为零、全偶数、全奇数、从开头开始的区间。变体：也可用至多 K 减至多 K-1 的滑动窗口。

\subsection{自测标准}

二刷时不要只确认自己见过这道题。请遮住答案，先讲题面模型，再讲暴力基线和瓶颈，然后说出优化结构保存了什么、不变量如何排除错误候选、边界实验如何打破错误实现。

\section{二分查找与答案空间：二刷复盘卡}
这一大节只围绕一个题型复盘，不把题号直接铺成目录。阅读时先看复盘目标，再读核心题卡，最后用自测标准检查自己是否能独立重讲。

\subsection{复盘目标}

追问通常会要求你讲清答案边界、可行性谓词或局部排除规则。请准备说明 left、right 的语义，为什么 mid 被排除或保留，以及返回值是否还需要二次验证。

\subsection{核心题卡}

\noindent\textbf{LeetCode 4 寻找两个正序数组的中位数。} 模型：两个有序数组的中位数可以看成左右两半切分后的边界值。推导重点：二分较短数组的切分位置，让左半元素数量正确且左半最大值不大于右半最小值。

\textbf{二刷读题。} 读题时先找候选空间和目标边界。本题可以压成“两个有序数组的中位数可以看成左右两半切分后的边界值”。然后把关键观察写成可执行规则：二分较短数组的切分位置，让左半元素数量正确且左半最大值不大于右半最小值。这条规则必须能安全排除一侧候选，或收缩到可检查的答案边界。先遮住答案，只保留题面，复述候选对象、合法性条件和输出目标。

\textbf{暴力回放。} 慢版本按顺序扫描下标、逐个试答案值，或直接检查候选直到遇到目标边界。它先保证候选空间覆盖完整，但每次只排除一个候选，浪费了题目给出的顺序、比较或可行性边界。二刷时先写慢版本或口述慢版本，用它确认不会漏答案。

\textbf{特例回放。} 拿 [1, 3] 和 [2] 手算，合并后中位数是 2；但不用真合并，只要把总长度切成左半两个数、右半一个数。再拿 [1, 2] 和 [3, 4]，如果短数组切 1 个、长数组切 1 个，左半最大 3 大于右半最小 2，切分错了，应把短数组切分往右调。规律是二分短数组切分点，使左半元素个数正确且左右边界有序。把这个特例继续拆开看：每次比较 mid 之后，必须能指出哪一侧整体不可能包含目标，或哪一侧整体已经满足可行性。二分的核心不是 mid 公式，而是被丢弃区间确实不含答案。复盘时把这个特例压成状态字段：left、right、mid、mid 处比较值或检查返回值、答案区间含义、返回值语义；每一轮更新都要保留目标位置或第一个可行值；再用边界 某个数组为空、总长度奇偶、切分在边界、哨兵值不能参与真实比较 反查初始化、更新顺序和退出条件。

\textbf{状态回放。} 手算路线：手算时列出 left、mid、right，以及 mid 处比较结果、可行性检查结果或局部趋势。每一步都要写出被丢弃的一侧为什么不含答案。状态字段：left、right、mid、mid 处比较值或检查返回值、答案区间含义、返回值语义；每一轮更新都要保留目标位置或第一个可行值。状态升级：把逐个试候选改成维护答案所在区间；边界谓词、可行性检查或局部比较给出分支方向，每个分支都必须能排除一侧候选。

\textbf{证明和边界。} 证明从排除一侧开始：答案空间题证明谓词单调，局部比较题证明保留下来的区间仍含目标；不能只靠经验猜测左右边界。边界：某个数组为空、总长度奇偶、切分在边界、哨兵值不能参与真实比较。变体：若只要求第 k 小元素，可以用递归丢弃较小前缀或二分切分的同类思想。

\noindent\textbf{LeetCode 33 搜索旋转排序数组。} 模型：旋转数组每次二分至少有一半保持有序。推导重点：判断哪一半有序，再看目标是否落在有序半边范围内。

\textbf{二刷读题。} 读题时先找候选空间和目标边界。本题可以压成“旋转数组每次二分至少有一半保持有序”。然后把关键观察写成可执行规则：判断哪一半有序，再看目标是否落在有序半边范围内。这条规则必须能安全排除一侧候选，或收缩到可检查的答案边界。先遮住答案，只保留题面，复述候选对象、合法性条件和输出目标。

\textbf{暴力回放。} 慢版本按顺序扫描下标、逐个试答案值，或直接检查候选直到遇到目标边界。它先保证候选空间覆盖完整，但每次只排除一个候选，浪费了题目给出的顺序、比较或可行性边界。二刷时先写慢版本或口述慢版本，用它确认不会漏答案。

\textbf{特例回放。} 拿 [4,5,6,7,0,1,2] 找 0 手算，mid=7 时左半 [4,5,6,7] 有序，但目标 0 不在这个范围，所以整段左半都能丢掉。下一轮目标落在右半。规律是每次先判断哪一半有序，再用目标是否落入有序半边来排除一整段；有重复元素时有序半边可能判断不出来。把这个特例继续拆开看：每次比较 mid 之后，必须能指出哪一侧整体不可能包含目标，或哪一侧整体已经满足可行性。二分的核心不是 mid 公式，而是被丢弃区间确实不含答案。复盘时把这个特例压成状态字段：left、right、mid、mid 处比较值或检查返回值、答案区间含义、返回值语义；每一轮更新都要保留目标位置或第一个可行值；再用边界 未旋转、长度一、目标不存在、边界等号、存在重复元素时退化 反查初始化、更新顺序和退出条件。

\textbf{状态回放。} 手算路线：手算时列出 left、mid、right，以及 mid 处比较结果、可行性检查结果或局部趋势。每一步都要写出被丢弃的一侧为什么不含答案。状态字段：left、right、mid、mid 处比较值或检查返回值、答案区间含义、返回值语义；每一轮更新都要保留目标位置或第一个可行值。状态升级：把逐个试候选改成维护答案所在区间；边界谓词、可行性检查或局部比较给出分支方向，每个分支都必须能排除一侧候选。

\textbf{证明和边界。} 证明从排除一侧开始：答案空间题证明谓词单调，局部比较题证明保留下来的区间仍含目标；不能只靠经验猜测左右边界。边界：未旋转、长度一、目标不存在、边界等号、存在重复元素时退化。变体：若重复元素很多，需要在无法判断有序半边时收缩边界。

\noindent\textbf{LeetCode 34 排序数组中查找首尾位置。} 模型：首尾位置可以拆成两个边界：第一个大于等于和第一个大于。推导重点：不要找到一个目标后向两边线性扩展，否则重复元素很多时退化。

\textbf{二刷读题。} 读题时先找候选空间和目标边界。本题可以压成“首尾位置可以拆成两个边界：第一个大于等于和第一个大于”。然后把关键观察写成可执行规则：不要找到一个目标后向两边线性扩展，否则重复元素很多时退化。这条规则必须能安全排除一侧候选，或收缩到可检查的答案边界。先遮住答案，只保留题面，复述候选对象、合法性条件和输出目标。

\textbf{暴力回放。} 慢版本按顺序扫描下标、逐个试答案值，或直接检查候选直到遇到目标边界。它先保证候选空间覆盖完整，但每次只排除一个候选，浪费了题目给出的顺序、比较或可行性边界。二刷时先写慢版本或口述慢版本，用它确认不会漏答案。

\textbf{特例回放。} 拿 [5,7,7,8,8,10] 找 8 手算，找到一个 8 后向两边扩展能得到答案，但全数组都是 8 时会退化成线性。改成两次边界二分：第一个 >=8 的位置是 3，第一个 >8 的位置是 5，区间就是 [3,4]。规律是首尾位置不是一次查找，而是两个单调边界。把这个特例继续拆开看：每次比较 mid 之后，必须能指出哪一侧整体不可能包含目标，或哪一侧整体已经满足可行性。二分的核心不是 mid 公式，而是被丢弃区间确实不含答案。复盘时把这个特例压成状态字段：left、right、mid、mid 处比较值或检查返回值、答案区间含义、返回值语义；每一轮更新都要保留目标位置或第一个可行值；再用边界 目标不存在、目标在开头或结尾、数组为空、全是目标 反查初始化、更新顺序和退出条件。

\textbf{状态回放。} 手算路线：手算时列出 left、mid、right，以及 mid 处比较结果、可行性检查结果或局部趋势。每一步都要写出被丢弃的一侧为什么不含答案。状态字段：left、right、mid、mid 处比较值或检查返回值、答案区间含义、返回值语义；每一轮更新都要保留目标位置或第一个可行值。状态升级：把逐个试候选改成维护答案所在区间；边界谓词、可行性检查或局部比较给出分支方向，每个分支都必须能排除一侧候选。

\textbf{证明和边界。} 证明从排除一侧开始：答案空间题证明谓词单调，局部比较题证明保留下来的区间仍含目标；不能只靠经验猜测左右边界。边界：目标不存在、目标在开头或结尾、数组为空、全是目标。变体：这个模型能迁移到计数某值出现次数和插入区间边界。

\noindent\textbf{LeetCode 35 搜索插入位置。} 模型：题目等价于寻找第一个大于等于目标的位置。推导重点：统一用 lower bound 语义，不在相等时提前返回也能保持边界一致。

\textbf{二刷读题。} 读题时先找候选空间和目标边界。本题可以压成“题目等价于寻找第一个大于等于目标的位置”。然后把关键观察写成可执行规则：统一用 lower bound 语义，不在相等时提前返回也能保持边界一致。这条规则必须能安全排除一侧候选，或收缩到可检查的答案边界。先遮住答案，只保留题面，复述候选对象、合法性条件和输出目标。

\textbf{暴力回放。} 慢版本按顺序扫描下标、逐个试答案值，或直接检查候选直到遇到目标边界。它先保证候选空间覆盖完整，但每次只排除一个候选，浪费了题目给出的顺序、比较或可行性边界。二刷时先写慢版本或口述慢版本，用它确认不会漏答案。

\textbf{特例回放。} 拿 [1,3,5,6] 插入 2 手算，线性看 1 还小、3 已经不小，所以答案是下标 1。二分时保持 left 是第一个可能不小于目标的位置，mid 值小于目标才丢左侧。规律是统一成 lower\_bound，不需要在相等时提前返回；目标大于所有元素时 left 会走到 n。把这个特例继续拆开看：每次比较 mid 之后，必须能指出哪一侧整体不可能包含目标，或哪一侧整体已经满足可行性。二分的核心不是 mid 公式，而是被丢弃区间确实不含答案。复盘时把这个特例压成状态字段：left、right、mid、mid 处比较值或检查返回值、答案区间含义、返回值语义；每一轮更新都要保留目标位置或第一个可行值；再用边界 目标小于所有元素、目标大于所有元素、重复值、空数组 反查初始化、更新顺序和退出条件。

\textbf{状态回放。} 手算路线：手算时列出 left、mid、right，以及 mid 处比较结果、可行性检查结果或局部趋势。每一步都要写出被丢弃的一侧为什么不含答案。状态字段：left、right、mid、mid 处比较值或检查返回值、答案区间含义、返回值语义；每一轮更新都要保留目标位置或第一个可行值。状态升级：把逐个试候选改成维护答案所在区间；边界谓词、可行性检查或局部比较给出分支方向，每个分支都必须能排除一侧候选。

\textbf{证明和边界。} 证明从排除一侧开始：答案空间题证明谓词单调，局部比较题证明保留下来的区间仍含目标；不能只靠经验猜测左右边界。边界：目标小于所有元素、目标大于所有元素、重复值、空数组。变体：手写二分时用左闭右开可直接返回 left。

\noindent\textbf{LeetCode 69 x 的平方根。} 模型：答案是最大的整数 r，使得 r 的平方不超过 x。推导重点：在整数范围上二分右边界，比较时用除法或 long long 避免平方溢出。

\textbf{二刷读题。} 读题时先找候选空间和目标边界。本题可以压成“答案是最大的整数 r，使得 r 的平方不超过 x”。然后把关键观察写成可执行规则：在整数范围上二分右边界，比较时用除法或 long long 避免平方溢出。这条规则必须能安全排除一侧候选，或收缩到可检查的答案边界。先遮住答案，只保留题面，复述候选对象、合法性条件和输出目标。

\textbf{暴力回放。} 慢版本按顺序扫描下标、逐个试答案值，或直接检查候选直到遇到目标边界。它先保证候选空间覆盖完整，但每次只排除一个候选，浪费了题目给出的顺序、比较或可行性边界。二刷时先写慢版本或口述慢版本，用它确认不会漏答案。

\textbf{特例回放。} 拿 x=8 手算，二的平方不超过 8，而三的平方已经超过 8，所以答案是 2。规律是找最后一个平方不超过 x 的整数，比较平方时用 long long 或除法避免溢出。把这个特例继续拆开看：每次比较 mid 之后，必须能指出哪一侧整体不可能包含目标，或哪一侧整体已经满足可行性。二分的核心不是 mid 公式，而是被丢弃区间确实不含答案。复盘时把这个特例压成状态字段：left、right、mid、mid 处比较值或检查返回值、答案区间含义、返回值语义；每一轮更新都要保留目标位置或第一个可行值；再用边界 x 为零、一、非完全平方、接近 int 最大值 反查初始化、更新顺序和退出条件。

\textbf{状态回放。} 手算路线：手算时列出 left、mid、right，以及 mid 处比较结果、可行性检查结果或局部趋势。每一步都要写出被丢弃的一侧为什么不含答案。状态字段：left、right、mid、mid 处比较值或检查返回值、答案区间含义、返回值语义；每一轮更新都要保留目标位置或第一个可行值。状态升级：把逐个试候选改成维护答案所在区间；边界谓词、可行性检查或局部比较给出分支方向，每个分支都必须能排除一侧候选。

\textbf{证明和边界。} 证明从排除一侧开始：答案空间题证明谓词单调，局部比较题证明保留下来的区间仍含目标；不能只靠经验猜测左右边界。边界：x 为零、一、非完全平方、接近 int 最大值。变体：若要求浮点精度，可以二分实数区间或用牛顿迭代。

\noindent\textbf{LeetCode 74 搜索二维矩阵。} 模型：矩阵每行有序且下一行首元素大于上一行末元素，可视为一维有序数组。推导重点：把下标 mid 映射到 row 和 col，再做标准二分。

\textbf{二刷读题。} 读题时先找候选空间和目标边界。本题可以压成“矩阵每行有序且下一行首元素大于上一行末元素，可视为一维有序数组”。然后把关键观察写成可执行规则：把下标 mid 映射到 row 和 col，再做标准二分。这条规则必须能安全排除一侧候选，或收缩到可检查的答案边界。先遮住答案，只保留题面，复述候选对象、合法性条件和输出目标。

\textbf{暴力回放。} 慢版本按顺序扫描下标、逐个试答案值，或直接检查候选直到遇到目标边界。它先保证候选空间覆盖完整，但每次只排除一个候选，浪费了题目给出的顺序、比较或可行性边界。二刷时先写慢版本或口述慢版本，用它确认不会漏答案。

\textbf{特例回放。} 拿矩阵 [[1,3,5],[7,9,11]] 找 9 手算，把它看成一维有序数组 [1,3,5,7,9,11]，下标 4 映射到 row=1、col=1。规律是行尾小于下一行行首时可以整体展平二分，空矩阵和单行要先处理。把这个特例继续拆开看：每次比较 mid 之后，必须能指出哪一侧整体不可能包含目标，或哪一侧整体已经满足可行性。二分的核心不是 mid 公式，而是被丢弃区间确实不含答案。复盘时把这个特例压成状态字段：left、right、mid、mid 处比较值或检查返回值、答案区间含义、返回值语义；每一轮更新都要保留目标位置或第一个可行值；再用边界 空矩阵、单行、单列、目标小于全部或大于全部 反查初始化、更新顺序和退出条件。

\textbf{状态回放。} 手算路线：手算时列出 left、mid、right，以及 mid 处比较结果、可行性检查结果或局部趋势。每一步都要写出被丢弃的一侧为什么不含答案。状态字段：left、right、mid、mid 处比较值或检查返回值、答案区间含义、返回值语义；每一轮更新都要保留目标位置或第一个可行值。状态升级：把逐个试候选改成维护答案所在区间；边界谓词、可行性检查或局部比较给出分支方向，每个分支都必须能排除一侧候选。

\textbf{证明和边界。} 证明从排除一侧开始：答案空间题证明谓词单调，局部比较题证明保留下来的区间仍含目标；不能只靠经验猜测左右边界。边界：空矩阵、单行、单列、目标小于全部或大于全部。变体：若只是行列分别有序但行之间不连续，应改用右上角排除法。

\noindent\textbf{LeetCode 81 搜索旋转排序数组 II。} 模型：旋转数组存在重复值时，二分可能无法判断哪一半有序。推导重点：当左右和中点值相同导致信息不足时，收缩边界；否则仍按有序半边排除。

\textbf{二刷读题。} 读题时先找候选空间和目标边界。本题可以压成“旋转数组存在重复值时，二分可能无法判断哪一半有序”。然后把关键观察写成可执行规则：当左右和中点值相同导致信息不足时，收缩边界；否则仍按有序半边排除。这条规则必须能安全排除一侧候选，或收缩到可检查的答案边界。先遮住答案，只保留题面，复述候选对象、合法性条件和输出目标。

\textbf{暴力回放。} 慢版本按顺序扫描下标、逐个试答案值，或直接检查候选直到遇到目标边界。它先保证候选空间覆盖完整，但每次只排除一个候选，浪费了题目给出的顺序、比较或可行性边界。二刷时先写慢版本或口述慢版本，用它确认不会漏答案。

\textbf{特例回放。} 拿 [2,5,6,0,0,1,2] 找 0 手算，可以判断一侧有序并排除；拿 [1,1,1,1] 时 mid、left、right 相同，无法判断，只能收缩边界。规律是重复元素让旋转二分可能退化，但每次丢掉一个相同边界不会漏目标。把这个特例继续拆开看：每次比较 mid 之后，必须能指出哪一侧整体不可能包含目标，或哪一侧整体已经满足可行性。二分的核心不是 mid 公式，而是被丢弃区间确实不含答案。复盘时把这个特例压成状态字段：left、right、mid、mid 处比较值或检查返回值、答案区间含义、返回值语义；每一轮更新都要保留目标位置或第一个可行值；再用边界 全相同元素、目标不存在、重复值跨过旋转点、退化到线性 反查初始化、更新顺序和退出条件。

\textbf{状态回放。} 手算路线：手算时列出 left、mid、right，以及 mid 处比较结果、可行性检查结果或局部趋势。每一步都要写出被丢弃的一侧为什么不含答案。状态字段：left、right、mid、mid 处比较值或检查返回值、答案区间含义、返回值语义；每一轮更新都要保留目标位置或第一个可行值。状态升级：把逐个试候选改成维护答案所在区间；边界谓词、可行性检查或局部比较给出分支方向，每个分支都必须能排除一侧候选。

\textbf{证明和边界。} 证明从排除一侧开始：答案空间题证明谓词单调，局部比较题证明保留下来的区间仍含目标；不能只靠经验猜测左右边界。边界：全相同元素、目标不存在、重复值跨过旋转点、退化到线性。变体：这题说明二分依赖信息量，重复值会削弱每次排除一半的能力。

\noindent\textbf{LeetCode 153 寻找旋转排序数组中的最小值。} 模型：最小值是旋转断点，右端点可帮助判断中点在哪一段。推导重点：若 mid 大于 right，最小值在右侧；否则在左侧含 mid。

\textbf{二刷读题。} 读题时先找候选空间和目标边界。本题可以压成“最小值是旋转断点，右端点可帮助判断中点在哪一段”。然后把关键观察写成可执行规则：若 mid 大于 right，最小值在右侧；否则在左侧含 mid。这条规则必须能安全排除一侧候选，或收缩到可检查的答案边界。先遮住答案，只保留题面，复述候选对象、合法性条件和输出目标。

\textbf{暴力回放。} 慢版本按顺序扫描下标、逐个试答案值，或直接检查候选直到遇到目标边界。它先保证候选空间覆盖完整，但每次只排除一个候选，浪费了题目给出的顺序、比较或可行性边界。二刷时先写慢版本或口述慢版本，用它确认不会漏答案。

\textbf{特例回放。} 拿 [3,4,5,1,2] 手算，mid=5 大于右端 2，说明最小值一定在 mid 右侧；若 mid 小于右端，最小值在左半含 mid。规律是和右端比较能判断旋转点所在侧，未旋转数组会一路收缩到第一个元素。把这个特例继续拆开看：每次比较 mid 之后，必须能指出哪一侧整体不可能包含目标，或哪一侧整体已经满足可行性。二分的核心不是 mid 公式，而是被丢弃区间确实不含答案。复盘时把这个特例压成状态字段：left、right、mid、mid 处比较值或检查返回值、答案区间含义、返回值语义；每一轮更新都要保留目标位置或第一个可行值；再用边界 未旋转、长度一、旋转一位、无重复条件 反查初始化、更新顺序和退出条件。

\textbf{状态回放。} 手算路线：手算时列出 left、mid、right，以及 mid 处比较结果、可行性检查结果或局部趋势。每一步都要写出被丢弃的一侧为什么不含答案。状态字段：left、right、mid、mid 处比较值或检查返回值、答案区间含义、返回值语义；每一轮更新都要保留目标位置或第一个可行值。状态升级：把逐个试候选改成维护答案所在区间；边界谓词、可行性检查或局部比较给出分支方向，每个分支都必须能排除一侧候选。

\textbf{证明和边界。} 证明从排除一侧开始：答案空间题证明谓词单调，局部比较题证明保留下来的区间仍含目标；不能只靠经验猜测左右边界。边界：未旋转、长度一、旋转一位、无重复条件。变体：有重复时遇到相等只能收缩 right，最坏退化。

\noindent\textbf{LeetCode 154 寻找旋转排序数组中的最小值 II。} 模型：重复元素会让 mid 和 right 相等时无法判断最小值方向。推导重点：相等时只能安全地丢掉一个 right，因为它不是唯一必要候选。

\textbf{二刷读题。} 读题时先找候选空间和目标边界。本题可以压成“重复元素会让 mid 和 right 相等时无法判断最小值方向”。然后把关键观察写成可执行规则：相等时只能安全地丢掉一个 right，因为它不是唯一必要候选。这条规则必须能安全排除一侧候选，或收缩到可检查的答案边界。先遮住答案，只保留题面，复述候选对象、合法性条件和输出目标。

\textbf{暴力回放。} 慢版本按顺序扫描下标、逐个试答案值，或直接检查候选直到遇到目标边界。它先保证候选空间覆盖完整，但每次只排除一个候选，浪费了题目给出的顺序、比较或可行性边界。二刷时先写慢版本或口述慢版本，用它确认不会漏答案。

\textbf{特例回放。} 拿 [2,2,2,0,1] 手算，mid 和 right 都是 2 时无法判断最小值在哪侧，只能 right-- 丢掉一个重复值。规律是重复元素会削弱二分信息，最坏会退化成线性，但不会丢失最小值。把这个特例继续拆开看：每次比较 mid 之后，必须能指出哪一侧整体不可能包含目标，或哪一侧整体已经满足可行性。二分的核心不是 mid 公式，而是被丢弃区间确实不含答案。复盘时把这个特例压成状态字段：left、right、mid、mid 处比较值或检查返回值、答案区间含义、返回值语义；每一轮更新都要保留目标位置或第一个可行值；再用边界 全相等、重复值跨断点、最小值出现多次 反查初始化、更新顺序和退出条件。

\textbf{状态回放。} 手算路线：手算时列出 left、mid、right，以及 mid 处比较结果、可行性检查结果或局部趋势。每一步都要写出被丢弃的一侧为什么不含答案。状态字段：left、right、mid、mid 处比较值或检查返回值、答案区间含义、返回值语义；每一轮更新都要保留目标位置或第一个可行值。状态升级：把逐个试候选改成维护答案所在区间；边界谓词、可行性检查或局部比较给出分支方向，每个分支都必须能排除一侧候选。

\textbf{证明和边界。} 证明从排除一侧开始：答案空间题证明谓词单调，局部比较题证明保留下来的区间仍含目标；不能只靠经验猜测左右边界。边界：全相等、重复值跨断点、最小值出现多次。变体：这题说明二分依赖的信息不足时会退化，不是所有旋转题都严格对数。

\noindent\textbf{LeetCode 162 寻找峰值。} 模型：比较 mid 和 mid 加一 的斜率能判断某侧一定存在峰值。推导重点：上升则右侧有峰，下降则左侧含 mid 有峰。

\textbf{二刷读题。} 读题时先找候选空间和目标边界。本题可以压成“比较 mid 和 mid 加一 的斜率能判断某侧一定存在峰值”。然后把关键观察写成可执行规则：上升则右侧有峰，下降则左侧含 mid 有峰。这条规则必须能安全排除一侧候选，或收缩到可检查的答案边界。先遮住答案，只保留题面，复述候选对象、合法性条件和输出目标。

\textbf{暴力回放。} 慢版本按顺序扫描下标、逐个试答案值，或直接检查候选直到遇到目标边界。它先保证候选空间覆盖完整，但每次只排除一个候选，浪费了题目给出的顺序、比较或可行性边界。二刷时先写慢版本或口述慢版本，用它确认不会漏答案。

\textbf{特例回放。} 拿 [1,2,3,1] 手算，mid=2 时 nums[mid] 小于右邻 3，说明右侧上坡方向必有峰值；若大于右邻，左侧含 mid 必有峰值。规律是比较 mid 和 mid+1 就能保留一侧峰值区间，边界外可视为负无穷。把这个特例继续拆开看：每次比较 mid 之后，必须能指出哪一侧整体不可能包含目标，或哪一侧整体已经满足可行性。二分的核心不是 mid 公式，而是被丢弃区间确实不含答案。复盘时把这个特例压成状态字段：left、right、mid、mid 处比较值或检查返回值、答案区间含义、返回值语义；每一轮更新都要保留目标位置或第一个可行值；再用边界 长度一、峰在边界、多个峰、相邻不相等假设 反查初始化、更新顺序和退出条件。

\textbf{状态回放。} 手算路线：手算时列出 left、mid、right，以及 mid 处比较结果、可行性检查结果或局部趋势。每一步都要写出被丢弃的一侧为什么不含答案。状态字段：left、right、mid、mid 处比较值或检查返回值、答案区间含义、返回值语义；每一轮更新都要保留目标位置或第一个可行值。状态升级：把逐个试候选改成维护答案所在区间；边界谓词、可行性检查或局部比较给出分支方向，每个分支都必须能排除一侧候选。

\textbf{证明和边界。} 证明从排除一侧开始：答案空间题证明谓词单调，局部比较题证明保留下来的区间仍含目标；不能只靠经验猜测左右边界。边界：长度一、峰在边界、多个峰、相邻不相等假设。变体：二维峰值需要按行列最大值进行更复杂二分。

\noindent\textbf{LeetCode 240 搜索二维矩阵 II。} 模型：右上角同时具有向左变小、向下变大的排除能力。推导重点：当前值大于目标就左移，小于目标就下移。

\textbf{二刷读题。} 读题时先找候选空间和目标边界。本题可以压成“右上角同时具有向左变小、向下变大的排除能力”。然后把关键观察写成可执行规则：当前值大于目标就左移，小于目标就下移。这条规则必须能安全排除一侧候选，或收缩到可检查的答案边界。先遮住答案，只保留题面，复述候选对象、合法性条件和输出目标。

\textbf{暴力回放。} 慢版本按顺序扫描下标、逐个试答案值，或直接检查候选直到遇到目标边界。它先保证候选空间覆盖完整，但每次只排除一个候选，浪费了题目给出的顺序、比较或可行性边界。二刷时先写慢版本或口述慢版本，用它确认不会漏答案。

\textbf{特例回放。} 拿矩阵 [[1,4,7],[2,5,8],[3,6,9]] 找 6 手算，从右上角 7 开始，7 太大就左移到 4，4 太小就下移到 5，再下移到 6。规律是右上角一列向下增大、一行向左减小，每步能排除一行或一列。把这个特例继续拆开看：每次比较 mid 之后，必须能指出哪一侧整体不可能包含目标，或哪一侧整体已经满足可行性。二分的核心不是 mid 公式，而是被丢弃区间确实不含答案。复盘时把这个特例压成状态字段：left、right、mid、mid 处比较值或检查返回值、答案区间含义、返回值语义；每一轮更新都要保留目标位置或第一个可行值；再用边界 空矩阵、单行、单列、目标在角落、重复值 反查初始化、更新顺序和退出条件。

\textbf{状态回放。} 手算路线：手算时列出 left、mid、right，以及 mid 处比较结果、可行性检查结果或局部趋势。每一步都要写出被丢弃的一侧为什么不含答案。状态字段：left、right、mid、mid 处比较值或检查返回值、答案区间含义、返回值语义；每一轮更新都要保留目标位置或第一个可行值。状态升级：把逐个试候选改成维护答案所在区间；边界谓词、可行性检查或局部比较给出分支方向，每个分支都必须能排除一侧候选。

\textbf{证明和边界。} 证明从排除一侧开始：答案空间题证明谓词单调，局部比较题证明保留下来的区间仍含目标；不能只靠经验猜测左右边界。边界：空矩阵、单行、单列、目标在角落、重复值。变体：从左上角无法一次排除一行或一列。

\noindent\textbf{LeetCode 410 分割数组的最大值。} 模型：连续分成 m 段后最大段和越大，越容易完成分割。推导重点：二分最大段和上限，检查函数贪心累计，超过上限就新开一段。

\textbf{二刷读题。} 读题时先找候选空间和目标边界。本题可以压成“连续分成 m 段后最大段和越大，越容易完成分割”。然后把关键观察写成可执行规则：二分最大段和上限，检查函数贪心累计，超过上限就新开一段。这条规则必须能安全排除一侧候选，或收缩到可检查的答案边界。先遮住答案，只保留题面，复述候选对象、合法性条件和输出目标。

\textbf{暴力回放。} 慢版本按顺序扫描下标、逐个试答案值，或直接检查候选直到遇到目标边界。它先保证候选空间覆盖完整，但每次只排除一个候选，浪费了题目给出的顺序、比较或可行性边界。二刷时先写慢版本或口述慢版本，用它确认不会漏答案。

\textbf{特例回放。} 拿 [7,2,5,10,8] 分成 2 段手算，若最大段和限制为 18，可以分成 [7,2,5] 和 [10,8]；限制为 17 就需要三段。规律是答案空间是最大段和，可行性随容量增大单调变真，所以二分最小可行容量。把这个特例继续拆开看：每次比较 mid 之后，必须能指出哪一侧整体不可能包含目标，或哪一侧整体已经满足可行性。二分的核心不是 mid 公式，而是被丢弃区间确实不含答案。复盘时把这个特例压成状态字段：left、right、mid、mid 处比较值或检查返回值、答案区间含义、返回值语义；每一轮更新都要保留目标位置或第一个可行值；再用边界 m 为一、m 等于数组长度、单个元素超过猜测上限、总和溢出 反查初始化、更新顺序和退出条件。

\textbf{状态回放。} 手算路线：手算时列出 left、mid、right，以及 mid 处比较结果、可行性检查结果或局部趋势。每一步都要写出被丢弃的一侧为什么不含答案。状态字段：left、right、mid、mid 处比较值或检查返回值、答案区间含义、返回值语义；每一轮更新都要保留目标位置或第一个可行值。状态升级：把逐个试候选改成维护答案所在区间；边界谓词、可行性检查或局部比较给出分支方向，每个分支都必须能排除一侧候选。

\textbf{证明和边界。} 证明从排除一侧开始：答案空间题证明谓词单调，局部比较题证明保留下来的区间仍含目标；不能只靠经验猜测左右边界。边界：m 为一、m 等于数组长度、单个元素超过猜测上限、总和溢出。变体：也可用 DP 求精确最优，但答案二分更适合大输入。

\noindent\textbf{LeetCode 540 有序数组中的单一元素。} 模型：除一个元素外其余元素都成对出现，单一元素会改变配对下标的奇偶规律。推导重点：把 mid 调整到偶数位，若 nums[mid] 等于 nums[mid+1]，单一元素在右侧，否则在左侧。

\textbf{二刷读题。} 读题时先找候选空间和目标边界。本题可以压成“除一个元素外其余元素都成对出现，单一元素会改变配对下标的奇偶规律”。然后把关键观察写成可执行规则：把 mid 调整到偶数位，若 nums[mid] 等于 nums[mid+1]，单一元素在右侧，否则在左侧。这条规则必须能安全排除一侧候选，或收缩到可检查的答案边界。先遮住答案，只保留题面，复述候选对象、合法性条件和输出目标。

\textbf{暴力回放。} 慢版本按顺序扫描下标、逐个试答案值，或直接检查候选直到遇到目标边界。它先保证候选空间覆盖完整，但每次只排除一个候选，浪费了题目给出的顺序、比较或可行性边界。二刷时先写慢版本或口述慢版本，用它确认不会漏答案。

\textbf{特例回放。} 拿 [1,1,2,3,3] 手算，单一元素 2 之前成对元素从偶数下标开始，之后配对位置会错位。规律是用 mid 的偶偶配对关系判断单一元素在哪边，保持答案区间长度为奇数。把这个特例继续拆开看：每次比较 mid 之后，必须能指出哪一侧整体不可能包含目标，或哪一侧整体已经满足可行性。二分的核心不是 mid 公式，而是被丢弃区间确实不含答案。复盘时把这个特例压成状态字段：left、right、mid、mid 处比较值或检查返回值、答案区间含义、返回值语义；每一轮更新都要保留目标位置或第一个可行值；再用边界 单一元素在开头或结尾、数组长度一、mid 加一越界 反查初始化、更新顺序和退出条件。

\textbf{状态回放。} 手算路线：手算时列出 left、mid、right，以及 mid 处比较结果、可行性检查结果或局部趋势。每一步都要写出被丢弃的一侧为什么不含答案。状态字段：left、right、mid、mid 处比较值或检查返回值、答案区间含义、返回值语义；每一轮更新都要保留目标位置或第一个可行值。状态升级：把逐个试候选改成维护答案所在区间；边界谓词、可行性检查或局部比较给出分支方向，每个分支都必须能排除一侧候选。

\textbf{证明和边界。} 证明从排除一侧开始：答案空间题证明谓词单调，局部比较题证明保留下来的区间仍含目标；不能只靠经验猜测左右边界。边界：单一元素在开头或结尾、数组长度一、mid 加一越界。变体：异或也能求值，但二分利用有序性达到对数时间。

\noindent\textbf{LeetCode 704 二分查找。} 模型：基础题的重点是固定区间语义。推导重点：左闭右开写法能统一 lower bound 和存在性检查。

\textbf{二刷读题。} 读题时先找候选空间和目标边界。本题可以压成“基础题的重点是固定区间语义”。然后把关键观察写成可执行规则：左闭右开写法能统一 lower bound 和存在性检查。这条规则必须能安全排除一侧候选，或收缩到可检查的答案边界。先遮住答案，只保留题面，复述候选对象、合法性条件和输出目标。

\textbf{暴力回放。} 慢版本按顺序扫描下标、逐个试答案值，或直接检查候选直到遇到目标边界。它先保证候选空间覆盖完整，但每次只排除一个候选，浪费了题目给出的顺序、比较或可行性边界。二刷时先写慢版本或口述慢版本，用它确认不会漏答案。

\textbf{特例回放。} 拿 [1,3,5] 找 3 手算，mid 正好命中返回；找 4 时，3 太小丢左侧，5 太大丢右侧，区间为空返回 -1。规律是标准二分每次排除一半，边界写法必须和退出条件匹配。把这个特例继续拆开看：每次比较 mid 之后，必须能指出哪一侧整体不可能包含目标，或哪一侧整体已经满足可行性。二分的核心不是 mid 公式，而是被丢弃区间确实不含答案。复盘时把这个特例压成状态字段：left、right、mid、mid 处比较值或检查返回值、答案区间含义、返回值语义；每一轮更新都要保留目标位置或第一个可行值；再用边界 空数组、长度一、目标在边界、目标不存在 反查初始化、更新顺序和退出条件。

\textbf{状态回放。} 手算路线：手算时列出 left、mid、right，以及 mid 处比较结果、可行性检查结果或局部趋势。每一步都要写出被丢弃的一侧为什么不含答案。状态字段：left、right、mid、mid 处比较值或检查返回值、答案区间含义、返回值语义；每一轮更新都要保留目标位置或第一个可行值。状态升级：把逐个试候选改成维护答案所在区间；边界谓词、可行性检查或局部比较给出分支方向，每个分支都必须能排除一侧候选。

\textbf{证明和边界。} 证明从排除一侧开始：答案空间题证明谓词单调，局部比较题证明保留下来的区间仍含目标；不能只靠经验猜测左右边界。边界：空数组、长度一、目标在边界、目标不存在。变体：所有复杂二分都应该能退回到这道题的边界不变量。

\noindent\textbf{LeetCode 875 爱吃香蕉的珂珂。} 模型：速度越大，总耗时越小，答案具有单调可行性。推导重点：二分速度，检查函数用向上取整累加小时数。

\textbf{二刷读题。} 读题时先找候选空间和目标边界。本题可以压成“速度越大，总耗时越小，答案具有单调可行性”。然后把关键观察写成可执行规则：二分速度，检查函数用向上取整累加小时数。这条规则必须能安全排除一侧候选，或收缩到可检查的答案边界。先遮住答案，只保留题面，复述候选对象、合法性条件和输出目标。

\textbf{暴力回放。} 慢版本按顺序扫描下标、逐个试答案值，或直接检查候选直到遇到目标边界。它先保证候选空间覆盖完整，但每次只排除一个候选，浪费了题目给出的顺序、比较或可行性边界。二刷时先写慢版本或口述慢版本，用它确认不会漏答案。

\textbf{特例回放。} 拿 piles=[3,6,7,11]、h=8 手算，速度 4 时四堆分别需要 1、2、2、3 小时，总计 8 小时，刚好可行；速度 3 会超过 8 小时。规律是速度越大耗时越少，可行性单调，二分最小可行速度。把这个特例继续拆开看：每次比较 mid 之后，必须能指出哪一侧整体不可能包含目标，或哪一侧整体已经满足可行性。二分的核心不是 mid 公式，而是被丢弃区间确实不含答案。复盘时把这个特例压成状态字段：left、right、mid、mid 处比较值或检查返回值、答案区间含义、返回值语义；每一轮更新都要保留目标位置或第一个可行值；再用边界 速度下界、堆很大、h 等于堆数、乘法加法溢出 反查初始化、更新顺序和退出条件。

\textbf{状态回放。} 手算路线：手算时列出 left、mid、right，以及 mid 处比较结果、可行性检查结果或局部趋势。每一步都要写出被丢弃的一侧为什么不含答案。状态字段：left、right、mid、mid 处比较值或检查返回值、答案区间含义、返回值语义；每一轮更新都要保留目标位置或第一个可行值。状态升级：把逐个试候选改成维护答案所在区间；边界谓词、可行性检查或局部比较给出分支方向，每个分支都必须能排除一侧候选。

\textbf{证明和边界。} 证明从排除一侧开始：答案空间题证明谓词单调，局部比较题证明保留下来的区间仍含目标；不能只靠经验猜测左右边界。边界：速度下界、堆很大、h 等于堆数、乘法加法溢出。变体：船运包裹和分割数组都是同类最小可行上限。

\noindent\textbf{LeetCode 981 基于时间的键值存储。} 模型：同一个 key 下的记录按时间戳递增，查询是不超过目标时间的最新值。推导重点：哈希表定位 key 后，在该 key 的时间列表中二分第一个大于目标时间的位置并回退一格。

\textbf{二刷读题。} 读题时先找候选空间和目标边界。本题可以压成“同一个 key 下的记录按时间戳递增，查询是不超过目标时间的最新值”。然后把关键观察写成可执行规则：哈希表定位 key 后，在该 key 的时间列表中二分第一个大于目标时间的位置并回退一格。这条规则必须能安全排除一侧候选，或收缩到可检查的答案边界。先遮住答案，只保留题面，复述候选对象、合法性条件和输出目标。

\textbf{暴力回放。} 慢版本按顺序扫描下标、逐个试答案值，或直接检查候选直到遇到目标边界。它先保证候选空间覆盖完整，但每次只排除一个候选，浪费了题目给出的顺序、比较或可行性边界。二刷时先写慢版本或口述慢版本，用它确认不会漏答案。

\textbf{特例回放。} 拿 key=foo 的记录 (1,bar)、(4,baz)，查询 t=3 手算，答案应是时间不超过 3 的最后一条，也就是 bar。查询 t=4 才是 baz。规律是每个 key 内部时间递增，二分第一个大于查询时间的位置，再回退一格。把这个特例继续拆开看：每次比较 mid 之后，必须能指出哪一侧整体不可能包含目标，或哪一侧整体已经满足可行性。二分的核心不是 mid 公式，而是被丢弃区间确实不含答案。复盘时把这个特例压成状态字段：left、right、mid、mid 处比较值或检查返回值、答案区间含义、返回值语义；每一轮更新都要保留目标位置或第一个可行值；再用边界 key 不存在、查询早于第一条记录、相同 key 多次 set、时间戳递增假设 反查初始化、更新顺序和退出条件。

\textbf{状态回放。} 手算路线：手算时列出 left、mid、right，以及 mid 处比较结果、可行性检查结果或局部趋势。每一步都要写出被丢弃的一侧为什么不含答案。状态字段：left、right、mid、mid 处比较值或检查返回值、答案区间含义、返回值语义；每一轮更新都要保留目标位置或第一个可行值。状态升级：把逐个试候选改成维护答案所在区间；边界谓词、可行性检查或局部比较给出分支方向，每个分支都必须能排除一侧候选。

\textbf{证明和边界。} 证明从排除一侧开始：答案空间题证明谓词单调，局部比较题证明保留下来的区间仍含目标；不能只靠经验猜测左右边界。边界：key 不存在、查询早于第一条记录、相同 key 多次 set、时间戳递增假设。变体：如果时间戳无序到达，必须先排序或改成有序表维护每个 key 的记录。

\noindent\textbf{LeetCode 1011 在 D 天内送达包裹。} 模型：容量越大，需要天数越少，存在最小可行容量。推导重点：下界是最大包裹，上界是总重量，检查函数按顺序贪心装船。

\textbf{二刷读题。} 读题时先找候选空间和目标边界。本题可以压成“容量越大，需要天数越少，存在最小可行容量”。然后把关键观察写成可执行规则：下界是最大包裹，上界是总重量，检查函数按顺序贪心装船。这条规则必须能安全排除一侧候选，或收缩到可检查的答案边界。先遮住答案，只保留题面，复述候选对象、合法性条件和输出目标。

\textbf{暴力回放。} 慢版本按顺序扫描下标、逐个试答案值，或直接检查候选直到遇到目标边界。它先保证候选空间覆盖完整，但每次只排除一个候选，浪费了题目给出的顺序、比较或可行性边界。二刷时先写慢版本或口述慢版本，用它确认不会漏答案。

\textbf{特例回放。} 拿 weights=[1,2,3,4,5]、D=3 手算，容量 6 可以分成 [1,2,3]、[4]、[5] 三天；容量 5 需要更多天。规律是容量越大所需天数越少，可行性单调，二分最小容量。把这个特例继续拆开看：每次比较 mid 之后，必须能指出哪一侧整体不可能包含目标，或哪一侧整体已经满足可行性。二分的核心不是 mid 公式，而是被丢弃区间确实不含答案。复盘时把这个特例压成状态字段：left、right、mid、mid 处比较值或检查返回值、答案区间含义、返回值语义；每一轮更新都要保留目标位置或第一个可行值；再用边界 D 为一、D 等于包裹数、单个包裹超过猜测容量 反查初始化、更新顺序和退出条件。

\textbf{状态回放。} 手算路线：手算时列出 left、mid、right，以及 mid 处比较结果、可行性检查结果或局部趋势。每一步都要写出被丢弃的一侧为什么不含答案。状态字段：left、right、mid、mid 处比较值或检查返回值、答案区间含义、返回值语义；每一轮更新都要保留目标位置或第一个可行值。状态升级：把逐个试候选改成维护答案所在区间；边界谓词、可行性检查或局部比较给出分支方向，每个分支都必须能排除一侧候选。

\textbf{证明和边界。} 证明从排除一侧开始：答案空间题证明谓词单调，局部比较题证明保留下来的区间仍含目标；不能只靠经验猜测左右边界。边界：D 为一、D 等于包裹数、单个包裹超过猜测容量。变体：它和分割数组最大值都是连续分段最大和模型。

\subsection{自测标准}

二刷时不要只确认自己见过这道题。请遮住答案，先讲题面模型，再讲暴力基线和瓶颈，然后说出优化结构保存了什么、不变量如何排除错误候选、边界实验如何打破错误实现。

\section{栈、单调栈与表达式：二刷复盘卡}
这一大节只围绕一个题型复盘，不把题号直接铺成目录。阅读时先看复盘目标，再读核心题卡，最后用自测标准检查自己是否能独立重讲。

\subsection{复盘目标}

追问通常会问栈里为什么存下标、相等元素如何处理、弹出时到底结算了谁。请准备用一个严格递增和一个严格递减样例解释入栈出栈次数。

\subsection{核心题卡}

\noindent\textbf{LeetCode 20 有效的括号。} 模型：括号匹配要求最近打开的左括号先被关闭。推导重点：栈保存尚未匹配的左括号，遇到右括号必须和栈顶类型一致。

\textbf{二刷读题。} 读题时先找未完成候选：哪些对象必须等到右侧、未来字符或后续运算符出现。本题可以压成“括号匹配要求最近打开的左括号先被关闭”，栈只保存这些尚未结算的对象。先遮住答案，只保留题面，复述候选对象、合法性条件和输出目标。

\textbf{暴力回放。} 慢版本让每个位置向左或向右线性寻找边界，或者让每个结构反复等待匹配。它能算出答案，但很多尚未完成的候选会被未来同一个事件一起解决。二刷时先写慢版本或口述慢版本，用它确认不会漏答案。

\textbf{特例回放。} 拿 ([)] 手算：只计数左右括号会误判，因为闭合顺序错了。栈顶必须保存最近的未匹配左括号；遇到 ] 时栈顶若不是 [，立即失败。规律是括号匹配不是数量问题，而是最近未闭合结构的后进先出问题。把这个特例继续拆开看：栈里的对象都是答案尚未确定的候选；一旦当前元素触发弹栈，被弹出的候选必须已经同时拿到了左边界和右边界，未来元素不再有资格改变它的答案。复盘时把这个特例压成状态字段：栈内对象的业务含义、当前输入、弹出时结算的对象、左边界和右边界；栈中存下标还是值要明确；再用边界 空串、以右括号开头、交叉匹配、结束后栈非空 反查初始化、更新顺序和退出条件。

\textbf{状态回放。} 手算路线：手算时记录栈内元素的业务含义，不只记录数值。入栈表示答案未定，弹栈表示当前事件已经让它的答案定下来。状态字段：栈内对象的业务含义、当前输入、弹出时结算的对象、左边界和右边界；栈中存下标还是值要明确。状态升级：把反复寻找最近边界改成维护未完成候选；新元素只和栈顶比较，连续弹出一批已经被当前元素解决的对象。

\textbf{证明和边界。} 证明从弹出时机开始：被弹出的对象在当前时刻答案已经确定，未来元素无法再改变它。边界：空串、以右括号开头、交叉匹配、结束后栈非空。变体：若加入通配符星号，需要额外维护可行区间或双栈。

\noindent\textbf{LeetCode 32 最长有效括号。} 模型：有效括号子串需要知道每段非法边界之后的最早可连接位置。推导重点：栈保存下标，先放入哨兵负一；匹配后用当前下标减栈顶得到长度。

\textbf{二刷读题。} 读题时先找未完成候选：哪些对象必须等到右侧、未来字符或后续运算符出现。本题可以压成“有效括号子串需要知道每段非法边界之后的最早可连接位置”，栈只保存这些尚未结算的对象。先遮住答案，只保留题面，复述候选对象、合法性条件和输出目标。

\textbf{暴力回放。} 慢版本让每个位置向左或向右线性寻找边界，或者让每个结构反复等待匹配。它能算出答案，但很多尚未完成的候选会被未来同一个事件一起解决。二刷时先写慢版本或口述慢版本，用它确认不会漏答案。

\textbf{特例回放。} 拿字符串 )()()) 手算。第一个 ) 不是可匹配对象，而是非法边界；后面每次匹配成功，当前有效长度由当前位置减去栈顶边界得出。若栈空，要把当前右括号作为新非法边界压入。规律是栈里既保存未匹配左括号下标，也保存最近非法边界。把这个特例继续拆开看：栈里的对象都是答案尚未确定的候选；一旦当前元素触发弹栈，被弹出的候选必须已经同时拿到了左边界和右边界，未来元素不再有资格改变它的答案。复盘时把这个特例压成状态字段：栈内对象的业务含义、当前输入、弹出时结算的对象、左边界和右边界；栈中存下标还是值要明确；再用边界 空串、全左括号、全右括号、多个断开的合法段 反查初始化、更新顺序和退出条件。

\textbf{状态回放。} 手算路线：手算时记录栈内元素的业务含义，不只记录数值。入栈表示答案未定，弹栈表示当前事件已经让它的答案定下来。状态字段：栈内对象的业务含义、当前输入、弹出时结算的对象、左边界和右边界；栈中存下标还是值要明确。状态升级：把反复寻找最近边界改成维护未完成候选；新元素只和栈顶比较，连续弹出一批已经被当前元素解决的对象。

\textbf{证明和边界。} 证明从弹出时机开始：被弹出的对象在当前时刻答案已经确定，未来元素无法再改变它。边界：空串、全左括号、全右括号、多个断开的合法段。变体：DP 写法也可做，但栈写法更直接表达最近非法边界。

\noindent\textbf{LeetCode 42 接雨水。} 模型：每个位置的水由左侧最高和右侧最高的较小者决定。推导重点：前后缀、双指针、单调栈都是不同状态维护方式；要能说明各自结算对象。

\textbf{二刷读题。} 读题时先找未完成候选：哪些对象必须等到右侧、未来字符或后续运算符出现。本题可以压成“每个位置的水由左侧最高和右侧最高的较小者决定”，栈只保存这些尚未结算的对象。先遮住答案，只保留题面，复述候选对象、合法性条件和输出目标。

\textbf{暴力回放。} 慢版本让每个位置向左或向右线性寻找边界，或者让每个结构反复等待匹配。它能算出答案，但很多尚未完成的候选会被未来同一个事件一起解决。二刷时先写慢版本或口述慢版本，用它确认不会漏答案。

\textbf{特例回放。} 拿高度 [0,2,0,3] 手算，中间的 0 左侧最高是 2，右侧最高是 3，所以能接 2 格水。若用单调栈，遇到 3 时弹出 0，左边界是 2，右边界是 3，宽度为 1，高度差为 2。规律是每个凹槽必须同时等到左右边界，弹栈那一刻才真正结算水量。把这个特例继续拆开看：栈里的对象都是答案尚未确定的候选；一旦当前元素触发弹栈，被弹出的候选必须已经同时拿到了左边界和右边界，未来元素不再有资格改变它的答案。复盘时把这个特例压成状态字段：栈内对象的业务含义、当前输入、弹出时结算的对象、左边界和右边界；栈中存下标还是值要明确；再用边界 单调递增、单调递减、平台高度、多个凹槽、数组长度不足三 反查初始化、更新顺序和退出条件。

\textbf{状态回放。} 手算路线：手算时记录栈内元素的业务含义，不只记录数值。入栈表示答案未定，弹栈表示当前事件已经让它的答案定下来。状态字段：栈内对象的业务含义、当前输入、弹出时结算的对象、左边界和右边界；栈中存下标还是值要明确。状态升级：把反复寻找最近边界改成维护未完成候选；新元素只和栈顶比较，连续弹出一批已经被当前元素解决的对象。

\textbf{证明和边界。} 证明从弹出时机开始：被弹出的对象在当前时刻答案已经确定，未来元素无法再改变它。边界：单调递增、单调递减、平台高度、多个凹槽、数组长度不足三。变体：若柱子变成二维网格，需要最小堆从边界向内扩展。

\noindent\textbf{LeetCode 84 柱状图中最大的矩形。} 模型：每根柱子作为最矮高度时，需要知道左右第一个更矮位置。推导重点：单调递增栈在遇到更矮柱时结算被弹出柱子的最大宽度。

\textbf{二刷读题。} 读题时先找未完成候选：哪些对象必须等到右侧、未来字符或后续运算符出现。本题可以压成“每根柱子作为最矮高度时，需要知道左右第一个更矮位置”，栈只保存这些尚未结算的对象。先遮住答案，只保留题面，复述候选对象、合法性条件和输出目标。

\textbf{暴力回放。} 慢版本让每个位置向左或向右线性寻找边界，或者让每个结构反复等待匹配。它能算出答案，但很多尚未完成的候选会被未来同一个事件一起解决。二刷时先写慢版本或口述慢版本，用它确认不会漏答案。

\textbf{特例回放。} 拿高度 [2,1,2] 手算。第一个 2 遇到更矮的 1 时，右边界已经确定在当前位置，左边界是弹出后新栈顶的后一位，所以面积可结算。若一直递增，最后需要哨兵触发结算。规律是栈保存递增高度下标，弹出时才知道该柱子的最大扩展宽度。把这个特例继续拆开看：栈里的对象都是答案尚未确定的候选；一旦当前元素触发弹栈，被弹出的候选必须已经同时拿到了左边界和右边界，未来元素不再有资格改变它的答案。复盘时把这个特例压成状态字段：栈内对象的业务含义、当前输入、弹出时结算的对象、左边界和右边界；栈中存下标还是值要明确；再用边界 递增、递减、相等高度、哨兵零、宽度计算 反查初始化、更新顺序和退出条件。

\textbf{状态回放。} 手算路线：手算时记录栈内元素的业务含义，不只记录数值。入栈表示答案未定，弹栈表示当前事件已经让它的答案定下来。状态字段：栈内对象的业务含义、当前输入、弹出时结算的对象、左边界和右边界；栈中存下标还是值要明确。状态升级：把反复寻找最近边界改成维护未完成候选；新元素只和栈顶比较，连续弹出一批已经被当前元素解决的对象。

\textbf{证明和边界。} 证明从弹出时机开始：被弹出的对象在当前时刻答案已经确定，未来元素无法再改变它。边界：递增、递减、相等高度、哨兵零、宽度计算。变体：最大矩形可以作为二维矩阵最大矩形的行压缩子问题。

\noindent\textbf{LeetCode 85 最大矩形。} 模型：每一行都可以把连续的一当成柱状图高度。推导重点：逐行更新每列高度，然后把当前行转成柱状图最大矩形问题。

\textbf{二刷读题。} 读题时先找未完成候选：哪些对象必须等到右侧、未来字符或后续运算符出现。本题可以压成“每一行都可以把连续的一当成柱状图高度”，栈只保存这些尚未结算的对象。先遮住答案，只保留题面，复述候选对象、合法性条件和输出目标。

\textbf{暴力回放。} 慢版本让每个位置向左或向右线性寻找边界，或者让每个结构反复等待匹配。它能算出答案，但很多尚未完成的候选会被未来同一个事件一起解决。二刷时先写慢版本或口述慢版本，用它确认不会漏答案。

\textbf{特例回放。} 拿矩阵每一行向上累计高度手算，某行高度数组若为 [2,1,2]，就退化成柱状图最大矩形。规律是逐行更新每列连续 1 的高度，再对每行高度数组跑单调栈；全零行会把高度清零。把这个特例继续拆开看：栈里的对象都是答案尚未确定的候选；一旦当前元素触发弹栈，被弹出的候选必须已经同时拿到了左边界和右边界，未来元素不再有资格改变它的答案。复盘时把这个特例压成状态字段：栈内对象的业务含义、当前输入、弹出时结算的对象、左边界和右边界；栈中存下标还是值要明确；再用边界 全零、全一、单行、单列、宽度和高度结算 反查初始化、更新顺序和退出条件。

\textbf{状态回放。} 手算路线：手算时记录栈内元素的业务含义，不只记录数值。入栈表示答案未定，弹栈表示当前事件已经让它的答案定下来。状态字段：栈内对象的业务含义、当前输入、弹出时结算的对象、左边界和右边界；栈中存下标还是值要明确。状态升级：把反复寻找最近边界改成维护未完成候选；新元素只和栈顶比较，连续弹出一批已经被当前元素解决的对象。

\textbf{证明和边界。} 证明从弹出时机开始：被弹出的对象在当前时刻答案已经确定，未来元素无法再改变它。边界：全零、全一、单行、单列、宽度和高度结算。变体：这是二维问题压成一维单调栈的典型做法。

\noindent\textbf{LeetCode 150 逆波兰表达式求值。} 模型：后缀表达式把优先级和括号都编码进 token 顺序。推导重点：遇到数字入栈，遇到运算符弹出右操作数和左操作数计算。

\textbf{二刷读题。} 读题时先找未完成候选：哪些对象必须等到右侧、未来字符或后续运算符出现。本题可以压成“后缀表达式把优先级和括号都编码进 token 顺序”，栈只保存这些尚未结算的对象。先遮住答案，只保留题面，复述候选对象、合法性条件和输出目标。

\textbf{暴力回放。} 慢版本让每个位置向左或向右线性寻找边界，或者让每个结构反复等待匹配。它能算出答案，但很多尚未完成的候选会被未来同一个事件一起解决。二刷时先写慢版本或口述慢版本，用它确认不会漏答案。

\textbf{特例回放。} 拿 tokens=[2,1,+,3,*] 手算，2 和 1 入栈，遇到 + 弹出右操作数 1 和左操作数 2 得 3，再与 3 相乘得 9。规律是后缀表达式遇到运算符时，栈顶两个数已经是它的完整左右操作数，减法和除法顺序不能反。把这个特例继续拆开看：栈里的对象都是答案尚未确定的候选；一旦当前元素触发弹栈，被弹出的候选必须已经同时拿到了左边界和右边界，未来元素不再有资格改变它的答案。复盘时把这个特例压成状态字段：栈内对象的业务含义、当前输入、弹出时结算的对象、左边界和右边界；栈中存下标还是值要明确；再用边界 负数 token、多位数、除法向零截断、减法顺序 反查初始化、更新顺序和退出条件。

\textbf{状态回放。} 手算路线：手算时记录栈内元素的业务含义，不只记录数值。入栈表示答案未定，弹栈表示当前事件已经让它的答案定下来。状态字段：栈内对象的业务含义、当前输入、弹出时结算的对象、左边界和右边界；栈中存下标还是值要明确。状态升级：把反复寻找最近边界改成维护未完成候选；新元素只和栈顶比较，连续弹出一批已经被当前元素解决的对象。

\textbf{证明和边界。} 证明从弹出时机开始：被弹出的对象在当前时刻答案已经确定，未来元素无法再改变它。边界：负数 token、多位数、除法向零截断、减法顺序。变体：中缀表达式求值需要另一个运算符栈处理优先级。

\noindent\textbf{LeetCode 155 最小栈。} 模型：普通栈只能看到栈顶，不能常数时间知道历史最小值。推导重点：辅助栈同步保存每个深度对应的最小值。

\textbf{二刷读题。} 读题时先找未完成候选：哪些对象必须等到右侧、未来字符或后续运算符出现。本题可以压成“普通栈只能看到栈顶，不能常数时间知道历史最小值”，栈只保存这些尚未结算的对象。先遮住答案，只保留题面，复述候选对象、合法性条件和输出目标。

\textbf{暴力回放。} 慢版本让每个位置向左或向右线性寻找边界，或者让每个结构反复等待匹配。它能算出答案，但很多尚未完成的候选会被未来同一个事件一起解决。二刷时先写慢版本或口述慢版本，用它确认不会漏答案。

\textbf{特例回放。} 拿 push 2、push 0、push 3、push 0 手算，当前最小值序列是 2、0、0、0；弹出最后一个 0 后，最小值仍是 0。规律是辅助栈要在每个深度同步保存当前最小值，重复最小值不能只保存一次。把这个特例继续拆开看：栈里的对象都是答案尚未确定的候选；一旦当前元素触发弹栈，被弹出的候选必须已经同时拿到了左边界和右边界，未来元素不再有资格改变它的答案。复盘时把这个特例压成状态字段：栈内对象的业务含义、当前输入、弹出时结算的对象、左边界和右边界；栈中存下标还是值要明确；再用边界 重复最小值、弹出最小值、空栈操作约定 反查初始化、更新顺序和退出条件。

\textbf{状态回放。} 手算路线：手算时记录栈内元素的业务含义，不只记录数值。入栈表示答案未定，弹栈表示当前事件已经让它的答案定下来。状态字段：栈内对象的业务含义、当前输入、弹出时结算的对象、左边界和右边界；栈中存下标还是值要明确。状态升级：把反复寻找最近边界改成维护未完成候选；新元素只和栈顶比较，连续弹出一批已经被当前元素解决的对象。

\textbf{证明和边界。} 证明从弹出时机开始：被弹出的对象在当前时刻答案已经确定，未来元素无法再改变它。边界：重复最小值、弹出最小值、空栈操作约定。变体：也可保存差值压缩空间，但可读性和溢出处理更复杂。

\noindent\textbf{LeetCode 224 基本计算器。} 模型：括号会开启新的表达式上下文，括号结束时要回到上一层。推导重点：栈保存进入括号前的结果和符号，当前层算完后乘以上层符号并合并。

\textbf{二刷读题。} 读题时先找未完成候选：哪些对象必须等到右侧、未来字符或后续运算符出现。本题可以压成“括号会开启新的表达式上下文，括号结束时要回到上一层”，栈只保存这些尚未结算的对象。先遮住答案，只保留题面，复述候选对象、合法性条件和输出目标。

\textbf{暴力回放。} 慢版本让每个位置向左或向右线性寻找边界，或者让每个结构反复等待匹配。它能算出答案，但很多尚未完成的候选会被未来同一个事件一起解决。二刷时先写慢版本或口述慢版本，用它确认不会漏答案。

\textbf{特例回放。} 拿 1-(2+3) 手算，进入括号前当前结果是 1，符号是 -；括号内算出 5 后要乘以上层符号并合并成 -4。规律是栈保存进入括号前的结果和符号，一元负号和空格要按字符流处理。把这个特例继续拆开看：栈里的对象都是答案尚未确定的候选；一旦当前元素触发弹栈，被弹出的候选必须已经同时拿到了左边界和右边界，未来元素不再有资格改变它的答案。复盘时把这个特例压成状态字段：栈内对象的业务含义、当前输入、弹出时结算的对象、左边界和右边界；栈中存下标还是值要明确；再用边界 多位数、前导空格、一元负号、嵌套括号 反查初始化、更新顺序和退出条件。

\textbf{状态回放。} 手算路线：手算时记录栈内元素的业务含义，不只记录数值。入栈表示答案未定，弹栈表示当前事件已经让它的答案定下来。状态字段：栈内对象的业务含义、当前输入、弹出时结算的对象、左边界和右边界；栈中存下标还是值要明确。状态升级：把反复寻找最近边界改成维护未完成候选；新元素只和栈顶比较，连续弹出一批已经被当前元素解决的对象。

\textbf{证明和边界。} 证明从弹出时机开始：被弹出的对象在当前时刻答案已经确定，未来元素无法再改变它。边界：多位数、前导空格、一元负号、嵌套括号。变体：基本计算器 II 没有括号但有乘除优先级，需要不同的栈语义。

\noindent\textbf{LeetCode 227 基本计算器 II。} 模型：乘除优先级高于加减，可以在读到下一个运算符时结算上一项。推导重点：栈保存已经确定符号的项，乘除直接修改栈顶，加减压入正负项。

\textbf{二刷读题。} 读题时先找未完成候选：哪些对象必须等到右侧、未来字符或后续运算符出现。本题可以压成“乘除优先级高于加减，可以在读到下一个运算符时结算上一项”，栈只保存这些尚未结算的对象。先遮住答案，只保留题面，复述候选对象、合法性条件和输出目标。

\textbf{暴力回放。} 慢版本让每个位置向左或向右线性寻找边界，或者让每个结构反复等待匹配。它能算出答案，但很多尚未完成的候选会被未来同一个事件一起解决。二刷时先写慢版本或口述慢版本，用它确认不会漏答案。

\textbf{特例回放。} 拿 3+2*2 手算，加号前的 3 可以先入栈，乘号遇到 2 时要立刻把栈顶 2 乘以下一个数 2，而不能等到最后按顺序相加。规律是栈保存已确定符号的项，乘除立即修改栈顶，加减压入正负项。把这个特例继续拆开看：栈里的对象都是答案尚未确定的候选；一旦当前元素触发弹栈，被弹出的候选必须已经同时拿到了左边界和右边界，未来元素不再有资格改变它的答案。复盘时把这个特例压成状态字段：栈内对象的业务含义、当前输入、弹出时结算的对象、左边界和右边界；栈中存下标还是值要明确；再用边界 多位数、空格、除法向零截断、最后一个数字需要结算 反查初始化、更新顺序和退出条件。

\textbf{状态回放。} 手算路线：手算时记录栈内元素的业务含义，不只记录数值。入栈表示答案未定，弹栈表示当前事件已经让它的答案定下来。状态字段：栈内对象的业务含义、当前输入、弹出时结算的对象、左边界和右边界；栈中存下标还是值要明确。状态升级：把反复寻找最近边界改成维护未完成候选；新元素只和栈顶比较，连续弹出一批已经被当前元素解决的对象。

\textbf{证明和边界。} 证明从弹出时机开始：被弹出的对象在当前时刻答案已经确定，未来元素无法再改变它。边界：多位数、空格、除法向零截断、最后一个数字需要结算。变体：若加入括号，需要再引入上下文栈或递归下降解析。

\noindent\textbf{LeetCode 232 用栈实现队列。} 模型：两个栈分别负责输入顺序和输出顺序。推导重点：只有输出栈为空时，才把输入栈整体倒过去，保证均摊常数。

\textbf{二刷读题。} 读题时先找未完成候选：哪些对象必须等到右侧、未来字符或后续运算符出现。本题可以压成“两个栈分别负责输入顺序和输出顺序”，栈只保存这些尚未结算的对象。先遮住答案，只保留题面，复述候选对象、合法性条件和输出目标。

\textbf{暴力回放。} 慢版本让每个位置向左或向右线性寻找边界，或者让每个结构反复等待匹配。它能算出答案，但很多尚未完成的候选会被未来同一个事件一起解决。二刷时先写慢版本或口述慢版本，用它确认不会漏答案。

\textbf{特例回放。} 拿 push 1、push 2、pop 手算，如果输出栈为空，把输入栈整体倒过去得到栈顶 1，正好符合队列先进先出。规律是只有输出栈为空才搬运，单个元素最多进出两个栈各一次，所以均摊常数。把这个特例继续拆开看：栈里的对象都是答案尚未确定的候选；一旦当前元素触发弹栈，被弹出的候选必须已经同时拿到了左边界和右边界，未来元素不再有资格改变它的答案。复盘时把这个特例压成状态字段：栈内对象的业务含义、当前输入、弹出时结算的对象、左边界和右边界；栈中存下标还是值要明确；再用边界 连续 peek、空队列、交替 push pop 反查初始化、更新顺序和退出条件。

\textbf{状态回放。} 手算路线：手算时记录栈内元素的业务含义，不只记录数值。入栈表示答案未定，弹栈表示当前事件已经让它的答案定下来。状态字段：栈内对象的业务含义、当前输入、弹出时结算的对象、左边界和右边界；栈中存下标还是值要明确。状态升级：把反复寻找最近边界改成维护未完成候选；新元素只和栈顶比较，连续弹出一批已经被当前元素解决的对象。

\textbf{证明和边界。} 证明从弹出时机开始：被弹出的对象在当前时刻答案已经确定，未来元素无法再改变它。边界：连续 peek、空队列、交替 push pop。变体：用队列实现栈则要通过轮转保持最近元素在队首。

\noindent\textbf{LeetCode 239 滑动窗口最大值。} 模型：每个窗口最大值来自仍未过期且没有被新元素支配的候选。推导重点：单调队列保存下标，队尾小于等于新值的候选被弹出。

\textbf{二刷读题。} 读题时先找未完成候选：哪些对象必须等到右侧、未来字符或后续运算符出现。本题可以压成“每个窗口最大值来自仍未过期且没有被新元素支配的候选”，栈只保存这些尚未结算的对象。先遮住答案，只保留题面，复述候选对象、合法性条件和输出目标。

\textbf{暴力回放。} 慢版本让每个位置向左或向右线性寻找边界，或者让每个结构反复等待匹配。它能算出答案，但很多尚未完成的候选会被未来同一个事件一起解决。二刷时先写慢版本或口述慢版本，用它确认不会漏答案。

\textbf{特例回放。} 拿 [1,3,-1]、窗口 3 手算，队尾的 1 在 3 进入后永远不可能成为最大值，因为 3 更新且更靠右。于是 1 可以弹出。规律是单调队列保存可能成为当前或未来窗口最大值的下标，过期下标从队首清理。把这个特例继续拆开看：栈里的对象都是答案尚未确定的候选；一旦当前元素触发弹栈，被弹出的候选必须已经同时拿到了左边界和右边界，未来元素不再有资格改变它的答案。复盘时把这个特例压成状态字段：栈内对象的业务含义、当前输入、弹出时结算的对象、左边界和右边界；栈中存下标还是值要明确；再用边界 k 为一、k 等于数组长度、重复最大值、队首过期 反查初始化、更新顺序和退出条件。

\textbf{状态回放。} 手算路线：手算时记录栈内元素的业务含义，不只记录数值。入栈表示答案未定，弹栈表示当前事件已经让它的答案定下来。状态字段：栈内对象的业务含义、当前输入、弹出时结算的对象、左边界和右边界；栈中存下标还是值要明确。状态升级：把反复寻找最近边界改成维护未完成候选；新元素只和栈顶比较，连续弹出一批已经被当前元素解决的对象。

\textbf{证明和边界。} 证明从弹出时机开始：被弹出的对象在当前时刻答案已经确定，未来元素无法再改变它。边界：k 为一、k 等于数组长度、重复最大值、队首过期。变体：受限子序列和使用同样队列维护最近 k 个 DP 最大值。

\noindent\textbf{LeetCode 394 字符串解码。} 模型：嵌套结构需要保存上一层字符串和重复次数。推导重点：遇到左括号入栈当前状态，右括号弹出并拼接展开。

\textbf{二刷读题。} 读题时先找未完成候选：哪些对象必须等到右侧、未来字符或后续运算符出现。本题可以压成“嵌套结构需要保存上一层字符串和重复次数”，栈只保存这些尚未结算的对象。先遮住答案，只保留题面，复述候选对象、合法性条件和输出目标。

\textbf{暴力回放。} 慢版本让每个位置向左或向右线性寻找边界，或者让每个结构反复等待匹配。它能算出答案，但很多尚未完成的候选会被未来同一个事件一起解决。二刷时先写慢版本或口述慢版本，用它确认不会漏答案。

\textbf{特例回放。} 拿 3[a2[c]] 手算，遇到左括号前保存当前字符串和重复次数；遇到右括号时弹出上一层，把当前 cc 重复 2 次拼成 acc，再被外层重复 3 次。规律是栈保存进入括号前的上下文，右括号触发一次完整展开。把这个特例继续拆开看：栈里的对象都是答案尚未确定的候选；一旦当前元素触发弹栈，被弹出的候选必须已经同时拿到了左边界和右边界，未来元素不再有资格改变它的答案。复盘时把这个特例压成状态字段：栈内对象的业务含义、当前输入、弹出时结算的对象、左边界和右边界；栈中存下标还是值要明确；再用边界 多位数字、嵌套、空片段、数字后必须跟括号 反查初始化、更新顺序和退出条件。

\textbf{状态回放。} 手算路线：手算时记录栈内元素的业务含义，不只记录数值。入栈表示答案未定，弹栈表示当前事件已经让它的答案定下来。状态字段：栈内对象的业务含义、当前输入、弹出时结算的对象、左边界和右边界；栈中存下标还是值要明确。状态升级：把反复寻找最近边界改成维护未完成候选；新元素只和栈顶比较，连续弹出一批已经被当前元素解决的对象。

\textbf{证明和边界。} 证明从弹出时机开始：被弹出的对象在当前时刻答案已经确定，未来元素无法再改变它。边界：多位数字、嵌套、空片段、数字后必须跟括号。变体：递归下降解析更通用，但迭代栈更符合工程栈控制。

\noindent\textbf{LeetCode 402 移掉 K 位数字。} 模型：要删除 k 位使剩余数字最小，本质是让高位尽量小。推导重点：单调递增栈中遇到更小数字时，弹出前面较大的数字并消耗删除次数。

\textbf{二刷读题。} 读题时先找未完成候选：哪些对象必须等到右侧、未来字符或后续运算符出现。本题可以压成“要删除 k 位使剩余数字最小，本质是让高位尽量小”，栈只保存这些尚未结算的对象。先遮住答案，只保留题面，复述候选对象、合法性条件和输出目标。

\textbf{暴力回放。} 慢版本让每个位置向左或向右线性寻找边界，或者让每个结构反复等待匹配。它能算出答案，但很多尚未完成的候选会被未来同一个事件一起解决。二刷时先写慢版本或口述慢版本，用它确认不会漏答案。

\textbf{特例回放。} 拿 1432219、k=3 手算，遇到 3 后面的 2 时，前面的 4 和 3 都比 2 大且可删，弹出能让高位更小；最后去掉前导零。规律是单调递增栈让更小数字尽量靠前，k 未用完时从尾部继续删。把这个特例继续拆开看：栈里的对象都是答案尚未确定的候选；一旦当前元素触发弹栈，被弹出的候选必须已经同时拿到了左边界和右边界，未来元素不再有资格改变它的答案。复盘时把这个特例压成状态字段：栈内对象的业务含义、当前输入、弹出时结算的对象、左边界和右边界；栈中存下标还是值要明确；再用边界 前导零、k 未用完、k 等于长度、数字已经递增 反查初始化、更新顺序和退出条件。

\textbf{状态回放。} 手算路线：手算时记录栈内元素的业务含义，不只记录数值。入栈表示答案未定，弹栈表示当前事件已经让它的答案定下来。状态字段：栈内对象的业务含义、当前输入、弹出时结算的对象、左边界和右边界；栈中存下标还是值要明确。状态升级：把反复寻找最近边界改成维护未完成候选；新元素只和栈顶比较，连续弹出一批已经被当前元素解决的对象。

\textbf{证明和边界。} 证明从弹出时机开始：被弹出的对象在当前时刻答案已经确定，未来元素无法再改变它。边界：前导零、k 未用完、k 等于长度、数字已经递增。变体：去除重复字母也用单调栈，但还要保证每个字符保留一次。

\noindent\textbf{LeetCode 496 下一个更大元素 I。} 模型：第二个数组提供每个值右侧第一个更大值的全局关系。推导重点：单调栈预处理值到下一个更大值映射，再回答第一个数组查询。

\textbf{二刷读题。} 读题时先找未完成候选：哪些对象必须等到右侧、未来字符或后续运算符出现。本题可以压成“第二个数组提供每个值右侧第一个更大值的全局关系”，栈只保存这些尚未结算的对象。先遮住答案，只保留题面，复述候选对象、合法性条件和输出目标。

\textbf{暴力回放。} 慢版本让每个位置向左或向右线性寻找边界，或者让每个结构反复等待匹配。它能算出答案，但很多尚未完成的候选会被未来同一个事件一起解决。二刷时先写慢版本或口述慢版本，用它确认不会漏答案。

\textbf{特例回放。} 拿 nums2=[1,3,4,2] 手算，1 的下一个更大是 3，3 的下一个更大是 4，4 没有。规律是单调栈在扫描 nums2 时建立值到下一个更大值的映射，再回答 nums1 查询。把这个特例继续拆开看：栈里的对象都是答案尚未确定的候选；一旦当前元素触发弹栈，被弹出的候选必须已经同时拿到了左边界和右边界，未来元素不再有资格改变它的答案。复盘时把这个特例压成状态字段：栈内对象的业务含义、当前输入、弹出时结算的对象、左边界和右边界；栈中存下标还是值要明确；再用边界 不存在更大值、递减数组、值唯一假设 反查初始化、更新顺序和退出条件。

\textbf{状态回放。} 手算路线：手算时记录栈内元素的业务含义，不只记录数值。入栈表示答案未定，弹栈表示当前事件已经让它的答案定下来。状态字段：栈内对象的业务含义、当前输入、弹出时结算的对象、左边界和右边界；栈中存下标还是值要明确。状态升级：把反复寻找最近边界改成维护未完成候选；新元素只和栈顶比较，连续弹出一批已经被当前元素解决的对象。

\textbf{证明和边界。} 证明从弹出时机开始：被弹出的对象在当前时刻答案已经确定，未来元素无法再改变它。边界：不存在更大值、递减数组、值唯一假设。变体：若值不唯一，应按下标而不是值建映射。

\noindent\textbf{LeetCode 503 下一个更大元素 II。} 模型：环形数组可以通过遍历两倍下标模拟绕回。推导重点：第一遍入栈，第二遍只帮助未解决下标找到答案。

\textbf{二刷读题。} 读题时先找未完成候选：哪些对象必须等到右侧、未来字符或后续运算符出现。本题可以压成“环形数组可以通过遍历两倍下标模拟绕回”，栈只保存这些尚未结算的对象。先遮住答案，只保留题面，复述候选对象、合法性条件和输出目标。

\textbf{暴力回放。} 慢版本让每个位置向左或向右线性寻找边界，或者让每个结构反复等待匹配。它能算出答案，但很多尚未完成的候选会被未来同一个事件一起解决。二刷时先写慢版本或口述慢版本，用它确认不会漏答案。

\textbf{特例回放。} 拿 [1,2,1] 手算，第一个 1 的下一个更大是 2，最后一个 1 环形回到前面也能找到 2。规律是扫描两遍模拟环形，第一遍入栈，第二遍只帮助未解决下标结算。把这个特例继续拆开看：栈里的对象都是答案尚未确定的候选；一旦当前元素触发弹栈，被弹出的候选必须已经同时拿到了左边界和右边界，未来元素不再有资格改变它的答案。复盘时把这个特例压成状态字段：栈内对象的业务含义、当前输入、弹出时结算的对象、左边界和右边界；栈中存下标还是值要明确；再用边界 全递减、全相等、长度一、取模下标 反查初始化、更新顺序和退出条件。

\textbf{状态回放。} 手算路线：手算时记录栈内元素的业务含义，不只记录数值。入栈表示答案未定，弹栈表示当前事件已经让它的答案定下来。状态字段：栈内对象的业务含义、当前输入、弹出时结算的对象、左边界和右边界；栈中存下标还是值要明确。状态升级：把反复寻找最近边界改成维护未完成候选；新元素只和栈顶比较，连续弹出一批已经被当前元素解决的对象。

\textbf{证明和边界。} 证明从弹出时机开始：被弹出的对象在当前时刻答案已经确定，未来元素无法再改变它。边界：全递减、全相等、长度一、取模下标。变体：不要在第二遍重复入栈，否则可能死循环或重复计算。

\noindent\textbf{LeetCode 739 每日温度。} 模型：每一天等待右侧第一个更高温度。推导重点：单调递减栈保存尚未找到答案的下标。

\textbf{二刷读题。} 读题时先找未完成候选：哪些对象必须等到右侧、未来字符或后续运算符出现。本题可以压成“每一天等待右侧第一个更高温度”，栈只保存这些尚未结算的对象。先遮住答案，只保留题面，复述候选对象、合法性条件和输出目标。

\textbf{暴力回放。} 慢版本让每个位置向左或向右线性寻找边界，或者让每个结构反复等待匹配。它能算出答案，但很多尚未完成的候选会被未来同一个事件一起解决。二刷时先写慢版本或口述慢版本，用它确认不会漏答案。

\textbf{特例回放。} 拿 [73,74,75,71,69,72,76,73] 手算，73 遇到 74 时等待 1 天；71 要等到 72 才解决。规律是单调递减栈保存尚未等到更高温的下标，遇到更高温时连续弹栈结算天数。把这个特例继续拆开看：栈里的对象都是答案尚未确定的候选；一旦当前元素触发弹栈，被弹出的候选必须已经同时拿到了左边界和右边界，未来元素不再有资格改变它的答案。复盘时把这个特例压成状态字段：栈内对象的业务含义、当前输入、弹出时结算的对象、左边界和右边界；栈中存下标还是值要明确；再用边界 没有更高温、相等温度、递增、递减 反查初始化、更新顺序和退出条件。

\textbf{状态回放。} 手算路线：手算时记录栈内元素的业务含义，不只记录数值。入栈表示答案未定，弹栈表示当前事件已经让它的答案定下来。状态字段：栈内对象的业务含义、当前输入、弹出时结算的对象、左边界和右边界；栈中存下标还是值要明确。状态升级：把反复寻找最近边界改成维护未完成候选；新元素只和栈顶比较，连续弹出一批已经被当前元素解决的对象。

\textbf{证明和边界。} 证明从弹出时机开始：被弹出的对象在当前时刻答案已经确定，未来元素无法再改变它。边界：没有更高温、相等温度、递增、递减。变体：下一个更大元素与它同型，只是输出值或距离不同。

\noindent\textbf{LeetCode 895 最大频率栈。} 模型：弹出时既要频率最高，又要在同频中最近入栈。推导重点：哈希表记录值频率，频率到栈记录达到该频率的入栈顺序，维护当前最大频率。

\textbf{二刷读题。} 读题时先找未完成候选：哪些对象必须等到右侧、未来字符或后续运算符出现。本题可以压成“弹出时既要频率最高，又要在同频中最近入栈”，栈只保存这些尚未结算的对象。先遮住答案，只保留题面，复述候选对象、合法性条件和输出目标。

\textbf{暴力回放。} 慢版本让每个位置向左或向右线性寻找边界，或者让每个结构反复等待匹配。它能算出答案，但很多尚未完成的候选会被未来同一个事件一起解决。二刷时先写慢版本或口述慢版本，用它确认不会漏答案。

\textbf{特例回放。} 拿 push 5,7,5,7,4,5 后 pop 手算，频率最高是 5，先弹 5；下一次 5 和 7 同频，弹最近达到该频率的 7。规律是值到频率、频率到栈、当前最大频率三者共同维护。把这个特例继续拆开看：栈里的对象都是答案尚未确定的候选；一旦当前元素触发弹栈，被弹出的候选必须已经同时拿到了左边界和右边界，未来元素不再有资格改变它的答案。复盘时把这个特例压成状态字段：栈内对象的业务含义、当前输入、弹出时结算的对象、左边界和右边界；栈中存下标还是值要明确；再用边界 重复 push、pop 后最大频率下降、同频最近性、空栈不在题意内 反查初始化、更新顺序和退出条件。

\textbf{状态回放。} 手算路线：手算时记录栈内元素的业务含义，不只记录数值。入栈表示答案未定，弹栈表示当前事件已经让它的答案定下来。状态字段：栈内对象的业务含义、当前输入、弹出时结算的对象、左边界和右边界；栈中存下标还是值要明确。状态升级：把反复寻找最近边界改成维护未完成候选；新元素只和栈顶比较，连续弹出一批已经被当前元素解决的对象。

\textbf{证明和边界。} 证明从弹出时机开始：被弹出的对象在当前时刻答案已经确定，未来元素无法再改变它。边界：重复 push、pop 后最大频率下降、同频最近性、空栈不在题意内。变体：它是栈和哈希表按频率维度组合的设计题。

\subsection{自测标准}

二刷时不要只确认自己见过这道题。请遮住答案，先讲题面模型，再讲暴力基线和瓶颈，然后说出优化结构保存了什么、不变量如何排除错误候选、边界实验如何打破错误实现。

\section{堆与动态极值：二刷复盘卡}
这一大节只围绕一个题型复盘，不把题号直接铺成目录。阅读时先看复盘目标，再读核心题卡，最后用自测标准检查自己是否能独立重讲。

\subsection{复盘目标}

追问通常会问为什么不完整排序、堆顶是否可能过期、比较器方向是否写反。请准备说明堆里元素的业务含义，而不是只说 priority\_queue。

\subsection{核心题卡}

\noindent\textbf{LeetCode 23 合并 K 个升序链表。} 模型：每条链表当前头节点都是可能的下一个最小元素。推导重点：小顶堆保存每条非空链的当前头，弹出最小节点后把它的后继入堆。

\textbf{二刷读题。} 读题时先找动态候选集合以及每一步需要的极值。本题可以压成“每条链表当前头节点都是可能的下一个最小元素”，堆顶必须正好回答下一步选择。先遮住答案，只保留题面，复述候选对象、合法性条件和输出目标。

\textbf{暴力回放。} 慢版本每轮扫描全部候选来找最大或最小，再更新候选集合。它的答案选择过程清楚，但动态集合每变化一次就重新找极值，代价太高。二刷时先写慢版本或口述慢版本，用它确认不会漏答案。

\textbf{特例回放。} 拿三条链表头 1、1、2 手算，每次答案只需要当前最小头节点；弹出某条链的头后，只有它的后继会成为新候选。暴力每轮扫 k 个头，堆把“找当前最小头”变成对数操作。规律是堆里始终最多保存每条非空链的一个当前头。把这个特例继续拆开看：堆顶必须正好代表下一步最需要处理的候选；若堆里可能有过期对象，读取堆顶前要先清理。堆只解决动态极值，不自动证明被弹出或被丢弃的元素无用。复盘时把这个特例压成状态字段：堆元素字段、堆顶比较规则、候选是否过期、答案读取时机；如果需要懒删除，还要保存删除计数；再用边界 空链表数组、某条链为空、重复值、堆里保存节点指针还是值 反查初始化、更新顺序和退出条件。

\textbf{状态回放。} 手算路线：手算时记录堆顶、候选集合和是否有过期元素。每次使用堆顶前都要确认它仍然属于当前问题。状态字段：堆元素字段、堆顶比较规则、候选是否过期、答案读取时机；如果需要懒删除，还要保存删除计数。状态升级：把每轮全量扫描改成堆顶查询；插入、弹出和过期清理共同维护动态候选集。

\textbf{证明和边界。} 证明从候选覆盖开始：所有可能成为下一步答案的对象都在堆中，堆顶过期时不会被使用。边界：空链表数组、某条链为空、重复值、堆里保存节点指针还是值。变体：也可以分治两两合并，堆更适合 K 较大且链长不均的场景。

\noindent\textbf{LeetCode 215 数组中的第 K 个最大元素。} 模型：只关心第 K 大，不需要完整排序。推导重点：大小为 K 的小顶堆保存当前前 K 大边界；快速选择可做到平均线性。

\textbf{二刷读题。} 读题时先找动态候选集合以及每一步需要的极值。本题可以压成“只关心第 K 大，不需要完整排序”，堆顶必须正好回答下一步选择。先遮住答案，只保留题面，复述候选对象、合法性条件和输出目标。

\textbf{暴力回放。} 慢版本每轮扫描全部候选来找最大或最小，再更新候选集合。它的答案选择过程清楚，但动态集合每变化一次就重新找极值，代价太高。二刷时先写慢版本或口述慢版本，用它确认不会漏答案。

\textbf{特例回放。} 拿 [3,2,1,5,6,4]、k=2 手算，固定大小为 2 的小顶堆最终保存最大的两个数，堆顶就是第 2 大。规律是堆里只保留当前前 k 大候选，比堆顶小的数不可能进入答案。把这个特例继续拆开看：堆顶必须正好代表下一步最需要处理的候选；若堆里可能有过期对象，读取堆顶前要先清理。堆只解决动态极值，不自动证明被弹出或被丢弃的元素无用。复盘时把这个特例压成状态字段：堆元素字段、堆顶比较规则、候选是否过期、答案读取时机；如果需要懒删除，还要保存删除计数；再用边界 K 为一、K 为 n、重复值、随机 pivot 反查初始化、更新顺序和退出条件。

\textbf{状态回放。} 手算路线：手算时记录堆顶、候选集合和是否有过期元素。每次使用堆顶前都要确认它仍然属于当前问题。状态字段：堆元素字段、堆顶比较规则、候选是否过期、答案读取时机；如果需要懒删除，还要保存删除计数。状态升级：把每轮全量扫描改成堆顶查询；插入、弹出和过期清理共同维护动态候选集。

\textbf{证明和边界。} 证明从候选覆盖开始：所有可能成为下一步答案的对象都在堆中，堆顶过期时不会被使用。边界：K 为一、K 为 n、重复值、随机 pivot。变体：若数据流持续到来，堆方案比一次性快速选择更自然。

\noindent\textbf{LeetCode 295 数据流的中位数。} 模型：数据流需要动态维护较小一半和较大一半。推导重点：大顶堆保存较小一半，小顶堆保存较大一半，大小差不超过一。

\textbf{二刷读题。} 读题时先找动态候选集合以及每一步需要的极值。本题可以压成“数据流需要动态维护较小一半和较大一半”，堆顶必须正好回答下一步选择。先遮住答案，只保留题面，复述候选对象、合法性条件和输出目标。

\textbf{暴力回放。} 慢版本每轮扫描全部候选来找最大或最小，再更新候选集合。它的答案选择过程清楚，但动态集合每变化一次就重新找极值，代价太高。二刷时先写慢版本或口述慢版本，用它确认不会漏答案。

\textbf{特例回放。} 拿数据流 1、2、3 手算，插入 1 后中位数 1；插入 2 后两半是 [1] 和 [2]，中位数 1.5；插入 3 后右半多一个，要平衡成左半 [2,1]、右半 [3]。规律是大顶堆保存较小一半，小顶堆保存较大一半，大小差最多 1。把这个特例继续拆开看：堆顶必须正好代表下一步最需要处理的候选；若堆里可能有过期对象，读取堆顶前要先清理。堆只解决动态极值，不自动证明被弹出或被丢弃的元素无用。复盘时把这个特例压成状态字段：堆元素字段、堆顶比较规则、候选是否过期、答案读取时机；如果需要懒删除，还要保存删除计数；再用边界 偶数个元素、负数、重复值、平均值溢出 反查初始化、更新顺序和退出条件。

\textbf{状态回放。} 手算路线：手算时记录堆顶、候选集合和是否有过期元素。每次使用堆顶前都要确认它仍然属于当前问题。状态字段：堆元素字段、堆顶比较规则、候选是否过期、答案读取时机；如果需要懒删除，还要保存删除计数。状态升级：把每轮全量扫描改成堆顶查询；插入、弹出和过期清理共同维护动态候选集。

\textbf{证明和边界。} 证明从候选覆盖开始：所有可能成为下一步答案的对象都在堆中，堆顶过期时不会被使用。边界：偶数个元素、负数、重复值、平均值溢出。变体：若还要删除任意元素，需要延迟删除哈希表。

\noindent\textbf{LeetCode 347 前 K 个高频元素。} 模型：先统计频次，再只维护前 K 个候选。推导重点：小顶堆大小超过 K 就弹出最低频；桶排序利用频次上界。

\textbf{二刷读题。} 读题时先找动态候选集合以及每一步需要的极值。本题可以压成“先统计频次，再只维护前 K 个候选”，堆顶必须正好回答下一步选择。先遮住答案，只保留题面，复述候选对象、合法性条件和输出目标。

\textbf{暴力回放。} 慢版本每轮扫描全部候选来找最大或最小，再更新候选集合。它的答案选择过程清楚，但动态集合每变化一次就重新找极值，代价太高。二刷时先写慢版本或口述慢版本，用它确认不会漏答案。

\textbf{特例回放。} 拿 [1,1,1,2,2,3]、k=2 手算，频次是 1->3、2->2、3->1，前两个高频是 1 和 2。规律是先用哈希计数，再用大小为 k 的小顶堆或桶按频次取 Top K；频次相同不影响题目集合答案。把这个特例继续拆开看：堆顶必须正好代表下一步最需要处理的候选；若堆里可能有过期对象，读取堆顶前要先清理。堆只解决动态极值，不自动证明被弹出或被丢弃的元素无用。复盘时把这个特例压成状态字段：堆元素字段、堆顶比较规则、候选是否过期、答案读取时机；如果需要懒删除，还要保存删除计数；再用边界 K 等于不同元素数、频次相同、结果顺序不要求 反查初始化、更新顺序和退出条件。

\textbf{状态回放。} 手算路线：手算时记录堆顶、候选集合和是否有过期元素。每次使用堆顶前都要确认它仍然属于当前问题。状态字段：堆元素字段、堆顶比较规则、候选是否过期、答案读取时机；如果需要懒删除，还要保存删除计数。状态升级：把每轮全量扫描改成堆顶查询；插入、弹出和过期清理共同维护动态候选集。

\textbf{证明和边界。} 证明从候选覆盖开始：所有可能成为下一步答案的对象都在堆中，堆顶过期时不会被使用。边界：K 等于不同元素数、频次相同、结果顺序不要求。变体：若要按频次和字典序排序，比较器要同时处理两个键。

\noindent\textbf{LeetCode 480 滑动窗口中位数。} 模型：窗口移动时既要加入新元素，也要删除滑出的旧元素并查询中位数。推导重点：双堆维护左右两半，延迟删除表记录已离窗但尚未到堆顶的元素。

\textbf{二刷读题。} 读题时先找动态候选集合以及每一步需要的极值。本题可以压成“窗口移动时既要加入新元素，也要删除滑出的旧元素并查询中位数”，堆顶必须正好回答下一步选择。先遮住答案，只保留题面，复述候选对象、合法性条件和输出目标。

\textbf{暴力回放。} 慢版本每轮扫描全部候选来找最大或最小，再更新候选集合。它的答案选择过程清楚，但动态集合每变化一次就重新找极值，代价太高。二刷时先写慢版本或口述慢版本，用它确认不会漏答案。

\textbf{特例回放。} 拿 [1,3,-1]、k=3 手算，中位数是 1；窗口右移时要删除 1、加入 -3，普通堆不能任意删除。规律是双堆维护两半，再用延迟删除清理过期元素，堆顶使用前必须有效。把这个特例继续拆开看：堆顶必须正好代表下一步最需要处理的候选；若堆里可能有过期对象，读取堆顶前要先清理。堆只解决动态极值，不自动证明被弹出或被丢弃的元素无用。复盘时把这个特例压成状态字段：堆元素字段、堆顶比较规则、候选是否过期、答案读取时机；如果需要懒删除，还要保存删除计数；再用边界 重复值、偶数窗口、平均值溢出、堆顶过期 反查初始化、更新顺序和退出条件。

\textbf{状态回放。} 手算路线：手算时记录堆顶、候选集合和是否有过期元素。每次使用堆顶前都要确认它仍然属于当前问题。状态字段：堆元素字段、堆顶比较规则、候选是否过期、答案读取时机；如果需要懒删除，还要保存删除计数。状态升级：把每轮全量扫描改成堆顶查询；插入、弹出和过期清理共同维护动态候选集。

\textbf{证明和边界。} 证明从候选覆盖开始：所有可能成为下一步答案的对象都在堆中，堆顶过期时不会被使用。边界：重复值、偶数窗口、平均值溢出、堆顶过期。变体：普通 priority\_queue 不能任意删除，所以必须解释懒删除。

\noindent\textbf{LeetCode 621 任务调度器。} 模型：相同任务之间有冷却间隔，核心受最高频任务骨架限制。推导重点：可用大顶堆模拟，也可用最高频公式计算最短时间。

\textbf{二刷读题。} 读题时先找动态候选集合以及每一步需要的极值。本题可以压成“相同任务之间有冷却间隔，核心受最高频任务骨架限制”，堆顶必须正好回答下一步选择。先遮住答案，只保留题面，复述候选对象、合法性条件和输出目标。

\textbf{暴力回放。} 慢版本每轮扫描全部候选来找最大或最小，再更新候选集合。它的答案选择过程清楚，但动态集合每变化一次就重新找极值，代价太高。二刷时先写慢版本或口述慢版本，用它确认不会漏答案。

\textbf{特例回放。} 拿 A,A,A,B,B,B、n=2 手算，若只按顺序执行 A 会违反冷却，必须在相同任务之间间隔两个时间单位。规律是最高频任务决定骨架长度，也可用堆按剩余次数模拟，每轮安排最多 n+1 个不同任务。把这个特例继续拆开看：堆顶必须正好代表下一步最需要处理的候选；若堆里可能有过期对象，读取堆顶前要先清理。堆只解决动态极值，不自动证明被弹出或被丢弃的元素无用。复盘时把这个特例压成状态字段：堆元素字段、堆顶比较规则、候选是否过期、答案读取时机；如果需要懒删除，还要保存删除计数；再用边界 冷却为零、多个最高频、空闲被其他任务填满 反查初始化、更新顺序和退出条件。

\textbf{状态回放。} 手算路线：手算时记录堆顶、候选集合和是否有过期元素。每次使用堆顶前都要确认它仍然属于当前问题。状态字段：堆元素字段、堆顶比较规则、候选是否过期、答案读取时机；如果需要懒删除，还要保存删除计数。状态升级：把每轮全量扫描改成堆顶查询；插入、弹出和过期清理共同维护动态候选集。

\textbf{证明和边界。} 证明从候选覆盖开始：所有可能成为下一步答案的对象都在堆中，堆顶过期时不会被使用。边界：冷却为零、多个最高频、空闲被其他任务填满。变体：若任务有不同释放时间或执行时长，就变成更一般的调度问题。

\noindent\textbf{LeetCode 632 最小区间覆盖 K 个列表。} 模型：每个列表至少取一个元素时，当前区间由最小当前值和最大当前值决定。推导重点：小顶堆保存各列表当前元素，同时维护当前最大值；弹出最小所在列表并推进。

\textbf{二刷读题。} 读题时先找动态候选集合以及每一步需要的极值。本题可以压成“每个列表至少取一个元素时，当前区间由最小当前值和最大当前值决定”，堆顶必须正好回答下一步选择。先遮住答案，只保留题面，复述候选对象、合法性条件和输出目标。

\textbf{暴力回放。} 慢版本每轮扫描全部候选来找最大或最小，再更新候选集合。它的答案选择过程清楚，但动态集合每变化一次就重新找极值，代价太高。二刷时先写慢版本或口述慢版本，用它确认不会漏答案。

\textbf{特例回放。} 拿三组 [4,10]、[0,9]、[5,18] 手算，当前头是 4、0、5，区间 [0,5] 覆盖三组；弹出最小头 0 后推进该组。规律是堆保存每组当前元素，同时维护当前最大值，最小头决定下一步推进哪一组。把这个特例继续拆开看：堆顶必须正好代表下一步最需要处理的候选；若堆里可能有过期对象，读取堆顶前要先清理。堆只解决动态极值，不自动证明被弹出或被丢弃的元素无用。复盘时把这个特例压成状态字段：堆元素字段、堆顶比较规则、候选是否过期、答案读取时机；如果需要懒删除，还要保存删除计数；再用边界 某个列表为空、重复值、区间长度相同的 tie-breaker、列表耗尽 反查初始化、更新顺序和退出条件。

\textbf{状态回放。} 手算路线：手算时记录堆顶、候选集合和是否有过期元素。每次使用堆顶前都要确认它仍然属于当前问题。状态字段：堆元素字段、堆顶比较规则、候选是否过期、答案读取时机；如果需要懒删除，还要保存删除计数。状态升级：把每轮全量扫描改成堆顶查询；插入、弹出和过期清理共同维护动态候选集。

\textbf{证明和边界。} 证明从候选覆盖开始：所有可能成为下一步答案的对象都在堆中，堆顶过期时不会被使用。边界：某个列表为空、重复值、区间长度相同的 tie-breaker、列表耗尽。变体：这是多路归并和覆盖窗口的结合。

\noindent\textbf{LeetCode 703 数据流中的第 K 大元素。} 模型：数据不断到达，只关心当前第 K 大而非完整排序。推导重点：维护大小为 K 的小顶堆，堆顶就是当前第 K 大元素。

\textbf{二刷读题。} 读题时先找动态候选集合以及每一步需要的极值。本题可以压成“数据不断到达，只关心当前第 K 大而非完整排序”，堆顶必须正好回答下一步选择。先遮住答案，只保留题面，复述候选对象、合法性条件和输出目标。

\textbf{暴力回放。} 慢版本每轮扫描全部候选来找最大或最小，再更新候选集合。它的答案选择过程清楚，但动态集合每变化一次就重新找极值，代价太高。二刷时先写慢版本或口述慢版本，用它确认不会漏答案。

\textbf{特例回放。} 拿 k=3，初始 4,5,8,2 手算，小顶堆保存 4,5,8，堆顶 4 是第 3 大；add 3 不进入堆，add 5 会替换 4。规律是数据流只保留前 k 大，堆顶就是当前第 k 大。把这个特例继续拆开看：堆顶必须正好代表下一步最需要处理的候选；若堆里可能有过期对象，读取堆顶前要先清理。堆只解决动态极值，不自动证明被弹出或被丢弃的元素无用。复盘时把这个特例压成状态字段：堆元素字段、堆顶比较规则、候选是否过期、答案读取时机；如果需要懒删除，还要保存删除计数；再用边界 初始元素少于 K、重复值、K 为一、连续 add 反查初始化、更新顺序和退出条件。

\textbf{状态回放。} 手算路线：手算时记录堆顶、候选集合和是否有过期元素。每次使用堆顶前都要确认它仍然属于当前问题。状态字段：堆元素字段、堆顶比较规则、候选是否过期、答案读取时机；如果需要懒删除，还要保存删除计数。状态升级：把每轮全量扫描改成堆顶查询；插入、弹出和过期清理共同维护动态候选集。

\textbf{证明和边界。} 证明从候选覆盖开始：所有可能成为下一步答案的对象都在堆中，堆顶过期时不会被使用。边界：初始元素少于 K、重复值、K 为一、连续 add。变体：这题是 Top K 的在线版本，和一次性排序的适用场景不同。

\noindent\textbf{LeetCode 857 雇佣 K 名工人的最低成本。} 模型：工资比例由当前队伍中最高工资质量比决定，总成本等于比例乘质量和。推导重点：按比例升序扫描，用大顶堆维护当前 K 个工人的最小质量和。

\textbf{二刷读题。} 读题时先找动态候选集合以及每一步需要的极值。本题可以压成“工资比例由当前队伍中最高工资质量比决定，总成本等于比例乘质量和”，堆顶必须正好回答下一步选择。先遮住答案，只保留题面，复述候选对象、合法性条件和输出目标。

\textbf{暴力回放。} 慢版本每轮扫描全部候选来找最大或最小，再更新候选集合。它的答案选择过程清楚，但动态集合每变化一次就重新找极值，代价太高。二刷时先写慢版本或口述慢版本，用它确认不会漏答案。

\textbf{特例回放。} 拿 quality=[10,20,5]、wage=[70,50,30] 手算，按工资质量比选组时，当前最大比率决定组内所有人的单位工资；再用堆保留质量和最小的 k 人。规律是排序枚举最高比率，堆优化质量总和。把这个特例继续拆开看：堆顶必须正好代表下一步最需要处理的候选；若堆里可能有过期对象，读取堆顶前要先清理。堆只解决动态极值，不自动证明被弹出或被丢弃的元素无用。复盘时把这个特例压成状态字段：堆元素字段、堆顶比较规则、候选是否过期、答案读取时机；如果需要懒删除，还要保存删除计数；再用边界 K 为一、比例相同、浮点精度、质量和更新 反查初始化、更新顺序和退出条件。

\textbf{状态回放。} 手算路线：手算时记录堆顶、候选集合和是否有过期元素。每次使用堆顶前都要确认它仍然属于当前问题。状态字段：堆元素字段、堆顶比较规则、候选是否过期、答案读取时机；如果需要懒删除，还要保存删除计数。状态升级：把每轮全量扫描改成堆顶查询；插入、弹出和过期清理共同维护动态候选集。

\textbf{证明和边界。} 证明从候选覆盖开始：所有可能成为下一步答案的对象都在堆中，堆顶过期时不会被使用。边界：K 为一、比例相同、浮点精度、质量和更新。变体：这是排序确定价格比例，堆维护可替换候选的组合题。

\noindent\textbf{LeetCode 871 最低加油次数。} 模型：行驶过程中，经过的加油站形成可反悔的燃料候选。推导重点：按位置扫描，油量不足到达下一点时，从已路过站点中取最大燃料补充。

\textbf{二刷读题。} 读题时先找动态候选集合以及每一步需要的极值。本题可以压成“行驶过程中，经过的加油站形成可反悔的燃料候选”，堆顶必须正好回答下一步选择。先遮住答案，只保留题面，复述候选对象、合法性条件和输出目标。

\textbf{暴力回放。} 慢版本每轮扫描全部候选来找最大或最小，再更新候选集合。它的答案选择过程清楚，但动态集合每变化一次就重新找极值，代价太高。二刷时先写慢版本或口述慢版本，用它确认不会漏答案。

\textbf{特例回放。} 拿 target=100、startFuel=10、站点 (10,60),(20,30),(30,30),(60,40) 手算，到达 10 后把 60 放入堆，燃料不够去下一站时先加此前可达站里最大油量。规律是能到的站都作为候选，加油时选最大油量，保证次数最少。把这个特例继续拆开看：堆顶必须正好代表下一步最需要处理的候选；若堆里可能有过期对象，读取堆顶前要先清理。堆只解决动态极值，不自动证明被弹出或被丢弃的元素无用。复盘时把这个特例压成状态字段：堆元素字段、堆顶比较规则、候选是否过期、答案读取时机；如果需要懒删除，还要保存删除计数；再用边界 起点油足够、无法到达第一个站、多个站同位置、目标作为终点事件 反查初始化、更新顺序和退出条件。

\textbf{状态回放。} 手算路线：手算时记录堆顶、候选集合和是否有过期元素。每次使用堆顶前都要确认它仍然属于当前问题。状态字段：堆元素字段、堆顶比较规则、候选是否过期、答案读取时机；如果需要懒删除，还要保存删除计数。状态升级：把每轮全量扫描改成堆顶查询；插入、弹出和过期清理共同维护动态候选集。

\textbf{证明和边界。} 证明从候选覆盖开始：所有可能成为下一步答案的对象都在堆中，堆顶过期时不会被使用。边界：起点油足够、无法到达第一个站、多个站同位置、目标作为终点事件。变体：这是反悔贪心，堆里保存的是已经经过但尚未使用的资源。

\noindent\textbf{LeetCode 973 最接近原点的 K 个点。} 模型：距离只需比较平方，避免开方误差和成本。推导重点：大小为 K 的大顶堆保存当前最近 K 个点中最远者。

\textbf{二刷读题。} 读题时先找动态候选集合以及每一步需要的极值。本题可以压成“距离只需比较平方，避免开方误差和成本”，堆顶必须正好回答下一步选择。先遮住答案，只保留题面，复述候选对象、合法性条件和输出目标。

\textbf{暴力回放。} 慢版本每轮扫描全部候选来找最大或最小，再更新候选集合。它的答案选择过程清楚，但动态集合每变化一次就重新找极值，代价太高。二刷时先写慢版本或口述慢版本，用它确认不会漏答案。

\textbf{特例回放。} 拿点 (1,3) 和 (-2,2) 手算，距离平方分别是 10 和 8，所以 k=1 时选 (-2,2)。规律是比较距离平方避免开方，用大小为 k 的大顶堆保留当前最近点。把这个特例继续拆开看：堆顶必须正好代表下一步最需要处理的候选；若堆里可能有过期对象，读取堆顶前要先清理。堆只解决动态极值，不自动证明被弹出或被丢弃的元素无用。复盘时把这个特例压成状态字段：堆元素字段、堆顶比较规则、候选是否过期、答案读取时机；如果需要懒删除，还要保存删除计数；再用边界 K 为一、K 为全部、距离相同、坐标平方溢出 反查初始化、更新顺序和退出条件。

\textbf{状态回放。} 手算路线：手算时记录堆顶、候选集合和是否有过期元素。每次使用堆顶前都要确认它仍然属于当前问题。状态字段：堆元素字段、堆顶比较规则、候选是否过期、答案读取时机；如果需要懒删除，还要保存删除计数。状态升级：把每轮全量扫描改成堆顶查询；插入、弹出和过期清理共同维护动态候选集。

\textbf{证明和边界。} 证明从候选覆盖开始：所有可能成为下一步答案的对象都在堆中，堆顶过期时不会被使用。边界：K 为一、K 为全部、距离相同、坐标平方溢出。变体：快速选择可平均线性，但堆更适合数据流。

\noindent\textbf{LeetCode 1046 最后一块石头的重量。} 模型：每次都要取当前最重的两块石头碰撞。推导重点：大顶堆保存所有重量，弹出两个最大值，若不同则把差值放回。

\textbf{二刷读题。} 读题时先找动态候选集合以及每一步需要的极值。本题可以压成“每次都要取当前最重的两块石头碰撞”，堆顶必须正好回答下一步选择。先遮住答案，只保留题面，复述候选对象、合法性条件和输出目标。

\textbf{暴力回放。} 慢版本每轮扫描全部候选来找最大或最小，再更新候选集合。它的答案选择过程清楚，但动态集合每变化一次就重新找极值，代价太高。二刷时先写慢版本或口述慢版本，用它确认不会漏答案。

\textbf{特例回放。} 拿 [2,7,4,1,8,1] 手算，每次取两块最大石头 8 和 7，撞剩 1 再放回。规律是动态最大值查询，用大顶堆保存当前石头重量，直到剩 0 或 1 块。把这个特例继续拆开看：堆顶必须正好代表下一步最需要处理的候选；若堆里可能有过期对象，读取堆顶前要先清理。堆只解决动态极值，不自动证明被弹出或被丢弃的元素无用。复盘时把这个特例压成状态字段：堆元素字段、堆顶比较规则、候选是否过期、答案读取时机；如果需要懒删除，还要保存删除计数；再用边界 没有石头、一块石头、两块相等、多次产生相同重量 反查初始化、更新顺序和退出条件。

\textbf{状态回放。} 手算路线：手算时记录堆顶、候选集合和是否有过期元素。每次使用堆顶前都要确认它仍然属于当前问题。状态字段：堆元素字段、堆顶比较规则、候选是否过期、答案读取时机；如果需要懒删除，还要保存删除计数。状态升级：把每轮全量扫描改成堆顶查询；插入、弹出和过期清理共同维护动态候选集。

\textbf{证明和边界。} 证明从候选覆盖开始：所有可能成为下一步答案的对象都在堆中，堆顶过期时不会被使用。边界：没有石头、一块石头、两块相等、多次产生相同重量。变体：最后一块石头 II 则是分组差值最小化，模型转为 0-1 背包。

\noindent\textbf{LeetCode 1353 最多可以参加的会议数目。} 模型：每天只能参加一个正在进行的会议，应优先选择最早结束的会议。推导重点：按开始日排序扫描日期，把当天可参加会议的结束日入小顶堆，弹出最早结束者。

\textbf{二刷读题。} 读题时先找动态候选集合以及每一步需要的极值。本题可以压成“每天只能参加一个正在进行的会议，应优先选择最早结束的会议”，堆顶必须正好回答下一步选择。先遮住答案，只保留题面，复述候选对象、合法性条件和输出目标。

\textbf{暴力回放。} 慢版本每轮扫描全部候选来找最大或最小，再更新候选集合。它的答案选择过程清楚，但动态集合每变化一次就重新找极值，代价太高。二刷时先写慢版本或口述慢版本，用它确认不会漏答案。

\textbf{特例回放。} 拿 events=[[1,2],[2,3],[3,4]] 手算，第 1 天把开始的会议入堆并参加最早结束的；第 2 天再清理过期并参加一个。规律是按天扫描，用小顶堆保存当前可参加会议的结束日，每天选最早结束者。把这个特例继续拆开看：堆顶必须正好代表下一步最需要处理的候选；若堆里可能有过期对象，读取堆顶前要先清理。堆只解决动态极值，不自动证明被弹出或被丢弃的元素无用。复盘时把这个特例压成状态字段：堆元素字段、堆顶比较规则、候选是否过期、答案读取时机；如果需要懒删除，还要保存删除计数；再用边界 同一天多个会议、过期会议清理、空闲日期跳跃、结束日相同 反查初始化、更新顺序和退出条件。

\textbf{状态回放。} 手算路线：手算时记录堆顶、候选集合和是否有过期元素。每次使用堆顶前都要确认它仍然属于当前问题。状态字段：堆元素字段、堆顶比较规则、候选是否过期、答案读取时机；如果需要懒删除，还要保存删除计数。状态升级：把每轮全量扫描改成堆顶查询；插入、弹出和过期清理共同维护动态候选集。

\textbf{证明和边界。} 证明从候选覆盖开始：所有可能成为下一步答案的对象都在堆中，堆顶过期时不会被使用。边界：同一天多个会议、过期会议清理、空闲日期跳跃、结束日相同。变体：这是扫描线和堆结合的调度题。

\noindent\textbf{LeetCode 1675 数组的最小偏移量。} 模型：奇数可以先乘二，偶数可以不断除二，目标是缩小当前最大最小差。推导重点：把所有数规范到可降低的偶数形态，用大顶堆维护当前最大值，同时记录当前最小值。

\textbf{二刷读题。} 读题时先找动态候选集合以及每一步需要的极值。本题可以压成“奇数可以先乘二，偶数可以不断除二，目标是缩小当前最大最小差”，堆顶必须正好回答下一步选择。先遮住答案，只保留题面，复述候选对象、合法性条件和输出目标。

\textbf{暴力回放。} 慢版本每轮扫描全部候选来找最大或最小，再更新候选集合。它的答案选择过程清楚，但动态集合每变化一次就重新找极值，代价太高。二刷时先写慢版本或口述慢版本，用它确认不会漏答案。

\textbf{特例回放。} 拿 [1,2,3,4] 手算，奇数可先翻倍成偶数，之后只能不断把当前最大偶数减半；每次记录最大值和当前最小值差。规律是把所有数变到可下降的最高状态，再用大顶堆逐步降低最大值。把这个特例继续拆开看：堆顶必须正好代表下一步最需要处理的候选；若堆里可能有过期对象，读取堆顶前要先清理。堆只解决动态极值，不自动证明被弹出或被丢弃的元素无用。复盘时把这个特例压成状态字段：堆元素字段、堆顶比较规则、候选是否过期、答案读取时机；如果需要懒删除，还要保存删除计数；再用边界 全奇数、全偶数、最大值变奇数停止、数值范围 反查初始化、更新顺序和退出条件。

\textbf{状态回放。} 手算路线：手算时记录堆顶、候选集合和是否有过期元素。每次使用堆顶前都要确认它仍然属于当前问题。状态字段：堆元素字段、堆顶比较规则、候选是否过期、答案读取时机；如果需要懒删除，还要保存删除计数。状态升级：把每轮全量扫描改成堆顶查询；插入、弹出和过期清理共同维护动态候选集。

\textbf{证明和边界。} 证明从候选覆盖开始：所有可能成为下一步答案的对象都在堆中，堆顶过期时不会被使用。边界：全奇数、全偶数、最大值变奇数停止、数值范围。变体：这是堆维护动态极值和状态空间规范化的结合。

\noindent\textbf{LeetCode 1851 包含每个查询的最小区间。} 模型：每个查询点需要找覆盖它的最短区间，区间和查询都可离线排序。推导重点：按查询值升序扫描，把左端不大于查询的区间入堆，堆顶过期区间弹出。

\textbf{二刷读题。} 读题时先找动态候选集合以及每一步需要的极值。本题可以压成“每个查询点需要找覆盖它的最短区间，区间和查询都可离线排序”，堆顶必须正好回答下一步选择。先遮住答案，只保留题面，复述候选对象、合法性条件和输出目标。

\textbf{暴力回放。} 慢版本每轮扫描全部候选来找最大或最小，再更新候选集合。它的答案选择过程清楚，但动态集合每变化一次就重新找极值，代价太高。二刷时先写慢版本或口述慢版本，用它确认不会漏答案。

\textbf{特例回放。} 拿 intervals=[[1,4],[2,4],[3,6]]、queries=[2,3,5] 手算，查询 2 时可选 [1,4] 和 [2,4]，较短的是 [2,4]。规律是离线按查询升序加入左端不大于查询的区间，堆顶清理右端已过期区间。把这个特例继续拆开看：堆顶必须正好代表下一步最需要处理的候选；若堆里可能有过期对象，读取堆顶前要先清理。堆只解决动态极值，不自动证明被弹出或被丢弃的元素无用。复盘时把这个特例压成状态字段：堆元素字段、堆顶比较规则、候选是否过期、答案读取时机；如果需要懒删除，还要保存删除计数；再用边界 没有覆盖区间、区间长度相同、查询重复、右端过期 反查初始化、更新顺序和退出条件。

\textbf{状态回放。} 手算路线：手算时记录堆顶、候选集合和是否有过期元素。每次使用堆顶前都要确认它仍然属于当前问题。状态字段：堆元素字段、堆顶比较规则、候选是否过期、答案读取时机；如果需要懒删除，还要保存删除计数。状态升级：把每轮全量扫描改成堆顶查询；插入、弹出和过期清理共同维护动态候选集。

\textbf{证明和边界。} 证明从候选覆盖开始：所有可能成为下一步答案的对象都在堆中，堆顶过期时不会被使用。边界：没有覆盖区间、区间长度相同、查询重复、右端过期。变体：这是离线扫描线、堆和区间覆盖的组合题。

\subsection{自测标准}

二刷时不要只确认自己见过这道题。请遮住答案，先讲题面模型，再讲暴力基线和瓶颈，然后说出优化结构保存了什么、不变量如何排除错误候选、边界实验如何打破错误实现。

\section{链表与指针重连：二刷复盘卡}
这一大节只围绕一个题型复盘，不把题号直接铺成目录。阅读时先看复盘目标，再读核心题卡，最后用自测标准检查自己是否能独立重讲。

\subsection{复盘目标}

追问通常会让你画指针。请准备说明上一段尾、当前段头、当前节点、后继节点分别是谁，以及哪一步保存 next 防止断链。

\subsection{核心题卡}

\noindent\textbf{LeetCode 2 两数相加。} 模型：链表按低位到高位保存数字，本质是竖式加法而不是整数转换。推导重点：维护两个游标和进位，任一链未结束或进位非零都继续生成节点。

\textbf{二刷读题。} 读题时先区分节点值、节点身份和连接关系。本题可以压成“链表按低位到高位保存数字，本质是竖式加法而不是整数转换”，后续推导要围绕哪些指针会断开、哪些节点仍要可达。先遮住答案，只保留题面，复述候选对象、合法性条件和输出目标。

\textbf{暴力回放。} 慢版本可以先把节点顺序抄到数组里，按数组推导目标顺序。回到链表后，难点不再是值的顺序，而是节点身份和 next 指针不能丢。二刷时先写慢版本或口述慢版本，用它确认不会漏答案。

\textbf{特例回放。} 拿 2->4->3 加 5->6->4 手算，个位 2+5 得到 7，十位 4+6 得到 10，所以当前位写 0 并把进位 1 带到下一位。再走百位 3+4+1 得到 8。规律是链表从低位开始，状态只需要两个游标、当前进位和结果尾节点；两条链长度不同或最后还有进位时，循环仍不能停。把这个特例继续拆开看：每个指针都要有角色，上一段尾、当前段头、当前节点和后继入口不能混。只要一次改 next 之前没有保存后继，就可能丢掉整段未处理链。复盘时把这个特例压成状态字段：虚拟头、上一段尾、当前段头、当前节点、后继节点；任何改 next 的操作之前都要保存后续入口；再用边界 两链长度不同、最后进位、输入为零、不能用内置整数承载超长数字 反查初始化、更新顺序和退出条件。

\textbf{状态回放。} 手算路线：手算时画出 prev、cur、next 和下一段头节点。每次改 next 前先保存后继，这是链表推导的安全线。状态字段：虚拟头、上一段尾、当前段头、当前节点、后继节点；任何改 next 的操作之前都要保存后续入口。状态升级：把数组辅助结构换成几个稳定指针角色；重连顺序必须保证已处理链合法、未处理链仍可达。

\textbf{证明和边界。} 证明从链结构开始：已处理链连通、未处理链可达、二者连接点明确，节点没有丢失也没有成环。边界：两链长度不同、最后进位、输入为零、不能用内置整数承载超长数字。变体：若链表高位在前，可以用栈反向处理或先反转链表。

\noindent\textbf{LeetCode 19 删除链表的倒数第 N 个结点。} 模型：倒数位置可以转成两个指针之间固定 n 步的距离。推导重点：快指针先走 n 步，再让快慢指针同步移动，慢指针停在待删节点前驱。

\textbf{二刷读题。} 读题时先区分节点值、节点身份和连接关系。本题可以压成“倒数位置可以转成两个指针之间固定 n 步的距离”，后续推导要围绕哪些指针会断开、哪些节点仍要可达。先遮住答案，只保留题面，复述候选对象、合法性条件和输出目标。

\textbf{暴力回放。} 慢版本可以先把节点顺序抄到数组里，按数组推导目标顺序。回到链表后，难点不再是值的顺序，而是节点身份和 next 指针不能丢。二刷时先写慢版本或口述慢版本，用它确认不会漏答案。

\textbf{特例回放。} 拿长度 5 删除倒数第 2 个手算。快指针先走 2 步，慢指针和快指针保持固定距离；当快指针到尾后，慢指针正好停在待删节点前驱。再试删除头节点，若没有虚拟头就缺前驱。规律是虚拟头统一头部删除，快慢指针固定距离定位前驱。把这个特例继续拆开看：每个指针都要有角色，上一段尾、当前段头、当前节点和后继入口不能混。只要一次改 next 之前没有保存后继，就可能丢掉整段未处理链。复盘时把这个特例压成状态字段：虚拟头、上一段尾、当前段头、当前节点、后继节点；任何改 next 的操作之前都要保存后续入口；再用边界 删除头节点、n 等于链表长度、只有一个节点、题目保证 n 合法 反查初始化、更新顺序和退出条件。

\textbf{状态回放。} 手算路线：手算时画出 prev、cur、next 和下一段头节点。每次改 next 前先保存后继，这是链表推导的安全线。状态字段：虚拟头、上一段尾、当前段头、当前节点、后继节点；任何改 next 的操作之前都要保存后续入口。状态升级：把数组辅助结构换成几个稳定指针角色；重连顺序必须保证已处理链合法、未处理链仍可达。

\textbf{证明和边界。} 证明从链结构开始：已处理链连通、未处理链可达、二者连接点明确，节点没有丢失也没有成环。边界：删除头节点、n 等于链表长度、只有一个节点、题目保证 n 合法。变体：若 n 不保证合法，需要在快指针预走阶段检测长度不足。

\noindent\textbf{LeetCode 21 合并两个有序链表。} 模型：两个链表已经各自有序，下一节点只能来自两个当前头中较小者。推导重点：虚拟头统一结果链表，尾指针不断接上较小节点并推进来源链。

\textbf{二刷读题。} 读题时先区分节点值、节点身份和连接关系。本题可以压成“两个链表已经各自有序，下一节点只能来自两个当前头中较小者”，后续推导要围绕哪些指针会断开、哪些节点仍要可达。先遮住答案，只保留题面，复述候选对象、合法性条件和输出目标。

\textbf{暴力回放。} 慢版本可以先把节点顺序抄到数组里，按数组推导目标顺序。回到链表后，难点不再是值的顺序，而是节点身份和 next 指针不能丢。二刷时先写慢版本或口述慢版本，用它确认不会漏答案。

\textbf{特例回放。} 拿 1->3 和 2->4 手算，先接 1，再比较 3 和 2 接 2，之后接 3 和 4。规律是结果尾指针每次接较小头节点，某条链耗尽后直接接另一条剩余部分；虚拟头统一处理首节点。把这个特例继续拆开看：每个指针都要有角色，上一段尾、当前段头、当前节点和后继入口不能混。只要一次改 next 之前没有保存后继，就可能丢掉整段未处理链。复盘时把这个特例压成状态字段：虚拟头、上一段尾、当前段头、当前节点、后继节点；任何改 next 的操作之前都要保存后续入口；再用边界 某一条链为空、重复值、剩余整段接回、是否复用原节点 反查初始化、更新顺序和退出条件。

\textbf{状态回放。} 手算路线：手算时画出 prev、cur、next 和下一段头节点。每次改 next 前先保存后继，这是链表推导的安全线。状态字段：虚拟头、上一段尾、当前段头、当前节点、后继节点；任何改 next 的操作之前都要保存后续入口。状态升级：把数组辅助结构换成几个稳定指针角色；重连顺序必须保证已处理链合法、未处理链仍可达。

\textbf{证明和边界。} 证明从链结构开始：已处理链连通、未处理链可达、二者连接点明确，节点没有丢失也没有成环。边界：某一条链为空、重复值、剩余整段接回、是否复用原节点。变体：合并 K 条链表可用分治或堆，核心仍是维护每条链的当前头。

\noindent\textbf{LeetCode 24 两两交换链表中的节点。} 模型：链表按固定长度为二的小组做局部重连。推导重点：使用组前驱、第一节点、第二节点和下一组头，按顺序改指针并推进前驱。

\textbf{二刷读题。} 读题时先区分节点值、节点身份和连接关系。本题可以压成“链表按固定长度为二的小组做局部重连”，后续推导要围绕哪些指针会断开、哪些节点仍要可达。先遮住答案，只保留题面，复述候选对象、合法性条件和输出目标。

\textbf{暴力回放。} 慢版本可以先把节点顺序抄到数组里，按数组推导目标顺序。回到链表后，难点不再是值的顺序，而是节点身份和 next 指针不能丢。二刷时先写慢版本或口述慢版本，用它确认不会漏答案。

\textbf{特例回放。} 拿 1->2->3->4 手算，第一组交换后应是 2->1，并且 1 要接到下一组交换后的头。规律是每次固定 pre、first、second、next 四个指针，先保存 next，再重连 second->first，first->next。把这个特例继续拆开看：每个指针都要有角色，上一段尾、当前段头、当前节点和后继入口不能混。只要一次改 next 之前没有保存后继，就可能丢掉整段未处理链。复盘时把这个特例压成状态字段：虚拟头、上一段尾、当前段头、当前节点、后继节点；任何改 next 的操作之前都要保存后续入口；再用边界 空链、单节点、奇数长度、交换后前驱更新到组尾 反查初始化、更新顺序和退出条件。

\textbf{状态回放。} 手算路线：手算时画出 prev、cur、next 和下一段头节点。每次改 next 前先保存后继，这是链表推导的安全线。状态字段：虚拟头、上一段尾、当前段头、当前节点、后继节点；任何改 next 的操作之前都要保存后续入口。状态升级：把数组辅助结构换成几个稳定指针角色；重连顺序必须保证已处理链合法、未处理链仍可达。

\textbf{证明和边界。} 证明从链结构开始：已处理链连通、未处理链可达、二者连接点明确，节点没有丢失也没有成环。边界：空链、单节点、奇数长度、交换后前驱更新到组尾。变体：K 个一组翻转是同一思想的推广，只是组内反转更长。

\noindent\textbf{LeetCode 25 K 个一组翻转链表。} 模型：链表被切成长度为 K 的连续组，只有完整组才反转。推导重点：每轮先探测 K 个节点确认完整，再保存下一组头并做局部反转接回。

\textbf{二刷读题。} 读题时先区分节点值、节点身份和连接关系。本题可以压成“链表被切成长度为 K 的连续组，只有完整组才反转”，后续推导要围绕哪些指针会断开、哪些节点仍要可达。先遮住答案，只保留题面，复述候选对象、合法性条件和输出目标。

\textbf{暴力回放。} 慢版本可以先把节点顺序抄到数组里，按数组推导目标顺序。回到链表后，难点不再是值的顺序，而是节点身份和 next 指针不能丢。二刷时先写慢版本或口述慢版本，用它确认不会漏答案。

\textbf{特例回放。} 拿 1->2->3->4、k=2 画图。第一组边界是 1 和 2，下一组入口是 3；反转前如果不保存 3，链会断。反转后上一组尾接 2，旧头 1 接 3。规律是每组先找 kth 和 group\_next，再局部反转并接回前后两段。把这个特例继续拆开看：每个指针都要有角色，上一段尾、当前段头、当前节点和后继入口不能混。只要一次改 next 之前没有保存后继，就可能丢掉整段未处理链。复盘时把这个特例压成状态字段：虚拟头、上一段尾、当前段头、当前节点、后继节点；任何改 next 的操作之前都要保存后续入口；再用边界 不足 K 个不反转、K 为一、刚好整除、反转后头尾接错 反查初始化、更新顺序和退出条件。

\textbf{状态回放。} 手算路线：手算时画出 prev、cur、next 和下一段头节点。每次改 next 前先保存后继，这是链表推导的安全线。状态字段：虚拟头、上一段尾、当前段头、当前节点、后继节点；任何改 next 的操作之前都要保存后续入口。状态升级：把数组辅助结构换成几个稳定指针角色；重连顺序必须保证已处理链合法、未处理链仍可达。

\textbf{证明和边界。} 证明从链结构开始：已处理链连通、未处理链可达、二者连接点明确，节点没有丢失也没有成环。边界：不足 K 个不反转、K 为一、刚好整除、反转后头尾接错。变体：递归写法短但可能有栈风险，迭代写法更接近工程实现。

\noindent\textbf{LeetCode 61 旋转链表。} 模型：右旋 k 位等价于把链表尾部一段切到头部。推导重点：先求长度并让 k 对长度取模，再找到新尾节点并重连成新头。

\textbf{二刷读题。} 读题时先区分节点值、节点身份和连接关系。本题可以压成“右旋 k 位等价于把链表尾部一段切到头部”，后续推导要围绕哪些指针会断开、哪些节点仍要可达。先遮住答案，只保留题面，复述候选对象、合法性条件和输出目标。

\textbf{暴力回放。} 慢版本可以先把节点顺序抄到数组里，按数组推导目标顺序。回到链表后，难点不再是值的顺序，而是节点身份和 next 指针不能丢。二刷时先写慢版本或口述慢版本，用它确认不会漏答案。

\textbf{特例回放。} 拿 1->2->3->4->5、k=2 手算，右移后应为 4->5->1->2->3；先连成环，再从新尾 3 处断开。规律是 k 先对长度取模，新头是第 n-k 个节点的后继，断环前必须保存新头。把这个特例继续拆开看：每个指针都要有角色，上一段尾、当前段头、当前节点和后继入口不能混。只要一次改 next 之前没有保存后继，就可能丢掉整段未处理链。复盘时把这个特例压成状态字段：虚拟头、上一段尾、当前段头、当前节点、后继节点；任何改 next 的操作之前都要保存后续入口；再用边界 空链、单节点、k 为零、k 大于长度、刚好整圈 反查初始化、更新顺序和退出条件。

\textbf{状态回放。} 手算路线：手算时画出 prev、cur、next 和下一段头节点。每次改 next 前先保存后继，这是链表推导的安全线。状态字段：虚拟头、上一段尾、当前段头、当前节点、后继节点；任何改 next 的操作之前都要保存后续入口。状态升级：把数组辅助结构换成几个稳定指针角色；重连顺序必须保证已处理链合法、未处理链仍可达。

\textbf{证明和边界。} 证明从链结构开始：已处理链连通、未处理链可达、二者连接点明确，节点没有丢失也没有成环。边界：空链、单节点、k 为零、k 大于长度、刚好整圈。变体：若本来把链表接成环，最后必须断开新尾的 next。

\noindent\textbf{LeetCode 82 删除排序链表中的重复元素 II。} 模型：有序链表中重复值连续出现，题目要求把所有重复值整组删除。推导重点：虚拟头保存结果前驱，遇到一段重复值时跳过整段，否则前驱前进。

\textbf{二刷读题。} 读题时先区分节点值、节点身份和连接关系。本题可以压成“有序链表中重复值连续出现，题目要求把所有重复值整组删除”，后续推导要围绕哪些指针会断开、哪些节点仍要可达。先遮住答案，只保留题面，复述候选对象、合法性条件和输出目标。

\textbf{暴力回放。} 慢版本可以先把节点顺序抄到数组里，按数组推导目标顺序。回到链表后，难点不再是值的顺序，而是节点身份和 next 指针不能丢。二刷时先写慢版本或口述慢版本，用它确认不会漏答案。

\textbf{特例回放。} 拿 1->2->3->3->4 手算，两个 3 都要删除，不是保留一个；pre 停在 2，cur 扫完整段重复 3 后接到 4。规律是虚拟头加 pre 指向已确认无重复前缀尾，遇到重复段时整段跳过。把这个特例继续拆开看：每个指针都要有角色，上一段尾、当前段头、当前节点和后继入口不能混。只要一次改 next 之前没有保存后继，就可能丢掉整段未处理链。复盘时把这个特例压成状态字段：虚拟头、上一段尾、当前段头、当前节点、后继节点；任何改 next 的操作之前都要保存后续入口；再用边界 头部重复、尾部重复、全部重复、没有重复 反查初始化、更新顺序和退出条件。

\textbf{状态回放。} 手算路线：手算时画出 prev、cur、next 和下一段头节点。每次改 next 前先保存后继，这是链表推导的安全线。状态字段：虚拟头、上一段尾、当前段头、当前节点、后继节点；任何改 next 的操作之前都要保存后续入口。状态升级：把数组辅助结构换成几个稳定指针角色；重连顺序必须保证已处理链合法、未处理链仍可达。

\textbf{证明和边界。} 证明从链结构开始：已处理链连通、未处理链可达、二者连接点明确，节点没有丢失也没有成环。边界：头部重复、尾部重复、全部重复、没有重复。变体：若只保留一个重复值，就是删除重复元素的较简单版本。

\noindent\textbf{LeetCode 83 删除排序链表中的重复元素。} 模型：有序链表让重复节点相邻，只需保留每段相同值的第一个节点。推导重点：当前节点和后继值相同就跳过后继，否则当前节点前进。

\textbf{二刷读题。} 读题时先区分节点值、节点身份和连接关系。本题可以压成“有序链表让重复节点相邻，只需保留每段相同值的第一个节点”，后续推导要围绕哪些指针会断开、哪些节点仍要可达。先遮住答案，只保留题面，复述候选对象、合法性条件和输出目标。

\textbf{暴力回放。} 慢版本可以先把节点顺序抄到数组里，按数组推导目标顺序。回到链表后，难点不再是值的顺序，而是节点身份和 next 指针不能丢。二刷时先写慢版本或口述慢版本，用它确认不会漏答案。

\textbf{特例回放。} 拿 1->1->2 手算，当前节点 1 的 next 值相同，就让当前节点跳过 next；遇到 2 时再前进。规律是有序链表让重复值相邻，只需比较当前节点和下一个节点，保留每组第一个。把这个特例继续拆开看：每个指针都要有角色，上一段尾、当前段头、当前节点和后继入口不能混。只要一次改 next 之前没有保存后继，就可能丢掉整段未处理链。复盘时把这个特例压成状态字段：虚拟头、上一段尾、当前段头、当前节点、后继节点；任何改 next 的操作之前都要保存后续入口；再用边界 空链、单节点、全重复、重复段在尾部 反查初始化、更新顺序和退出条件。

\textbf{状态回放。} 手算路线：手算时画出 prev、cur、next 和下一段头节点。每次改 next 前先保存后继，这是链表推导的安全线。状态字段：虚拟头、上一段尾、当前段头、当前节点、后继节点；任何改 next 的操作之前都要保存后续入口。状态升级：把数组辅助结构换成几个稳定指针角色；重连顺序必须保证已处理链合法、未处理链仍可达。

\textbf{证明和边界。} 证明从链结构开始：已处理链连通、未处理链可达、二者连接点明确，节点没有丢失也没有成环。边界：空链、单节点、全重复、重复段在尾部。变体：和删除重复元素 II 的区别是是否整段删除。

\noindent\textbf{LeetCode 92 反转链表 II。} 模型：只反转链表中一段连续区间，区间外节点必须保持原连接。推导重点：用虚拟头找到区间前驱，再用头插法把区间内部节点逐个移到前面。

\textbf{二刷读题。} 读题时先区分节点值、节点身份和连接关系。本题可以压成“只反转链表中一段连续区间，区间外节点必须保持原连接”，后续推导要围绕哪些指针会断开、哪些节点仍要可达。先遮住答案，只保留题面，复述候选对象、合法性条件和输出目标。

\textbf{暴力回放。} 慢版本可以先把节点顺序抄到数组里，按数组推导目标顺序。回到链表后，难点不再是值的顺序，而是节点身份和 next 指针不能丢。二刷时先写慢版本或口述慢版本，用它确认不会漏答案。

\textbf{特例回放。} 拿 1->2->3->4->5、left=2、right=4 手算，反转中段后是 1->4->3->2->5；反转前要保存 left 前驱和 right 后继。规律是虚拟头找到子段前驱，局部头插或三指针反转，最后接回前后两段。把这个特例继续拆开看：每个指针都要有角色，上一段尾、当前段头、当前节点和后继入口不能混。只要一次改 next 之前没有保存后继，就可能丢掉整段未处理链。复盘时把这个特例压成状态字段：虚拟头、上一段尾、当前段头、当前节点、后继节点；任何改 next 的操作之前都要保存后续入口；再用边界 从头开始反转、只反转一个节点、反转到尾部、区间边界 反查初始化、更新顺序和退出条件。

\textbf{状态回放。} 手算路线：手算时画出 prev、cur、next 和下一段头节点。每次改 next 前先保存后继，这是链表推导的安全线。状态字段：虚拟头、上一段尾、当前段头、当前节点、后继节点；任何改 next 的操作之前都要保存后续入口。状态升级：把数组辅助结构换成几个稳定指针角色；重连顺序必须保证已处理链合法、未处理链仍可达。

\textbf{证明和边界。} 证明从链结构开始：已处理链连通、未处理链可达、二者连接点明确，节点没有丢失也没有成环。边界：从头开始反转、只反转一个节点、反转到尾部、区间边界。变体：K 个一组反转是在多个固定长度区间上重复这一类局部重连。

\noindent\textbf{LeetCode 141 环形链表。} 模型：快慢指针在有环时必然相遇，无环时快指针先到空。推导重点：不用哈希表也能常数空间检测环。

\textbf{二刷读题。} 读题时先区分节点值、节点身份和连接关系。本题可以压成“快慢指针在有环时必然相遇，无环时快指针先到空”，后续推导要围绕哪些指针会断开、哪些节点仍要可达。先遮住答案，只保留题面，复述候选对象、合法性条件和输出目标。

\textbf{暴力回放。} 慢版本可以先把节点顺序抄到数组里，按数组推导目标顺序。回到链表后，难点不再是值的顺序，而是节点身份和 next 指针不能丢。二刷时先写慢版本或口述慢版本，用它确认不会漏答案。

\textbf{特例回放。} 拿 1->2->3 且 3 指回 2 手算，快指针每次走两步，慢指针每次走一步，进入环后二者距离会被不断缩小直到相遇。规律是快慢指针的相遇说明有环；无环时快指针会先到空。把这个特例继续拆开看：每个指针都要有角色，上一段尾、当前段头、当前节点和后继入口不能混。只要一次改 next 之前没有保存后继，就可能丢掉整段未处理链。复盘时把这个特例压成状态字段：虚拟头、上一段尾、当前段头、当前节点、后继节点；任何改 next 的操作之前都要保存后续入口；再用边界 空链、单节点自环、两节点环、尾部无环 反查初始化、更新顺序和退出条件。

\textbf{状态回放。} 手算路线：手算时画出 prev、cur、next 和下一段头节点。每次改 next 前先保存后继，这是链表推导的安全线。状态字段：虚拟头、上一段尾、当前段头、当前节点、后继节点；任何改 next 的操作之前都要保存后续入口。状态升级：把数组辅助结构换成几个稳定指针角色；重连顺序必须保证已处理链合法、未处理链仍可达。

\textbf{证明和边界。} 证明从链结构开始：已处理链连通、未处理链可达、二者连接点明确，节点没有丢失也没有成环。边界：空链、单节点自环、两节点环、尾部无环。变体：若要求环入口，用相遇后双指针同步走。

\noindent\textbf{LeetCode 142 环形链表 II。} 模型：相遇点到入口的距离关系来自快慢指针速度差。推导重点：相遇后一个指针回头，两个指针每次走一步，再次相遇就是入口。

\textbf{二刷读题。} 读题时先区分节点值、节点身份和连接关系。本题可以压成“相遇点到入口的距离关系来自快慢指针速度差”，后续推导要围绕哪些指针会断开、哪些节点仍要可达。先遮住答案，只保留题面，复述候选对象、合法性条件和输出目标。

\textbf{暴力回放。} 慢版本可以先把节点顺序抄到数组里，按数组推导目标顺序。回到链表后，难点不再是值的顺序，而是节点身份和 next 指针不能丢。二刷时先写慢版本或口述慢版本，用它确认不会漏答案。

\textbf{特例回放。} 拿 1->2->3->4 且 4 指回 2 手算，快慢指针在环内相遇后，把一个指针放回头节点，两个指针同速走，会在入环点 2 相遇。规律是相遇点到入环点的距离和头节点到入环点的距离同余。把这个特例继续拆开看：每个指针都要有角色，上一段尾、当前段头、当前节点和后继入口不能混。只要一次改 next 之前没有保存后继，就可能丢掉整段未处理链。复盘时把这个特例压成状态字段：虚拟头、上一段尾、当前段头、当前节点、后继节点；任何改 next 的操作之前都要保存后续入口；再用边界 入口在头节点、长尾短环、无环、单节点环 反查初始化、更新顺序和退出条件。

\textbf{状态回放。} 手算路线：手算时画出 prev、cur、next 和下一段头节点。每次改 next 前先保存后继，这是链表推导的安全线。状态字段：虚拟头、上一段尾、当前段头、当前节点、后继节点；任何改 next 的操作之前都要保存后续入口。状态升级：把数组辅助结构换成几个稳定指针角色；重连顺序必须保证已处理链合法、未处理链仍可达。

\textbf{证明和边界。} 证明从链结构开始：已处理链连通、未处理链可达、二者连接点明确，节点没有丢失也没有成环。边界：入口在头节点、长尾短环、无环、单节点环。变体：也可用哈希集合，但空间更高。

\noindent\textbf{LeetCode 146 LRU 缓存。} 模型：哈希表负责按 key 定位，双向链表负责维护新旧顺序。推导重点：访问和更新都把节点移动到头部，容量满时删除尾部最旧节点。

\textbf{二刷读题。} 读题时先区分节点值、节点身份和连接关系。本题可以压成“哈希表负责按 key 定位，双向链表负责维护新旧顺序”，后续推导要围绕哪些指针会断开、哪些节点仍要可达。先遮住答案，只保留题面，复述候选对象、合法性条件和输出目标。

\textbf{暴力回放。} 慢版本可以先把节点顺序抄到数组里，按数组推导目标顺序。回到链表后，难点不再是值的顺序，而是节点身份和 next 指针不能丢。二刷时先写慢版本或口述慢版本，用它确认不会漏答案。

\textbf{特例回放。} 拿容量 2 依次 put(1)、put(2)、get(1)、put(3) 手算。get(1) 后 1 变成最新，淘汰时应删 2。数组维护顺序每次移动成本高；链表移动节点到头 O(1)，哈希表按 key 定位节点 O(1)。规律是哈希负责找节点，双向链表负责维护新旧顺序。把这个特例继续拆开看：每个指针都要有角色，上一段尾、当前段头、当前节点和后继入口不能混。只要一次改 next 之前没有保存后继，就可能丢掉整段未处理链。复盘时把这个特例压成状态字段：虚拟头、上一段尾、当前段头、当前节点、后继节点；任何改 next 的操作之前都要保存后续入口；再用边界 容量为一、重复 put、get 不存在、更新已有值 反查初始化、更新顺序和退出条件。

\textbf{状态回放。} 手算路线：手算时画出 prev、cur、next 和下一段头节点。每次改 next 前先保存后继，这是链表推导的安全线。状态字段：虚拟头、上一段尾、当前段头、当前节点、后继节点；任何改 next 的操作之前都要保存后续入口。状态升级：把数组辅助结构换成几个稳定指针角色；重连顺序必须保证已处理链合法、未处理链仍可达。

\textbf{证明和边界。} 证明从链结构开始：已处理链连通、未处理链可达、二者连接点明确，节点没有丢失也没有成环。边界：容量为一、重复 put、get 不存在、更新已有值。变体：C++ 可用 \code{std::list} 加迭代器，也可手写双向链表理解所有权。

\noindent\textbf{LeetCode 148 排序链表。} 模型：链表不能随机访问，适合归并排序而不是快速排序数组 partition。推导重点：自底向上迭代归并能避免递归栈，同时保持 O(n log n)。

\textbf{二刷读题。} 读题时先区分节点值、节点身份和连接关系。本题可以压成“链表不能随机访问，适合归并排序而不是快速排序数组 partition”，后续推导要围绕哪些指针会断开、哪些节点仍要可达。先遮住答案，只保留题面，复述候选对象、合法性条件和输出目标。

\textbf{暴力回放。} 慢版本可以先把节点顺序抄到数组里，按数组推导目标顺序。回到链表后，难点不再是值的顺序，而是节点身份和 next 指针不能丢。二刷时先写慢版本或口述慢版本，用它确认不会漏答案。

\textbf{特例回放。} 拿 4->2->1->3 手算，快慢指针把链表切成 4->2 和 1->3，分别排好后再归并。规律是链表排序不能随机访问，归并排序只需要切断中点和有序合并，复杂度稳定在 O(n log n)。把这个特例继续拆开看：每个指针都要有角色，上一段尾、当前段头、当前节点和后继入口不能混。只要一次改 next 之前没有保存后继，就可能丢掉整段未处理链。复盘时把这个特例压成状态字段：虚拟头、上一段尾、当前段头、当前节点、后继节点；任何改 next 的操作之前都要保存后续入口；再用边界 空链、单节点、重复值、奇数长度、切断子链 反查初始化、更新顺序和退出条件。

\textbf{状态回放。} 手算路线：手算时画出 prev、cur、next 和下一段头节点。每次改 next 前先保存后继，这是链表推导的安全线。状态字段：虚拟头、上一段尾、当前段头、当前节点、后继节点；任何改 next 的操作之前都要保存后续入口。状态升级：把数组辅助结构换成几个稳定指针角色；重连顺序必须保证已处理链合法、未处理链仍可达。

\textbf{证明和边界。} 证明从链结构开始：已处理链连通、未处理链可达、二者连接点明确，节点没有丢失也没有成环。边界：空链、单节点、重复值、奇数长度、切断子链。变体：若把值复制到数组排序，会改变节点身份语义。

\noindent\textbf{LeetCode 160 相交链表。} 模型：相交是节点地址相同，不是节点值相同。推导重点：两个指针走完各自链表后切换到另一条链，走过总长度相同后相遇。

\textbf{二刷读题。} 读题时先区分节点值、节点身份和连接关系。本题可以压成“相交是节点地址相同，不是节点值相同”，后续推导要围绕哪些指针会断开、哪些节点仍要可达。先遮住答案，只保留题面，复述候选对象、合法性条件和输出目标。

\textbf{暴力回放。} 慢版本可以先把节点顺序抄到数组里，按数组推导目标顺序。回到链表后，难点不再是值的顺序，而是节点身份和 next 指针不能丢。二刷时先写慢版本或口述慢版本，用它确认不会漏答案。

\textbf{特例回放。} 拿 A: a1->a2->c1->c2 和 B: b1->c1->c2 手算，交点是节点地址 c1，不是值相等。两个指针走到尾后切换到对方头部，会在总路程相同后于 c1 相遇。规律是比较节点身份，长度差由切换链表自动抵消。把这个特例继续拆开看：每个指针都要有角色，上一段尾、当前段头、当前节点和后继入口不能混。只要一次改 next 之前没有保存后继，就可能丢掉整段未处理链。复盘时把这个特例压成状态字段：虚拟头、上一段尾、当前段头、当前节点、后继节点；任何改 next 的操作之前都要保存后续入口；再用边界 不相交、头节点相交、尾节点相交、长度差很大 反查初始化、更新顺序和退出条件。

\textbf{状态回放。} 手算路线：手算时画出 prev、cur、next 和下一段头节点。每次改 next 前先保存后继，这是链表推导的安全线。状态字段：虚拟头、上一段尾、当前段头、当前节点、后继节点；任何改 next 的操作之前都要保存后续入口。状态升级：把数组辅助结构换成几个稳定指针角色；重连顺序必须保证已处理链合法、未处理链仍可达。

\textbf{证明和边界。} 证明从链结构开始：已处理链连通、未处理链可达、二者连接点明确，节点没有丢失也没有成环。边界：不相交、头节点相交、尾节点相交、长度差很大。变体：哈希集合更直观但空间更高。

\noindent\textbf{LeetCode 206 反转链表。} 模型：每次把当前节点的 next 指向已经反转好的前缀头。推导重点：先保存后继，再改指针，再推进 previous 和 current。

\textbf{二刷读题。} 读题时先区分节点值、节点身份和连接关系。本题可以压成“每次把当前节点的 next 指向已经反转好的前缀头”，后续推导要围绕哪些指针会断开、哪些节点仍要可达。先遮住答案，只保留题面，复述候选对象、合法性条件和输出目标。

\textbf{暴力回放。} 慢版本可以先把节点顺序抄到数组里，按数组推导目标顺序。回到链表后，难点不再是值的顺序，而是节点身份和 next 指针不能丢。二刷时先写慢版本或口述慢版本，用它确认不会漏答案。

\textbf{特例回放。} 拿 1->2->3 手算，处理 1 时先保存 next=2，再让 1->null；处理 2 时保存 3，再让 2->1。规律是 prev 保存已反转前缀头，cur 保存待处理节点，每次改 next 前必须保存后继。把这个特例继续拆开看：每个指针都要有角色，上一段尾、当前段头、当前节点和后继入口不能混。只要一次改 next 之前没有保存后继，就可能丢掉整段未处理链。复盘时把这个特例压成状态字段：虚拟头、上一段尾、当前段头、当前节点、后继节点；任何改 next 的操作之前都要保存后续入口；再用边界 空链、单节点、两个节点、不要丢失剩余链 反查初始化、更新顺序和退出条件。

\textbf{状态回放。} 手算路线：手算时画出 prev、cur、next 和下一段头节点。每次改 next 前先保存后继，这是链表推导的安全线。状态字段：虚拟头、上一段尾、当前段头、当前节点、后继节点；任何改 next 的操作之前都要保存后续入口。状态升级：把数组辅助结构换成几个稳定指针角色；重连顺序必须保证已处理链合法、未处理链仍可达。

\textbf{证明和边界。} 证明从链结构开始：已处理链连通、未处理链可达、二者连接点明确，节点没有丢失也没有成环。边界：空链、单节点、两个节点、不要丢失剩余链。变体：K 组反转是在这段局部反转外增加分组边界。

\noindent\textbf{LeetCode 234 回文链表。} 模型：链表回文要求前半段和后半段逆序后逐一相等。推导重点：快慢指针找中点，反转后半段，再和前半段同步比较。

\textbf{二刷读题。} 读题时先区分节点值、节点身份和连接关系。本题可以压成“链表回文要求前半段和后半段逆序后逐一相等”，后续推导要围绕哪些指针会断开、哪些节点仍要可达。先遮住答案，只保留题面，复述候选对象、合法性条件和输出目标。

\textbf{暴力回放。} 慢版本可以先把节点顺序抄到数组里，按数组推导目标顺序。回到链表后，难点不再是值的顺序，而是节点身份和 next 指针不能丢。二刷时先写慢版本或口述慢版本，用它确认不会漏答案。

\textbf{特例回放。} 拿 1->2->2->1 手算，快慢指针找到后半段起点，再反转后半段得到 1->2，与前半段逐个比较。规律是链表不能随机访问，先找中点、反转后半、比较，再按需要恢复结构。把这个特例继续拆开看：每个指针都要有角色，上一段尾、当前段头、当前节点和后继入口不能混。只要一次改 next 之前没有保存后继，就可能丢掉整段未处理链。复盘时把这个特例压成状态字段：虚拟头、上一段尾、当前段头、当前节点、后继节点；任何改 next 的操作之前都要保存后续入口；再用边界 空链、单节点、奇数长度、偶数长度、是否恢复链表 反查初始化、更新顺序和退出条件。

\textbf{状态回放。} 手算路线：手算时画出 prev、cur、next 和下一段头节点。每次改 next 前先保存后继，这是链表推导的安全线。状态字段：虚拟头、上一段尾、当前段头、当前节点、后继节点；任何改 next 的操作之前都要保存后续入口。状态升级：把数组辅助结构换成几个稳定指针角色；重连顺序必须保证已处理链合法、未处理链仍可达。

\textbf{证明和边界。} 证明从链结构开始：已处理链连通、未处理链可达、二者连接点明确，节点没有丢失也没有成环。边界：空链、单节点、奇数长度、偶数长度、是否恢复链表。变体：如果不允许修改链表，可用栈保存前半段但空间变为线性。

\noindent\textbf{LeetCode 430 扁平化多级双向链表。} 模型：多级链表按深度优先顺序展开成一条双向链表。推导重点：显式栈保存后续 next，遇到 child 先接 child，再在子链结束后恢复原 next。

\textbf{二刷读题。} 读题时先区分节点值、节点身份和连接关系。本题可以压成“多级链表按深度优先顺序展开成一条双向链表”，后续推导要围绕哪些指针会断开、哪些节点仍要可达。先遮住答案，只保留题面，复述候选对象、合法性条件和输出目标。

\textbf{暴力回放。} 慢版本可以先把节点顺序抄到数组里，按数组推导目标顺序。回到链表后，难点不再是值的顺序，而是节点身份和 next 指针不能丢。二刷时先写慢版本或口述慢版本，用它确认不会漏答案。

\textbf{特例回放。} 拿 1-2-3 且 2 有子链 4-5 手算，扁平化顺序应为 1-2-4-5-3；若不保存 2 原来的 next=3，子链接完后会丢失 3。规律是遇到 child 时先保存 next，再把 child 接入，并把 child 尾部接回原 next。把这个特例继续拆开看：每个指针都要有角色，上一段尾、当前段头、当前节点和后继入口不能混。只要一次改 next 之前没有保存后继，就可能丢掉整段未处理链。复盘时把这个特例压成状态字段：虚拟头、上一段尾、当前段头、当前节点、后继节点；任何改 next 的操作之前都要保存后续入口；再用边界 child 为空、next 和 child 同时存在、多层嵌套、prev 指针恢复 反查初始化、更新顺序和退出条件。

\textbf{状态回放。} 手算路线：手算时画出 prev、cur、next 和下一段头节点。每次改 next 前先保存后继，这是链表推导的安全线。状态字段：虚拟头、上一段尾、当前段头、当前节点、后继节点；任何改 next 的操作之前都要保存后续入口。状态升级：把数组辅助结构换成几个稳定指针角色；重连顺序必须保证已处理链合法、未处理链仍可达。

\textbf{证明和边界。} 证明从链结构开始：已处理链连通、未处理链可达、二者连接点明确，节点没有丢失也没有成环。边界：child 为空、next 和 child 同时存在、多层嵌套、prev 指针恢复。变体：它和二叉树前序展开类似，核心是保存未处理分支。

\noindent\textbf{LeetCode 460 LFU 缓存。} 模型：缓存淘汰同时按访问频次和同频最近性排序，需要维护两个维度的索引。推导重点：哈希表按 key 定位节点，频次到双向链表保存同频节点顺序，并维护当前最小频次。

\textbf{二刷读题。} 读题时先区分节点值、节点身份和连接关系。本题可以压成“缓存淘汰同时按访问频次和同频最近性排序，需要维护两个维度的索引”，后续推导要围绕哪些指针会断开、哪些节点仍要可达。先遮住答案，只保留题面，复述候选对象、合法性条件和输出目标。

\textbf{暴力回放。} 慢版本可以先把节点顺序抄到数组里，按数组推导目标顺序。回到链表后，难点不再是值的顺序，而是节点身份和 next 指针不能丢。二刷时先写慢版本或口述慢版本，用它确认不会漏答案。

\textbf{特例回放。} 拿容量 2，put(1)、put(2)、get(1)、put(3) 手算。此时 1 的频次高于 2，应淘汰 2；若两个 key 频次相同，还要淘汰同频里最旧的。规律是 key 哈希定位节点，频次哈希定位同频链表，minFreq 指向当前最低频次。把这个特例继续拆开看：每个指针都要有角色，上一段尾、当前段头、当前节点和后继入口不能混。只要一次改 next 之前没有保存后继，就可能丢掉整段未处理链。复盘时把这个特例压成状态字段：虚拟头、上一段尾、当前段头、当前节点、后继节点；任何改 next 的操作之前都要保存后续入口；再用边界 容量为零、更新已有 key、频次链表变空时最小频次更新、同频淘汰最旧节点 反查初始化、更新顺序和退出条件。

\textbf{状态回放。} 手算路线：手算时画出 prev、cur、next 和下一段头节点。每次改 next 前先保存后继，这是链表推导的安全线。状态字段：虚拟头、上一段尾、当前段头、当前节点、后继节点；任何改 next 的操作之前都要保存后续入口。状态升级：把数组辅助结构换成几个稳定指针角色；重连顺序必须保证已处理链合法、未处理链仍可达。

\textbf{证明和边界。} 证明从链结构开始：已处理链连通、未处理链可达、二者连接点明确，节点没有丢失也没有成环。边界：容量为零、更新已有 key、频次链表变空时最小频次更新、同频淘汰最旧节点。变体：LRU 只有时间顺序一个维度，LFU 多了频次维度，结构维护明显更复杂。

\subsection{自测标准}

二刷时不要只确认自己见过这道题。请遮住答案，先讲题面模型，再讲暴力基线和瓶颈，然后说出优化结构保存了什么、不变量如何排除错误候选、边界实验如何打破错误实现。

\section{二叉树与树形状态：二刷复盘卡}
这一大节只围绕一个题型复盘，不把题号直接铺成目录。阅读时先看复盘目标，再读核心题卡，最后用自测标准检查自己是否能独立重讲。

\subsection{复盘目标}

追问通常会问返回值语义、空节点初值、全负数和极端深树。请准备把递归思想改写成显式栈或队列的口头方案。

\subsection{核心题卡}

\noindent\textbf{LeetCode 94 二叉树中序遍历。} 模型：中序访问顺序是左子树、根、右子树，BST 中会得到有序序列。推导重点：用显式栈模拟一路向左，再回退访问根并转向右子树。

\textbf{二刷读题。} 读题时先判断状态方向：父到子、子到父，还是按层传播。本题可以压成“中序访问顺序是左子树、根、右子树，BST 中会得到有序序列”，每个节点的输入和输出状态必须写清。先遮住答案，只保留题面，复述候选对象、合法性条件和输出目标。

\textbf{暴力回放。} 慢版本在每个节点重新扫描子树或路径，先求出局部答案。重复扫描暴露了真正状态：每个节点只需要从孩子或父亲那里拿到少量信息。二刷时先写慢版本或口述慢版本，用它确认不会漏答案。

\textbf{特例回放。} 拿根 2、左 1、右 3 手算，中序结果必须是 1,2,3。显式栈做法先一路压到最左，弹出 1 后访问，再回到 2。规律是栈保存尚未访问但左侧已经处理完的祖先，BST 题里这个顺序会产生递增序列。把这个特例继续拆开看：节点向父亲返回的量和全局答案通常不是同一个量。先写叶子或空节点的返回值，再写父节点如何合并左右孩子，递归式才有边界基础。复盘时把这个特例压成状态字段：节点指针、传入约束或返回值、当前答案、层号或路径状态；空节点的返回值必须和状态定义匹配；再用边界 空树、单节点、链式左树、链式右树 反查初始化、更新顺序和退出条件。

\textbf{状态回放。} 手算路线：手算时从叶子开始写返回值，或者从根开始写传入约束。不要只写访问顺序，要写每条边上传递的信息。状态字段：节点指针、传入约束或返回值、当前答案、层号或路径状态；空节点的返回值必须和状态定义匹配。状态升级：把对子树的重复扫描压成返回值或队列状态；每个节点只合并孩子给出的稳定信息。

\textbf{证明和边界。} 证明从子树归纳开始：孩子状态正确时，父节点按题目规则合并后也正确。边界：空树、单节点、链式左树、链式右树。变体：Morris 遍历可做到常数额外空间，但会临时改树结构。

\noindent\textbf{LeetCode 98 验证二叉搜索树。} 模型：BST 约束是全局上下界，不是只比较父子节点。推导重点：中序严格递增或显式传递上下界都可以；中序版本要保存前一个访问值。

\textbf{二刷读题。} 读题时先判断状态方向：父到子、子到父，还是按层传播。本题可以压成“BST 约束是全局上下界，不是只比较父子节点”，每个节点的输入和输出状态必须写清。先遮住答案，只保留题面，复述候选对象、合法性条件和输出目标。

\textbf{暴力回放。} 慢版本在每个节点重新扫描子树或路径，先求出局部答案。重复扫描暴露了真正状态：每个节点只需要从孩子或父亲那里拿到少量信息。二刷时先写慢版本或口述慢版本，用它确认不会漏答案。

\textbf{特例回放。} 拿 [5,1,4,null,null,3,6] 手算，节点 3 小于根 5 却在右子树里，只比较父子 4 和 3 会误判。规律是 BST 约束来自整条祖先上下界，或中序序列必须严格递增；重复值和 int 边界要按题意单独处理。把这个特例继续拆开看：节点向父亲返回的量和全局答案通常不是同一个量。先写叶子或空节点的返回值，再写父节点如何合并左右孩子，递归式才有边界基础。复盘时把这个特例压成状态字段：节点指针、传入约束或返回值、当前答案、层号或路径状态；空节点的返回值必须和状态定义匹配；再用边界 重复值、int 最小最大值、局部合法但全局非法 反查初始化、更新顺序和退出条件。

\textbf{状态回放。} 手算路线：手算时从叶子开始写返回值，或者从根开始写传入约束。不要只写访问顺序，要写每条边上传递的信息。状态字段：节点指针、传入约束或返回值、当前答案、层号或路径状态；空节点的返回值必须和状态定义匹配。状态升级：把对子树的重复扫描压成返回值或队列状态；每个节点只合并孩子给出的稳定信息。

\textbf{证明和边界。} 证明从子树归纳开始：孩子状态正确时，父节点按题目规则合并后也正确。边界：重复值、int 最小最大值、局部合法但全局非法。变体：若允许重复值，需要明确重复放左还是放右。

\noindent\textbf{LeetCode 102 二叉树层序遍历。} 模型：队列在每轮开始时保存当前层所有节点。推导重点：记录 level size 固定层边界，避免下一层节点混入当前层。

\textbf{二刷读题。} 读题时先判断状态方向：父到子、子到父，还是按层传播。本题可以压成“队列在每轮开始时保存当前层所有节点”，每个节点的输入和输出状态必须写清。先遮住答案，只保留题面，复述候选对象、合法性条件和输出目标。

\textbf{暴力回放。} 慢版本在每个节点重新扫描子树或路径，先求出局部答案。重复扫描暴露了真正状态：每个节点只需要从孩子或父亲那里拿到少量信息。二刷时先写慢版本或口述慢版本，用它确认不会漏答案。

\textbf{特例回放。} 拿根 3、左 9、右 20 手算，队列第一轮只有 3，第二轮才是 9 和 20。若不固定本层 size，新入队的下一层节点会混进当前层。规律是每轮开始记录队列长度，这个长度就是当前层边界。把这个特例继续拆开看：节点向父亲返回的量和全局答案通常不是同一个量。先写叶子或空节点的返回值，再写父节点如何合并左右孩子，递归式才有边界基础。复盘时把这个特例压成状态字段：节点指针、传入约束或返回值、当前答案、层号或路径状态；空节点的返回值必须和状态定义匹配；再用边界 空树、只有根、最后一层不满、左右顺序 反查初始化、更新顺序和退出条件。

\textbf{状态回放。} 手算路线：手算时从叶子开始写返回值，或者从根开始写传入约束。不要只写访问顺序，要写每条边上传递的信息。状态字段：节点指针、传入约束或返回值、当前答案、层号或路径状态；空节点的返回值必须和状态定义匹配。状态升级：把对子树的重复扫描压成返回值或队列状态；每个节点只合并孩子给出的稳定信息。

\textbf{证明和边界。} 证明从子树归纳开始：孩子状态正确时，父节点按题目规则合并后也正确。边界：空树、只有根、最后一层不满、左右顺序。变体：锯齿形层序只是在每层输出顺序上增加方向控制。

\noindent\textbf{LeetCode 103 二叉树锯齿形层序遍历。} 模型：BFS 层边界不变，变化的是当前层写入结果的方向。推导重点：可以层内反转，也可以预分配数组按方向写入目标下标。

\textbf{二刷读题。} 读题时先判断状态方向：父到子、子到父，还是按层传播。本题可以压成“BFS 层边界不变，变化的是当前层写入结果的方向”，每个节点的输入和输出状态必须写清。先遮住答案，只保留题面，复述候选对象、合法性条件和输出目标。

\textbf{暴力回放。} 慢版本在每个节点重新扫描子树或路径，先求出局部答案。重复扫描暴露了真正状态：每个节点只需要从孩子或父亲那里拿到少量信息。二刷时先写慢版本或口述慢版本，用它确认不会漏答案。

\textbf{特例回放。} 拿三层树手算，第一层从左到右，第二层从右到左，第三层再从左到右；BFS 层边界不变，只是写入结果的位置反向。规律是层序遍历负责分层，方向标志只影响当前层数组的填充顺序。把这个特例继续拆开看：节点向父亲返回的量和全局答案通常不是同一个量。先写叶子或空节点的返回值，再写父节点如何合并左右孩子，递归式才有边界基础。复盘时把这个特例压成状态字段：节点指针、传入约束或返回值、当前答案、层号或路径状态；空节点的返回值必须和状态定义匹配；再用边界 单层、偶数层、奇数层、空树 反查初始化、更新顺序和退出条件。

\textbf{状态回放。} 手算路线：手算时从叶子开始写返回值，或者从根开始写传入约束。不要只写访问顺序，要写每条边上传递的信息。状态字段：节点指针、传入约束或返回值、当前答案、层号或路径状态；空节点的返回值必须和状态定义匹配。状态升级：把对子树的重复扫描压成返回值或队列状态；每个节点只合并孩子给出的稳定信息。

\textbf{证明和边界。} 证明从子树归纳开始：孩子状态正确时，父节点按题目规则合并后也正确。边界：单层、偶数层、奇数层、空树。变体：若要求从底向上输出，收集后反转层列表即可。

\noindent\textbf{LeetCode 104 二叉树最大深度。} 模型：深度可以看成层数，也可以作为栈中携带的状态。推导重点：显式栈保存节点和深度，避免极端链式树递归栈过深。

\textbf{二刷读题。} 读题时先判断状态方向：父到子、子到父，还是按层传播。本题可以压成“深度可以看成层数，也可以作为栈中携带的状态”，每个节点的输入和输出状态必须写清。先遮住答案，只保留题面，复述候选对象、合法性条件和输出目标。

\textbf{暴力回放。} 慢版本在每个节点重新扫描子树或路径，先求出局部答案。重复扫描暴露了真正状态：每个节点只需要从孩子或父亲那里拿到少量信息。二刷时先写慢版本或口述慢版本，用它确认不会漏答案。

\textbf{特例回放。} 拿链式树 1->2->3 手算，最大深度是 3；空树是 0。递归返回左右子树深度的较大值加一，显式栈则把节点和当前深度一起压栈。规律是状态必须携带深度，不能只记录访问过哪些节点。把这个特例继续拆开看：节点向父亲返回的量和全局答案通常不是同一个量。先写叶子或空节点的返回值，再写父节点如何合并左右孩子，递归式才有边界基础。复盘时把这个特例压成状态字段：节点指针、传入约束或返回值、当前答案、层号或路径状态；空节点的返回值必须和状态定义匹配；再用边界 空树返回零、根深度是一、链式树 反查初始化、更新顺序和退出条件。

\textbf{状态回放。} 手算路线：手算时从叶子开始写返回值，或者从根开始写传入约束。不要只写访问顺序，要写每条边上传递的信息。状态字段：节点指针、传入约束或返回值、当前答案、层号或路径状态；空节点的返回值必须和状态定义匹配。状态升级：把对子树的重复扫描压成返回值或队列状态；每个节点只合并孩子给出的稳定信息。

\textbf{证明和边界。} 证明从子树归纳开始：孩子状态正确时，父节点按题目规则合并后也正确。边界：空树返回零、根深度是一、链式树。变体：最小深度不能简单取左右最小，因为空子树不构成叶子路径。

\noindent\textbf{LeetCode 105 前序中序构造二叉树。} 模型：前序给根，中序把左右子树切开。推导重点：用哈希表记录中序下标，避免每次线性寻找根位置。

\textbf{二刷读题。} 读题时先判断状态方向：父到子、子到父，还是按层传播。本题可以压成“前序给根，中序把左右子树切开”，每个节点的输入和输出状态必须写清。先遮住答案，只保留题面，复述候选对象、合法性条件和输出目标。

\textbf{暴力回放。} 慢版本在每个节点重新扫描子树或路径，先求出局部答案。重复扫描暴露了真正状态：每个节点只需要从孩子或父亲那里拿到少量信息。二刷时先写慢版本或口述慢版本，用它确认不会漏答案。

\textbf{特例回放。} 拿 preorder=[3,9,20,15,7]、inorder=[9,3,15,20,7] 手算，前序第一个 3 是根，中序中 3 左边是左子树，右边是右子树。规律是用哈希表快速找根在中序的位置，再用左右子树长度切分前序区间。把这个特例继续拆开看：节点向父亲返回的量和全局答案通常不是同一个量。先写叶子或空节点的返回值，再写父节点如何合并左右孩子，递归式才有边界基础。复盘时把这个特例压成状态字段：节点指针、传入约束或返回值、当前答案、层号或路径状态；空节点的返回值必须和状态定义匹配；再用边界 节点值唯一、空子树、左右子树长度计算、输入不合法 反查初始化、更新顺序和退出条件。

\textbf{状态回放。} 手算路线：手算时从叶子开始写返回值，或者从根开始写传入约束。不要只写访问顺序，要写每条边上传递的信息。状态字段：节点指针、传入约束或返回值、当前答案、层号或路径状态；空节点的返回值必须和状态定义匹配。状态升级：把对子树的重复扫描压成返回值或队列状态；每个节点只合并孩子给出的稳定信息。

\textbf{证明和边界。} 证明从子树归纳开始：孩子状态正确时，父节点按题目规则合并后也正确。边界：节点值唯一、空子树、左右子树长度计算、输入不合法。变体：后序中序构造类似，只是根来自后序末尾。

\noindent\textbf{LeetCode 106 中序后序构造二叉树。} 模型：后序末尾是根，中序位置确定左右子树规模。推导重点：构造时要小心后序左右区间边界，避免切分错位。

\textbf{二刷读题。} 读题时先判断状态方向：父到子、子到父，还是按层传播。本题可以压成“后序末尾是根，中序位置确定左右子树规模”，每个节点的输入和输出状态必须写清。先遮住答案，只保留题面，复述候选对象、合法性条件和输出目标。

\textbf{暴力回放。} 慢版本在每个节点重新扫描子树或路径，先求出局部答案。重复扫描暴露了真正状态：每个节点只需要从孩子或父亲那里拿到少量信息。二刷时先写慢版本或口述慢版本，用它确认不会漏答案。

\textbf{特例回放。} 拿 inorder=[9,3,15,20,7]、postorder=[9,15,7,20,3] 手算，后序最后一个 3 是根，中序把左右子树切开。规律是根来自后序末尾，左右区间长度必须同步更新；全左链和全右链最容易暴露边界错位。把这个特例继续拆开看：节点向父亲返回的量和全局答案通常不是同一个量。先写叶子或空节点的返回值，再写父节点如何合并左右孩子，递归式才有边界基础。复盘时把这个特例压成状态字段：节点指针、传入约束或返回值、当前答案、层号或路径状态；空节点的返回值必须和状态定义匹配；再用边界 空树、单节点、全左链、全右链 反查初始化、更新顺序和退出条件。

\textbf{状态回放。} 手算路线：手算时从叶子开始写返回值，或者从根开始写传入约束。不要只写访问顺序，要写每条边上传递的信息。状态字段：节点指针、传入约束或返回值、当前答案、层号或路径状态；空节点的返回值必须和状态定义匹配。状态升级：把对子树的重复扫描压成返回值或队列状态；每个节点只合并孩子给出的稳定信息。

\textbf{证明和边界。} 证明从子树归纳开始：孩子状态正确时，父节点按题目规则合并后也正确。边界：空树、单节点、全左链、全右链。变体：若给前序和后序，二叉树通常不能唯一确定，除非额外限制。

\noindent\textbf{LeetCode 108 有序数组转二叉搜索树。} 模型：有序数组中点作为根能让左右规模接近。推导重点：每次选择区间中点，左右子区间构造左右子树。

\textbf{二刷读题。} 读题时先判断状态方向：父到子、子到父，还是按层传播。本题可以压成“有序数组中点作为根能让左右规模接近”，每个节点的输入和输出状态必须写清。先遮住答案，只保留题面，复述候选对象、合法性条件和输出目标。

\textbf{暴力回放。} 慢版本在每个节点重新扫描子树或路径，先求出局部答案。重复扫描暴露了真正状态：每个节点只需要从孩子或父亲那里拿到少量信息。二刷时先写慢版本或口述慢版本，用它确认不会漏答案。

\textbf{特例回放。} 拿 [-10,-3,0,5,9] 手算，选 0 做根，左边两个数构成左子树，右边两个数构成右子树，高度接近平衡。规律是有序数组的中点天然把值域分成 BST 左右两半，空区间返回空节点。把这个特例继续拆开看：节点向父亲返回的量和全局答案通常不是同一个量。先写叶子或空节点的返回值，再写父节点如何合并左右孩子，递归式才有边界基础。复盘时把这个特例压成状态字段：节点指针、传入约束或返回值、当前答案、层号或路径状态；空节点的返回值必须和状态定义匹配；再用边界 空区间、偶数长度选左中还是右中、高度平衡定义 反查初始化、更新顺序和退出条件。

\textbf{状态回放。} 手算路线：手算时从叶子开始写返回值，或者从根开始写传入约束。不要只写访问顺序，要写每条边上传递的信息。状态字段：节点指针、传入约束或返回值、当前答案、层号或路径状态；空节点的返回值必须和状态定义匹配。状态升级：把对子树的重复扫描压成返回值或队列状态；每个节点只合并孩子给出的稳定信息。

\textbf{证明和边界。} 证明从子树归纳开始：孩子状态正确时，父节点按题目规则合并后也正确。边界：空区间、偶数长度选左中还是右中、高度平衡定义。变体：工程上递归可读，若严格避免递归可用队列保存区间和父指针。

\noindent\textbf{LeetCode 110 平衡二叉树。} 模型：每个节点需要知道子树高度以及是否已经失衡。推导重点：后序汇总时一旦子树失衡就向上返回失败标记，避免重复算高度。

\textbf{二刷读题。} 读题时先判断状态方向：父到子、子到父，还是按层传播。本题可以压成“每个节点需要知道子树高度以及是否已经失衡”，每个节点的输入和输出状态必须写清。先遮住答案，只保留题面，复述候选对象、合法性条件和输出目标。

\textbf{暴力回放。} 慢版本在每个节点重新扫描子树或路径，先求出局部答案。重复扫描暴露了真正状态：每个节点只需要从孩子或父亲那里拿到少量信息。二刷时先写慢版本或口述慢版本，用它确认不会漏答案。

\textbf{特例回放。} 拿链式树 1->2->3 手算，根的左右高度差为 2，已经不平衡；若子树已经失衡，父节点没必要继续精算高度。规律是后序返回高度或失败标记，空树高度为 0，左右差正好 1 仍合法。把这个特例继续拆开看：节点向父亲返回的量和全局答案通常不是同一个量。先写叶子或空节点的返回值，再写父节点如何合并左右孩子，递归式才有边界基础。复盘时把这个特例压成状态字段：节点指针、传入约束或返回值、当前答案、层号或路径状态；空节点的返回值必须和状态定义匹配；再用边界 空树、单节点、左右高度差正好一、链式树 反查初始化、更新顺序和退出条件。

\textbf{状态回放。} 手算路线：手算时从叶子开始写返回值，或者从根开始写传入约束。不要只写访问顺序，要写每条边上传递的信息。状态字段：节点指针、传入约束或返回值、当前答案、层号或路径状态；空节点的返回值必须和状态定义匹配。状态升级：把对子树的重复扫描压成返回值或队列状态；每个节点只合并孩子给出的稳定信息。

\textbf{证明和边界。} 证明从子树归纳开始：孩子状态正确时，父节点按题目规则合并后也正确。边界：空树、单节点、左右高度差正好一、链式树。变体：可以把返回值设计成 optional 高度，空表示不平衡。

\noindent\textbf{LeetCode 111 二叉树最小深度。} 模型：最小深度是到最近叶子的节点数，不是左右深度简单取最小。推导重点：BFS 第一次遇到叶子就是最小深度，因为按层扩展。

\textbf{二刷读题。} 读题时先判断状态方向：父到子、子到父，还是按层传播。本题可以压成“最小深度是到最近叶子的节点数，不是左右深度简单取最小”，每个节点的输入和输出状态必须写清。先遮住答案，只保留题面，复述候选对象、合法性条件和输出目标。

\textbf{暴力回放。} 慢版本在每个节点重新扫描子树或路径，先求出局部答案。重复扫描暴露了真正状态：每个节点只需要从孩子或父亲那里拿到少量信息。二刷时先写慢版本或口述慢版本，用它确认不会漏答案。

\textbf{特例回放。} 拿根 1 只有左孩子 2 手算，最小深度是 2，不是 min(左深度1,右深度0)+1。规律是空子树不能当作叶子路径；BFS 第一次遇到叶子最直接，因为层次递增。把这个特例继续拆开看：节点向父亲返回的量和全局答案通常不是同一个量。先写叶子或空节点的返回值，再写父节点如何合并左右孩子，递归式才有边界基础。复盘时把这个特例压成状态字段：节点指针、传入约束或返回值、当前答案、层号或路径状态；空节点的返回值必须和状态定义匹配；再用边界 只有左子树、只有右子树、根为空、根是叶子 反查初始化、更新顺序和退出条件。

\textbf{状态回放。} 手算路线：手算时从叶子开始写返回值，或者从根开始写传入约束。不要只写访问顺序，要写每条边上传递的信息。状态字段：节点指针、传入约束或返回值、当前答案、层号或路径状态；空节点的返回值必须和状态定义匹配。状态升级：把对子树的重复扫描压成返回值或队列状态；每个节点只合并孩子给出的稳定信息。

\textbf{证明和边界。} 证明从子树归纳开始：孩子状态正确时，父节点按题目规则合并后也正确。边界：只有左子树、只有右子树、根为空、根是叶子。变体：DFS 写法必须特殊处理空子树不能作为最小路径。

\noindent\textbf{LeetCode 112 路径总和。} 模型：路径必须从根到叶子，状态是剩余目标值。推导重点：沿路径减去当前节点值，到叶子时检查剩余是否等于叶子值。

\textbf{二刷读题。} 读题时先判断状态方向：父到子、子到父，还是按层传播。本题可以压成“路径必须从根到叶子，状态是剩余目标值”，每个节点的输入和输出状态必须写清。先遮住答案，只保留题面，复述候选对象、合法性条件和输出目标。

\textbf{暴力回放。} 慢版本在每个节点重新扫描子树或路径，先求出局部答案。重复扫描暴露了真正状态：每个节点只需要从孩子或父亲那里拿到少量信息。二刷时先写慢版本或口述慢版本，用它确认不会漏答案。

\textbf{特例回放。} 拿路径 5->4->11、target=20 手算，沿途把目标减成 15、11、0，到叶子时正好为 0 才成功；非叶子提前等于目标不算。规律是状态是剩余目标值，终止条件必须同时满足叶子节点。把这个特例继续拆开看：节点向父亲返回的量和全局答案通常不是同一个量。先写叶子或空节点的返回值，再写父节点如何合并左右孩子，递归式才有边界基础。复盘时把这个特例压成状态字段：节点指针、传入约束或返回值、当前答案、层号或路径状态；空节点的返回值必须和状态定义匹配；再用边界 负数节点、目标为零、非叶子提前达到目标、空树 反查初始化、更新顺序和退出条件。

\textbf{状态回放。} 手算路线：手算时从叶子开始写返回值，或者从根开始写传入约束。不要只写访问顺序，要写每条边上传递的信息。状态字段：节点指针、传入约束或返回值、当前答案、层号或路径状态；空节点的返回值必须和状态定义匹配。状态升级：把对子树的重复扫描压成返回值或队列状态；每个节点只合并孩子给出的稳定信息。

\textbf{证明和边界。} 证明从子树归纳开始：孩子状态正确时，父节点按题目规则合并后也正确。边界：负数节点、目标为零、非叶子提前达到目标、空树。变体：若要求列出所有路径，需要保存路径数组并在回退时弹出。

\noindent\textbf{LeetCode 113 路径总和 II。} 模型：相比判定题，状态还要保存当前路径。推导重点：到叶子且和满足时复制路径；回溯时恢复路径。

\textbf{二刷读题。} 读题时先判断状态方向：父到子、子到父，还是按层传播。本题可以压成“相比判定题，状态还要保存当前路径”，每个节点的输入和输出状态必须写清。先遮住答案，只保留题面，复述候选对象、合法性条件和输出目标。

\textbf{暴力回放。} 慢版本在每个节点重新扫描子树或路径，先求出局部答案。重复扫描暴露了真正状态：每个节点只需要从孩子或父亲那里拿到少量信息。二刷时先写慢版本或口述慢版本，用它确认不会漏答案。

\textbf{特例回放。} 拿 target=22 和路径 5->4->11->2 手算，到叶子时剩余为 0，需要复制当前路径；回退后要弹出 2，继续找别的分支。规律是判定题多了路径数组，保存答案必须复制，不能保存同一个可变容器引用。把这个特例继续拆开看：节点向父亲返回的量和全局答案通常不是同一个量。先写叶子或空节点的返回值，再写父节点如何合并左右孩子，递归式才有边界基础。复盘时把这个特例压成状态字段：节点指针、传入约束或返回值、当前答案、层号或路径状态；空节点的返回值必须和状态定义匹配；再用边界 多个答案、负数、空树、路径数组复制时机 反查初始化、更新顺序和退出条件。

\textbf{状态回放。} 手算路线：手算时从叶子开始写返回值，或者从根开始写传入约束。不要只写访问顺序，要写每条边上传递的信息。状态字段：节点指针、传入约束或返回值、当前答案、层号或路径状态；空节点的返回值必须和状态定义匹配。状态升级：把对子树的重复扫描压成返回值或队列状态；每个节点只合并孩子给出的稳定信息。

\textbf{证明和边界。} 证明从子树归纳开始：孩子状态正确时，父节点按题目规则合并后也正确。边界：多个答案、负数、空树、路径数组复制时机。变体：若路径不要求从根到叶，模型会变成前缀和加哈希。

\noindent\textbf{LeetCode 114 二叉树展开为链表。} 模型：目标是按前序顺序把树原地改成右指针链。推导重点：显式栈按右后左先入顺序处理，逐个接到前驱右侧。

\textbf{二刷读题。} 读题时先判断状态方向：父到子、子到父，还是按层传播。本题可以压成“目标是按前序顺序把树原地改成右指针链”，每个节点的输入和输出状态必须写清。先遮住答案，只保留题面，复述候选对象、合法性条件和输出目标。

\textbf{暴力回放。} 慢版本在每个节点重新扫描子树或路径，先求出局部答案。重复扫描暴露了真正状态：每个节点只需要从孩子或父亲那里拿到少量信息。二刷时先写慢版本或口述慢版本，用它确认不会漏答案。

\textbf{特例回放。} 拿 1 的左子树 2 和右子树 5 手算，前序链应该是 1->2->...->5；若直接把左子树接到右边而不保存原右子树，5 会丢失。规律是按前序顺序重连，左指针置空，原右子树必须接到左子树展开后的尾部。把这个特例继续拆开看：节点向父亲返回的量和全局答案通常不是同一个量。先写叶子或空节点的返回值，再写父节点如何合并左右孩子，递归式才有边界基础。复盘时把这个特例压成状态字段：节点指针、传入约束或返回值、当前答案、层号或路径状态；空节点的返回值必须和状态定义匹配；再用边界 空树、单节点、原有右子树不能丢、左指针置空 反查初始化、更新顺序和退出条件。

\textbf{状态回放。} 手算路线：手算时从叶子开始写返回值，或者从根开始写传入约束。不要只写访问顺序，要写每条边上传递的信息。状态字段：节点指针、传入约束或返回值、当前答案、层号或路径状态；空节点的返回值必须和状态定义匹配。状态升级：把对子树的重复扫描压成返回值或队列状态；每个节点只合并孩子给出的稳定信息。

\textbf{证明和边界。} 证明从子树归纳开始：孩子状态正确时，父节点按题目规则合并后也正确。边界：空树、单节点、原有右子树不能丢、左指针置空。变体：也可用寻找左子树最右节点的原地方法。

\noindent\textbf{LeetCode 124 二叉树最大路径和。} 模型：路径可以从任意节点到任意节点，但不能分叉向父节点贡献两边。推导重点：每个节点向父亲贡献单边最大增益，全局答案可用左右两边加当前节点更新。

\textbf{二刷读题。} 读题时先判断状态方向：父到子、子到父，还是按层传播。本题可以压成“路径可以从任意节点到任意节点，但不能分叉向父节点贡献两边”，每个节点的输入和输出状态必须写清。先遮住答案，只保留题面，复述候选对象、合法性条件和输出目标。

\textbf{暴力回放。} 慢版本在每个节点重新扫描子树或路径，先求出局部答案。重复扫描暴露了真正状态：每个节点只需要从孩子或父亲那里拿到少量信息。二刷时先写慢版本或口述慢版本，用它确认不会漏答案。

\textbf{特例回放。} 拿根 -10、左 9、右 20 手算，经过 20 的路径可以同时吃左增益 15 和右增益 7 更新全局答案，但向父亲只能贡献一条单边路径。规律是返回给父亲的是单边最大增益，全局答案才允许左右两边加当前节点。把这个特例继续拆开看：节点向父亲返回的量和全局答案通常不是同一个量。先写叶子或空节点的返回值，再写父节点如何合并左右孩子，递归式才有边界基础。复盘时把这个特例压成状态字段：节点指针、传入约束或返回值、当前答案、层号或路径状态；空节点的返回值必须和状态定义匹配；再用边界 全负数、单节点、左右增益为负时取零、答案初始值 反查初始化、更新顺序和退出条件。

\textbf{状态回放。} 手算路线：手算时从叶子开始写返回值，或者从根开始写传入约束。不要只写访问顺序，要写每条边上传递的信息。状态字段：节点指针、传入约束或返回值、当前答案、层号或路径状态；空节点的返回值必须和状态定义匹配。状态升级：把对子树的重复扫描压成返回值或队列状态；每个节点只合并孩子给出的稳定信息。

\textbf{证明和边界。} 证明从子树归纳开始：孩子状态正确时，父节点按题目规则合并后也正确。边界：全负数、单节点、左右增益为负时取零、答案初始值。变体：若路径必须从根到叶，状态和答案位置完全不同。

\noindent\textbf{LeetCode 199 二叉树右视图。} 模型：每层从右侧能看到的就是该层最后访问或最先访问的右侧节点。推导重点：BFS 固定层边界，取每层最后一个；或 DFS 按右优先首次到达深度。

\textbf{二刷读题。} 读题时先判断状态方向：父到子、子到父，还是按层传播。本题可以压成“每层从右侧能看到的就是该层最后访问或最先访问的右侧节点”，每个节点的输入和输出状态必须写清。先遮住答案，只保留题面，复述候选对象、合法性条件和输出目标。

\textbf{暴力回放。} 慢版本在每个节点重新扫描子树或路径，先求出局部答案。重复扫描暴露了真正状态：每个节点只需要从孩子或父亲那里拿到少量信息。二刷时先写慢版本或口述慢版本，用它确认不会漏答案。

\textbf{特例回放。} 拿根 1、左 2、右 3 手算，从右侧看第一层是 1，第二层是 3；若右边缺失，才会看到左侧节点。规律是 BFS 每层最后一个节点，或 DFS 优先右子树并第一次到达某层时记录。把这个特例继续拆开看：节点向父亲返回的量和全局答案通常不是同一个量。先写叶子或空节点的返回值，再写父节点如何合并左右孩子，递归式才有边界基础。复盘时把这个特例压成状态字段：节点指针、传入约束或返回值、当前答案、层号或路径状态；空节点的返回值必须和状态定义匹配；再用边界 空树、左链、右链、左右交错 反查初始化、更新顺序和退出条件。

\textbf{状态回放。} 手算路线：手算时从叶子开始写返回值，或者从根开始写传入约束。不要只写访问顺序，要写每条边上传递的信息。状态字段：节点指针、传入约束或返回值、当前答案、层号或路径状态；空节点的返回值必须和状态定义匹配。状态升级：把对子树的重复扫描压成返回值或队列状态；每个节点只合并孩子给出的稳定信息。

\textbf{证明和边界。} 证明从子树归纳开始：孩子状态正确时，父节点按题目规则合并后也正确。边界：空树、左链、右链、左右交错。变体：若求左视图，只改变层内取值方向。

\noindent\textbf{LeetCode 208 实现 Trie。} 模型：前缀树把字符串公共前缀共享成路径。推导重点：插入和查询都逐字符沿边移动，不存在的边按需要创建。

\textbf{二刷读题。} 读题时先判断状态方向：父到子、子到父，还是按层传播。本题可以压成“前缀树把字符串公共前缀共享成路径”，每个节点的输入和输出状态必须写清。先遮住答案，只保留题面，复述候选对象、合法性条件和输出目标。

\textbf{暴力回放。} 慢版本在每个节点重新扫描子树或路径，先求出局部答案。重复扫描暴露了真正状态：每个节点只需要从孩子或父亲那里拿到少量信息。二刷时先写慢版本或口述慢版本，用它确认不会漏答案。

\textbf{特例回放。} 拿插入 apple 再查 app 手算，app 的路径存在，但如果 p 节点没有终止标记，search(app) 应为假，startsWith(app) 才为真。规律是 Trie 节点既要有孩子边，也要有单词终止标记。把这个特例继续拆开看：节点向父亲返回的量和全局答案通常不是同一个量。先写叶子或空节点的返回值，再写父节点如何合并左右孩子，递归式才有边界基础。复盘时把这个特例压成状态字段：节点指针、传入约束或返回值、当前答案、层号或路径状态；空节点的返回值必须和状态定义匹配；再用边界 空字符串约定、单词结束标记、前缀存在不等于单词存在 反查初始化、更新顺序和退出条件。

\textbf{状态回放。} 手算路线：手算时从叶子开始写返回值，或者从根开始写传入约束。不要只写访问顺序，要写每条边上传递的信息。状态字段：节点指针、传入约束或返回值、当前答案、层号或路径状态；空节点的返回值必须和状态定义匹配。状态升级：把对子树的重复扫描压成返回值或队列状态；每个节点只合并孩子给出的稳定信息。

\textbf{证明和边界。} 证明从子树归纳开始：孩子状态正确时，父节点按题目规则合并后也正确。边界：空字符串约定、单词结束标记、前缀存在不等于单词存在。变体：若字符集很大，用哈希子节点；小写字母用数组更快。

\noindent\textbf{LeetCode 211 添加与搜索单词。} 模型：点号通配符让查询可能分叉，但前缀树仍能剪掉不匹配前缀。推导重点：普通字符沿唯一边走，点号枚举当前节点所有孩子。

\textbf{二刷读题。} 读题时先判断状态方向：父到子、子到父，还是按层传播。本题可以压成“点号通配符让查询可能分叉，但前缀树仍能剪掉不匹配前缀”，每个节点的输入和输出状态必须写清。先遮住答案，只保留题面，复述候选对象、合法性条件和输出目标。

\textbf{暴力回放。} 慢版本在每个节点重新扫描子树或路径，先求出局部答案。重复扫描暴露了真正状态：每个节点只需要从孩子或父亲那里拿到少量信息。二刷时先写慢版本或口述慢版本，用它确认不会漏答案。

\textbf{特例回放。} 拿 add bad、dad、mad 后查 .ad 手算，点号可以走 b、d、m 三条边；查 b.. 时后两个点继续在子树里展开。规律是普通字符沿唯一边走，点号触发 DFS 枚举孩子，终点仍要检查单词结束标记。把这个特例继续拆开看：节点向父亲返回的量和全局答案通常不是同一个量。先写叶子或空节点的返回值，再写父节点如何合并左右孩子，递归式才有边界基础。复盘时把这个特例压成状态字段：节点指针、传入约束或返回值、当前答案、层号或路径状态；空节点的返回值必须和状态定义匹配；再用边界 连续通配符、空结果、单词结束标记、字符集大小 反查初始化、更新顺序和退出条件。

\textbf{状态回放。} 手算路线：手算时从叶子开始写返回值，或者从根开始写传入约束。不要只写访问顺序，要写每条边上传递的信息。状态字段：节点指针、传入约束或返回值、当前答案、层号或路径状态；空节点的返回值必须和状态定义匹配。状态升级：把对子树的重复扫描压成返回值或队列状态；每个节点只合并孩子给出的稳定信息。

\textbf{证明和边界。} 证明从子树归纳开始：孩子状态正确时，父节点按题目规则合并后也正确。边界：连续通配符、空结果、单词结束标记、字符集大小。变体：大量通配符会接近遍历整棵 Trie，需要规模约束。

\noindent\textbf{LeetCode 226 翻转二叉树。} 模型：每个节点交换左右孩子即可，遍历顺序不影响最终结构。推导重点：显式队列或栈逐个节点交换，避免递归深度。

\textbf{二刷读题。} 读题时先判断状态方向：父到子、子到父，还是按层传播。本题可以压成“每个节点交换左右孩子即可，遍历顺序不影响最终结构”，每个节点的输入和输出状态必须写清。先遮住答案，只保留题面，复述候选对象、合法性条件和输出目标。

\textbf{暴力回放。} 慢版本在每个节点重新扫描子树或路径，先求出局部答案。重复扫描暴露了真正状态：每个节点只需要从孩子或父亲那里拿到少量信息。二刷时先写慢版本或口述慢版本，用它确认不会漏答案。

\textbf{特例回放。} 拿根 4、左 2、右 7 手算，翻转后左右孩子交换，2 和 7 的子树也要递归交换。规律是每个节点只负责交换自己的左右指针，子树正确性由递归返回；空节点直接返回。把这个特例继续拆开看：节点向父亲返回的量和全局答案通常不是同一个量。先写叶子或空节点的返回值，再写父节点如何合并左右孩子，递归式才有边界基础。复盘时把这个特例压成状态字段：节点指针、传入约束或返回值、当前答案、层号或路径状态；空节点的返回值必须和状态定义匹配；再用边界 空树、单节点、只有一侧子树、重复值 反查初始化、更新顺序和退出条件。

\textbf{状态回放。} 手算路线：手算时从叶子开始写返回值，或者从根开始写传入约束。不要只写访问顺序，要写每条边上传递的信息。状态字段：节点指针、传入约束或返回值、当前答案、层号或路径状态；空节点的返回值必须和状态定义匹配。状态升级：把对子树的重复扫描压成返回值或队列状态；每个节点只合并孩子给出的稳定信息。

\textbf{证明和边界。} 证明从子树归纳开始：孩子状态正确时，父节点按题目规则合并后也正确。边界：空树、单节点、只有一侧子树、重复值。变体：这题简单但适合训练树节点指针修改。

\noindent\textbf{LeetCode 230 二叉搜索树中第 K 小。} 模型：BST 中序遍历按升序访问节点。推导重点：显式栈中序，访问第 k 个节点时返回。

\textbf{二刷读题。} 读题时先判断状态方向：父到子、子到父，还是按层传播。本题可以压成“BST 中序遍历按升序访问节点”，每个节点的输入和输出状态必须写清。先遮住答案，只保留题面，复述候选对象、合法性条件和输出目标。

\textbf{暴力回放。} 慢版本在每个节点重新扫描子树或路径，先求出局部答案。重复扫描暴露了真正状态：每个节点只需要从孩子或父亲那里拿到少量信息。二刷时先写慢版本或口述慢版本，用它确认不会漏答案。

\textbf{特例回放。} 拿 BST [3,1,4,null,2]、k=1 手算，中序访问顺序是 1,2,3,4，第一个就是答案。规律是 BST 的中序严格递增，计数到第 k 个即可停止；重复值是否允许要按题意定义。把这个特例继续拆开看：节点向父亲返回的量和全局答案通常不是同一个量。先写叶子或空节点的返回值，再写父节点如何合并左右孩子，递归式才有边界基础。复盘时把这个特例压成状态字段：节点指针、传入约束或返回值、当前答案、层号或路径状态；空节点的返回值必须和状态定义匹配；再用边界 k 为一、k 为节点数、树不平衡、是否有重复值 反查初始化、更新顺序和退出条件。

\textbf{状态回放。} 手算路线：手算时从叶子开始写返回值，或者从根开始写传入约束。不要只写访问顺序，要写每条边上传递的信息。状态字段：节点指针、传入约束或返回值、当前答案、层号或路径状态；空节点的返回值必须和状态定义匹配。状态升级：把对子树的重复扫描压成返回值或队列状态；每个节点只合并孩子给出的稳定信息。

\textbf{证明和边界。} 证明从子树归纳开始：孩子状态正确时，父节点按题目规则合并后也正确。边界：k 为一、k 为节点数、树不平衡、是否有重复值。变体：若频繁查询，可在节点维护子树大小成为顺序统计树。

\noindent\textbf{LeetCode 235 二叉搜索树最近公共祖先。} 模型：BST 的值域关系能决定两个目标在当前节点的哪一侧。推导重点：两个值都小去左，都大去右，否则当前节点就是分叉点。

\textbf{二刷读题。} 读题时先判断状态方向：父到子、子到父，还是按层传播。本题可以压成“BST 的值域关系能决定两个目标在当前节点的哪一侧”，每个节点的输入和输出状态必须写清。先遮住答案，只保留题面，复述候选对象、合法性条件和输出目标。

\textbf{暴力回放。} 慢版本在每个节点重新扫描子树或路径，先求出局部答案。重复扫描暴露了真正状态：每个节点只需要从孩子或父亲那里拿到少量信息。二刷时先写慢版本或口述慢版本，用它确认不会漏答案。

\textbf{特例回放。} 拿 BST 根 6，p=2，q=8 手算，2 在左侧、8 在右侧，根 6 就是最近公共祖先；若 p 和 q 都小于根，就继续去左子树。规律是 BST 的值域能直接决定搜索方向。把这个特例继续拆开看：节点向父亲返回的量和全局答案通常不是同一个量。先写叶子或空节点的返回值，再写父节点如何合并左右孩子，递归式才有边界基础。复盘时把这个特例压成状态字段：节点指针、传入约束或返回值、当前答案、层号或路径状态；空节点的返回值必须和状态定义匹配；再用边界 一个节点是另一个祖先、目标顺序、重复值约定 反查初始化、更新顺序和退出条件。

\textbf{状态回放。} 手算路线：手算时从叶子开始写返回值，或者从根开始写传入约束。不要只写访问顺序，要写每条边上传递的信息。状态字段：节点指针、传入约束或返回值、当前答案、层号或路径状态；空节点的返回值必须和状态定义匹配。状态升级：把对子树的重复扫描压成返回值或队列状态；每个节点只合并孩子给出的稳定信息。

\textbf{证明和边界。} 证明从子树归纳开始：孩子状态正确时，父节点按题目规则合并后也正确。边界：一个节点是另一个祖先、目标顺序、重复值约定。变体：普通二叉树不能用值域剪枝，需要父指针或后序状态。

\noindent\textbf{LeetCode 236 二叉树最近公共祖先。} 模型：普通树没有有序性，只能依靠祖先关系。推导重点：父指针集合或后序返回目标命中状态都可行。

\textbf{二刷读题。} 读题时先判断状态方向：父到子、子到父，还是按层传播。本题可以压成“普通树没有有序性，只能依靠祖先关系”，每个节点的输入和输出状态必须写清。先遮住答案，只保留题面，复述候选对象、合法性条件和输出目标。

\textbf{暴力回放。} 慢版本在每个节点重新扫描子树或路径，先求出局部答案。重复扫描暴露了真正状态：每个节点只需要从孩子或父亲那里拿到少量信息。二刷时先写慢版本或口述慢版本，用它确认不会漏答案。

\textbf{特例回放。} 拿普通二叉树根 3，p=5、q=1 手算，左右子树分别找到一个目标，根 3 就是最近公共祖先；若两个目标都在左子树，答案应由左子树返回。规律是后序返回是否找到目标或祖先，不能用 BST 的大小关系。把这个特例继续拆开看：节点向父亲返回的量和全局答案通常不是同一个量。先写叶子或空节点的返回值，再写父节点如何合并左右孩子，递归式才有边界基础。复盘时把这个特例压成状态字段：节点指针、传入约束或返回值、当前答案、层号或路径状态；空节点的返回值必须和状态定义匹配；再用边界 目标不存在、p 等于 q、一个是另一个祖先、比较节点地址 反查初始化、更新顺序和退出条件。

\textbf{状态回放。} 手算路线：手算时从叶子开始写返回值，或者从根开始写传入约束。不要只写访问顺序，要写每条边上传递的信息。状态字段：节点指针、传入约束或返回值、当前答案、层号或路径状态；空节点的返回值必须和状态定义匹配。状态升级：把对子树的重复扫描压成返回值或队列状态；每个节点只合并孩子给出的稳定信息。

\textbf{证明和边界。} 证明从子树归纳开始：孩子状态正确时，父节点按题目规则合并后也正确。边界：目标不存在、p 等于 q、一个是另一个祖先、比较节点地址。变体：多次 LCA 查询可预处理倍增祖先表。

\noindent\textbf{LeetCode 337 打家劫舍 III。} 模型：树上相邻节点不能同时选，每个节点需要偷和不偷两个状态。推导重点：后序汇总：偷当前则孩子不能偷，不偷当前则孩子可偷可不偷。

\textbf{二刷读题。} 读题时先判断状态方向：父到子、子到父，还是按层传播。本题可以压成“树上相邻节点不能同时选，每个节点需要偷和不偷两个状态”，每个节点的输入和输出状态必须写清。先遮住答案，只保留题面，复述候选对象、合法性条件和输出目标。

\textbf{暴力回放。} 慢版本在每个节点重新扫描子树或路径，先求出局部答案。重复扫描暴露了真正状态：每个节点只需要从孩子或父亲那里拿到少量信息。二刷时先写慢版本或口述慢版本，用它确认不会漏答案。

\textbf{特例回放。} 拿树 [3,2,3,null,3,null,1] 手算，偷根 3 时不能偷两个孩子，只能接孙子；不偷根时可以分别取左右孩子的最优。规律是每个节点返回两个状态：偷当前和不偷当前，父节点再组合。把这个特例继续拆开看：节点向父亲返回的量和全局答案通常不是同一个量。先写叶子或空节点的返回值，再写父节点如何合并左右孩子，递归式才有边界基础。复盘时把这个特例压成状态字段：节点指针、传入约束或返回值、当前答案、层号或路径状态；空节点的返回值必须和状态定义匹配；再用边界 空树、负值是否可能、链式树、递归深度 反查初始化、更新顺序和退出条件。

\textbf{状态回放。} 手算路线：手算时从叶子开始写返回值，或者从根开始写传入约束。不要只写访问顺序，要写每条边上传递的信息。状态字段：节点指针、传入约束或返回值、当前答案、层号或路径状态；空节点的返回值必须和状态定义匹配。状态升级：把对子树的重复扫描压成返回值或队列状态；每个节点只合并孩子给出的稳定信息。

\textbf{证明和边界。} 证明从子树归纳开始：孩子状态正确时，父节点按题目规则合并后也正确。边界：空树、负值是否可能、链式树、递归深度。变体：可用显式栈后序保存节点状态，避免调用栈。

\noindent\textbf{LeetCode 450 删除二叉搜索树中的节点。} 模型：删除节点后仍要保持 BST 中序有序。推导重点：若有两个孩子，用后继或前驱替换当前值，再删除那个后继节点。

\textbf{二刷读题。} 读题时先判断状态方向：父到子、子到父，还是按层传播。本题可以压成“删除节点后仍要保持 BST 中序有序”，每个节点的输入和输出状态必须写清。先遮住答案，只保留题面，复述候选对象、合法性条件和输出目标。

\textbf{暴力回放。} 慢版本在每个节点重新扫描子树或路径，先求出局部答案。重复扫描暴露了真正状态：每个节点只需要从孩子或父亲那里拿到少量信息。二刷时先写慢版本或口述慢版本，用它确认不会漏答案。

\textbf{特例回放。} 拿 BST 删除 3 且 3 有左右孩子手算，不能直接丢掉整棵子树；可用右子树最小节点接替 3，再在右子树删除这个后继。规律是删除分三类：叶子、单孩子、双孩子，双孩子要用前驱或后继保持 BST 顺序。把这个特例继续拆开看：节点向父亲返回的量和全局答案通常不是同一个量。先写叶子或空节点的返回值，再写父节点如何合并左右孩子，递归式才有边界基础。复盘时把这个特例压成状态字段：节点指针、传入约束或返回值、当前答案、层号或路径状态；空节点的返回值必须和状态定义匹配；再用边界 删除叶子、一个孩子、两个孩子、删除根节点 反查初始化、更新顺序和退出条件。

\textbf{状态回放。} 手算路线：手算时从叶子开始写返回值，或者从根开始写传入约束。不要只写访问顺序，要写每条边上传递的信息。状态字段：节点指针、传入约束或返回值、当前答案、层号或路径状态；空节点的返回值必须和状态定义匹配。状态升级：把对子树的重复扫描压成返回值或队列状态；每个节点只合并孩子给出的稳定信息。

\textbf{证明和边界。} 证明从子树归纳开始：孩子状态正确时，父节点按题目规则合并后也正确。边界：删除叶子、一个孩子、两个孩子、删除根节点。变体：工程实现要清楚是移动值还是重连节点。

\noindent\textbf{LeetCode 543 二叉树直径。} 模型：直径经过某节点时等于左高加右高，但向父亲只能贡献单边高度。推导重点：后序遍历同时更新全局直径并返回高度。

\textbf{二刷读题。} 读题时先判断状态方向：父到子、子到父，还是按层传播。本题可以压成“直径经过某节点时等于左高加右高，但向父亲只能贡献单边高度”，每个节点的输入和输出状态必须写清。先遮住答案，只保留题面，复述候选对象、合法性条件和输出目标。

\textbf{暴力回放。} 慢版本在每个节点重新扫描子树或路径，先求出局部答案。重复扫描暴露了真正状态：每个节点只需要从孩子或父亲那里拿到少量信息。二刷时先写慢版本或口述慢版本，用它确认不会漏答案。

\textbf{特例回放。} 拿树 1 的左右子树高度分别为 2 和 1 手算，经过根的直径是 3 条边；但向父节点只能返回单侧高度。规律是后序返回高度，全局答案用 left\_height+right\_height 更新。把这个特例继续拆开看：节点向父亲返回的量和全局答案通常不是同一个量。先写叶子或空节点的返回值，再写父节点如何合并左右孩子，递归式才有边界基础。复盘时把这个特例压成状态字段：节点指针、传入约束或返回值、当前答案、层号或路径状态；空节点的返回值必须和状态定义匹配；再用边界 空树、单节点、链式树、边数还是节点数定义 反查初始化、更新顺序和退出条件。

\textbf{状态回放。} 手算路线：手算时从叶子开始写返回值，或者从根开始写传入约束。不要只写访问顺序，要写每条边上传递的信息。状态字段：节点指针、传入约束或返回值、当前答案、层号或路径状态；空节点的返回值必须和状态定义匹配。状态升级：把对子树的重复扫描压成返回值或队列状态；每个节点只合并孩子给出的稳定信息。

\textbf{证明和边界。} 证明从子树归纳开始：孩子状态正确时，父节点按题目规则合并后也正确。边界：空树、单节点、链式树、边数还是节点数定义。变体：最大路径和是同型题，但贡献值和全局更新有负数处理。

\subsection{自测标准}

二刷时不要只确认自己见过这道题。请遮住答案，先讲题面模型，再讲暴力基线和瓶颈，然后说出优化结构保存了什么、不变量如何排除错误候选、边界实验如何打破错误实现。

\section{图建模、BFS 与 DFS：二刷复盘卡}
这一大节只围绕一个题型复盘，不把题号直接铺成目录。阅读时先看复盘目标，再读核心题卡，最后用自测标准检查自己是否能独立重讲。

\subsection{复盘目标}

追问通常会问节点和边的建模、访问标记时机、BFS 为什么最短。请准备说明状态是否还包含资源，为什么不能只按位置去重。

\subsection{核心题卡}

\noindent\textbf{LeetCode 126 单词接龙 II。} 模型：不仅要最短转换长度，还要输出所有最短转换路径。推导重点：BFS 只建立最短层级和前驱列表，同一层的多个前驱都保留，最后再从终点回溯路径。

\textbf{二刷读题。} 读题时先定义状态图。本题可以压成“不仅要最短转换长度，还要输出所有最短转换路径”，状态节点和题目原始位置不一定相同。先遮住答案，只保留题面，复述候选对象、合法性条件和输出目标。

\textbf{暴力回放。} 慢版本先按题意画出完整状态图，用普通搜索跑通可达性。这个过程用于确认不仅要最短转换长度，还要输出所有最短转换路径的节点和边，之后再讨论访问标记、层次或代价顺序。二刷时先写慢版本或口述慢版本，用它确认不会漏答案。

\textbf{特例回放。} 拿 hit 到 cog 的小词表手算，同一层 hot 可以通向 dot 和 lot，后续 dog/log 都可能到 cog。若某个单词在本层第一次见到就立刻删掉，会丢掉另一条同长度父路径。规律是 BFS 建最短层级，前驱列表保留同层多父节点，最后再回溯所有路径。把这个特例继续拆开看：状态节点不一定等于原始位置，还可能包含钥匙、剩余资源、时间或访问集合。visited 只能合并未来能力完全相同的状态，否则会把合法路径剪掉。复盘时把这个特例压成状态字段：状态编号、邻接生成方式、访问标记、距离或层数、队列内容；若状态含资源，visited 不能只按位置记录；再用边界 终点不在词表、同层多父节点、本层访问不能提前删除等长前驱、输出规模可能很大 反查初始化、更新顺序和退出条件。

\textbf{状态回放。} 手算路线：手算时画队列或栈，状态必须包含题目需要的所有资源。入队时标记还是出队时标记，要和去重语义一致。状态字段：状态编号、邻接生成方式、访问标记、距离或层数、队列内容；若状态含资源，visited 不能只按位置记录。状态升级：把重复搜索压成访问标记、距离数组或拓扑状态；邻居生成也要从全量扫描变成结构化枚举。

\textbf{证明和边界。} 证明从可达性开始：搜索覆盖所有合法边能到达的状态，访问标记只删除重复状态而不删除不同未来能力。边界：终点不在词表、同层多父节点、本层访问不能提前删除等长前驱、输出规模可能很大。变体：若只问最短长度，普通 BFS 或双向 BFS 更轻；若要路径，就必须保存前驱关系。

\noindent\textbf{LeetCode 127 单词接龙。} 模型：每个单词是节点，一次修改一个字符形成边。推导重点：通配模式桶加速邻居查询，BFS 第一次到达终点就是最短转换长度。

\textbf{二刷读题。} 读题时先定义状态图。本题可以压成“每个单词是节点，一次修改一个字符形成边”，状态节点和题目原始位置不一定相同。先遮住答案，只保留题面，复述候选对象、合法性条件和输出目标。

\textbf{暴力回放。} 慢版本先按题意画出完整状态图，用普通搜索跑通可达性。这个过程用于确认每个单词是节点，一次修改一个字符形成边的节点和边，之后再讨论访问标记、层次或代价顺序。二刷时先写慢版本或口述慢版本，用它确认不会漏答案。

\textbf{特例回放。} 拿 hit -> hot -> dot -> dog -> cog 手算，每次只能改一个字符，所以 BFS 第一层是 hot，第二层才到 dot 或 lot。规律是单词作为节点，通配桶如 h*t 把邻居快速找出；第一次到达终点就是最短转换长度。把这个特例继续拆开看：状态节点不一定等于原始位置，还可能包含钥匙、剩余资源、时间或访问集合。visited 只能合并未来能力完全相同的状态，否则会把合法路径剪掉。复盘时把这个特例压成状态字段：状态编号、邻接生成方式、访问标记、距离或层数、队列内容；若状态含资源，visited 不能只按位置记录；再用边界 终点不在词表、起点和终点相同、桶重复扫描、单词长度一致 反查初始化、更新顺序和退出条件。

\textbf{状态回放。} 手算路线：手算时画队列或栈，状态必须包含题目需要的所有资源。入队时标记还是出队时标记，要和去重语义一致。状态字段：状态编号、邻接生成方式、访问标记、距离或层数、队列内容；若状态含资源，visited 不能只按位置记录。状态升级：把重复搜索压成访问标记、距离数组或拓扑状态；邻居生成也要从全量扫描变成结构化枚举。

\textbf{证明和边界。} 证明从可达性开始：搜索覆盖所有合法边能到达的状态，访问标记只删除重复状态而不删除不同未来能力。边界：终点不在词表、起点和终点相同、桶重复扫描、单词长度一致。变体：双向 BFS 可以进一步降低搜索空间。

\noindent\textbf{LeetCode 130 被围绕的区域。} 模型：只有不连通到边界的 O 才会被翻转。推导重点：从边界 O 出发标记安全区域，再翻转剩余 O。

\textbf{二刷读题。} 读题时先定义状态图。本题可以压成“只有不连通到边界的 O 才会被翻转”，状态节点和题目原始位置不一定相同。先遮住答案，只保留题面，复述候选对象、合法性条件和输出目标。

\textbf{暴力回放。} 慢版本先按题意画出完整状态图，用普通搜索跑通可达性。这个过程用于确认只有不连通到边界的 O 才会被翻转的节点和边，之后再讨论访问标记、层次或代价顺序。二刷时先写慢版本或口述慢版本，用它确认不会漏答案。

\textbf{特例回放。} 拿边界上有 O、内部也有 O 的棋盘手算，边界 O 以及和它连通的 O 不能翻转，真正要翻的是不接触边界的内部块。规律是反向从边界 O 出发标记安全区，再翻转剩下的 O；单行单列全部接触边界。把这个特例继续拆开看：状态节点不一定等于原始位置，还可能包含钥匙、剩余资源、时间或访问集合。visited 只能合并未来能力完全相同的状态，否则会把合法路径剪掉。复盘时把这个特例压成状态字段：状态编号、邻接生成方式、访问标记、距离或层数、队列内容；若状态含资源，visited 不能只按位置记录；再用边界 全边界、无 O、单行单列、内部岛屿 反查初始化、更新顺序和退出条件。

\textbf{状态回放。} 手算路线：手算时画队列或栈，状态必须包含题目需要的所有资源。入队时标记还是出队时标记，要和去重语义一致。状态字段：状态编号、邻接生成方式、访问标记、距离或层数、队列内容；若状态含资源，visited 不能只按位置记录。状态升级：把重复搜索压成访问标记、距离数组或拓扑状态；邻居生成也要从全量扫描变成结构化枚举。

\textbf{证明和边界。} 证明从可达性开始：搜索覆盖所有合法边能到达的状态，访问标记只删除重复状态而不删除不同未来能力。边界：全边界、无 O、单行单列、内部岛屿。变体：这是反向思维：找不能被翻转的区域比直接判断每块区域更简单。

\noindent\textbf{LeetCode 133 克隆图。} 模型：每个原节点需要唯一对应一个新节点，边按原图连接。推导重点：哈希表保存原节点到克隆节点的映射，BFS 或 DFS 遍历图。

\textbf{二刷读题。} 读题时先定义状态图。本题可以压成“每个原节点需要唯一对应一个新节点，边按原图连接”，状态节点和题目原始位置不一定相同。先遮住答案，只保留题面，复述候选对象、合法性条件和输出目标。

\textbf{暴力回放。} 慢版本先按题意画出完整状态图，用普通搜索跑通可达性。这个过程用于确认每个原节点需要唯一对应一个新节点，边按原图连接的节点和边，之后再讨论访问标记、层次或代价顺序。二刷时先写慢版本或口述慢版本，用它确认不会漏答案。

\textbf{特例回放。} 拿三角形 1-2-3-1 手算，克隆 1 时会遇到 2 和 3；克隆 2 时又会遇到 1，如果没有哈希表会无限复制。规律是哈希表保存原节点到新节点的映射，先建节点再连邻居，环和自环都依赖这个映射。把这个特例继续拆开看：状态节点不一定等于原始位置，还可能包含钥匙、剩余资源、时间或访问集合。visited 只能合并未来能力完全相同的状态，否则会把合法路径剪掉。复盘时把这个特例压成状态字段：状态编号、邻接生成方式、访问标记、距离或层数、队列内容；若状态含资源，visited 不能只按位置记录；再用边界 空图、自环、重边、环结构、不能按节点值唯一假设 反查初始化、更新顺序和退出条件。

\textbf{状态回放。} 手算路线：手算时画队列或栈，状态必须包含题目需要的所有资源。入队时标记还是出队时标记，要和去重语义一致。状态字段：状态编号、邻接生成方式、访问标记、距离或层数、队列内容；若状态含资源，visited 不能只按位置记录。状态升级：把重复搜索压成访问标记、距离数组或拓扑状态；邻居生成也要从全量扫描变成结构化枚举。

\textbf{证明和边界。} 证明从可达性开始：搜索覆盖所有合法边能到达的状态，访问标记只删除重复状态而不删除不同未来能力。边界：空图、自环、重边、环结构、不能按节点值唯一假设。变体：若图节点值不唯一，必须用地址作为键。

\noindent\textbf{LeetCode 200 岛屿数量。} 模型：陆地格子是节点，四方向相邻是边，岛屿是连通块。推导重点：扫描遇到未访问陆地就启动搜索并标记整块。

\textbf{二刷读题。} 读题时先定义状态图。本题可以压成“陆地格子是节点，四方向相邻是边，岛屿是连通块”，状态节点和题目原始位置不一定相同。先遮住答案，只保留题面，复述候选对象、合法性条件和输出目标。

\textbf{暴力回放。} 慢版本先按题意画出完整状态图，用普通搜索跑通可达性。这个过程用于确认陆地格子是节点，四方向相邻是边，岛屿是连通块的节点和边，之后再讨论访问标记、层次或代价顺序。二刷时先写慢版本或口述慢版本，用它确认不会漏答案。

\textbf{特例回放。} 拿一个 2x2 小岛手算，从一个陆地出发能沿上下左右标记同一块岛。若不标记访问，会重复计数；若对角线也连通，会误合并。规律是每发现一个未访问陆地，答案加一，并用 DFS/BFS 标记与它四连通的全部陆地。把这个特例继续拆开看：状态节点不一定等于原始位置，还可能包含钥匙、剩余资源、时间或访问集合。visited 只能合并未来能力完全相同的状态，否则会把合法路径剪掉。复盘时把这个特例压成状态字段：状态编号、邻接生成方式、访问标记、距离或层数、队列内容；若状态含资源，visited 不能只按位置记录；再用边界 空网格、全水、全陆、斜对角不连通、是否允许修改输入 反查初始化、更新顺序和退出条件。

\textbf{状态回放。} 手算路线：手算时画队列或栈，状态必须包含题目需要的所有资源。入队时标记还是出队时标记，要和去重语义一致。状态字段：状态编号、邻接生成方式、访问标记、距离或层数、队列内容；若状态含资源，visited 不能只按位置记录。状态升级：把重复搜索压成访问标记、距离数组或拓扑状态；邻居生成也要从全量扫描变成结构化枚举。

\textbf{证明和边界。} 证明从可达性开始：搜索覆盖所有合法边能到达的状态，访问标记只删除重复状态而不删除不同未来能力。边界：空网格、全水、全陆、斜对角不连通、是否允许修改输入。变体：也可用并查集合并相邻陆地，适合动态岛屿扩展。

\noindent\textbf{LeetCode 207 课程表。} 模型：课程是节点，先修关系是有向边，问题是判断是否有环。推导重点：入度拓扑排序不断学习当前无依赖课程。

\textbf{二刷读题。} 读题时先定义状态图。本题可以压成“课程是节点，先修关系是有向边，问题是判断是否有环”，状态节点和题目原始位置不一定相同。先遮住答案，只保留题面，复述候选对象、合法性条件和输出目标。

\textbf{暴力回放。} 慢版本先按题意画出完整状态图，用普通搜索跑通可达性。这个过程用于确认课程是节点，先修关系是有向边，问题是判断是否有环的节点和边，之后再讨论访问标记、层次或代价顺序。二刷时先写慢版本或口述慢版本，用它确认不会漏答案。

\textbf{特例回放。} 拿课程 0->1 手算，入度为 0 的课程 0 先入队，学完后 1 的入度降为 0；若有 0->1 和 1->0，队列一开始为空。规律是拓扑排序用入度为 0 的节点推进，处理数量小于课程数说明有环。把这个特例继续拆开看：状态节点不一定等于原始位置，还可能包含钥匙、剩余资源、时间或访问集合。visited 只能合并未来能力完全相同的状态，否则会把合法路径剪掉。复盘时把这个特例压成状态字段：状态编号、邻接生成方式、访问标记、距离或层数、队列内容；若状态含资源，visited 不能只按位置记录；再用边界 边方向、孤立课程、环、自依赖、重复边 反查初始化、更新顺序和退出条件。

\textbf{状态回放。} 手算路线：手算时画队列或栈，状态必须包含题目需要的所有资源。入队时标记还是出队时标记，要和去重语义一致。状态字段：状态编号、邻接生成方式、访问标记、距离或层数、队列内容；若状态含资源，visited 不能只按位置记录。状态升级：把重复搜索压成访问标记、距离数组或拓扑状态；邻居生成也要从全量扫描变成结构化枚举。

\textbf{证明和边界。} 证明从可达性开始：搜索覆盖所有合法边能到达的状态，访问标记只删除重复状态而不删除不同未来能力。边界：边方向、孤立课程、环、自依赖、重复边。变体：DFS 三色标记也能判环，但要明确访问状态。

\noindent\textbf{LeetCode 210 课程表 II。} 模型：在判环基础上还要输出一个可行拓扑序。推导重点：每弹出一个入度为零课程，就把它加入顺序并删除出边。

\textbf{二刷读题。} 读题时先定义状态图。本题可以压成“在判环基础上还要输出一个可行拓扑序”，状态节点和题目原始位置不一定相同。先遮住答案，只保留题面，复述候选对象、合法性条件和输出目标。

\textbf{暴力回放。} 慢版本先按题意画出完整状态图，用普通搜索跑通可达性。这个过程用于确认在判环基础上还要输出一个可行拓扑序的节点和边，之后再讨论访问标记、层次或代价顺序。二刷时先写慢版本或口述慢版本，用它确认不会漏答案。

\textbf{特例回放。} 拿 0->1、0->2 手算，课程 0 入度为 0，先输出 0，再把 1 和 2 的入度减到 0。若存在环，输出数量会少于课程数。规律是课程表 II 比判定题多了输出顺序，队列弹出的顺序就是一种可行拓扑序。把这个特例继续拆开看：状态节点不一定等于原始位置，还可能包含钥匙、剩余资源、时间或访问集合。visited 只能合并未来能力完全相同的状态，否则会把合法路径剪掉。复盘时把这个特例压成状态字段：状态编号、邻接生成方式、访问标记、距离或层数、队列内容；若状态含资源，visited 不能只按位置记录；再用边界 有环返回空数组、多个可行顺序、边方向 反查初始化、更新顺序和退出条件。

\textbf{状态回放。} 手算路线：手算时画队列或栈，状态必须包含题目需要的所有资源。入队时标记还是出队时标记，要和去重语义一致。状态字段：状态编号、邻接生成方式、访问标记、距离或层数、队列内容；若状态含资源，visited 不能只按位置记录。状态升级：把重复搜索压成访问标记、距离数组或拓扑状态；邻居生成也要从全量扫描变成结构化枚举。

\textbf{证明和边界。} 证明从可达性开始：搜索覆盖所有合法边能到达的状态，访问标记只删除重复状态而不删除不同未来能力。边界：有环返回空数组、多个可行顺序、边方向。变体：若需要字典序最小拓扑序，把队列换成小顶堆。

\noindent\textbf{LeetCode 417 太平洋大西洋水流问题。} 模型：从每个格子向海搜索很慢，反向从海边沿非递减高度搜索。推导重点：分别标记能到太平洋和大西洋的格子，交集就是答案。

\textbf{二刷读题。} 读题时先定义状态图。本题可以压成“从每个格子向海搜索很慢，反向从海边沿非递减高度搜索”，状态节点和题目原始位置不一定相同。先遮住答案，只保留题面，复述候选对象、合法性条件和输出目标。

\textbf{暴力回放。} 慢版本先按题意画出完整状态图，用普通搜索跑通可达性。这个过程用于确认从每个格子向海搜索很慢，反向从海边沿非递减高度搜索的节点和边，之后再讨论访问标记、层次或代价顺序。二刷时先写慢版本或口述慢版本，用它确认不会漏答案。

\textbf{特例回放。} 拿高度矩阵 [[1,2],[4,3]] 手算，水从高处流向低处；反向从太平洋边界向内，只能走到更高或相等的格子。规律是分别从两个海洋边界反向 BFS/DFS，两个访问集合的交集就是答案。把这个特例继续拆开看：状态节点不一定等于原始位置，还可能包含钥匙、剩余资源、时间或访问集合。visited 只能合并未来能力完全相同的状态，否则会把合法路径剪掉。复盘时把这个特例压成状态字段：状态编号、邻接生成方式、访问标记、距离或层数、队列内容；若状态含资源，visited 不能只按位置记录；再用边界 单行单列、平台高度、边界重复入队、方向条件反向 反查初始化、更新顺序和退出条件。

\textbf{状态回放。} 手算路线：手算时画队列或栈，状态必须包含题目需要的所有资源。入队时标记还是出队时标记，要和去重语义一致。状态字段：状态编号、邻接生成方式、访问标记、距离或层数、队列内容；若状态含资源，visited 不能只按位置记录。状态升级：把重复搜索压成访问标记、距离数组或拓扑状态；邻居生成也要从全量扫描变成结构化枚举。

\textbf{证明和边界。} 证明从可达性开始：搜索覆盖所有合法边能到达的状态，访问标记只删除重复状态而不删除不同未来能力。边界：单行单列、平台高度、边界重复入队、方向条件反向。变体：反向搜索是很多可达性题的重要技巧。

\noindent\textbf{LeetCode 695 岛屿的最大面积。} 模型：面积就是一个连通块中陆地格子的数量。推导重点：每发现新陆地就搜索完整连通块并计数，取最大值。

\textbf{二刷读题。} 读题时先定义状态图。本题可以压成“面积就是一个连通块中陆地格子的数量”，状态节点和题目原始位置不一定相同。先遮住答案，只保留题面，复述候选对象、合法性条件和输出目标。

\textbf{暴力回放。} 慢版本先按题意画出完整状态图，用普通搜索跑通可达性。这个过程用于确认面积就是一个连通块中陆地格子的数量的节点和边，之后再讨论访问标记、层次或代价顺序。二刷时先写慢版本或口述慢版本，用它确认不会漏答案。

\textbf{特例回放。} 拿网格 [[1,1,0],[0,1,0]] 手算，从左上陆地出发能访问 3 个四连通格子；对角线不算连通。规律是每发现未访问陆地就 DFS/BFS 统计面积并标记，最大值随每个连通块更新。把这个特例继续拆开看：状态节点不一定等于原始位置，还可能包含钥匙、剩余资源、时间或访问集合。visited 只能合并未来能力完全相同的状态，否则会把合法路径剪掉。复盘时把这个特例压成状态字段：状态编号、邻接生成方式、访问标记、距离或层数、队列内容；若状态含资源，visited 不能只按位置记录；再用边界 全水、全陆、细长岛、是否允许修改输入 反查初始化、更新顺序和退出条件。

\textbf{状态回放。} 手算路线：手算时画队列或栈，状态必须包含题目需要的所有资源。入队时标记还是出队时标记，要和去重语义一致。状态字段：状态编号、邻接生成方式、访问标记、距离或层数、队列内容；若状态含资源，visited 不能只按位置记录。状态升级：把重复搜索压成访问标记、距离数组或拓扑状态；邻居生成也要从全量扫描变成结构化枚举。

\textbf{证明和边界。} 证明从可达性开始：搜索覆盖所有合法边能到达的状态，访问标记只删除重复状态而不删除不同未来能力。边界：全水、全陆、细长岛、是否允许修改输入。变体：与岛屿数量相比，只是连通块聚合值不同。

\noindent\textbf{LeetCode 752 打开转盘锁。} 模型：四位密码是状态，每次转动一位形成边。推导重点：BFS 从起点扩展，死亡状态不能入队，访问状态避免重复。

\textbf{二刷读题。} 读题时先定义状态图。本题可以压成“四位密码是状态，每次转动一位形成边”，状态节点和题目原始位置不一定相同。先遮住答案，只保留题面，复述候选对象、合法性条件和输出目标。

\textbf{暴力回放。} 慢版本先按题意画出完整状态图，用普通搜索跑通可达性。这个过程用于确认四位密码是状态，每次转动一位形成边的节点和边，之后再讨论访问标记、层次或代价顺序。二刷时先写慢版本或口述慢版本，用它确认不会漏答案。

\textbf{特例回放。} 拿锁 0000 到 0001 手算，每一步只能拨动一位加一或减一，0000 的邻居有 1000、9000、0100 等。规律是状态是四位字符串，BFS 按步数扩展，deadend 入队前就要排除。把这个特例继续拆开看：状态节点不一定等于原始位置，还可能包含钥匙、剩余资源、时间或访问集合。visited 只能合并未来能力完全相同的状态，否则会把合法路径剪掉。复盘时把这个特例压成状态字段：状态编号、邻接生成方式、访问标记、距离或层数、队列内容；若状态含资源，visited 不能只按位置记录；再用边界 起点死亡、目标就是起点、数字环绕、死亡集合很大 反查初始化、更新顺序和退出条件。

\textbf{状态回放。} 手算路线：手算时画队列或栈，状态必须包含题目需要的所有资源。入队时标记还是出队时标记，要和去重语义一致。状态字段：状态编号、邻接生成方式、访问标记、距离或层数、队列内容；若状态含资源，visited 不能只按位置记录。状态升级：把重复搜索压成访问标记、距离数组或拓扑状态；邻居生成也要从全量扫描变成结构化枚举。

\textbf{证明和边界。} 证明从可达性开始：搜索覆盖所有合法边能到达的状态，访问标记只删除重复状态而不删除不同未来能力。边界：起点死亡、目标就是起点、数字环绕、死亡集合很大。变体：双向 BFS 能明显减少状态展开。

\noindent\textbf{LeetCode 815 公交路线。} 模型：公交路线和站点构成二分关系，换乘次数比站点步数更重要。推导重点：用站点到路线列表的哈希表找可乘路线，以路线或乘车层数做 BFS，访问过的路线不再重复展开。

\textbf{二刷读题。} 读题时先定义状态图。本题可以压成“公交路线和站点构成二分关系，换乘次数比站点步数更重要”，状态节点和题目原始位置不一定相同。先遮住答案，只保留题面，复述候选对象、合法性条件和输出目标。

\textbf{暴力回放。} 慢版本先按题意画出完整状态图，用普通搜索跑通可达性。这个过程用于确认公交路线和站点构成二分关系，换乘次数比站点步数更重要的节点和边，之后再讨论访问标记、层次或代价顺序。二刷时先写慢版本或口述慢版本，用它确认不会漏答案。

\textbf{特例回放。} 拿起点站可坐路线 A，A 和 B 在某站相交，B 到终点手算。按站点一步步 BFS 会反复展开同一条路线的所有站；按路线 BFS，一次乘车层数就是换乘次数。规律是访问过的路线不再展开，站点到路线表只用于找到下一批路线。把这个特例继续拆开看：状态节点不一定等于原始位置，还可能包含钥匙、剩余资源、时间或访问集合。visited 只能合并未来能力完全相同的状态，否则会把合法路径剪掉。复盘时把这个特例压成状态字段：状态编号、邻接生成方式、访问标记、距离或层数、队列内容；若状态含资源，visited 不能只按位置记录；再用边界 起点等于终点、某站经过路线很多、路线重复访问、目标站不可达 反查初始化、更新顺序和退出条件。

\textbf{状态回放。} 手算路线：手算时画队列或栈，状态必须包含题目需要的所有资源。入队时标记还是出队时标记，要和去重语义一致。状态字段：状态编号、邻接生成方式、访问标记、距离或层数、队列内容；若状态含资源，visited 不能只按位置记录。状态升级：把重复搜索压成访问标记、距离数组或拓扑状态；邻居生成也要从全量扫描变成结构化枚举。

\textbf{证明和边界。} 证明从可达性开始：搜索覆盖所有合法边能到达的状态，访问标记只删除重复状态而不删除不同未来能力。边界：起点等于终点、某站经过路线很多、路线重复访问、目标站不可达。变体：若把每个站点之间完全建边会爆炸，应利用路线这一层压缩邻居生成。

\noindent\textbf{LeetCode 864 获取所有钥匙的最短路径。} 模型：位置相同但钥匙集合不同，未来能开的门不同。推导重点：BFS 状态包含行、列和钥匙掩码，访问标记也必须包含掩码。

\textbf{二刷读题。} 读题时先定义状态图。本题可以压成“位置相同但钥匙集合不同，未来能开的门不同”，状态节点和题目原始位置不一定相同。先遮住答案，只保留题面，复述候选对象、合法性条件和输出目标。

\textbf{暴力回放。} 慢版本先按题意画出完整状态图，用普通搜索跑通可达性。这个过程用于确认位置相同但钥匙集合不同，未来能开的门不同的节点和边，之后再讨论访问标记、层次或代价顺序。二刷时先写慢版本或口述慢版本，用它确认不会漏答案。

\textbf{特例回放。} 拿网格 @.a / \#.\# / b.A 手算，走到 a 后状态不是位置变了而是钥匙集合多了一把；同一位置带不同钥匙未来能力不同。规律是 BFS 的 visited 必须按 (row,col,keymask) 记录，拿齐所有钥匙第一次出队就是最短步数。把这个特例继续拆开看：状态节点不一定等于原始位置，还可能包含钥匙、剩余资源、时间或访问集合。visited 只能合并未来能力完全相同的状态，否则会把合法路径剪掉。复盘时把这个特例压成状态字段：状态编号、邻接生成方式、访问标记、距离或层数、队列内容；若状态含资源，visited 不能只按位置记录；再用边界 起点旁边有门、钥匙顺序、不可达、钥匙数量决定掩码大小 反查初始化、更新顺序和退出条件。

\textbf{状态回放。} 手算路线：手算时画队列或栈，状态必须包含题目需要的所有资源。入队时标记还是出队时标记，要和去重语义一致。状态字段：状态编号、邻接生成方式、访问标记、距离或层数、队列内容；若状态含资源，visited 不能只按位置记录。状态升级：把重复搜索压成访问标记、距离数组或拓扑状态；邻居生成也要从全量扫描变成结构化枚举。

\textbf{证明和边界。} 证明从可达性开始：搜索覆盖所有合法边能到达的状态，访问标记只删除重复状态而不删除不同未来能力。边界：起点旁边有门、钥匙顺序、不可达、钥匙数量决定掩码大小。变体：这是状态图维度不能只按位置去重的典型题。

\noindent\textbf{LeetCode 994 腐烂的橘子。} 模型：所有初始腐烂橘子同时作为 BFS 源点。推导重点：按层扩散，分钟数等于最后一个新鲜橘子被首次访问的层数。

\textbf{二刷读题。} 读题时先定义状态图。本题可以压成“所有初始腐烂橘子同时作为 BFS 源点”，状态节点和题目原始位置不一定相同。先遮住答案，只保留题面，复述候选对象、合法性条件和输出目标。

\textbf{暴力回放。} 慢版本先按题意画出完整状态图，用普通搜索跑通可达性。这个过程用于确认所有初始腐烂橘子同时作为 BFS 源点的节点和边，之后再讨论访问标记、层次或代价顺序。二刷时先写慢版本或口述慢版本，用它确认不会漏答案。

\textbf{特例回放。} 拿 [[2,1,1],[1,1,0],[0,1,1]] 手算，所有腐烂橘子同时向四周扩散，第一分钟会感染相邻新鲜橘子。规律是多源 BFS，初始所有腐烂橘子同时入队，层数就是分钟数。把这个特例继续拆开看：状态节点不一定等于原始位置，还可能包含钥匙、剩余资源、时间或访问集合。visited 只能合并未来能力完全相同的状态，否则会把合法路径剪掉。复盘时把这个特例压成状态字段：状态编号、邻接生成方式、访问标记、距离或层数、队列内容；若状态含资源，visited 不能只按位置记录；再用边界 没有新鲜橘子、没有腐烂源、隔离新鲜橘子、空格阻断 反查初始化、更新顺序和退出条件。

\textbf{状态回放。} 手算路线：手算时画队列或栈，状态必须包含题目需要的所有资源。入队时标记还是出队时标记，要和去重语义一致。状态字段：状态编号、邻接生成方式、访问标记、距离或层数、队列内容；若状态含资源，visited 不能只按位置记录。状态升级：把重复搜索压成访问标记、距离数组或拓扑状态；邻居生成也要从全量扫描变成结构化枚举。

\textbf{证明和边界。} 证明从可达性开始：搜索覆盖所有合法边能到达的状态，访问标记只删除重复状态而不删除不同未来能力。边界：没有新鲜橘子、没有腐烂源、隔离新鲜橘子、空格阻断。变体：多源 BFS 也适用于最近零距离、地图扩散等题。

\noindent\textbf{LeetCode 1091 二进制矩阵中的最短路径。} 模型：八方向网格无权边，答案是从左上到右下的最少格子数。推导重点：BFS 入队时标记访问，层数或距离随状态保存。

\textbf{二刷读题。} 读题时先定义状态图。本题可以压成“八方向网格无权边，答案是从左上到右下的最少格子数”，状态节点和题目原始位置不一定相同。先遮住答案，只保留题面，复述候选对象、合法性条件和输出目标。

\textbf{暴力回放。} 慢版本先按题意画出完整状态图，用普通搜索跑通可达性。这个过程用于确认八方向网格无权边，答案是从左上到右下的最少格子数的节点和边，之后再讨论访问标记、层次或代价顺序。二刷时先写慢版本或口述慢版本，用它确认不会漏答案。

\textbf{特例回放。} 拿 [[0,1],[1,0]] 手算，从左上可以斜着一步到右下，路径长度是 2；若起点或终点为 1 直接不可达。规律是八方向 BFS，第一次到达终点就是最短路径长度。把这个特例继续拆开看：状态节点不一定等于原始位置，还可能包含钥匙、剩余资源、时间或访问集合。visited 只能合并未来能力完全相同的状态，否则会把合法路径剪掉。复盘时把这个特例压成状态字段：状态编号、邻接生成方式、访问标记、距离或层数、队列内容；若状态含资源，visited 不能只按位置记录；再用边界 起点或终点阻塞、单格、八方向越界、无路可达 反查初始化、更新顺序和退出条件。

\textbf{状态回放。} 手算路线：手算时画队列或栈，状态必须包含题目需要的所有资源。入队时标记还是出队时标记，要和去重语义一致。状态字段：状态编号、邻接生成方式、访问标记、距离或层数、队列内容；若状态含资源，visited 不能只按位置记录。状态升级：把重复搜索压成访问标记、距离数组或拓扑状态；邻居生成也要从全量扫描变成结构化枚举。

\textbf{证明和边界。} 证明从可达性开始：搜索覆盖所有合法边能到达的状态，访问标记只删除重复状态而不删除不同未来能力。边界：起点或终点阻塞、单格、八方向越界、无路可达。变体：边权相同才能用 BFS，若每格有代价就要最短路。

\noindent\textbf{LeetCode 1293 网格中的最短路径。} 模型：走到同一格时，剩余可消除障碍数量不同会影响未来可达性。推导重点：BFS 状态包含位置和剩余消除次数，visited 记录该维度或记录到达该格的最大剩余次数。

\textbf{二刷读题。} 读题时先定义状态图。本题可以压成“走到同一格时，剩余可消除障碍数量不同会影响未来可达性”，状态节点和题目原始位置不一定相同。先遮住答案，只保留题面，复述候选对象、合法性条件和输出目标。

\textbf{暴力回放。} 慢版本先按题意画出完整状态图，用普通搜索跑通可达性。这个过程用于确认走到同一格时，剩余可消除障碍数量不同会影响未来可达性的节点和边，之后再讨论访问标记、层次或代价顺序。二刷时先写慢版本或口述慢版本，用它确认不会漏答案。

\textbf{特例回放。} 拿网格三列中间有障碍且 k=1 手算，走到同一格时，如果剩余消障次数不同，未来可走路径也不同。规律是 BFS 的状态必须包含剩余 k，visited 不能只按位置记录。把这个特例继续拆开看：状态节点不一定等于原始位置，还可能包含钥匙、剩余资源、时间或访问集合。visited 只能合并未来能力完全相同的状态，否则会把合法路径剪掉。复盘时把这个特例压成状态字段：状态编号、邻接生成方式、访问标记、距离或层数、队列内容；若状态含资源，visited 不能只按位置记录；再用边界 起点终点、k 很大、障碍密集、二维 visited 错误剪枝 反查初始化、更新顺序和退出条件。

\textbf{状态回放。} 手算路线：手算时画队列或栈，状态必须包含题目需要的所有资源。入队时标记还是出队时标记，要和去重语义一致。状态字段：状态编号、邻接生成方式、访问标记、距离或层数、队列内容；若状态含资源，visited 不能只按位置记录。状态升级：把重复搜索压成访问标记、距离数组或拓扑状态；邻居生成也要从全量扫描变成结构化枚举。

\textbf{证明和边界。} 证明从可达性开始：搜索覆盖所有合法边能到达的状态，访问标记只删除重复状态而不删除不同未来能力。边界：起点终点、k 很大、障碍密集、二维 visited 错误剪枝。变体：它和钥匙题一样说明状态维度不能丢。

\noindent\textbf{LeetCode 1462 课程表 IV。} 模型：大量先修关系查询需要预处理课程之间的可达性。推导重点：可以用 Floyd 传递闭包、每个点 BFS，或拓扑 DP 合并祖先集合，查询时直接查可达表。

\textbf{二刷读题。} 读题时先定义状态图。本题可以压成“大量先修关系查询需要预处理课程之间的可达性”，状态节点和题目原始位置不一定相同。先遮住答案，只保留题面，复述候选对象、合法性条件和输出目标。

\textbf{暴力回放。} 慢版本先按题意画出完整状态图，用普通搜索跑通可达性。这个过程用于确认大量先修关系查询需要预处理课程之间的可达性的节点和边，之后再讨论访问标记、层次或代价顺序。二刷时先写慢版本或口述慢版本，用它确认不会漏答案。

\textbf{特例回放。} 拿 0->1->2 和查询 0 是否先修 2 手算，直接边只能回答 0->1、1->2，不能回答传递关系。若查询很多，每次 DFS 会重复走相同路径。规律是先做可达性闭包或拓扑合并祖先集合，再 O(1) 回答查询。把这个特例继续拆开看：状态节点不一定等于原始位置，还可能包含钥匙、剩余资源、时间或访问集合。visited 只能合并未来能力完全相同的状态，否则会把合法路径剪掉。复盘时把这个特例压成状态字段：状态编号、邻接生成方式、访问标记、距离或层数、队列内容；若状态含资源，visited 不能只按位置记录；再用边界 课程之间没有关系、自身是否算先修、查询很多、图中按题意应为 DAG 反查初始化、更新顺序和退出条件。

\textbf{状态回放。} 手算路线：手算时画队列或栈，状态必须包含题目需要的所有资源。入队时标记还是出队时标记，要和去重语义一致。状态字段：状态编号、邻接生成方式、访问标记、距离或层数、队列内容；若状态含资源，visited 不能只按位置记录。状态升级：把重复搜索压成访问标记、距离数组或拓扑状态；邻居生成也要从全量扫描变成结构化枚举。

\textbf{证明和边界。} 证明从可达性开始：搜索覆盖所有合法边能到达的状态，访问标记只删除重复状态而不删除不同未来能力。边界：课程之间没有关系、自身是否算先修、查询很多、图中按题意应为 DAG。变体：如果查询很少，逐次搜索可能更省预处理成本；多查询时闭包更稳。

\subsection{自测标准}

二刷时不要只确认自己见过这道题。请遮住答案，先讲题面模型，再讲暴力基线和瓶颈，然后说出优化结构保存了什么、不变量如何排除错误候选、边界实验如何打破错误实现。

\section{最短路与带权图：二刷复盘卡}
这一大节只围绕一个题型复盘，不把题号直接铺成目录。阅读时先看复盘目标，再读核心题卡，最后用自测标准检查自己是否能独立重讲。

\subsection{复盘目标}

追问通常会问为什么普通 BFS 不行、Dijkstra 的非负边条件、堆中旧状态如何处理。请准备一个不同边权反例。

\subsection{核心题卡}

\noindent\textbf{LeetCode 505 迷宫 II。} 模型：小球一旦选择方向就会滚到墙前停止，边权是滚动经过的格子数。推导重点：邻居生成不是走一步，而是沿四个方向模拟滚动到停点，再用 Dijkstra 按累计距离松弛。

\textbf{二刷读题。} 读题时先区分可达、最少边数和最小代价。本题可以压成“小球一旦选择方向就会滚到墙前停止，边权是滚动经过的格子数”，邻居生成不是走一步，而是沿四个方向模拟滚动到停点，再用 Dijkstra 按累计距离松弛决定扩展顺序。先遮住答案，只保留题面，复述候选对象、合法性条件和输出目标。

\textbf{暴力回放。} 慢版本可以枚举路径，或先用普通队列试跑一个小图。它会暴露步数最少和代价最小是否一致；一旦不一致，就必须围绕代价组织状态。二刷时先写慢版本或口述慢版本，用它确认不会漏答案。

\textbf{特例回放。} 拿迷宫中球从 (0,4) 滚向左手算，它不会每格停下，只会直到撞墙才停。普通 BFS 按格走会错，边权是滚动距离。规律是状态是停球位置，邻居是四个方向滚到的终点，距离要用 Dijkstra 更新。把这个特例继续拆开看：队列或堆弹出的顺序必须和代价规则一致；旧状态被跳过，是因为已经有更小代价到达同一状态。只要边权或代价合成方式改变，就要重新证明扩展顺序。复盘时把这个特例压成状态字段：距离数组、优先队列状态、松弛候选、旧状态判定、不可达哨兵；资源约束题还要把资源维度放进状态；再用边界 起点就是终点、目标无法停下、绕远路更短、同一停点多次被松弛 反查初始化、更新顺序和退出条件。

\textbf{状态回放。} 手算路线：手算时记录每次弹出的代价、当前距离数组和松弛结果。旧状态被弹出时为什么跳过，也要写在表里。状态字段：距离数组、优先队列状态、松弛候选、旧状态判定、不可达哨兵；资源约束题还要把资源维度放进状态。状态升级：把路径枚举压成距离数组和优先队列；每次只扩展当前最有希望的未过期状态。

\textbf{证明和边界。} 证明从扩展顺序开始：当前弹出的状态在代价规则下已经不会被未来路径改得更优。边界：起点就是终点、目标无法停下、绕远路更短、同一停点多次被松弛。变体：若只问能否到达，可用 BFS 或 DFS；若问最少距离，就必须处理带权边。

\noindent\textbf{LeetCode 743 网络延迟时间。} 模型：有向正权图从起点到所有节点的最短路最大值。推导重点：Dijkstra 求每个节点最短距离，最后取最大，若有无穷则不可达。

\textbf{二刷读题。} 读题时先区分可达、最少边数和最小代价。本题可以压成“有向正权图从起点到所有节点的最短路最大值”，Dijkstra 求每个节点最短距离，最后取最大，若有无穷则不可达决定扩展顺序。先遮住答案，只保留题面，复述候选对象、合法性条件和输出目标。

\textbf{暴力回放。} 慢版本可以枚举路径，或先用普通队列试跑一个小图。它会暴露步数最少和代价最小是否一致；一旦不一致，就必须围绕代价组织状态。二刷时先写慢版本或口述慢版本，用它确认不会漏答案。

\textbf{特例回放。} 拿 times=[[2,1,1],[2,3,1],[3,4,1]]、起点 2 手算，2 到 1 和 3 距离都是 1，到 4 距离是 2。规律是非负边最短路用 Dijkstra，最后取所有最短距离的最大值；不可达节点导致答案 -1。把这个特例继续拆开看：队列或堆弹出的顺序必须和代价规则一致；旧状态被跳过，是因为已经有更小代价到达同一状态。只要边权或代价合成方式改变，就要重新证明扩展顺序。复盘时把这个特例压成状态字段：距离数组、优先队列状态、松弛候选、旧状态判定、不可达哨兵；资源约束题还要把资源维度放进状态；再用边界 节点编号从一开始、不可达、重边、过期堆元素 反查初始化、更新顺序和退出条件。

\textbf{状态回放。} 手算路线：手算时记录每次弹出的代价、当前距离数组和松弛结果。旧状态被弹出时为什么跳过，也要写在表里。状态字段：距离数组、优先队列状态、松弛候选、旧状态判定、不可达哨兵；资源约束题还要把资源维度放进状态。状态升级：把路径枚举压成距离数组和优先队列；每次只扩展当前最有希望的未过期状态。

\textbf{证明和边界。} 证明从扩展顺序开始：当前弹出的状态在代价规则下已经不会被未来路径改得更优。边界：节点编号从一开始、不可达、重边、过期堆元素。变体：边权相同才可用普通 BFS。

\noindent\textbf{LeetCode 778 水位上升的泳池中游泳。} 模型：到达某格子的代价是路径上最大高度，而不是高度和。推导重点：用 Dijkstra 按当前路径最大高度扩展，松弛时取 max 当前代价和邻格高度。

\textbf{二刷读题。} 读题时先区分可达、最少边数和最小代价。本题可以压成“到达某格子的代价是路径上最大高度，而不是高度和”，用 Dijkstra 按当前路径最大高度扩展，松弛时取 max 当前代价和邻格高度决定扩展顺序。先遮住答案，只保留题面，复述候选对象、合法性条件和输出目标。

\textbf{暴力回放。} 慢版本可以枚举路径，或先用普通队列试跑一个小图。它会暴露步数最少和代价最小是否一致；一旦不一致，就必须围绕代价组织状态。二刷时先写慢版本或口述慢版本，用它确认不会漏答案。

\textbf{特例回放。} 拿 [[0,2],[1,3]] 手算，时间必须至少到 3 才能进入右下角；路径代价不是边权和，而是沿路最大高度。规律是 Dijkstra 的距离定义为到达某格所需的最小最大高度，或二分时间做可达性。把这个特例继续拆开看：队列或堆弹出的顺序必须和代价规则一致；旧状态被跳过，是因为已经有更小代价到达同一状态。只要边权或代价合成方式改变，就要重新证明扩展顺序。复盘时把这个特例压成状态字段：距离数组、优先队列状态、松弛候选、旧状态判定、不可达哨兵；资源约束题还要把资源维度放进状态；再用边界 单格、起点高度很高、需要绕路、多个相同高度 反查初始化、更新顺序和退出条件。

\textbf{状态回放。} 手算路线：手算时记录每次弹出的代价、当前距离数组和松弛结果。旧状态被弹出时为什么跳过，也要写在表里。状态字段：距离数组、优先队列状态、松弛候选、旧状态判定、不可达哨兵；资源约束题还要把资源维度放进状态。状态升级：把路径枚举压成距离数组和优先队列；每次只扩展当前最有希望的未过期状态。

\textbf{证明和边界。} 证明从扩展顺序开始：当前弹出的状态在代价规则下已经不会被未来路径改得更优。边界：单格、起点高度很高、需要绕路、多个相同高度。变体：也可以按水位二分可达，或按高度排序并查集合并。

\noindent\textbf{LeetCode 787 K 站中转内最便宜的航班。} 模型：路径代价受中转次数限制，不能只按费用最小弹出节点后永久确定。推导重点：用按边数分层的 Bellman-Ford 或带层数状态的优先队列，状态包含城市和已用边数。

\textbf{二刷读题。} 读题时先区分可达、最少边数和最小代价。本题可以压成“路径代价受中转次数限制，不能只按费用最小弹出节点后永久确定”，用按边数分层的 Bellman-Ford 或带层数状态的优先队列，状态包含城市和已用边数决定扩展顺序。先遮住答案，只保留题面，复述候选对象、合法性条件和输出目标。

\textbf{暴力回放。} 慢版本可以枚举路径，或先用普通队列试跑一个小图。它会暴露步数最少和代价最小是否一致；一旦不一致，就必须围绕代价组织状态。二刷时先写慢版本或口述慢版本，用它确认不会漏答案。

\textbf{特例回放。} 拿航班 0->1 价格 100、1->2 价格 100、0->2 价格 500，K=1 手算，允许一次中转时 0->1->2 更便宜。规律是状态必须包含已用中转次数或边数，不能只保存每个城市一个最低价格。把这个特例继续拆开看：队列或堆弹出的顺序必须和代价规则一致；旧状态被跳过，是因为已经有更小代价到达同一状态。只要边权或代价合成方式改变，就要重新证明扩展顺序。复盘时把这个特例压成状态字段：距离数组、优先队列状态、松弛候选、旧状态判定、不可达哨兵；资源约束题还要把资源维度放进状态；再用边界 起点终点相同、不可达、K 为零、便宜路径中转过多 反查初始化、更新顺序和退出条件。

\textbf{状态回放。} 手算路线：手算时记录每次弹出的代价、当前距离数组和松弛结果。旧状态被弹出时为什么跳过，也要写在表里。状态字段：距离数组、优先队列状态、松弛候选、旧状态判定、不可达哨兵；资源约束题还要把资源维度放进状态。状态升级：把路径枚举压成距离数组和优先队列；每次只扩展当前最有希望的未过期状态。

\textbf{证明和边界。} 证明从扩展顺序开始：当前弹出的状态在代价规则下已经不会被未来路径改得更优。边界：起点终点相同、不可达、K 为零、便宜路径中转过多。变体：若去掉中转限制，普通 Dijkstra 才能直接按费用组织扩展。

\noindent\textbf{LeetCode 1102 得分最高的路径。} 模型：路径得分是沿途最小格子值，要最大化这个最小值。推导重点：用最大堆按当前路径得分扩展，进入邻格时取当前得分和邻格值的较小者。

\textbf{二刷读题。} 读题时先区分可达、最少边数和最小代价。本题可以压成“路径得分是沿途最小格子值，要最大化这个最小值”，用最大堆按当前路径得分扩展，进入邻格时取当前得分和邻格值的较小者决定扩展顺序。先遮住答案，只保留题面，复述候选对象、合法性条件和输出目标。

\textbf{暴力回放。} 慢版本可以枚举路径，或先用普通队列试跑一个小图。它会暴露步数最少和代价最小是否一致；一旦不一致，就必须围绕代价组织状态。二刷时先写慢版本或口述慢版本，用它确认不会漏答案。

\textbf{特例回放。} 拿 [[5,4,5],[1,2,6],[7,4,6]] 手算，路径分数是沿路最小值，要最大化这个最小值。规律是优先扩展当前路径分数最高的格子，或二分阈值判断能否只走不低于阈值的格子。把这个特例继续拆开看：队列或堆弹出的顺序必须和代价规则一致；旧状态被跳过，是因为已经有更小代价到达同一状态。只要边权或代价合成方式改变，就要重新证明扩展顺序。复盘时把这个特例压成状态字段：距离数组、优先队列状态、松弛候选、旧状态判定、不可达哨兵；资源约束题还要把资源维度放进状态；再用边界 单格、起点或终点很小、必须绕路、多个相同得分路径 反查初始化、更新顺序和退出条件。

\textbf{状态回放。} 手算路线：手算时记录每次弹出的代价、当前距离数组和松弛结果。旧状态被弹出时为什么跳过，也要写在表里。状态字段：距离数组、优先队列状态、松弛候选、旧状态判定、不可达哨兵；资源约束题还要把资源维度放进状态。状态升级：把路径枚举压成距离数组和优先队列；每次只扩展当前最有希望的未过期状态。

\textbf{证明和边界。} 证明从扩展顺序开始：当前弹出的状态在代价规则下已经不会被未来路径改得更优。边界：单格、起点或终点很小、必须绕路、多个相同得分路径。变体：也可把阈值二分后检查连通性，或按值从高到低并查集合并。

\noindent\textbf{LeetCode 1334 阈值距离内邻居最少的城市。} 模型：每个城市都需要知道到其他城市的最短距离是否不超过阈值。推导重点：节点数较小时可用 Floyd，边稀疏时也可从每个点跑 Dijkstra。

\textbf{二刷读题。} 读题时先区分可达、最少边数和最小代价。本题可以压成“每个城市都需要知道到其他城市的最短距离是否不超过阈值”，节点数较小时可用 Floyd，边稀疏时也可从每个点跑 Dijkstra决定扩展顺序。先遮住答案，只保留题面，复述候选对象、合法性条件和输出目标。

\textbf{暴力回放。} 慢版本可以枚举路径，或先用普通队列试跑一个小图。它会暴露步数最少和代价最小是否一致；一旦不一致，就必须围绕代价组织状态。二刷时先写慢版本或口述慢版本，用它确认不会漏答案。

\textbf{特例回放。} 拿四个城市和距离阈值 4 手算，每个城市需要统计最短距离不超过 4 的其他城市数；并列时选编号更大的城市。规律是先求全源最短路，再按邻居数量和编号规则选答案。把这个特例继续拆开看：队列或堆弹出的顺序必须和代价规则一致；旧状态被跳过，是因为已经有更小代价到达同一状态。只要边权或代价合成方式改变，就要重新证明扩展顺序。复盘时把这个特例压成状态字段：距离数组、优先队列状态、松弛候选、旧状态判定、不可达哨兵；资源约束题还要把资源维度放进状态；再用边界 距离等于阈值、并列时选编号最大、不可达、无向边 反查初始化、更新顺序和退出条件。

\textbf{状态回放。} 手算路线：手算时记录每次弹出的代价、当前距离数组和松弛结果。旧状态被弹出时为什么跳过，也要写在表里。状态字段：距离数组、优先队列状态、松弛候选、旧状态判定、不可达哨兵；资源约束题还要把资源维度放进状态。状态升级：把路径枚举压成距离数组和优先队列；每次只扩展当前最有希望的未过期状态。

\textbf{证明和边界。} 证明从扩展顺序开始：当前弹出的状态在代价规则下已经不会被未来路径改得更优。边界：距离等于阈值、并列时选编号最大、不可达、无向边。变体：这是多源最短路选择问题，算法取决于 n、边数和查询密度。

\noindent\textbf{LeetCode 1368 使网格图至少有一条有效路径的最小代价。} 模型：沿着格子箭头走代价为零，改方向走代价为一。推导重点：边权只有零和一，用 0-1 BFS：零代价邻居入队首，一代价邻居入队尾。

\textbf{二刷读题。} 读题时先区分可达、最少边数和最小代价。本题可以压成“沿着格子箭头走代价为零，改方向走代价为一”，边权只有零和一，用 0-1 BFS：零代价邻居入队首，一代价邻居入队尾决定扩展顺序。先遮住答案，只保留题面，复述候选对象、合法性条件和输出目标。

\textbf{暴力回放。} 慢版本可以枚举路径，或先用普通队列试跑一个小图。它会暴露步数最少和代价最小是否一致；一旦不一致，就必须围绕代价组织状态。二刷时先写慢版本或口述慢版本，用它确认不会漏答案。

\textbf{特例回放。} 拿网格箭头指向右和下手算，沿箭头走的代价是 0，改方向才付 1。规律是边权只有 0 和 1，可用 0-1 BFS；状态是格子，松弛代价取决于当前箭头是否指向邻居。把这个特例继续拆开看：队列或堆弹出的顺序必须和代价规则一致；旧状态被跳过，是因为已经有更小代价到达同一状态。只要边权或代价合成方式改变，就要重新证明扩展顺序。复盘时把这个特例压成状态字段：距离数组、优先队列状态、松弛候选、旧状态判定、不可达哨兵；资源约束题还要把资源维度放进状态；再用边界 单格、箭头指向越界、多个零代价路径、过期状态 反查初始化、更新顺序和退出条件。

\textbf{状态回放。} 手算路线：手算时记录每次弹出的代价、当前距离数组和松弛结果。旧状态被弹出时为什么跳过，也要写在表里。状态字段：距离数组、优先队列状态、松弛候选、旧状态判定、不可达哨兵；资源约束题还要把资源维度放进状态。状态升级：把路径枚举压成距离数组和优先队列；每次只扩展当前最有希望的未过期状态。

\textbf{证明和边界。} 证明从扩展顺序开始：当前弹出的状态在代价规则下已经不会被未来路径改得更优。边界：单格、箭头指向越界、多个零代价路径、过期状态。变体：普通 BFS 按步数扩展，不等于总代价最小。

\noindent\textbf{LeetCode 1514 概率最大的路径。} 模型：路径概率是边概率乘积，目标是最大乘积路径。推导重点：用大顶堆每次扩展当前概率最高节点，乘以边概率得到候选概率。

\textbf{二刷读题。} 读题时先区分可达、最少边数和最小代价。本题可以压成“路径概率是边概率乘积，目标是最大乘积路径”，用大顶堆每次扩展当前概率最高节点，乘以边概率得到候选概率决定扩展顺序。先遮住答案，只保留题面，复述候选对象、合法性条件和输出目标。

\textbf{暴力回放。} 慢版本可以枚举路径，或先用普通队列试跑一个小图。它会暴露步数最少和代价最小是否一致；一旦不一致，就必须围绕代价组织状态。二刷时先写慢版本或口述慢版本，用它确认不会漏答案。

\textbf{特例回放。} 拿 0->1 概率 0.5、1->2 概率 0.5、0->2 概率 0.2 手算，经 1 到 2 的概率 0.25 更大。规律是把距离改成最大概率，优先队列每次弹出当前概率最高的节点并做乘法松弛。把这个特例继续拆开看：队列或堆弹出的顺序必须和代价规则一致；旧状态被跳过，是因为已经有更小代价到达同一状态。只要边权或代价合成方式改变，就要重新证明扩展顺序。复盘时把这个特例压成状态字段：距离数组、优先队列状态、松弛候选、旧状态判定、不可达哨兵；资源约束题还要把资源维度放进状态；再用边界 不可达、概率为零、起点等于终点、浮点比较 反查初始化、更新顺序和退出条件。

\textbf{状态回放。} 手算路线：手算时记录每次弹出的代价、当前距离数组和松弛结果。旧状态被弹出时为什么跳过，也要写在表里。状态字段：距离数组、优先队列状态、松弛候选、旧状态判定、不可达哨兵；资源约束题还要把资源维度放进状态。状态升级：把路径枚举压成距离数组和优先队列；每次只扩展当前最有希望的未过期状态。

\textbf{证明和边界。} 证明从扩展顺序开始：当前弹出的状态在代价规则下已经不会被未来路径改得更优。边界：不可达、概率为零、起点等于终点、浮点比较。变体：取负对数后可转成最短路，但直接最大堆更直观。

\noindent\textbf{LeetCode 1631 最小体力消耗路径。} 模型：路径代价是沿途最大高度差，不是边权和。推导重点：Dijkstra 的松弛改成取当前代价和新边代价的最大值。

\textbf{二刷读题。} 读题时先区分可达、最少边数和最小代价。本题可以压成“路径代价是沿途最大高度差，不是边权和”，Dijkstra 的松弛改成取当前代价和新边代价的最大值决定扩展顺序。先遮住答案，只保留题面，复述候选对象、合法性条件和输出目标。

\textbf{暴力回放。} 慢版本可以枚举路径，或先用普通队列试跑一个小图。它会暴露步数最少和代价最小是否一致；一旦不一致，就必须围绕代价组织状态。二刷时先写慢版本或口述慢版本，用它确认不会漏答案。

\textbf{特例回放。} 拿 heights=[[1,2,2],[3,8,2],[5,3,5]] 手算，路径代价是相邻高度差的最大值，不是差值之和。规律是 Dijkstra 的距离定义为到达某格的最小最大边权，或二分最大体力做可达性。把这个特例继续拆开看：队列或堆弹出的顺序必须和代价规则一致；旧状态被跳过，是因为已经有更小代价到达同一状态。只要边权或代价合成方式改变，就要重新证明扩展顺序。复盘时把这个特例压成状态字段：距离数组、优先队列状态、松弛候选、旧状态判定、不可达哨兵；资源约束题还要把资源维度放进状态；再用边界 单格、平坦网格、必须绕路、多个相同代价 反查初始化、更新顺序和退出条件。

\textbf{状态回放。} 手算路线：手算时记录每次弹出的代价、当前距离数组和松弛结果。旧状态被弹出时为什么跳过，也要写在表里。状态字段：距离数组、优先队列状态、松弛候选、旧状态判定、不可达哨兵；资源约束题还要把资源维度放进状态。状态升级：把路径枚举压成距离数组和优先队列；每次只扩展当前最有希望的未过期状态。

\textbf{证明和边界。} 证明从扩展顺序开始：当前弹出的状态在代价规则下已经不会被未来路径改得更优。边界：单格、平坦网格、必须绕路、多个相同代价。变体：也可二分体力上限后 BFS 检查可达。

\noindent\textbf{LeetCode 1928 规定时间内到达终点的最小花费。} 模型：路径同时有时间约束和费用目标，单一距离无法描述状态。推导重点：状态包含节点和已用时间，DP 或 Dijkstra 在时间维度内松弛最小费用。

\textbf{二刷读题。} 读题时先区分可达、最少边数和最小代价。本题可以压成“路径同时有时间约束和费用目标，单一距离无法描述状态”，状态包含节点和已用时间，DP 或 Dijkstra 在时间维度内松弛最小费用决定扩展顺序。先遮住答案，只保留题面，复述候选对象、合法性条件和输出目标。

\textbf{暴力回放。} 慢版本可以枚举路径，或先用普通队列试跑一个小图。它会暴露步数最少和代价最小是否一致；一旦不一致，就必须围绕代价组织状态。二刷时先写慢版本或口述慢版本，用它确认不会漏答案。

\textbf{特例回放。} 拿 0->1 用时 10 费用 5，1->2 用时 10 费用 5，0->2 用时 25 费用 100、maxTime=30 手算，费用最小的是经 1 到 2。规律是每个节点不能只保存一个最小费用，状态要包含已用时间。把这个特例继续拆开看：队列或堆弹出的顺序必须和代价规则一致；旧状态被跳过，是因为已经有更小代价到达同一状态。只要边权或代价合成方式改变，就要重新证明扩展顺序。复盘时把这个特例压成状态字段：距离数组、优先队列状态、松弛候选、旧状态判定、不可达哨兵；资源约束题还要把资源维度放进状态；再用边界 时间刚好到达、费用高但时间短、不可达、状态数量由 maxTime 决定 反查初始化、更新顺序和退出条件。

\textbf{状态回放。} 手算路线：手算时记录每次弹出的代价、当前距离数组和松弛结果。旧状态被弹出时为什么跳过，也要写在表里。状态字段：距离数组、优先队列状态、松弛候选、旧状态判定、不可达哨兵；资源约束题还要把资源维度放进状态。状态升级：把路径枚举压成距离数组和优先队列；每次只扩展当前最有希望的未过期状态。

\textbf{证明和边界。} 证明从扩展顺序开始：当前弹出的状态在代价规则下已经不会被未来路径改得更优。边界：时间刚好到达、费用高但时间短、不可达、状态数量由 maxTime 决定。变体：这是资源约束最短路，不能只保存每个节点一个最小费用。

\noindent\textbf{LeetCode 1976 到达目的地的方案数。} 模型：需要在最短距离之外统计有多少条最短路径。推导重点：Dijkstra 松弛时若得到更短距离就重置方案数，若等于最短距离就累加方案数。

\textbf{二刷读题。} 读题时先区分可达、最少边数和最小代价。本题可以压成“需要在最短距离之外统计有多少条最短路径”，Dijkstra 松弛时若得到更短距离就重置方案数，若等于最短距离就累加方案数决定扩展顺序。先遮住答案，只保留题面，复述候选对象、合法性条件和输出目标。

\textbf{暴力回放。} 慢版本可以枚举路径，或先用普通队列试跑一个小图。它会暴露步数最少和代价最小是否一致；一旦不一致，就必须围绕代价组织状态。二刷时先写慢版本或口述慢版本，用它确认不会漏答案。

\textbf{特例回放。} 拿 0->1->3 和 0->2->3 两条长度都为 2 的路手算，最短距离相同，所以到 3 的方案数应累加为 2。规律是 Dijkstra 松弛到更短距离时重置方案数，等长时累加方案数。把这个特例继续拆开看：队列或堆弹出的顺序必须和代价规则一致；旧状态被跳过，是因为已经有更小代价到达同一状态。只要边权或代价合成方式改变，就要重新证明扩展顺序。复盘时把这个特例压成状态字段：距离数组、优先队列状态、松弛候选、旧状态判定、不可达哨兵；资源约束题还要把资源维度放进状态；再用边界 多条等长路径、取模、旧堆状态、不可达 反查初始化、更新顺序和退出条件。

\textbf{状态回放。} 手算路线：手算时记录每次弹出的代价、当前距离数组和松弛结果。旧状态被弹出时为什么跳过，也要写在表里。状态字段：距离数组、优先队列状态、松弛候选、旧状态判定、不可达哨兵；资源约束题还要把资源维度放进状态。状态升级：把路径枚举压成距离数组和优先队列；每次只扩展当前最有希望的未过期状态。

\textbf{证明和边界。} 证明从扩展顺序开始：当前弹出的状态在代价规则下已经不会被未来路径改得更优。边界：多条等长路径、取模、旧堆状态、不可达。变体：这是最短路和计数 DP 的结合，更新顺序必须和距离松弛一致。

\noindent\textbf{LeetCode 2045 到达目的地的第二短时间。} 模型：每个节点需要保留最短和次短到达步数，再考虑红绿灯等待时间。推导重点：BFS 式扩展边数或时间，只有候选时间不同且能进入前两名时才入队。

\textbf{二刷读题。} 读题时先区分可达、最少边数和最小代价。本题可以压成“每个节点需要保留最短和次短到达步数，再考虑红绿灯等待时间”，BFS 式扩展边数或时间，只有候选时间不同且能进入前两名时才入队决定扩展顺序。先遮住答案，只保留题面，复述候选对象、合法性条件和输出目标。

\textbf{暴力回放。} 慢版本可以枚举路径，或先用普通队列试跑一个小图。它会暴露步数最少和代价最小是否一致；一旦不一致，就必须围绕代价组织状态。二刷时先写慢版本或口述慢版本，用它确认不会漏答案。

\textbf{特例回放。} 拿 1--2、time=3、change=2 手算，第一次到 2 是最短；第二短必须从 2 回 1 再到 2，且出发时可能遇到红灯等待。规律是每个节点保存两个不同到达时间，不能把次短等同于最短。把这个特例继续拆开看：队列或堆弹出的顺序必须和代价规则一致；旧状态被跳过，是因为已经有更小代价到达同一状态。只要边权或代价合成方式改变，就要重新证明扩展顺序。复盘时把这个特例压成状态字段：距离数组、优先队列状态、松弛候选、旧状态判定、不可达哨兵；资源约束题还要把资源维度放进状态；再用边界 第二短不能等于最短、红灯等待、无向图往返、到达终点后继续找次短 反查初始化、更新顺序和退出条件。

\textbf{状态回放。} 手算路线：手算时记录每次弹出的代价、当前距离数组和松弛结果。旧状态被弹出时为什么跳过，也要写在表里。状态字段：距离数组、优先队列状态、松弛候选、旧状态判定、不可达哨兵；资源约束题还要把资源维度放进状态。状态升级：把路径枚举压成距离数组和优先队列；每次只扩展当前最有希望的未过期状态。

\textbf{证明和边界。} 证明从扩展顺序开始：当前弹出的状态在代价规则下已经不会被未来路径改得更优。边界：第二短不能等于最短、红灯等待、无向图往返、到达终点后继续找次短。变体：它说明最短路状态有时不是一个最优值，而是前两个不同值。

\subsection{自测标准}

二刷时不要只确认自己见过这道题。请遮住答案，先讲题面模型，再讲暴力基线和瓶颈，然后说出优化结构保存了什么、不变量如何排除错误候选、边界实验如何打破错误实现。

\section{并查集与动态连通性：二刷复盘卡}
这一大节只围绕一个题型复盘，不把题号直接铺成目录。阅读时先看复盘目标，再读核心题卡，最后用自测标准检查自己是否能独立重讲。

\subsection{复盘目标}

追问通常会问并查集能否回答路径长度、路径压缩是否改变集合语义、字符串如何编号。请准备说明合并失败为什么意味着环。

\subsection{核心题卡}

\noindent\textbf{LeetCode 305 岛屿数量 II。} 模型：陆地逐步加入，每次加入只可能影响相邻陆地所在集合。推导重点：新增陆地先让岛屿数加一，再和四邻已存在陆地合并，合并成功则岛屿数减一。

\textbf{二刷读题。} 读题时先确认关系是否只增加不删除，且问题只关心同组关系。本题可以压成“陆地逐步加入，每次加入只可能影响相邻陆地所在集合”，并查集的代表元就是集合身份。先遮住答案，只保留题面，复述候选对象、合法性条件和输出目标。

\textbf{暴力回放。} 慢版本每次合并后用 DFS 或 BFS 重新判断连通性。它正确但重复遍历已有连接；当关系只会合并不会拆开时，可以压成集合代表。二刷时先写慢版本或口述慢版本，用它确认不会漏答案。

\textbf{特例回放。} 拿 3x3 网格依次加 (0,0)、(0,1)、(1,2) 手算，前两个陆地相邻要合并，岛屿数从 2 降回 1；第三个不相邻，岛屿数加 1。规律是每次加陆地先新增一个集合，再和四邻陆地 union，重复加点不能重复计数。把这个特例继续拆开看：union 只是在维护等价类，不是在保存具体路径。两个节点代表元相同只能说明连通或关系一致；如果题目还要距离、比例或奇偶性，必须把附加信息挂到根关系上。复盘时把这个特例压成状态字段：parent、size 或 rank、find 返回的代表元、合并时的附加信息；带权或奇偶并查集还要维护到根的关系；再用边界 重复加入同一格、边界位置、四邻属于同一岛、空操作列表 反查初始化、更新顺序和退出条件。

\textbf{状态回放。} 手算路线：手算时写 parent 表和集合代表。每次 union 只改变根之间的关系，find 后代表相同才说明连通。状态字段：parent、size 或 rank、find 返回的代表元、合并时的附加信息；带权或奇偶并查集还要维护到根的关系。状态升级：把连通分量压成代表元；查询只比较代表，合并只发生在两个根之间。

\textbf{证明和边界。} 证明从等价关系开始：合并维护连通性的传递性，find 返回的代表元就是当前集合身份。边界：重复加入同一格、边界位置、四邻属于同一岛、空操作列表。变体：这是动态连通性，比每次重新 BFS 更适合并查集。

\noindent\textbf{LeetCode 399 除法求值。} 模型：变量之间的除法关系可以看成带权连通分量。推导重点：并查集维护节点到根的乘积权值，合并时根据等式调整根之间权重。

\textbf{二刷读题。} 读题时先确认关系是否只增加不删除，且问题只关心同组关系。本题可以压成“变量之间的除法关系可以看成带权连通分量”，并查集的代表元就是集合身份。先遮住答案，只保留题面，复述候选对象、合法性条件和输出目标。

\textbf{暴力回放。} 慢版本每次合并后用 DFS 或 BFS 重新判断连通性。它正确但重复遍历已有连接；当关系只会合并不会拆开时，可以压成集合代表。二刷时先写慢版本或口述慢版本，用它确认不会漏答案。

\textbf{特例回放。} 拿 a/b=2、b/c=3 手算，a 到 c 的比值是沿边相乘 6；若查询 c/a，则沿反向边乘 1/3 和 1/2。规律是并查集或图搜索都要维护到代表元的倍率关系，合并时同步更新权值。把这个特例继续拆开看：union 只是在维护等价类，不是在保存具体路径。两个节点代表元相同只能说明连通或关系一致；如果题目还要距离、比例或奇偶性，必须把附加信息挂到根关系上。复盘时把这个特例压成状态字段：parent、size 或 rank、find 返回的代表元、合并时的附加信息；带权或奇偶并查集还要维护到根的关系；再用边界 查询未知变量、自除、矛盾数据、路径压缩时权值同步更新 反查初始化、更新顺序和退出条件。

\textbf{状态回放。} 手算路线：手算时写 parent 表和集合代表。每次 union 只改变根之间的关系，find 后代表相同才说明连通。状态字段：parent、size 或 rank、find 返回的代表元、合并时的附加信息；带权或奇偶并查集还要维护到根的关系。状态升级：把连通分量压成代表元；查询只比较代表，合并只发生在两个根之间。

\textbf{证明和边界。} 证明从等价关系开始：合并维护连通性的传递性，find 返回的代表元就是当前集合身份。边界：查询未知变量、自除、矛盾数据、路径压缩时权值同步更新。变体：也可建图 DFS，但多次查询时带权并查集更适合。

\noindent\textbf{LeetCode 547 省份数量。} 模型：城市连接关系形成无向图，省份就是连通分量。推导重点：并查集合并所有直接连接的城市，最后统计根数量。

\textbf{二刷读题。} 读题时先确认关系是否只增加不删除，且问题只关心同组关系。本题可以压成“城市连接关系形成无向图，省份就是连通分量”，并查集的代表元就是集合身份。先遮住答案，只保留题面，复述候选对象、合法性条件和输出目标。

\textbf{暴力回放。} 慢版本每次合并后用 DFS 或 BFS 重新判断连通性。它正确但重复遍历已有连接；当关系只会合并不会拆开时，可以压成集合代表。二刷时先写慢版本或口述慢版本，用它确认不会漏答案。

\textbf{特例回放。} 拿 isConnected 中 0 和 1 相连、2 独立手算，0 和 1 union 后属于同一代表，2 保持自己，省份数为 2。规律是矩阵中的一条连接边对应集合合并，最后代表元数量就是连通分量数。把这个特例继续拆开看：union 只是在维护等价类，不是在保存具体路径。两个节点代表元相同只能说明连通或关系一致；如果题目还要距离、比例或奇偶性，必须把附加信息挂到根关系上。复盘时把这个特例压成状态字段：parent、size 或 rank、find 返回的代表元、合并时的附加信息；带权或奇偶并查集还要维护到根的关系；再用边界 邻接矩阵对称、自连接、孤立城市、只扫描上三角 反查初始化、更新顺序和退出条件。

\textbf{状态回放。} 手算路线：手算时写 parent 表和集合代表。每次 union 只改变根之间的关系，find 后代表相同才说明连通。状态字段：parent、size 或 rank、find 返回的代表元、合并时的附加信息；带权或奇偶并查集还要维护到根的关系。状态升级：把连通分量压成代表元；查询只比较代表，合并只发生在两个根之间。

\textbf{证明和边界。} 证明从等价关系开始：合并维护连通性的传递性，find 返回的代表元就是当前集合身份。边界：邻接矩阵对称、自连接、孤立城市、只扫描上三角。变体：BFS 也可做；并查集更突出动态合并。

\noindent\textbf{LeetCode 684 冗余连接。} 模型：树加一条边会形成唯一环。推导重点：按顺序加边，若一条边两端已连通，则它就是冗余边。

\textbf{二刷读题。} 读题时先确认关系是否只增加不删除，且问题只关心同组关系。本题可以压成“树加一条边会形成唯一环”，并查集的代表元就是集合身份。先遮住答案，只保留题面，复述候选对象、合法性条件和输出目标。

\textbf{暴力回放。} 慢版本每次合并后用 DFS 或 BFS 重新判断连通性。它正确但重复遍历已有连接；当关系只会合并不会拆开时，可以压成集合代表。二刷时先写慢版本或口述慢版本，用它确认不会漏答案。

\textbf{特例回放。} 拿 edges=[[1,2],[1,3],[2,3]] 手算，前两条边把 1、2、3 连成一棵树，第三条边连接的两个点已经同组，所以它就是冗余边。规律是无向图加边时，第一次 union 失败的边形成环。把这个特例继续拆开看：union 只是在维护等价类，不是在保存具体路径。两个节点代表元相同只能说明连通或关系一致；如果题目还要距离、比例或奇偶性，必须把附加信息挂到根关系上。复盘时把这个特例压成状态字段：parent、size 或 rank、find 返回的代表元、合并时的附加信息；带权或奇偶并查集还要维护到根的关系；再用边界 节点编号、最后一条成环边、重复边、并查集初始化大小 反查初始化、更新顺序和退出条件。

\textbf{状态回放。} 手算路线：手算时写 parent 表和集合代表。每次 union 只改变根之间的关系，find 后代表相同才说明连通。状态字段：parent、size 或 rank、find 返回的代表元、合并时的附加信息；带权或奇偶并查集还要维护到根的关系。状态升级：把连通分量压成代表元；查询只比较代表，合并只发生在两个根之间。

\textbf{证明和边界。} 证明从等价关系开始：合并维护连通性的传递性，find 返回的代表元就是当前集合身份。边界：节点编号、最后一条成环边、重复边、并查集初始化大小。变体：有向冗余连接需要同时处理入度为二和环，复杂度更高。

\noindent\textbf{LeetCode 685 冗余连接 II。} 模型：有向图中额外边可能导致某节点入度为二，也可能导致环。推导重点：先记录入度为二的两条候选边，再用并查集排除其中一条测试是否成树。

\textbf{二刷读题。} 读题时先确认关系是否只增加不删除，且问题只关心同组关系。本题可以压成“有向图中额外边可能导致某节点入度为二，也可能导致环”，并查集的代表元就是集合身份。先遮住答案，只保留题面，复述候选对象、合法性条件和输出目标。

\textbf{暴力回放。} 慢版本每次合并后用 DFS 或 BFS 重新判断连通性。它正确但重复遍历已有连接；当关系只会合并不会拆开时，可以压成集合代表。二刷时先写慢版本或口述慢版本，用它确认不会漏答案。

\textbf{特例回放。} 拿有向边 [[1,2],[1,3],[2,3]] 手算，节点 3 有两个父亲 1 和 2；需要判断去掉哪条边后剩余图是有根树。规律是先处理双父节点候选，再用并查集检测环，双父和成环组合时要按题意返回正确候选边。把这个特例继续拆开看：union 只是在维护等价类，不是在保存具体路径。两个节点代表元相同只能说明连通或关系一致；如果题目还要距离、比例或奇偶性，必须把附加信息挂到根关系上。复盘时把这个特例压成状态字段：parent、size 或 rank、find 返回的代表元、合并时的附加信息；带权或奇偶并查集还要维护到根的关系；再用边界 只有环、只有入度二、两者同时存在、返回最后出现的边 反查初始化、更新顺序和退出条件。

\textbf{状态回放。} 手算路线：手算时写 parent 表和集合代表。每次 union 只改变根之间的关系，find 后代表相同才说明连通。状态字段：parent、size 或 rank、find 返回的代表元、合并时的附加信息；带权或奇偶并查集还要维护到根的关系。状态升级：把连通分量压成代表元；查询只比较代表，合并只发生在两个根之间。

\textbf{证明和边界。} 证明从等价关系开始：合并维护连通性的传递性，find 返回的代表元就是当前集合身份。边界：只有环、只有入度二、两者同时存在、返回最后出现的边。变体：它比无向冗余连接多了有向树入度约束。

\noindent\textbf{LeetCode 721 账户合并。} 模型：共享邮箱说明账户属于同一人，邮箱之间形成连通分量。推导重点：给邮箱编号并查集合并，同根邮箱收集后排序输出。

\textbf{二刷读题。} 读题时先确认关系是否只增加不删除，且问题只关心同组关系。本题可以压成“共享邮箱说明账户属于同一人，邮箱之间形成连通分量”，并查集的代表元就是集合身份。先遮住答案，只保留题面，复述候选对象、合法性条件和输出目标。

\textbf{暴力回放。} 慢版本每次合并后用 DFS 或 BFS 重新判断连通性。它正确但重复遍历已有连接；当关系只会合并不会拆开时，可以压成集合代表。二刷时先写慢版本或口述慢版本，用它确认不会漏答案。

\textbf{特例回放。} 拿 John 的账号 [a,b] 和另一个 [b,c] 手算，邮箱 b 重叠，所以 a、b、c 属于同一人；名字只是输出标签，合并依据是邮箱。规律是邮箱映射到编号后并查集合并，最后按代表收集并排序邮箱。把这个特例继续拆开看：union 只是在维护等价类，不是在保存具体路径。两个节点代表元相同只能说明连通或关系一致；如果题目还要距离、比例或奇偶性，必须把附加信息挂到根关系上。复盘时把这个特例压成状态字段：parent、size 或 rank、find 返回的代表元、合并时的附加信息；带权或奇偶并查集还要维护到根的关系；再用边界 同名不同人、一个账户多个邮箱、重复邮箱、输出排序 反查初始化、更新顺序和退出条件。

\textbf{状态回放。} 手算路线：手算时写 parent 表和集合代表。每次 union 只改变根之间的关系，find 后代表相同才说明连通。状态字段：parent、size 或 rank、find 返回的代表元、合并时的附加信息；带权或奇偶并查集还要维护到根的关系。状态升级：把连通分量压成代表元；查询只比较代表，合并只发生在两个根之间。

\textbf{证明和边界。} 证明从等价关系开始：合并维护连通性的传递性，find 返回的代表元就是当前集合身份。边界：同名不同人、一个账户多个邮箱、重复邮箱、输出排序。变体：姓名不能作为合并依据，只能作为输出字段。

\noindent\textbf{LeetCode 947 移除最多的同行或同列石头。} 模型：同一行或同一列的石头属于同一连通分量，每个分量至少留一块。推导重点：把行和列编号后合并，答案是石头数减去连通分量数。

\textbf{二刷读题。} 读题时先确认关系是否只增加不删除，且问题只关心同组关系。本题可以压成“同一行或同一列的石头属于同一连通分量，每个分量至少留一块”，并查集的代表元就是集合身份。先遮住答案，只保留题面，复述候选对象、合法性条件和输出目标。

\textbf{暴力回放。} 慢版本每次合并后用 DFS 或 BFS 重新判断连通性。它正确但重复遍历已有连接；当关系只会合并不会拆开时，可以压成集合代表。二刷时先写慢版本或口述慢版本，用它确认不会漏答案。

\textbf{特例回放。} 拿石头 (0,0),(0,1),(1,0) 手算，它们通过同一行或同一列连成一个分量，三个石头最多移除 2 个。规律是每个连通分量至少留一个，答案是石头数减分量数。把这个特例继续拆开看：union 只是在维护等价类，不是在保存具体路径。两个节点代表元相同只能说明连通或关系一致；如果题目还要距离、比例或奇偶性，必须把附加信息挂到根关系上。复盘时把这个特例压成状态字段：parent、size 或 rank、find 返回的代表元、合并时的附加信息；带权或奇偶并查集还要维护到根的关系；再用边界 单块石头、行列编号冲突、多个分量、重复坐标 反查初始化、更新顺序和退出条件。

\textbf{状态回放。} 手算路线：手算时写 parent 表和集合代表。每次 union 只改变根之间的关系，find 后代表相同才说明连通。状态字段：parent、size 或 rank、find 返回的代表元、合并时的附加信息；带权或奇偶并查集还要维护到根的关系。状态升级：把连通分量压成代表元；查询只比较代表，合并只发生在两个根之间。

\textbf{证明和边界。} 证明从等价关系开始：合并维护连通性的传递性，find 返回的代表元就是当前集合身份。边界：单块石头、行列编号冲突、多个分量、重复坐标。变体：也可把石头当节点建图，但行列并查集更省边。

\noindent\textbf{LeetCode 959 由斜杠划分区域。} 模型：每个网格可拆成四个小三角，斜杠决定内部哪些三角相连。推导重点：给每格四个区域编号，先按字符合并内部，再合并相邻格边界区域。

\textbf{二刷读题。} 读题时先确认关系是否只增加不删除，且问题只关心同组关系。本题可以压成“每个网格可拆成四个小三角，斜杠决定内部哪些三角相连”，并查集的代表元就是集合身份。先遮住答案，只保留题面，复述候选对象、合法性条件和输出目标。

\textbf{暴力回放。} 慢版本每次合并后用 DFS 或 BFS 重新判断连通性。它正确但重复遍历已有连接；当关系只会合并不会拆开时，可以压成集合代表。二刷时先写慢版本或口述慢版本，用它确认不会漏答案。

\textbf{特例回放。} 拿 1x1 格子里一条斜杠手算，它把小方格分成两个区域；把每格拆成 4 个三角形后，斜杠决定哪些三角形不能合并。规律是格内按字符合并，格间相邻边合并，最终代表元数量就是区域数。把这个特例继续拆开看：union 只是在维护等价类，不是在保存具体路径。两个节点代表元相同只能说明连通或关系一致；如果题目还要距离、比例或奇偶性，必须把附加信息挂到根关系上。复盘时把这个特例压成状态字段：parent、size 或 rank、find 返回的代表元、合并时的附加信息；带权或奇偶并查集还要维护到根的关系；再用边界 空格、正斜杠、反斜杠、边界合并、区域计数 反查初始化、更新顺序和退出条件。

\textbf{状态回放。} 手算路线：手算时写 parent 表和集合代表。每次 union 只改变根之间的关系，find 后代表相同才说明连通。状态字段：parent、size 或 rank、find 返回的代表元、合并时的附加信息；带权或奇偶并查集还要维护到根的关系。状态升级：把连通分量压成代表元；查询只比较代表，合并只发生在两个根之间。

\textbf{证明和边界。} 证明从等价关系开始：合并维护连通性的传递性，find 返回的代表元就是当前集合身份。边界：空格、正斜杠、反斜杠、边界合并、区域计数。变体：这是把几何连通性离散成并查集元素的典型题。

\noindent\textbf{LeetCode 990 等式方程的可满足性。} 模型：相等关系形成集合，不等关系要求两端不在同一集合。推导重点：先合并所有等式，再检查每个不等式是否冲突。

\textbf{二刷读题。} 读题时先确认关系是否只增加不删除，且问题只关心同组关系。本题可以压成“相等关系形成集合，不等关系要求两端不在同一集合”，并查集的代表元就是集合身份。先遮住答案，只保留题面，复述候选对象、合法性条件和输出目标。

\textbf{暴力回放。} 慢版本每次合并后用 DFS 或 BFS 重新判断连通性。它正确但重复遍历已有连接；当关系只会合并不会拆开时，可以压成集合代表。二刷时先写慢版本或口述慢版本，用它确认不会漏答案。

\textbf{特例回放。} 拿 a==b、b==c、a!=c 手算，前两条等式把 a,b,c 合并，最后不等式发现二者同组，矛盾。规律是先处理所有等式建立集合，再检查不等式；顺序反过来会漏掉后续合并造成的冲突。把这个特例继续拆开看：union 只是在维护等价类，不是在保存具体路径。两个节点代表元相同只能说明连通或关系一致；如果题目还要距离、比例或奇偶性，必须把附加信息挂到根关系上。复盘时把这个特例压成状态字段：parent、size 或 rank、find 返回的代表元、合并时的附加信息；带权或奇偶并查集还要维护到根的关系；再用边界 自等式、自不等式、传递相等、多组变量 反查初始化、更新顺序和退出条件。

\textbf{状态回放。} 手算路线：手算时写 parent 表和集合代表。每次 union 只改变根之间的关系，find 后代表相同才说明连通。状态字段：parent、size 或 rank、find 返回的代表元、合并时的附加信息；带权或奇偶并查集还要维护到根的关系。状态升级：把连通分量压成代表元；查询只比较代表，合并只发生在两个根之间。

\textbf{证明和边界。} 证明从等价关系开始：合并维护连通性的传递性，find 返回的代表元就是当前集合身份。边界：自等式、自不等式、传递相等、多组变量。变体：先处理不等式会漏掉后续等式带来的传递冲突。

\noindent\textbf{LeetCode 1061 按字典序排列最小的等效字符串。} 模型：等价字符形成集合，每个集合应选择字典序最小代表。推导重点：合并两个字符集合时让较小根成为代表，扫描目标串时替换为根。

\textbf{二刷读题。} 读题时先确认关系是否只增加不删除，且问题只关心同组关系。本题可以压成“等价字符形成集合，每个集合应选择字典序最小代表”，并查集的代表元就是集合身份。先遮住答案，只保留题面，复述候选对象、合法性条件和输出目标。

\textbf{暴力回放。} 慢版本每次合并后用 DFS 或 BFS 重新判断连通性。它正确但重复遍历已有连接；当关系只会合并不会拆开时，可以压成集合代表。二刷时先写慢版本或口述慢版本，用它确认不会漏答案。

\textbf{特例回放。} 拿 s1=parker、s2=morris、base=parser 手算，p 和 m 等价，a 和 o 等价；每个等价类输出字典序最小字符。规律是并查集合并字符时让较小根做代表，转换 base 时查代表。把这个特例继续拆开看：union 只是在维护等价类，不是在保存具体路径。两个节点代表元相同只能说明连通或关系一致；如果题目还要距离、比例或奇偶性，必须把附加信息挂到根关系上。复盘时把这个特例压成状态字段：parent、size 或 rank、find 返回的代表元、合并时的附加信息；带权或奇偶并查集还要维护到根的关系；再用边界 传递等价、重复等价关系、自等价、未出现字符 反查初始化、更新顺序和退出条件。

\textbf{状态回放。} 手算路线：手算时写 parent 表和集合代表。每次 union 只改变根之间的关系，find 后代表相同才说明连通。状态字段：parent、size 或 rank、find 返回的代表元、合并时的附加信息；带权或奇偶并查集还要维护到根的关系。状态升级：把连通分量压成代表元；查询只比较代表，合并只发生在两个根之间。

\textbf{证明和边界。} 证明从等价关系开始：合并维护连通性的传递性，find 返回的代表元就是当前集合身份。边界：传递等价、重复等价关系、自等价、未出现字符。变体：如果集合代表还要保存其他权值，合并规则需要同步维护。

\noindent\textbf{LeetCode 1202 交换字符串中的元素。} 模型：可交换下标形成连通分量，同一分量内字符可以任意重排。推导重点：并查集合并所有可交换下标，对每个分量收集字符并按字典序分配回最小下标。

\textbf{二刷读题。} 读题时先确认关系是否只增加不删除，且问题只关心同组关系。本题可以压成“可交换下标形成连通分量，同一分量内字符可以任意重排”，并查集的代表元就是集合身份。先遮住答案，只保留题面，复述候选对象、合法性条件和输出目标。

\textbf{暴力回放。} 慢版本每次合并后用 DFS 或 BFS 重新判断连通性。它正确但重复遍历已有连接；当关系只会合并不会拆开时，可以压成集合代表。二刷时先写慢版本或口述慢版本，用它确认不会漏答案。

\textbf{特例回放。} 拿 s=dcab、pairs=[[0,3],[1,2]] 手算，0 和 3 可交换成一组，1 和 2 可交换成一组；每组内部字符排序后放回最小位置。规律是并查集找可交换连通分量，每个分量的字符可任意重排。把这个特例继续拆开看：union 只是在维护等价类，不是在保存具体路径。两个节点代表元相同只能说明连通或关系一致；如果题目还要距离、比例或奇偶性，必须把附加信息挂到根关系上。复盘时把这个特例压成状态字段：parent、size 或 rank、find 返回的代表元、合并时的附加信息；带权或奇偶并查集还要维护到根的关系；再用边界 没有交换、多个分量、重复字符、分量下标排序 反查初始化、更新顺序和退出条件。

\textbf{状态回放。} 手算路线：手算时写 parent 表和集合代表。每次 union 只改变根之间的关系，find 后代表相同才说明连通。状态字段：parent、size 或 rank、find 返回的代表元、合并时的附加信息；带权或奇偶并查集还要维护到根的关系。状态升级：把连通分量压成代表元；查询只比较代表，合并只发生在两个根之间。

\textbf{证明和边界。} 证明从等价关系开始：合并维护连通性的传递性，find 返回的代表元就是当前集合身份。边界：没有交换、多个分量、重复字符、分量下标排序。变体：它展示并查集不仅能回答连通性，还能把分量内资源重新组织。

\noindent\textbf{LeetCode 1319 连通网络的操作次数。} 模型：多余网线可以挪用来连接不同连通分量。推导重点：先判断边数至少为 n 减一，再用并查集统计连通分量，答案是分量数减一。

\textbf{二刷读题。} 读题时先确认关系是否只增加不删除，且问题只关心同组关系。本题可以压成“多余网线可以挪用来连接不同连通分量”，并查集的代表元就是集合身份。先遮住答案，只保留题面，复述候选对象、合法性条件和输出目标。

\textbf{暴力回放。} 慢版本每次合并后用 DFS 或 BFS 重新判断连通性。它正确但重复遍历已有连接；当关系只会合并不会拆开时，可以压成集合代表。二刷时先写慢版本或口述慢版本，用它确认不会漏答案。

\textbf{特例回放。} 拿 n=4、连接 [[0,1],[0,2],[1,2]] 手算，前三个点已有一条冗余边，但节点 3 孤立；总边数只有 3，刚好可以拿冗余边连上 3。规律是边数少于 n-1 直接失败，否则答案是连通分量数减一。把这个特例继续拆开看：union 只是在维护等价类，不是在保存具体路径。两个节点代表元相同只能说明连通或关系一致；如果题目还要距离、比例或奇偶性，必须把附加信息挂到根关系上。复盘时把这个特例压成状态字段：parent、size 或 rank、find 返回的代表元、合并时的附加信息；带权或奇偶并查集还要维护到根的关系；再用边界 边数不足、已经连通、重复边、多分量 反查初始化、更新顺序和退出条件。

\textbf{状态回放。} 手算路线：手算时写 parent 表和集合代表。每次 union 只改变根之间的关系，find 后代表相同才说明连通。状态字段：parent、size 或 rank、find 返回的代表元、合并时的附加信息；带权或奇偶并查集还要维护到根的关系。状态升级：把连通分量压成代表元；查询只比较代表，合并只发生在两个根之间。

\textbf{证明和边界。} 证明从等价关系开始：合并维护连通性的传递性，find 返回的代表元就是当前集合身份。边界：边数不足、已经连通、重复边、多分量。变体：合并失败的边代表冗余资源，但只要总边数足够即可。

\noindent\textbf{LeetCode 1579 保证图可完全遍历。} 模型：Alice 和 Bob 各自需要图连通，公共边应优先服务两人。推导重点：先处理类型三公共边同步合并两个并查集，再分别处理各自边，最后检查两图连通。

\textbf{二刷读题。} 读题时先确认关系是否只增加不删除，且问题只关心同组关系。本题可以压成“Alice 和 Bob 各自需要图连通，公共边应优先服务两人”，并查集的代表元就是集合身份。先遮住答案，只保留题面，复述候选对象、合法性条件和输出目标。

\textbf{暴力回放。} 慢版本每次合并后用 DFS 或 BFS 重新判断连通性。它正确但重复遍历已有连接；当关系只会合并不会拆开时，可以压成集合代表。二刷时先写慢版本或口述慢版本，用它确认不会漏答案。

\textbf{特例回放。} 拿公共边能同时连接 Alice 和 Bob 的两个分量手算，先用类型 3 边可以让两个人都受益；若先处理私有边，公共边可能变成冗余。规律是先处理公共边，再分别处理 Alice 和 Bob 的边，最后两套并查集都必须连通。把这个特例继续拆开看：union 只是在维护等价类，不是在保存具体路径。两个节点代表元相同只能说明连通或关系一致；如果题目还要距离、比例或奇偶性，必须把附加信息挂到根关系上。复盘时把这个特例压成状态字段：parent、size 或 rank、find 返回的代表元、合并时的附加信息；带权或奇偶并查集还要维护到根的关系；再用边界 公共边冗余、某一方不连通、边数很多、节点从一开始编号 反查初始化、更新顺序和退出条件。

\textbf{状态回放。} 手算路线：手算时写 parent 表和集合代表。每次 union 只改变根之间的关系，find 后代表相同才说明连通。状态字段：parent、size 或 rank、find 返回的代表元、合并时的附加信息；带权或奇偶并查集还要维护到根的关系。状态升级：把连通分量压成代表元；查询只比较代表，合并只发生在两个根之间。

\textbf{证明和边界。} 证明从等价关系开始：合并维护连通性的传递性，find 返回的代表元就是当前集合身份。边界：公共边冗余、某一方不连通、边数很多、节点从一开始编号。变体：这是双并查集和资源优先级的组合。

\subsection{自测标准}

二刷时不要只确认自己见过这道题。请遮住答案，先讲题面模型，再讲暴力基线和瓶颈，然后说出优化结构保存了什么、不变量如何排除错误候选、边界实验如何打破错误实现。

\section{回溯、枚举与剪枝：二刷复盘卡}
这一大节只围绕一个题型复盘，不把题号直接铺成目录。阅读时先看复盘目标，再读核心题卡，最后用自测标准检查自己是否能独立重讲。

\subsection{复盘目标}

追问通常会问去重发生在同层还是路径、剪枝是否会删掉合法解、输出规模是多少。请准备画出前三层搜索树。

\subsection{核心题卡}

\noindent\textbf{LeetCode 17 电话号码的字母组合。} 模型：每个数字对应若干字母，答案是所有按位选择形成的字符串。推导重点：暴力就是完整笛卡尔积，优化重点是路径长度和下标同步推进，避免在每层重新拼接大量中间字符串。

\textbf{二刷读题。} 读题时先画搜索树：层数代表什么，路径保存什么，终止条件是什么。本题可以压成“每个数字对应若干字母，答案是所有按位选择形成的字符串”，剪枝只能在这个搜索树上说明。先遮住答案，只保留题面，复述候选对象、合法性条件和输出目标。

\textbf{暴力回放。} 慢版本就是完整搜索树：每层列出选择列表，路径到达终止条件时检查答案。它先保证覆盖所有可能，再把明显无效或重复的分支剪掉。二刷时先写慢版本或口述慢版本，用它确认不会漏答案。

\textbf{特例回放。} 拿 digits=23 手算，2 对应 a/b/c，3 对应 d/e/f，第一层选 a 后第二层依次选 d/e/f，得到 ad、ae、af。规律是层数对应输入数字位置，路径保存已经选择的字母；空输入应返回空结果而不是包含空串。把这个特例继续拆开看：一层递归对应一个明确决策，路径保存已经做出的选择，选择列表保存仍可能扩展的分支。剪枝前必须能说明被剪掉的整棵子树没有合法答案，而不是只因为它看起来不优。复盘时把这个特例压成状态字段：路径、起点或使用标记、剩余目标、答案集合、剪枝条件；进入递归和退出递归必须成对恢复；再用边界 空输入、包含 7 或 9、输入中不应出现 0 和 1、输出规模指数增长 反查初始化、更新顺序和退出条件。

\textbf{状态回放。} 手算路线：手算时给每层标出选择列表和路径。剪枝发生前要指出被剪分支为什么已经不可能得到合法答案。状态字段：路径、起点或使用标记、剩余目标、答案集合、剪枝条件；进入递归和退出递归必须成对恢复。状态升级：把完整搜索树加上合法性检查、剩余资源剪枝和同层去重；路径状态进入和退出要成对恢复。

\textbf{证明和边界。} 证明从搜索树覆盖开始：每个合法答案有且只有一条路径，剪枝只删掉不可能成功的分支。边界：空输入、包含 7 或 9、输入中不应出现 0 和 1、输出规模指数增长。变体：若字符映射来自键盘以外的字典，只需要替换候选表，搜索骨架不变。

\noindent\textbf{LeetCode 22 括号生成。} 模型：合法括号串的前缀中右括号数量不能超过左括号数量。推导重点：状态保存已用左括号和右括号数量，只有满足约束的选择才继续递归。

\textbf{二刷读题。} 读题时先画搜索树：层数代表什么，路径保存什么，终止条件是什么。本题可以压成“合法括号串的前缀中右括号数量不能超过左括号数量”，剪枝只能在这个搜索树上说明。先遮住答案，只保留题面，复述候选对象、合法性条件和输出目标。

\textbf{暴力回放。} 慢版本就是完整搜索树：每层列出选择列表，路径到达终止条件时检查答案。它先保证覆盖所有可能，再把明显无效或重复的分支剪掉。二刷时先写慢版本或口述慢版本，用它确认不会漏答案。

\textbf{特例回放。} 拿 n=2 手算，第一步只能放左括号，之后可以放左括号得到 (())，也可以在已有左括号数量大于右括号时放右括号得到 ()()。规律是状态包含已放左括号数和右括号数，右括号不能超过左括号，左括号不能超过 n。把这个特例继续拆开看：一层递归对应一个明确决策，路径保存已经做出的选择，选择列表保存仍可能扩展的分支。剪枝前必须能说明被剪掉的整棵子树没有合法答案，而不是只因为它看起来不优。复盘时把这个特例压成状态字段：路径、起点或使用标记、剩余目标、答案集合、剪枝条件；进入递归和退出递归必须成对恢复；再用边界 n 为零或一、右括号提前超过左括号、输出规模是卡特兰数 反查初始化、更新顺序和退出条件。

\textbf{状态回放。} 手算路线：手算时给每层标出选择列表和路径。剪枝发生前要指出被剪分支为什么已经不可能得到合法答案。状态字段：路径、起点或使用标记、剩余目标、答案集合、剪枝条件；进入递归和退出递归必须成对恢复。状态升级：把完整搜索树加上合法性检查、剩余资源剪枝和同层去重；路径状态进入和退出要成对恢复。

\textbf{证明和边界。} 证明从搜索树覆盖开始：每个合法答案有且只有一条路径，剪枝只删掉不可能成功的分支。边界：n 为零或一、右括号提前超过左括号、输出规模是卡特兰数。变体：若有多种括号类型，状态还需要栈保存类型匹配关系。

\noindent\textbf{LeetCode 39 组合总和。} 模型：候选数字可重复使用，答案不关心排列顺序。推导重点：排序后按起点枚举，递归时仍传当前下标表示可以继续使用同一数字。

\textbf{二刷读题。} 读题时先画搜索树：层数代表什么，路径保存什么，终止条件是什么。本题可以压成“候选数字可重复使用，答案不关心排列顺序”，剪枝只能在这个搜索树上说明。先遮住答案，只保留题面，复述候选对象、合法性条件和输出目标。

\textbf{暴力回放。} 慢版本就是完整搜索树：每层列出选择列表，路径到达终止条件时检查答案。它先保证覆盖所有可能，再把明显无效或重复的分支剪掉。二刷时先写慢版本或口述慢版本，用它确认不会漏答案。

\textbf{特例回放。} 拿 candidates=[2,3,6,7]、target=7 手算，选择 2 后仍可以继续选择 2，因为每个数可重复；路径 2,2,3 成功，7 单独也成功。规律是回溯起点不回退，但递归下一层仍传当前下标以允许重复使用。把这个特例继续拆开看：一层递归对应一个明确决策，路径保存已经做出的选择，选择列表保存仍可能扩展的分支。剪枝前必须能说明被剪掉的整棵子树没有合法答案，而不是只因为它看起来不优。复盘时把这个特例压成状态字段：路径、起点或使用标记、剩余目标、答案集合、剪枝条件；进入递归和退出递归必须成对恢复；再用边界 目标小于最小数、刚好等于、候选有重复时题意如何处理、输出规模 反查初始化、更新顺序和退出条件。

\textbf{状态回放。} 手算路线：手算时给每层标出选择列表和路径。剪枝发生前要指出被剪分支为什么已经不可能得到合法答案。状态字段：路径、起点或使用标记、剩余目标、答案集合、剪枝条件；进入递归和退出递归必须成对恢复。状态升级：把完整搜索树加上合法性检查、剩余资源剪枝和同层去重；路径状态进入和退出要成对恢复。

\textbf{证明和边界。} 证明从搜索树覆盖开始：每个合法答案有且只有一条路径，剪枝只删掉不可能成功的分支。边界：目标小于最小数、刚好等于、候选有重复时题意如何处理、输出规模。变体：组合总和 II 每个数只能用一次，起点推进和去重条件都不同。

\noindent\textbf{LeetCode 40 组合总和 II。} 模型：每个候选只能使用一次，输入中可能存在重复值。推导重点：排序后递归下一层从 i 加一开始，同层遇到相同值要跳过。

\textbf{二刷读题。} 读题时先画搜索树：层数代表什么，路径保存什么，终止条件是什么。本题可以压成“每个候选只能使用一次，输入中可能存在重复值”，剪枝只能在这个搜索树上说明。先遮住答案，只保留题面，复述候选对象、合法性条件和输出目标。

\textbf{暴力回放。} 慢版本就是完整搜索树：每层列出选择列表，路径到达终止条件时检查答案。它先保证覆盖所有可能，再把明显无效或重复的分支剪掉。二刷时先写慢版本或口述慢版本，用它确认不会漏答案。

\textbf{特例回放。} 拿 candidates=[10,1,2,7,6,1,5]、target=8 手算，排序后两个 1 在同一层只能选第一个作为起点，否则会生成重复组合；但选了第一个 1 后，下一层可以选第二个 1。规律是每个数只能用一次，递归下一层传 i+1，并做同层去重。把这个特例继续拆开看：一层递归对应一个明确决策，路径保存已经做出的选择，选择列表保存仍可能扩展的分支。剪枝前必须能说明被剪掉的整棵子树没有合法答案，而不是只因为它看起来不优。复盘时把这个特例压成状态字段：路径、起点或使用标记、剩余目标、答案集合、剪枝条件；进入递归和退出递归必须成对恢复；再用边界 全重复、目标为零、同层去重写成跨层去重、答案顺序不要求 反查初始化、更新顺序和退出条件。

\textbf{状态回放。} 手算路线：手算时给每层标出选择列表和路径。剪枝发生前要指出被剪分支为什么已经不可能得到合法答案。状态字段：路径、起点或使用标记、剩余目标、答案集合、剪枝条件；进入递归和退出递归必须成对恢复。状态升级：把完整搜索树加上合法性检查、剩余资源剪枝和同层去重；路径状态进入和退出要成对恢复。

\textbf{证明和边界。} 证明从搜索树覆盖开始：每个合法答案有且只有一条路径，剪枝只删掉不可能成功的分支。边界：全重复、目标为零、同层去重写成跨层去重、答案顺序不要求。变体：子集 II 使用同样的同层去重，只是每个路径节点都可能成为答案。

\noindent\textbf{LeetCode 46 全排列。} 模型：排列关心元素顺序，每个位置从所有未使用元素中选择一个。推导重点：路径长度达到 n 时产生答案，used 数组保证同一元素不重复使用。

\textbf{二刷读题。} 读题时先画搜索树：层数代表什么，路径保存什么，终止条件是什么。本题可以压成“排列关心元素顺序，每个位置从所有未使用元素中选择一个”，剪枝只能在这个搜索树上说明。先遮住答案，只保留题面，复述候选对象、合法性条件和输出目标。

\textbf{暴力回放。} 慢版本就是完整搜索树：每层列出选择列表，路径到达终止条件时检查答案。它先保证覆盖所有可能，再把明显无效或重复的分支剪掉。二刷时先写慢版本或口述慢版本，用它确认不会漏答案。

\textbf{特例回放。} 拿 [1,2,3] 手算，第一位可选 1、2、3；选 1 后，第二位只能从未使用的 2、3 中选。规律是路径长度等于 n 时得到一个排列，used 数组表示哪些下标已经在当前路径里。把这个特例继续拆开看：一层递归对应一个明确决策，路径保存已经做出的选择，选择列表保存仍可能扩展的分支。剪枝前必须能说明被剪掉的整棵子树没有合法答案，而不是只因为它看起来不优。复盘时把这个特例压成状态字段：路径、起点或使用标记、剩余目标、答案集合、剪枝条件；进入递归和退出递归必须成对恢复；再用边界 空数组、单元素、输出规模为 n 阶乘、元素值唯一假设 反查初始化、更新顺序和退出条件。

\textbf{状态回放。} 手算路线：手算时给每层标出选择列表和路径。剪枝发生前要指出被剪分支为什么已经不可能得到合法答案。状态字段：路径、起点或使用标记、剩余目标、答案集合、剪枝条件；进入递归和退出递归必须成对恢复。状态升级：把完整搜索树加上合法性检查、剩余资源剪枝和同层去重；路径状态进入和退出要成对恢复。

\textbf{证明和边界。} 证明从搜索树覆盖开始：每个合法答案有且只有一条路径，剪枝只删掉不可能成功的分支。边界：空数组、单元素、输出规模为 n 阶乘、元素值唯一假设。变体：若有重复元素，需要排序并在同层跳过等价选择。

\noindent\textbf{LeetCode 47 全排列 II。} 模型：重复元素会让不同下标选择生成相同排列。推导重点：排序后使用 used 数组，同层遇到相同值且前一个相同值未使用时跳过。

\textbf{二刷读题。} 读题时先画搜索树：层数代表什么，路径保存什么，终止条件是什么。本题可以压成“重复元素会让不同下标选择生成相同排列”，剪枝只能在这个搜索树上说明。先遮住答案，只保留题面，复述候选对象、合法性条件和输出目标。

\textbf{暴力回放。} 慢版本就是完整搜索树：每层列出选择列表，路径到达终止条件时检查答案。它先保证覆盖所有可能，再把明显无效或重复的分支剪掉。二刷时先写慢版本或口述慢版本，用它确认不会漏答案。

\textbf{特例回放。} 拿 [1,1,2] 手算，第一位选第一个 1 和第二个 1 会生成同样分支，所以同层必须跳过后一个 1；但第一层选 1 后，第二层仍可选另一个 1。规律是排序后用 used 区分路径内使用，同层重复值只保留第一个未使用代表。把这个特例继续拆开看：一层递归对应一个明确决策，路径保存已经做出的选择，选择列表保存仍可能扩展的分支。剪枝前必须能说明被剪掉的整棵子树没有合法答案，而不是只因为它看起来不优。复盘时把这个特例压成状态字段：路径、起点或使用标记、剩余目标、答案集合、剪枝条件；进入递归和退出递归必须成对恢复；再用边界 全重复、部分重复、同层去重条件写反、输出顺序不要求 反查初始化、更新顺序和退出条件。

\textbf{状态回放。} 手算路线：手算时给每层标出选择列表和路径。剪枝发生前要指出被剪分支为什么已经不可能得到合法答案。状态字段：路径、起点或使用标记、剩余目标、答案集合、剪枝条件；进入递归和退出递归必须成对恢复。状态升级：把完整搜索树加上合法性检查、剩余资源剪枝和同层去重；路径状态进入和退出要成对恢复。

\textbf{证明和边界。} 证明从搜索树覆盖开始：每个合法答案有且只有一条路径，剪枝只删掉不可能成功的分支。边界：全重复、部分重复、同层去重条件写反、输出顺序不要求。变体：子集 II 和组合总和 II 也是同层去重，但起点推进方式不同。

\noindent\textbf{LeetCode 51 N 皇后。} 模型：每一行放一个皇后，列和两条对角线不能冲突。推导重点：逐行搜索，用列集合、主对角线集合和副对角线集合快速判断合法性。

\textbf{二刷读题。} 读题时先画搜索树：层数代表什么，路径保存什么，终止条件是什么。本题可以压成“每一行放一个皇后，列和两条对角线不能冲突”，剪枝只能在这个搜索树上说明。先遮住答案，只保留题面，复述候选对象、合法性条件和输出目标。

\textbf{暴力回放。} 慢版本就是完整搜索树：每层列出选择列表，路径到达终止条件时检查答案。它先保证覆盖所有可能，再把明显无效或重复的分支剪掉。二刷时先写慢版本或口述慢版本，用它确认不会漏答案。

\textbf{特例回放。} 拿 n=4 手算，第一行放某列后，下一行不能放同列、主对角线和副对角线冲突的位置。规律是按行递归，每行放一个皇后，三个集合分别记录列、row-col、row+col 是否被占用。把这个特例继续拆开看：一层递归对应一个明确决策，路径保存已经做出的选择，选择列表保存仍可能扩展的分支。剪枝前必须能说明被剪掉的整棵子树没有合法答案，而不是只因为它看起来不优。复盘时把这个特例压成状态字段：路径、起点或使用标记、剩余目标、答案集合、剪枝条件；进入递归和退出递归必须成对恢复；再用边界 n 为一、n 为二或三无解、对角线编号、回退时状态恢复 反查初始化、更新顺序和退出条件。

\textbf{状态回放。} 手算路线：手算时给每层标出选择列表和路径。剪枝发生前要指出被剪分支为什么已经不可能得到合法答案。状态字段：路径、起点或使用标记、剩余目标、答案集合、剪枝条件；进入递归和退出递归必须成对恢复。状态升级：把完整搜索树加上合法性检查、剩余资源剪枝和同层去重；路径状态进入和退出要成对恢复。

\textbf{证明和边界。} 证明从搜索树覆盖开始：每个合法答案有且只有一条路径，剪枝只删掉不可能成功的分支。边界：n 为一、n 为二或三无解、对角线编号、回退时状态恢复。变体：可用位掩码压缩三个约束集合，提高常数性能。

\noindent\textbf{LeetCode 78 子集。} 模型：每个元素都有选或不选两种选择，答案是所有路径节点。推导重点：回溯按起点向后枚举，避免不同顺序生成同一集合。

\textbf{二刷读题。} 读题时先画搜索树：层数代表什么，路径保存什么，终止条件是什么。本题可以压成“每个元素都有选或不选两种选择，答案是所有路径节点”，剪枝只能在这个搜索树上说明。先遮住答案，只保留题面，复述候选对象、合法性条件和输出目标。

\textbf{暴力回放。} 慢版本就是完整搜索树：每层列出选择列表，路径到达终止条件时检查答案。它先保证覆盖所有可能，再把明显无效或重复的分支剪掉。二刷时先写慢版本或口述慢版本，用它确认不会漏答案。

\textbf{特例回放。} 拿 [1,2] 画搜索树，第一层可以不选 1 或选 1；在选 1 的分支里再决定 2，得到 [1] 和 [1,2]。规律是每个元素对应选或不选，或用起点向后枚举，路径上的每个节点都是一个子集；空数组也要输出空集。把这个特例继续拆开看：一层递归对应一个明确决策，路径保存已经做出的选择，选择列表保存仍可能扩展的分支。剪枝前必须能说明被剪掉的整棵子树没有合法答案，而不是只因为它看起来不优。复盘时把这个特例压成状态字段：路径、起点或使用标记、剩余目标、答案集合、剪枝条件；进入递归和退出递归必须成对恢复；再用边界 空数组、单元素、输出数量为二的 n 次方、顺序不要求 反查初始化、更新顺序和退出条件。

\textbf{状态回放。} 手算路线：手算时给每层标出选择列表和路径。剪枝发生前要指出被剪分支为什么已经不可能得到合法答案。状态字段：路径、起点或使用标记、剩余目标、答案集合、剪枝条件；进入递归和退出递归必须成对恢复。状态升级：把完整搜索树加上合法性检查、剩余资源剪枝和同层去重；路径状态进入和退出要成对恢复。

\textbf{证明和边界。} 证明从搜索树覆盖开始：每个合法答案有且只有一条路径，剪枝只删掉不可能成功的分支。边界：空数组、单元素、输出数量为二的 n 次方、顺序不要求。变体：也可用位掩码枚举所有子集，适合 n 较小的场景。

\noindent\textbf{LeetCode 79 单词搜索。} 模型：网格路径不能重复使用同一格，状态包含位置、匹配下标和访问标记。推导重点：剪枝包括首字符不匹配、剩余长度不足、字符频次不足和越界访问。

\textbf{二刷读题。} 读题时先画搜索树：层数代表什么，路径保存什么，终止条件是什么。本题可以压成“网格路径不能重复使用同一格，状态包含位置、匹配下标和访问标记”，剪枝只能在这个搜索树上说明。先遮住答案，只保留题面，复述候选对象、合法性条件和输出目标。

\textbf{暴力回放。} 慢版本就是完整搜索树：每层列出选择列表，路径到达终止条件时检查答案。它先保证覆盖所有可能，再把明显无效或重复的分支剪掉。二刷时先写慢版本或口述慢版本，用它确认不会漏答案。

\textbf{特例回放。} 拿 board=[A B; C D]、word=AB 手算，从 A 出发只能走到相邻的 B，不能跳到 D；若搜索 ABA，第二个 A 不能复用起点格。规律是状态包含位置、匹配下标和访问标记，递归退出时要恢复访问标记。把这个特例继续拆开看：一层递归对应一个明确决策，路径保存已经做出的选择，选择列表保存仍可能扩展的分支。剪枝前必须能说明被剪掉的整棵子树没有合法答案，而不是只因为它看起来不优。复盘时把这个特例压成状态字段：路径、起点或使用标记、剩余目标、答案集合、剪枝条件；进入递归和退出递归必须成对恢复；再用边界 单字符单词、路径需要回退、重复字母、单元格不能复用 反查初始化、更新顺序和退出条件。

\textbf{状态回放。} 手算路线：手算时给每层标出选择列表和路径。剪枝发生前要指出被剪分支为什么已经不可能得到合法答案。状态字段：路径、起点或使用标记、剩余目标、答案集合、剪枝条件；进入递归和退出递归必须成对恢复。状态升级：把完整搜索树加上合法性检查、剩余资源剪枝和同层去重；路径状态进入和退出要成对恢复。

\textbf{证明和边界。} 证明从搜索树覆盖开始：每个合法答案有且只有一条路径，剪枝只删掉不可能成功的分支。边界：单字符单词、路径需要回退、重复字母、单元格不能复用。变体：若要查很多单词，应使用 Trie 合并搜索前缀。

\noindent\textbf{LeetCode 90 子集 II。} 模型：重复元素会让不同下标路径生成相同集合。推导重点：排序后在同一层跳过相同值，保留不同层使用重复值的可能。

\textbf{二刷读题。} 读题时先画搜索树：层数代表什么，路径保存什么，终止条件是什么。本题可以压成“重复元素会让不同下标路径生成相同集合”，剪枝只能在这个搜索树上说明。先遮住答案，只保留题面，复述候选对象、合法性条件和输出目标。

\textbf{暴力回放。} 慢版本就是完整搜索树：每层列出选择列表，路径到达终止条件时检查答案。它先保证覆盖所有可能，再把明显无效或重复的分支剪掉。二刷时先写慢版本或口述慢版本，用它确认不会漏答案。

\textbf{特例回放。} 拿 [1,2,2] 手算，排序后第二个 2 如果在同一层再次作为起点，会生成和第一个 2 同样的子集；但在已经选择第一个 2 的下一层，第二个 2 又可以被选出 [2,2]。规律是同层跳过重复值，不是全局禁止重复值。把这个特例继续拆开看：一层递归对应一个明确决策，路径保存已经做出的选择，选择列表保存仍可能扩展的分支。剪枝前必须能说明被剪掉的整棵子树没有合法答案，而不是只因为它看起来不优。复盘时把这个特例压成状态字段：路径、起点或使用标记、剩余目标、答案集合、剪枝条件；进入递归和退出递归必须成对恢复；再用边界 全重复、空集、重复值相邻、去重条件写成同层而不是全局 反查初始化、更新顺序和退出条件。

\textbf{状态回放。} 手算路线：手算时给每层标出选择列表和路径。剪枝发生前要指出被剪分支为什么已经不可能得到合法答案。状态字段：路径、起点或使用标记、剩余目标、答案集合、剪枝条件；进入递归和退出递归必须成对恢复。状态升级：把完整搜索树加上合法性检查、剩余资源剪枝和同层去重；路径状态进入和退出要成对恢复。

\textbf{证明和边界。} 证明从搜索树覆盖开始：每个合法答案有且只有一条路径，剪枝只删掉不可能成功的分支。边界：全重复、空集、重复值相邻、去重条件写成同层而不是全局。变体：组合总和 II 使用同样的同层去重思想。

\noindent\textbf{LeetCode 93 复原 IP 地址。} 模型：字符串要切成四段，每段必须是 0 到 255 且无非法前导零。推导重点：按段数和剩余长度剪枝，每次尝试长度一到三的片段。

\textbf{二刷读题。} 读题时先画搜索树：层数代表什么，路径保存什么，终止条件是什么。本题可以压成“字符串要切成四段，每段必须是 0 到 255 且无非法前导零”，剪枝只能在这个搜索树上说明。先遮住答案，只保留题面，复述候选对象、合法性条件和输出目标。

\textbf{暴力回放。} 慢版本就是完整搜索树：每层列出选择列表，路径到达终止条件时检查答案。它先保证覆盖所有可能，再把明显无效或重复的分支剪掉。二刷时先写慢版本或口述慢版本，用它确认不会漏答案。

\textbf{特例回放。} 拿 25525511135 手算，第一段可以是 255，但不能是 2552；段值必须在 0..255，且 01 这种前导零非法。规律是回溯选四段，每段长度 1 到 3，剩余字符数必须能填满剩余段数。把这个特例继续拆开看：一层递归对应一个明确决策，路径保存已经做出的选择，选择列表保存仍可能扩展的分支。剪枝前必须能说明被剪掉的整棵子树没有合法答案，而不是只因为它看起来不优。复盘时把这个特例压成状态字段：路径、起点或使用标记、剩余目标、答案集合、剪枝条件；进入递归和退出递归必须成对恢复；再用边界 前导零、数值超过 255、剩余字符太多或太少、刚好四段 反查初始化、更新顺序和退出条件。

\textbf{状态回放。} 手算路线：手算时给每层标出选择列表和路径。剪枝发生前要指出被剪分支为什么已经不可能得到合法答案。状态字段：路径、起点或使用标记、剩余目标、答案集合、剪枝条件；进入递归和退出递归必须成对恢复。状态升级：把完整搜索树加上合法性检查、剩余资源剪枝和同层去重；路径状态进入和退出要成对恢复。

\textbf{证明和边界。} 证明从搜索树覆盖开始：每个合法答案有且只有一条路径，剪枝只删掉不可能成功的分支。边界：前导零、数值超过 255、剩余字符太多或太少、刚好四段。变体：分割回文串也是前缀分割，但合法性检查不同。

\noindent\textbf{LeetCode 131 分割回文串。} 模型：每次选择一个从当前位置开始的回文前缀。推导重点：预处理回文 DP 表，让回溯合法性检查为常数。

\textbf{二刷读题。} 读题时先画搜索树：层数代表什么，路径保存什么，终止条件是什么。本题可以压成“每次选择一个从当前位置开始的回文前缀”，剪枝只能在这个搜索树上说明。先遮住答案，只保留题面，复述候选对象、合法性条件和输出目标。

\textbf{暴力回放。} 慢版本就是完整搜索树：每层列出选择列表，路径到达终止条件时检查答案。它先保证覆盖所有可能，再把明显无效或重复的分支剪掉。二刷时先写慢版本或口述慢版本，用它确认不会漏答案。

\textbf{特例回放。} 拿 aab 手算，第一刀可以切 a，后面 ab 不全是回文；也可以切 aa，后面 b 是回文，得到 [aa,b]。规律是回溯枚举切分点，但每段必须先通过回文检查；预处理回文表可以避免重复检查子串。把这个特例继续拆开看：一层递归对应一个明确决策，路径保存已经做出的选择，选择列表保存仍可能扩展的分支。剪枝前必须能说明被剪掉的整棵子树没有合法答案，而不是只因为它看起来不优。复盘时把这个特例压成状态字段：路径、起点或使用标记、剩余目标、答案集合、剪枝条件；进入递归和退出递归必须成对恢复；再用边界 空串、单字符、全相同字符、重复分割路径 反查初始化、更新顺序和退出条件。

\textbf{状态回放。} 手算路线：手算时给每层标出选择列表和路径。剪枝发生前要指出被剪分支为什么已经不可能得到合法答案。状态字段：路径、起点或使用标记、剩余目标、答案集合、剪枝条件；进入递归和退出递归必须成对恢复。状态升级：把完整搜索树加上合法性检查、剩余资源剪枝和同层去重；路径状态进入和退出要成对恢复。

\textbf{证明和边界。} 证明从搜索树覆盖开始：每个合法答案有且只有一条路径，剪枝只删掉不可能成功的分支。边界：空串、单字符、全相同字符、重复分割路径。变体：若只求最少切割次数，可改成 DP 最短切分。

\noindent\textbf{LeetCode 140 单词拆分 II。} 模型：需要输出所有由字典单词拼出的句子，输出规模本身可能指数级。推导重点：先用 DP 或记忆化判断后缀是否可拆，再 DFS 枚举可行前缀并在末尾拼接句子。

\textbf{二刷读题。} 读题时先画搜索树：层数代表什么，路径保存什么，终止条件是什么。本题可以压成“需要输出所有由字典单词拼出的句子，输出规模本身可能指数级”，剪枝只能在这个搜索树上说明。先遮住答案，只保留题面，复述候选对象、合法性条件和输出目标。

\textbf{暴力回放。} 慢版本就是完整搜索树：每层列出选择列表，路径到达终止条件时检查答案。它先保证覆盖所有可能，再把明显无效或重复的分支剪掉。二刷时先写慢版本或口述慢版本，用它确认不会漏答案。

\textbf{特例回放。} 拿 catsanddog 手算，前缀 cat 和 cats 都能走，但某些后缀继续切会失败。若直接 DFS，每次都会重复探索同一个后缀。规律是先用记忆化记录某个起点之后能生成哪些句子，或先用 DP 剪掉不可拆后缀，再枚举路径。把这个特例继续拆开看：一层递归对应一个明确决策，路径保存已经做出的选择，选择列表保存仍可能扩展的分支。剪枝前必须能说明被剪掉的整棵子树没有合法答案，而不是只因为它看起来不优。复盘时把这个特例压成状态字段：路径、起点或使用标记、剩余目标、答案集合、剪枝条件；进入递归和退出递归必须成对恢复；再用边界 无解后缀、重复单词、路径拼接成本、答案数量主导复杂度 反查初始化、更新顺序和退出条件。

\textbf{状态回放。} 手算路线：手算时给每层标出选择列表和路径。剪枝发生前要指出被剪分支为什么已经不可能得到合法答案。状态字段：路径、起点或使用标记、剩余目标、答案集合、剪枝条件；进入递归和退出递归必须成对恢复。状态升级：把完整搜索树加上合法性检查、剩余资源剪枝和同层去重；路径状态进入和退出要成对恢复。

\textbf{证明和边界。} 证明从搜索树覆盖开始：每个合法答案有且只有一条路径，剪枝只删掉不可能成功的分支。边界：无解后缀、重复单词、路径拼接成本、答案数量主导复杂度。变体：只问能否拆分时用布尔 DP 即可，不要承担枚举所有句子的成本。

\noindent\textbf{LeetCode 212 单词搜索 II。} 模型：多个单词在同一棋盘上搜索，公共前缀可以共享。推导重点：先建 Trie，再从每个格子回溯；路径到达单词结束节点就收集答案。

\textbf{二刷读题。} 读题时先画搜索树：层数代表什么，路径保存什么，终止条件是什么。本题可以压成“多个单词在同一棋盘上搜索，公共前缀可以共享”，剪枝只能在这个搜索树上说明。先遮住答案，只保留题面，复述候选对象、合法性条件和输出目标。

\textbf{暴力回放。} 慢版本就是完整搜索树：每层列出选择列表，路径到达终止条件时检查答案。它先保证覆盖所有可能，再把明显无效或重复的分支剪掉。二刷时先写慢版本或口述慢版本，用它确认不会漏答案。

\textbf{特例回放。} 拿 board 上有 oath、pea、eat 手算，单词共享前缀时不应对每个词重复全图搜索；从 o 出发沿 Trie 前缀 o-a-t-h 找到 oath。规律是 Trie 合并所有单词前缀，DFS 走棋盘同时走 Trie，命中单词后记录并可去重。把这个特例继续拆开看：一层递归对应一个明确决策，路径保存已经做出的选择，选择列表保存仍可能扩展的分支。剪枝前必须能说明被剪掉的整棵子树没有合法答案，而不是只因为它看起来不优。复盘时把这个特例压成状态字段：路径、起点或使用标记、剩余目标、答案集合、剪枝条件；进入递归和退出递归必须成对恢复；再用边界 重复单词、同一格不能复用、前缀是完整单词、剪枝删除空子树 反查初始化、更新顺序和退出条件。

\textbf{状态回放。} 手算路线：手算时给每层标出选择列表和路径。剪枝发生前要指出被剪分支为什么已经不可能得到合法答案。状态字段：路径、起点或使用标记、剩余目标、答案集合、剪枝条件；进入递归和退出递归必须成对恢复。状态升级：把完整搜索树加上合法性检查、剩余资源剪枝和同层去重；路径状态进入和退出要成对恢复。

\textbf{证明和边界。} 证明从搜索树覆盖开始：每个合法答案有且只有一条路径，剪枝只删掉不可能成功的分支。边界：重复单词、同一格不能复用、前缀是完整单词、剪枝删除空子树。变体：这是单词搜索从单目标扩展到多目标后的结构升级。

\noindent\textbf{LeetCode 216 组合总和 III。} 模型：从一到九中选 k 个数和为 n，每个数最多用一次。推导重点：按递增起点枚举，利用剩余个数和剩余和剪枝。

\textbf{二刷读题。} 读题时先画搜索树：层数代表什么，路径保存什么，终止条件是什么。本题可以压成“从一到九中选 k 个数和为 n，每个数最多用一次”，剪枝只能在这个搜索树上说明。先遮住答案，只保留题面，复述候选对象、合法性条件和输出目标。

\textbf{暴力回放。} 慢版本就是完整搜索树：每层列出选择列表，路径到达终止条件时检查答案。它先保证覆盖所有可能，再把明显无效或重复的分支剪掉。二刷时先写慢版本或口述慢版本，用它确认不会漏答案。

\textbf{特例回放。} 拿 k=3、n=7 手算，选择 1 后还需两个数凑 6，可以选 2 和 4；若当前和已经超过 7，就不用再往下扩展。规律是路径长度、剩余和和起点共同定义状态，数字只能向后选避免重复组合。把这个特例继续拆开看：一层递归对应一个明确决策，路径保存已经做出的选择，选择列表保存仍可能扩展的分支。剪枝前必须能说明被剪掉的整棵子树没有合法答案，而不是只因为它看起来不优。复盘时把这个特例压成状态字段：路径、起点或使用标记、剩余目标、答案集合、剪枝条件；进入递归和退出递归必须成对恢复；再用边界 n 太小、n 太大、k 为零、路径长度已经超过 k 反查初始化、更新顺序和退出条件。

\textbf{状态回放。} 手算路线：手算时给每层标出选择列表和路径。剪枝发生前要指出被剪分支为什么已经不可能得到合法答案。状态字段：路径、起点或使用标记、剩余目标、答案集合、剪枝条件；进入递归和退出递归必须成对恢复。状态升级：把完整搜索树加上合法性检查、剩余资源剪枝和同层去重；路径状态进入和退出要成对恢复。

\textbf{证明和边界。} 证明从搜索树覆盖开始：每个合法答案有且只有一条路径，剪枝只删掉不可能成功的分支。边界：n 太小、n 太大、k 为零、路径长度已经超过 k。变体：可预估剩余最小和最大和进一步剪枝。

\noindent\textbf{LeetCode 698 划分为 k 个相等的子集。} 模型：每个元素要分配到 k 个桶中，所有桶目标和相同。推导重点：先判断总和可整除并排序降序，再用桶剩余容量回溯放置元素。

\textbf{二刷读题。} 读题时先画搜索树：层数代表什么，路径保存什么，终止条件是什么。本题可以压成“每个元素要分配到 k 个桶中，所有桶目标和相同”，剪枝只能在这个搜索树上说明。先遮住答案，只保留题面，复述候选对象、合法性条件和输出目标。

\textbf{暴力回放。} 慢版本就是完整搜索树：每层列出选择列表，路径到达终止条件时检查答案。它先保证覆盖所有可能，再把明显无效或重复的分支剪掉。二刷时先写慢版本或口述慢版本，用它确认不会漏答案。

\textbf{特例回放。} 拿 [4,3,2,3,5,2,1]、k=4 手算，总和 20，每桶目标 5；先放 5 单独成桶，再尝试 4+1、3+2。规律是回溯状态是桶剩余容量和已用元素，排序后先放大数能更早剪枝。把这个特例继续拆开看：一层递归对应一个明确决策，路径保存已经做出的选择，选择列表保存仍可能扩展的分支。剪枝前必须能说明被剪掉的整棵子树没有合法答案，而不是只因为它看起来不优。复盘时把这个特例压成状态字段：路径、起点或使用标记、剩余目标、答案集合、剪枝条件；进入递归和退出递归必须成对恢复；再用边界 存在元素大于目标、重复元素、空桶对称、k 为一 反查初始化、更新顺序和退出条件。

\textbf{状态回放。} 手算路线：手算时给每层标出选择列表和路径。剪枝发生前要指出被剪分支为什么已经不可能得到合法答案。状态字段：路径、起点或使用标记、剩余目标、答案集合、剪枝条件；进入递归和退出递归必须成对恢复。状态升级：把完整搜索树加上合法性检查、剩余资源剪枝和同层去重；路径状态进入和退出要成对恢复。

\textbf{证明和边界。} 证明从搜索树覆盖开始：每个合法答案有且只有一条路径，剪枝只删掉不可能成功的分支。边界：存在元素大于目标、重复元素、空桶对称、k 为一。变体：状态压缩 DP 也能做，适合用掩码表示已选元素集合。

\subsection{自测标准}

二刷时不要只确认自己见过这道题。请遮住答案，先讲题面模型，再讲暴力基线和瓶颈，然后说出优化结构保存了什么、不变量如何排除错误候选、边界实验如何打破错误实现。

\section{动态规划与状态定义：二刷复盘卡}
这一大节只围绕一个题型复盘，不把题号直接铺成目录。阅读时先看复盘目标，再读核心题卡，最后用自测标准检查自己是否能独立重讲。

\subsection{复盘目标}

追问通常会问状态是否足够、初始化为什么这样、空间压缩读到的是旧值还是新值。请准备从暴力搜索树指出重复子问题。

\subsection{核心题卡}

\noindent\textbf{LeetCode 10 正则表达式匹配。} 模型：模式中的点号匹配单字符，星号修饰前一个字符并允许出现零次或多次。推导重点：二维 DP 表示两个前缀是否匹配，星号转移同时覆盖零次和消耗一个字符后继续匹配。

\textbf{二刷读题。} 读题时先找重复子问题，而不是直接背公式。本题可以压成“模式中的点号匹配单字符，星号修饰前一个字符并允许出现零次或多次”，二维 DP 表示两个前缀是否匹配，星号转移同时覆盖零次和消耗一个字符后继续匹配就是状态之间的依赖关系。先遮住答案，只保留题面，复述候选对象、合法性条件和输出目标。

\textbf{暴力回放。} 慢版本先写递归搜索树：节点表示处理位置、剩余资源和当前选择。只要不同路径反复到达同一个节点，就说明需要把这个节点命名成 DP 状态。二刷时先写慢版本或口述慢版本，用它确认不会漏答案。

\textbf{特例回放。} 拿 s=a、p=a* 手算，星号可以让 a 出现一次；拿 s=空、p=a*，星号又可以让 a 出现零次。再拿 s=aa、p=a*，消耗一个 a 后模式仍停在 a*，还能继续匹配。规律是星号有两条分叉：当零次用时看 dp[i][j-2]，当消耗当前字符时看 dp[i-1][j]。把这个特例继续拆开看：先问暴力搜索里哪些不同路径到达同一个后续问题，再给这个后续问题命名为状态。状态必须包含未来决策需要的全部信息，转移只从更小且已经算清的状态来。复盘时把这个特例压成状态字段：状态下标、状态值语义、初始化、转移来源、遍历顺序；空间压缩时还要指出读到的是旧值还是新值；再用边界 空串、能匹配空串的星号链、星号前必须有字符、点号和普通字符的匹配条件 反查初始化、更新顺序和退出条件。

\textbf{状态回放。} 手算路线：手算时先写一个很小的 DP 表或记忆化搜索记录。每个格子旁边写它来自哪些更小状态。状态字段：状态下标、状态值语义、初始化、转移来源、遍历顺序；空间压缩时还要指出读到的是旧值还是新值。状态升级：把重复递归节点命名为状态；遍历顺序只是在保证依赖状态已经算完。

\textbf{证明和边界。} 证明从最后一步开始：任意最优解的最后一步都在转移枚举中，转移来源已经由更小状态正确计算。边界：空串、能匹配空串的星号链、星号前必须有字符、点号和普通字符的匹配条件。变体：通配符匹配中的星号语义是任意长度字符串，不能直接搬用本题转移。

\noindent\textbf{LeetCode 44 通配符匹配。} 模型：问号匹配一个字符，星号匹配任意长度字符串，状态是两个前缀能否匹配。推导重点：DP 中星号可匹配空串或继续吞掉当前字符；贪心写法记录最近星号并在失配时回退扩展星号覆盖范围。

\textbf{二刷读题。} 读题时先找重复子问题，而不是直接背公式。本题可以压成“问号匹配一个字符，星号匹配任意长度字符串，状态是两个前缀能否匹配”，DP 中星号可匹配空串或继续吞掉当前字符；贪心写法记录最近星号并在失配时回退扩展星号覆盖范围就是状态之间的依赖关系。先遮住答案，只保留题面，复述候选对象、合法性条件和输出目标。

\textbf{暴力回放。} 慢版本先写递归搜索树：节点表示处理位置、剩余资源和当前选择。只要不同路径反复到达同一个节点，就说明需要把这个节点命名成 DP 状态。二刷时先写慢版本或口述慢版本，用它确认不会漏答案。

\textbf{特例回放。} 拿 s=ab、p=a* 手算，星号可以吞掉 b；拿 s=空、p=*，星号也可以代表空。DP 中星号的两条路是匹配空串或继续吞当前字符。贪心中失配时回到最近星号，让它多吞一个字符。规律是本题星号独立代表任意串，和正则里修饰前一字符的星号不同。把这个特例继续拆开看：先问暴力搜索里哪些不同路径到达同一个后续问题，再给这个后续问题命名为状态。状态必须包含未来决策需要的全部信息，转移只从更小且已经算清的状态来。复盘时把这个特例压成状态字段：状态下标、状态值语义、初始化、转移来源、遍历顺序；空间压缩时还要指出读到的是旧值还是新值；再用边界 连续星号、空串、模式只含星号、贪心回退时字符串位置不能倒退错 反查初始化、更新顺序和退出条件。

\textbf{状态回放。} 手算路线：手算时先写一个很小的 DP 表或记忆化搜索记录。每个格子旁边写它来自哪些更小状态。状态字段：状态下标、状态值语义、初始化、转移来源、遍历顺序；空间压缩时还要指出读到的是旧值还是新值。状态升级：把重复递归节点命名为状态；遍历顺序只是在保证依赖状态已经算完。

\textbf{证明和边界。} 证明从最后一步开始：任意最优解的最后一步都在转移枚举中，转移来源已经由更小状态正确计算。边界：连续星号、空串、模式只含星号、贪心回退时字符串位置不能倒退错。变体：正则表达式匹配的星号修饰前一个字符，和本题星号独立匹配任意串不同。

\noindent\textbf{LeetCode 53 最大子数组和。} 模型：状态是必须以当前位置结尾的最大连续和。推导重点：用 [-2,3] 和 [2,3] 对比，推出当前位置只在单独开始与接上旧后缀之间取较大值。

\textbf{二刷读题。} 读题时先找重复子问题，而不是直接背公式。本题可以压成“状态是必须以当前位置结尾的最大连续和”，用 [-2,3] 和 [2,3] 对比，推出当前位置只在单独开始与接上旧后缀之间取较大值就是状态之间的依赖关系。先遮住答案，只保留题面，复述候选对象、合法性条件和输出目标。

\textbf{暴力回放。} 慢版本先写递归搜索树：节点表示处理位置、剩余资源和当前选择。只要不同路径反复到达同一个节点，就说明需要把这个节点命名成 DP 状态。二刷时先写慢版本或口述慢版本，用它确认不会漏答案。

\textbf{特例回放。} 拿 [-2, 3] 手算。以 3 结尾的连续段只有两个候选：单独 [3]，或者接上前一个最佳后缀 [-2,3]，后者和为 1，比 3 更差，所以负后缀必须丢。再拿 [2, 3]，候选是 [3] 与 [2,3]，后者为 5，更好，所以正后缀要接。再拿 [-1, 0]，接不接都不增益，可以从当前重新开始而不影响最大和。规律是当前位置只需要比较 nums[i] 与 bestEndingHere + nums[i]，全局答案再和 bestEndingHere 比较；全负数组要求答案初始化为第一个数，不能初始化为 0。把这个特例继续拆开看：先问暴力搜索里哪些不同路径到达同一个后续问题，再给这个后续问题命名为状态。状态必须包含未来决策需要的全部信息，转移只从更小且已经算清的状态来。复盘时把这个特例压成状态字段：状态下标、状态值语义、初始化、转移来源、遍历顺序；空间压缩时还要指出读到的是旧值还是新值；再用边界 全负数组、单元素、和溢出、答案不能初始化为零 反查初始化、更新顺序和退出条件。

\textbf{状态回放。} 手算路线：手算时先写一个很小的 DP 表或记忆化搜索记录。每个格子旁边写它来自哪些更小状态。状态字段：状态下标、状态值语义、初始化、转移来源、遍历顺序；空间压缩时还要指出读到的是旧值还是新值。状态升级：把重复递归节点命名为状态；遍历顺序只是在保证依赖状态已经算完。

\textbf{证明和边界。} 证明从最后一步开始：任意最优解的最后一步都在转移枚举中，转移来源已经由更小状态正确计算。边界：全负数组、单元素、和溢出、答案不能初始化为零。变体：若要求返回区间下标，同步记录重新开始位置和最优左右端。

\noindent\textbf{LeetCode 62 不同路径。} 模型：网格状态由上方和左方两种最后一步转移而来。推导重点：二维表可压缩为一维，因为当前行只依赖上一行同列和当前行左侧。

\textbf{二刷读题。} 读题时先找重复子问题，而不是直接背公式。本题可以压成“网格状态由上方和左方两种最后一步转移而来”，二维表可压缩为一维，因为当前行只依赖上一行同列和当前行左侧就是状态之间的依赖关系。先遮住答案，只保留题面，复述候选对象、合法性条件和输出目标。

\textbf{暴力回放。} 慢版本先写递归搜索树：节点表示处理位置、剩余资源和当前选择。只要不同路径反复到达同一个节点，就说明需要把这个节点命名成 DP 状态。二刷时先写慢版本或口述慢版本，用它确认不会漏答案。

\textbf{特例回放。} 拿 2x3 网格手算，到每个格子的最后一步只能从上方或左方来；第一行和第一列都只有 1 条路，右下角等于左边 2 加上上边 1，得到 3。规律是 dp[row][col]=上方+左方，滚动数组只是在保存上一行同列和当前行左侧。把这个特例继续拆开看：先问暴力搜索里哪些不同路径到达同一个后续问题，再给这个后续问题命名为状态。状态必须包含未来决策需要的全部信息，转移只从更小且已经算清的状态来。复盘时把这个特例压成状态字段：状态下标、状态值语义、初始化、转移来源、遍历顺序；空间压缩时还要指出读到的是旧值还是新值；再用边界 一行、一列、较大网格结果溢出、障碍物版本边界 反查初始化、更新顺序和退出条件。

\textbf{状态回放。} 手算路线：手算时先写一个很小的 DP 表或记忆化搜索记录。每个格子旁边写它来自哪些更小状态。状态字段：状态下标、状态值语义、初始化、转移来源、遍历顺序；空间压缩时还要指出读到的是旧值还是新值。状态升级：把重复递归节点命名为状态；遍历顺序只是在保证依赖状态已经算完。

\textbf{证明和边界。} 证明从最后一步开始：任意最优解的最后一步都在转移枚举中，转移来源已经由更小状态正确计算。边界：一行、一列、较大网格结果溢出、障碍物版本边界。变体：带障碍物时遇到障碍状态清零，第一行第一列要按阻断处理。

\noindent\textbf{LeetCode 63 不同路径 II。} 模型：障碍物让某些格子的路径数变为零。推导重点：滚动数组中遇到障碍直接把当前位置清零，否则累加左侧。

\textbf{二刷读题。} 读题时先找重复子问题，而不是直接背公式。本题可以压成“障碍物让某些格子的路径数变为零”，滚动数组中遇到障碍直接把当前位置清零，否则累加左侧就是状态之间的依赖关系。先遮住答案，只保留题面，复述候选对象、合法性条件和输出目标。

\textbf{暴力回放。} 慢版本先写递归搜索树：节点表示处理位置、剩余资源和当前选择。只要不同路径反复到达同一个节点，就说明需要把这个节点命名成 DP 状态。二刷时先写慢版本或口述慢版本，用它确认不会漏答案。

\textbf{特例回放。} 拿 [[0,0],[1,0]] 手算，左下角是障碍，路径数必须清零；右下角只能从上方到达，所以答案是 1。规律是障碍格不是普通转移，而是把当前状态置零；起点或终点是障碍时也按同一清零语义处理。把这个特例继续拆开看：先问暴力搜索里哪些不同路径到达同一个后续问题，再给这个后续问题命名为状态。状态必须包含未来决策需要的全部信息，转移只从更小且已经算清的状态来。复盘时把这个特例压成状态字段：状态下标、状态值语义、初始化、转移来源、遍历顺序；空间压缩时还要指出读到的是旧值还是新值；再用边界 起点或终点是障碍、整行被截断、单行障碍 反查初始化、更新顺序和退出条件。

\textbf{状态回放。} 手算路线：手算时先写一个很小的 DP 表或记忆化搜索记录。每个格子旁边写它来自哪些更小状态。状态字段：状态下标、状态值语义、初始化、转移来源、遍历顺序；空间压缩时还要指出读到的是旧值还是新值。状态升级：把重复递归节点命名为状态；遍历顺序只是在保证依赖状态已经算完。

\textbf{证明和边界。} 证明从最后一步开始：任意最优解的最后一步都在转移枚举中，转移来源已经由更小状态正确计算。边界：起点或终点是障碍、整行被截断、单行障碍。变体：若允许向上或向左走，图会有环，普通填表不再成立。

\noindent\textbf{LeetCode 64 最小路径和。} 模型：状态是到达当前格子的最小代价，最后一步来自上方或左方。推导重点：先初始化第一行和第一列，再处理内部，主转移保持简单。

\textbf{二刷读题。} 读题时先找重复子问题，而不是直接背公式。本题可以压成“状态是到达当前格子的最小代价，最后一步来自上方或左方”，先初始化第一行和第一列，再处理内部，主转移保持简单就是状态之间的依赖关系。先遮住答案，只保留题面，复述候选对象、合法性条件和输出目标。

\textbf{暴力回放。} 慢版本先写递归搜索树：节点表示处理位置、剩余资源和当前选择。只要不同路径反复到达同一个节点，就说明需要把这个节点命名成 DP 状态。二刷时先写慢版本或口述慢版本，用它确认不会漏答案。

\textbf{特例回放。} 拿 [[1,3,1],[1,5,1],[4,2,1]] 手算，到 5 的最小代价来自左侧 4 或上方 4，再加 5；右下角只需要比较上方和左方的已知最小代价。规律是最后一步来自上或左，先初始化第一行第一列，内部统一取 min。把这个特例继续拆开看：先问暴力搜索里哪些不同路径到达同一个后续问题，再给这个后续问题命名为状态。状态必须包含未来决策需要的全部信息，转移只从更小且已经算清的状态来。复盘时把这个特例压成状态字段：状态下标、状态值语义、初始化、转移来源、遍历顺序；空间压缩时还要指出读到的是旧值还是新值；再用边界 单行、单列、权值为零、路径和溢出 反查初始化、更新顺序和退出条件。

\textbf{状态回放。} 手算路线：手算时先写一个很小的 DP 表或记忆化搜索记录。每个格子旁边写它来自哪些更小状态。状态字段：状态下标、状态值语义、初始化、转移来源、遍历顺序；空间压缩时还要指出读到的是旧值还是新值。状态升级：把重复递归节点命名为状态；遍历顺序只是在保证依赖状态已经算完。

\textbf{证明和边界。} 证明从最后一步开始：任意最优解的最后一步都在转移枚举中，转移来源已经由更小状态正确计算。边界：单行、单列、权值为零、路径和溢出。变体：若允许四方向移动，变成最短路而不是普通网格 DP。

\noindent\textbf{LeetCode 70 爬楼梯。} 模型：最后一步来自前一阶或前两阶，是最小的重叠子问题模型。推导重点：滚动变量保存最近两个状态即可，不需要完整数组。

\textbf{二刷读题。} 读题时先找重复子问题，而不是直接背公式。本题可以压成“最后一步来自前一阶或前两阶，是最小的重叠子问题模型”，滚动变量保存最近两个状态即可，不需要完整数组就是状态之间的依赖关系。先遮住答案，只保留题面，复述候选对象、合法性条件和输出目标。

\textbf{暴力回放。} 慢版本先写递归搜索树：节点表示处理位置、剩余资源和当前选择。只要不同路径反复到达同一个节点，就说明需要把这个节点命名成 DP 状态。二刷时先写慢版本或口述慢版本，用它确认不会漏答案。

\textbf{特例回放。} 拿 n=4 手算，最后一步要么从第 3 阶跨 1 阶，要么从第 2 阶跨 2 阶，这两类互斥且覆盖所有走法。于是 ways[4]=ways[3]+ways[2]。规律不是背斐波那契，而是最后一步只有两个来源，状态只需保存前两项。把这个特例继续拆开看：先问暴力搜索里哪些不同路径到达同一个后续问题，再给这个后续问题命名为状态。状态必须包含未来决策需要的全部信息，转移只从更小且已经算清的状态来。复盘时把这个特例压成状态字段：状态下标、状态值语义、初始化、转移来源、遍历顺序；空间压缩时还要指出读到的是旧值还是新值；再用边界 n 为一、n 为二、题目是否允许零阶、结果范围 反查初始化、更新顺序和退出条件。

\textbf{状态回放。} 手算路线：手算时先写一个很小的 DP 表或记忆化搜索记录。每个格子旁边写它来自哪些更小状态。状态字段：状态下标、状态值语义、初始化、转移来源、遍历顺序；空间压缩时还要指出读到的是旧值还是新值。状态升级：把重复递归节点命名为状态；遍历顺序只是在保证依赖状态已经算完。

\textbf{证明和边界。} 证明从最后一步开始：任意最优解的最后一步都在转移枚举中，转移来源已经由更小状态正确计算。边界：n 为一、n 为二、题目是否允许零阶、结果范围。变体：若每次可走一到 k 阶，转移变成固定宽度窗口求和。

\noindent\textbf{LeetCode 72 编辑距离。} 模型：二维状态表示两个前缀之间的最少编辑操作。推导重点：末尾字符相同继承左上；不同则枚举插入、删除、替换。

\textbf{二刷读题。} 读题时先找重复子问题，而不是直接背公式。本题可以压成“二维状态表示两个前缀之间的最少编辑操作”，末尾字符相同继承左上；不同则枚举插入、删除、替换就是状态之间的依赖关系。先遮住答案，只保留题面，复述候选对象、合法性条件和输出目标。

\textbf{暴力回放。} 慢版本先写递归搜索树：节点表示处理位置、剩余资源和当前选择。只要不同路径反复到达同一个节点，就说明需要把这个节点命名成 DP 状态。二刷时先写慢版本或口述慢版本，用它确认不会漏答案。

\textbf{特例回放。} 拿 word1=ab、word2=ac 手算最后一个字符。若 b 改成 c，是 dp[1][1]+1；若删掉 b，是 dp[1][2]+1；若插入 c，是 dp[2][1]+1。字符相等时才直接继承左上。规律是每个格子枚举最后一次编辑操作，而不是凭感觉填三邻居最小。把这个特例继续拆开看：先问暴力搜索里哪些不同路径到达同一个后续问题，再给这个后续问题命名为状态。状态必须包含未来决策需要的全部信息，转移只从更小且已经算清的状态来。复盘时把这个特例压成状态字段：状态下标、状态值语义、初始化、转移来源、遍历顺序；空间压缩时还要指出读到的是旧值还是新值；再用边界 空串、完全相同、完全不同、操作代价不同 反查初始化、更新顺序和退出条件。

\textbf{状态回放。} 手算路线：手算时先写一个很小的 DP 表或记忆化搜索记录。每个格子旁边写它来自哪些更小状态。状态字段：状态下标、状态值语义、初始化、转移来源、遍历顺序；空间压缩时还要指出读到的是旧值还是新值。状态升级：把重复递归节点命名为状态；遍历顺序只是在保证依赖状态已经算完。

\textbf{证明和边界。} 证明从最后一步开始：任意最优解的最后一步都在转移枚举中，转移来源已经由更小状态正确计算。边界：空串、完全相同、完全不同、操作代价不同。变体：若只允许插入删除，问题退化到和最长公共子序列相关。

\noindent\textbf{LeetCode 97 交错字符串。} 模型：目标串前缀由两个来源字符串的前缀交错组成。推导重点：状态 dp[i][j] 表示 s1 前 i 个字符和 s2 前 j 个字符能否组成 s3 前 i 加 j 个字符。

\textbf{二刷读题。} 读题时先找重复子问题，而不是直接背公式。本题可以压成“目标串前缀由两个来源字符串的前缀交错组成”，状态 dp[i][j] 表示 s1 前 i 个字符和 s2 前 j 个字符能否组成 s3 前 i 加 j 个字符就是状态之间的依赖关系。先遮住答案，只保留题面，复述候选对象、合法性条件和输出目标。

\textbf{暴力回放。} 慢版本先写递归搜索树：节点表示处理位置、剩余资源和当前选择。只要不同路径反复到达同一个节点，就说明需要把这个节点命名成 DP 状态。二刷时先写慢版本或口述慢版本，用它确认不会漏答案。

\textbf{特例回放。} 拿 s1=a、s2=b、s3=ab 手算，最后一个字符 b 可以来自 s2；拿 s1=a、s2=a、s3=aa，两个来源都能匹配，不能贪心固定选一个。规律是 dp[i][j] 表示两个前缀能否组成目标前 i+j 个字符，转移同时尝试来自 s1 和 s2 的最后一步。把这个特例继续拆开看：先问暴力搜索里哪些不同路径到达同一个后续问题，再给这个后续问题命名为状态。状态必须包含未来决策需要的全部信息，转移只从更小且已经算清的状态来。复盘时把这个特例压成状态字段：状态下标、状态值语义、初始化、转移来源、遍历顺序；空间压缩时还要指出读到的是旧值还是新值；再用边界 总长度不等、第一行第一列初始化、两个来源字符都能匹配时不能贪心只选一个 反查初始化、更新顺序和退出条件。

\textbf{状态回放。} 手算路线：手算时先写一个很小的 DP 表或记忆化搜索记录。每个格子旁边写它来自哪些更小状态。状态字段：状态下标、状态值语义、初始化、转移来源、遍历顺序；空间压缩时还要指出读到的是旧值还是新值。状态升级：把重复递归节点命名为状态；遍历顺序只是在保证依赖状态已经算完。

\textbf{证明和边界。} 证明从最后一步开始：任意最优解的最后一步都在转移枚举中，转移来源已经由更小状态正确计算。边界：总长度不等、第一行第一列初始化、两个来源字符都能匹配时不能贪心只选一个。变体：若要统计交错方案数，布尔状态要改成计数状态并处理大数或取模。

\noindent\textbf{LeetCode 115 不同的子序列。} 模型：从源串中删除若干字符后形成目标串，本质是两个前缀之间的计数问题。推导重点：状态 dp[i][j] 表示源串前 i 个字符中有多少个子序列等于目标前 j 个字符。

\textbf{二刷读题。} 读题时先找重复子问题，而不是直接背公式。本题可以压成“从源串中删除若干字符后形成目标串，本质是两个前缀之间的计数问题”，状态 dp[i][j] 表示源串前 i 个字符中有多少个子序列等于目标前 j 个字符就是状态之间的依赖关系。先遮住答案，只保留题面，复述候选对象、合法性条件和输出目标。

\textbf{暴力回放。} 慢版本先写递归搜索树：节点表示处理位置、剩余资源和当前选择。只要不同路径反复到达同一个节点，就说明需要把这个节点命名成 DP 状态。二刷时先写慢版本或口述慢版本，用它确认不会漏答案。

\textbf{特例回放。} 拿 s=bab、t=ba 手算。扫描到最后一个 b 时，它不能匹配 t 的 a，只能继承不用它的方案；扫描到 a 时，有两类方案：用这个 a 匹配 t 的末尾，或不用它。规律是字符相等时计数相加，不等时只继承不用源字符的方案。把这个特例继续拆开看：先问暴力搜索里哪些不同路径到达同一个后续问题，再给这个后续问题命名为状态。状态必须包含未来决策需要的全部信息，转移只从更小且已经算清的状态来。复盘时把这个特例压成状态字段：状态下标、状态值语义、初始化、转移来源、遍历顺序；空间压缩时还要指出读到的是旧值还是新值；再用边界 空目标串计数为一、源串为空、重复字符导致计数增长、计数可能超过 int 反查初始化、更新顺序和退出条件。

\textbf{状态回放。} 手算路线：手算时先写一个很小的 DP 表或记忆化搜索记录。每个格子旁边写它来自哪些更小状态。状态字段：状态下标、状态值语义、初始化、转移来源、遍历顺序；空间压缩时还要指出读到的是旧值还是新值。状态升级：把重复递归节点命名为状态；遍历顺序只是在保证依赖状态已经算完。

\textbf{证明和边界。} 证明从最后一步开始：任意最优解的最后一步都在转移枚举中，转移来源已经由更小状态正确计算。边界：空目标串计数为一、源串为空、重复字符导致计数增长、计数可能超过 int。变体：若只问是否存在目标子序列，可以用双指针，不需要二维计数 DP。

\noindent\textbf{LeetCode 121 买卖股票的最佳时机。} 模型：卖出日固定时，最佳买入日是此前最低价格。推导重点：扫描维护历史最低价和最大利润，一次交易无需复杂状态机。

\textbf{二刷读题。} 读题时先找重复子问题，而不是直接背公式。本题可以压成“卖出日固定时，最佳买入日是此前最低价格”，扫描维护历史最低价和最大利润，一次交易无需复杂状态机就是状态之间的依赖关系。先遮住答案，只保留题面，复述候选对象、合法性条件和输出目标。

\textbf{暴力回放。} 慢版本先写递归搜索树：节点表示处理位置、剩余资源和当前选择。只要不同路径反复到达同一个节点，就说明需要把这个节点命名成 DP 状态。二刷时先写慢版本或口述慢版本，用它确认不会漏答案。

\textbf{特例回放。} 拿价格 [7,1,5] 手算。站在 7 不能卖出获利；站在 1 更新最低买入价；站在 5 时利润是 5-1。规律是固定卖出日后，最佳买入日只需要历史最低价，状态不是所有买卖对。把这个特例继续拆开看：先问暴力搜索里哪些不同路径到达同一个后续问题，再给这个后续问题命名为状态。状态必须包含未来决策需要的全部信息，转移只从更小且已经算清的状态来。复盘时把这个特例压成状态字段：状态下标、状态值语义、初始化、转移来源、遍历顺序；空间压缩时还要指出读到的是旧值还是新值；再用边界 单日价格、一直下跌、一直上涨、价格相同 反查初始化、更新顺序和退出条件。

\textbf{状态回放。} 手算路线：手算时先写一个很小的 DP 表或记忆化搜索记录。每个格子旁边写它来自哪些更小状态。状态字段：状态下标、状态值语义、初始化、转移来源、遍历顺序；空间压缩时还要指出读到的是旧值还是新值。状态升级：把重复递归节点命名为状态；遍历顺序只是在保证依赖状态已经算完。

\textbf{证明和边界。} 证明从最后一步开始：任意最优解的最后一步都在转移枚举中，转移来源已经由更小状态正确计算。边界：单日价格、一直下跌、一直上涨、价格相同。变体：多次交易、冷冻期和手续费版本都要扩展成状态机。

\noindent\textbf{LeetCode 123 买卖股票 III。} 模型：最多两次交易，状态需要包含交易次数和持有状态。推导重点：维护 buy1、sell1、buy2、sell2 四个状态，按天更新。

\textbf{二刷读题。} 读题时先找重复子问题，而不是直接背公式。本题可以压成“最多两次交易，状态需要包含交易次数和持有状态”，维护 buy1、sell1、buy2、sell2 四个状态，按天更新就是状态之间的依赖关系。先遮住答案，只保留题面，复述候选对象、合法性条件和输出目标。

\textbf{暴力回放。} 慢版本先写递归搜索树：节点表示处理位置、剩余资源和当前选择。只要不同路径反复到达同一个节点，就说明需要把这个节点命名成 DP 状态。二刷时先写慢版本或口述慢版本，用它确认不会漏答案。

\textbf{特例回放。} 拿 [3,3,5,0,0,3,1,4] 手算，只保存一个最大利润不够，因为第二次买入依赖第一次卖出后的余额。规律是维护 buy1、sell1、buy2、sell2 四个状态，每天都表示到今天为止对应状态的最优值；两次交易不能重叠。把这个特例继续拆开看：先问暴力搜索里哪些不同路径到达同一个后续问题，再给这个后续问题命名为状态。状态必须包含未来决策需要的全部信息，转移只从更小且已经算清的状态来。复盘时把这个特例压成状态字段：状态下标、状态值语义、初始化、转移来源、遍历顺序；空间压缩时还要指出读到的是旧值还是新值；再用边界 价格为空、只够一次交易、两次交易间不能重叠 反查初始化、更新顺序和退出条件。

\textbf{状态回放。} 手算路线：手算时先写一个很小的 DP 表或记忆化搜索记录。每个格子旁边写它来自哪些更小状态。状态字段：状态下标、状态值语义、初始化、转移来源、遍历顺序；空间压缩时还要指出读到的是旧值还是新值。状态升级：把重复递归节点命名为状态；遍历顺序只是在保证依赖状态已经算完。

\textbf{证明和边界。} 证明从最后一步开始：任意最优解的最后一步都在转移枚举中，转移来源已经由更小状态正确计算。边界：价格为空、只够一次交易、两次交易间不能重叠。变体：最多 k 次交易是它的泛化，k 很大时退化为无限次交易。

\noindent\textbf{LeetCode 139 单词拆分。} 模型：状态表示前缀能否由字典单词拼出。推导重点：枚举最后一个单词的起点，左侧前缀可达且子串在字典中则当前可达。

\textbf{二刷读题。} 读题时先找重复子问题，而不是直接背公式。本题可以压成“状态表示前缀能否由字典单词拼出”，枚举最后一个单词的起点，左侧前缀可达且子串在字典中则当前可达就是状态之间的依赖关系。先遮住答案，只保留题面，复述候选对象、合法性条件和输出目标。

\textbf{暴力回放。} 慢版本先写递归搜索树：节点表示处理位置、剩余资源和当前选择。只要不同路径反复到达同一个节点，就说明需要把这个节点命名成 DP 状态。二刷时先写慢版本或口述慢版本，用它确认不会漏答案。

\textbf{特例回放。} 拿 s=leetcode、字典 [leet,code] 手算，dp[4] 因为 leet 可拆而为真，dp[8] 又因为 dp[4] 且 code 在字典里为真。规律是 dp[i] 表示前 i 个字符能否拆分，转移枚举最后一个单词起点。把这个特例继续拆开看：先问暴力搜索里哪些不同路径到达同一个后续问题，再给这个后续问题命名为状态。状态必须包含未来决策需要的全部信息，转移只从更小且已经算清的状态来。复盘时把这个特例压成状态字段：状态下标、状态值语义、初始化、转移来源、遍历顺序；空间压缩时还要指出读到的是旧值还是新值；再用边界 空串、重复单词、长单词、substr 复制成本 反查初始化、更新顺序和退出条件。

\textbf{状态回放。} 手算路线：手算时先写一个很小的 DP 表或记忆化搜索记录。每个格子旁边写它来自哪些更小状态。状态字段：状态下标、状态值语义、初始化、转移来源、遍历顺序；空间压缩时还要指出读到的是旧值还是新值。状态升级：把重复递归节点命名为状态；遍历顺序只是在保证依赖状态已经算完。

\textbf{证明和边界。} 证明从最后一步开始：任意最优解的最后一步都在转移枚举中，转移来源已经由更小状态正确计算。边界：空串、重复单词、长单词、substr 复制成本。变体：若要输出所有句子，需要 DFS 加记忆化或前驱列表。

\noindent\textbf{LeetCode 152 乘积最大子数组。} 模型：负数会让最大乘积和最小乘积互换。推导重点：同时维护以当前位置结尾的最大和最小乘积。

\textbf{二刷读题。} 读题时先找重复子问题，而不是直接背公式。本题可以压成“负数会让最大乘积和最小乘积互换”，同时维护以当前位置结尾的最大和最小乘积就是状态之间的依赖关系。先遮住答案，只保留题面，复述候选对象、合法性条件和输出目标。

\textbf{暴力回放。} 慢版本先写递归搜索树：节点表示处理位置、剩余资源和当前选择。只要不同路径反复到达同一个节点，就说明需要把这个节点命名成 DP 状态。二刷时先写慢版本或口述慢版本，用它确认不会漏答案。

\textbf{特例回放。} 拿 [-2,3,-4] 手算。到 3 时最大乘积是 3、最小是 -6；再乘 -4 后，原来的最小 -6 反而变成最大 24。规律是负数会交换最大和最小，状态必须同时保存以当前位置结尾的最大乘积和最小乘积。把这个特例继续拆开看：先问暴力搜索里哪些不同路径到达同一个后续问题，再给这个后续问题命名为状态。状态必须包含未来决策需要的全部信息，转移只从更小且已经算清的状态来。复盘时把这个特例压成状态字段：状态下标、状态值语义、初始化、转移来源、遍历顺序；空间压缩时还要指出读到的是旧值还是新值；再用边界 零、负数个数奇偶、单元素、乘积溢出 反查初始化、更新顺序和退出条件。

\textbf{状态回放。} 手算路线：手算时先写一个很小的 DP 表或记忆化搜索记录。每个格子旁边写它来自哪些更小状态。状态字段：状态下标、状态值语义、初始化、转移来源、遍历顺序；空间压缩时还要指出读到的是旧值还是新值。状态升级：把重复递归节点命名为状态；遍历顺序只是在保证依赖状态已经算完。

\textbf{证明和边界。} 证明从最后一步开始：任意最优解的最后一步都在转移枚举中，转移来源已经由更小状态正确计算。边界：零、负数个数奇偶、单元素、乘积溢出。变体：如果要求返回区间，同步记录状态来源。

\noindent\textbf{LeetCode 188 买卖股票 IV。} 模型：最多 k 次交易让状态必须同时包含交易次数和是否持有股票。推导重点：维护每个交易次数下的 buy 和 sell 状态；当 k 足够大时可退化为无限次交易。

\textbf{二刷读题。} 读题时先找重复子问题，而不是直接背公式。本题可以压成“最多 k 次交易让状态必须同时包含交易次数和是否持有股票”，维护每个交易次数下的 buy 和 sell 状态；当 k 足够大时可退化为无限次交易就是状态之间的依赖关系。先遮住答案，只保留题面，复述候选对象、合法性条件和输出目标。

\textbf{暴力回放。} 慢版本先写递归搜索树：节点表示处理位置、剩余资源和当前选择。只要不同路径反复到达同一个节点，就说明需要把这个节点命名成 DP 状态。二刷时先写慢版本或口述慢版本，用它确认不会漏答案。

\textbf{特例回放。} 拿 k=2、价格 [3,2,6] 手算。某天结束时只说最大利润不够，因为未来能否卖出取决于是否持股，也取决于已经完成几次交易。规律是状态必须包含交易次数和持有状态；k 很大时才退化为无限次交易。把这个特例继续拆开看：先问暴力搜索里哪些不同路径到达同一个后续问题，再给这个后续问题命名为状态。状态必须包含未来决策需要的全部信息，转移只从更小且已经算清的状态来。复盘时把这个特例压成状态字段：状态下标、状态值语义、初始化、转移来源、遍历顺序；空间压缩时还要指出读到的是旧值还是新值；再用边界 k 为零、价格天数很少、k 大于等于 n/2、更新顺序造成同一天状态污染 反查初始化、更新顺序和退出条件。

\textbf{状态回放。} 手算路线：手算时先写一个很小的 DP 表或记忆化搜索记录。每个格子旁边写它来自哪些更小状态。状态字段：状态下标、状态值语义、初始化、转移来源、遍历顺序；空间压缩时还要指出读到的是旧值还是新值。状态升级：把重复递归节点命名为状态；遍历顺序只是在保证依赖状态已经算完。

\textbf{证明和边界。} 证明从最后一步开始：任意最优解的最后一步都在转移枚举中，转移来源已经由更小状态正确计算。边界：k 为零、价格天数很少、k 大于等于 n/2、更新顺序造成同一天状态污染。变体：手续费、冷冻期和最多两次交易都是同一状态机框架的不同约束。

\noindent\textbf{LeetCode 198 打家劫舍。} 模型：每间房只有偷和不偷两种结局，偷当前就不能偷前一间。推导重点：滚动保存前一状态和前两状态。

\textbf{二刷读题。} 读题时先找重复子问题，而不是直接背公式。本题可以压成“每间房只有偷和不偷两种结局，偷当前就不能偷前一间”，滚动保存前一状态和前两状态就是状态之间的依赖关系。先遮住答案，只保留题面，复述候选对象、合法性条件和输出目标。

\textbf{暴力回放。} 慢版本先写递归搜索树：节点表示处理位置、剩余资源和当前选择。只要不同路径反复到达同一个节点，就说明需要把这个节点命名成 DP 状态。二刷时先写慢版本或口述慢版本，用它确认不会漏答案。

\textbf{特例回放。} 拿 [2,7,9] 手算，到房屋 9 时只有两类合法结局：不偷它，答案沿用前一间；偷它，就只能接前两间的最优。规律是 dp[i]=max(dp[i-1], dp[i-2]+nums[i])，相邻限制来自题面而不是公式本身。把这个特例继续拆开看：先问暴力搜索里哪些不同路径到达同一个后续问题，再给这个后续问题命名为状态。状态必须包含未来决策需要的全部信息，转移只从更小且已经算清的状态来。复盘时把这个特例压成状态字段：状态下标、状态值语义、初始化、转移来源、遍历顺序；空间压缩时还要指出读到的是旧值还是新值；再用边界 空数组、单间、全零、金额很大 反查初始化、更新顺序和退出条件。

\textbf{状态回放。} 手算路线：手算时先写一个很小的 DP 表或记忆化搜索记录。每个格子旁边写它来自哪些更小状态。状态字段：状态下标、状态值语义、初始化、转移来源、遍历顺序；空间压缩时还要指出读到的是旧值还是新值。状态升级：把重复递归节点命名为状态；遍历顺序只是在保证依赖状态已经算完。

\textbf{证明和边界。} 证明从最后一步开始：任意最优解的最后一步都在转移枚举中，转移来源已经由更小状态正确计算。边界：空数组、单间、全零、金额很大。变体：环形房屋要拆成不含首和不含尾两个线性问题。

\noindent\textbf{LeetCode 221 最大正方形。} 模型：以当前位置为右下角的最大正方形由上、左、左上三个状态限制。推导重点：当前位置为一时取三者最小加一，否则为零。

\textbf{二刷读题。} 读题时先找重复子问题，而不是直接背公式。本题可以压成“以当前位置为右下角的最大正方形由上、左、左上三个状态限制”，当前位置为一时取三者最小加一，否则为零就是状态之间的依赖关系。先遮住答案，只保留题面，复述候选对象、合法性条件和输出目标。

\textbf{暴力回放。} 慢版本先写递归搜索树：节点表示处理位置、剩余资源和当前选择。只要不同路径反复到达同一个节点，就说明需要把这个节点命名成 DP 状态。二刷时先写慢版本或口述慢版本，用它确认不会漏答案。

\textbf{特例回放。} 拿 2x2 全 1 方块手算，右下角能形成边长 2，因为上、左、左上三个格子都至少能形成边长 1。若其中任一方向为 0，只能形成边长 1。规律是 dp[row][col]=min(上,左,左上)+1，状态值是以当前格为右下角的最大正方形边长。把这个特例继续拆开看：先问暴力搜索里哪些不同路径到达同一个后续问题，再给这个后续问题命名为状态。状态必须包含未来决策需要的全部信息，转移只从更小且已经算清的状态来。复盘时把这个特例压成状态字段：状态下标、状态值语义、初始化、转移来源、遍历顺序；空间压缩时还要指出读到的是旧值还是新值；再用边界 全零、全一、单行单列、字符一和数字一不要混淆 反查初始化、更新顺序和退出条件。

\textbf{状态回放。} 手算路线：手算时先写一个很小的 DP 表或记忆化搜索记录。每个格子旁边写它来自哪些更小状态。状态字段：状态下标、状态值语义、初始化、转移来源、遍历顺序；空间压缩时还要指出读到的是旧值还是新值。状态升级：把重复递归节点命名为状态；遍历顺序只是在保证依赖状态已经算完。

\textbf{证明和边界。} 证明从最后一步开始：任意最优解的最后一步都在转移枚举中，转移来源已经由更小状态正确计算。边界：全零、全一、单行单列、字符一和数字一不要混淆。变体：最大矩形则需要按行构造柱状图再用单调栈。

\noindent\textbf{LeetCode 300 最长递增子序列。} 模型：二次 DP 固定结尾，二分优化维护每个长度的最小结尾。推导重点：tails 不是实际答案序列，而是未来扩展潜力表。

\textbf{二刷读题。} 读题时先找重复子问题，而不是直接背公式。本题可以压成“二次 DP 固定结尾，二分优化维护每个长度的最小结尾”，tails 不是实际答案序列，而是未来扩展潜力表就是状态之间的依赖关系。先遮住答案，只保留题面，复述候选对象、合法性条件和输出目标。

\textbf{暴力回放。} 慢版本先写递归搜索树：节点表示处理位置、剩余资源和当前选择。只要不同路径反复到达同一个节点，就说明需要把这个节点命名成 DP 状态。二刷时先写慢版本或口述慢版本，用它确认不会漏答案。

\textbf{特例回放。} 拿 [10,9,2,5,3] 手算。二次 DP 固定每个位置为结尾，5 可接在 2 后；优化时 tails[长度] 保存该长度递增子序列的最小尾值，尾值越小未来越容易接。规律是二分替换 tails，不代表真实序列完整内容，而是维护未来潜力。把这个特例继续拆开看：先问暴力搜索里哪些不同路径到达同一个后续问题，再给这个后续问题命名为状态。状态必须包含未来决策需要的全部信息，转移只从更小且已经算清的状态来。复盘时把这个特例压成状态字段：状态下标、状态值语义、初始化、转移来源、遍历顺序；空间压缩时还要指出读到的是旧值还是新值；再用边界 重复元素、严格递增和非递减区别、空数组 反查初始化、更新顺序和退出条件。

\textbf{状态回放。} 手算路线：手算时先写一个很小的 DP 表或记忆化搜索记录。每个格子旁边写它来自哪些更小状态。状态字段：状态下标、状态值语义、初始化、转移来源、遍历顺序；空间压缩时还要指出读到的是旧值还是新值。状态升级：把重复递归节点命名为状态；遍历顺序只是在保证依赖状态已经算完。

\textbf{证明和边界。} 证明从最后一步开始：任意最优解的最后一步都在转移枚举中，转移来源已经由更小状态正确计算。边界：重复元素、严格递增和非递减区别、空数组。变体：若要还原路径，需要额外前驱数组，不能只看 tails。

\noindent\textbf{LeetCode 309 最佳买卖股票含冷冻期。} 模型：冷冻期让买入来源受昨天卖出状态限制。推导重点：状态机维护持有、刚卖出、休息三种日终状态。

\textbf{二刷读题。} 读题时先找重复子问题，而不是直接背公式。本题可以压成“冷冻期让买入来源受昨天卖出状态限制”，状态机维护持有、刚卖出、休息三种日终状态就是状态之间的依赖关系。先遮住答案，只保留题面，复述候选对象、合法性条件和输出目标。

\textbf{暴力回放。} 慢版本先写递归搜索树：节点表示处理位置、剩余资源和当前选择。只要不同路径反复到达同一个节点，就说明需要把这个节点命名成 DP 状态。二刷时先写慢版本或口述慢版本，用它确认不会漏答案。

\textbf{特例回放。} 拿 [1,2,3,0,2] 手算，第 2 天卖出后第 3 天不能立刻买入，因为有冷冻期；状态必须区分持股、刚卖出、空仓可买。规律是冷冻期把普通买卖股票拆成多个日末状态，转移不能跨过限制。把这个特例继续拆开看：先问暴力搜索里哪些不同路径到达同一个后续问题，再给这个后续问题命名为状态。状态必须包含未来决策需要的全部信息，转移只从更小且已经算清的状态来。复盘时把这个特例压成状态字段：状态下标、状态值语义、初始化、转移来源、遍历顺序；空间压缩时还要指出读到的是旧值还是新值；再用边界 空价格、一天、连续上涨、卖出后不能次日买入 反查初始化、更新顺序和退出条件。

\textbf{状态回放。} 手算路线：手算时先写一个很小的 DP 表或记忆化搜索记录。每个格子旁边写它来自哪些更小状态。状态字段：状态下标、状态值语义、初始化、转移来源、遍历顺序；空间压缩时还要指出读到的是旧值还是新值。状态升级：把重复递归节点命名为状态；遍历顺序只是在保证依赖状态已经算完。

\textbf{证明和边界。} 证明从最后一步开始：任意最优解的最后一步都在转移枚举中，转移来源已经由更小状态正确计算。边界：空价格、一天、连续上涨、卖出后不能次日买入。变体：手续费版本可以在买入或卖出转移中扣费。

\noindent\textbf{LeetCode 312 戳气球。} 模型：先戳哪个会让邻居持续变化，改成最后戳某个气球后区间边界才稳定。推导重点：区间 DP 定义开区间内全部戳完的最大收益，枚举最后一个被戳的气球作为切分点。

\textbf{二刷读题。} 读题时先找重复子问题，而不是直接背公式。本题可以压成“先戳哪个会让邻居持续变化，改成最后戳某个气球后区间边界才稳定”，区间 DP 定义开区间内全部戳完的最大收益，枚举最后一个被戳的气球作为切分点就是状态之间的依赖关系。先遮住答案，只保留题面，复述候选对象、合法性条件和输出目标。

\textbf{暴力回放。} 慢版本先写递归搜索树：节点表示处理位置、剩余资源和当前选择。只要不同路径反复到达同一个节点，就说明需要把这个节点命名成 DP 状态。二刷时先写慢版本或口述慢版本，用它确认不会漏答案。

\textbf{特例回放。} 拿 [3,1,5] 手算，若先戳 1，左右邻居会变；若改问某个区间里最后戳谁，最后一刻左右边界固定，收益可拆成左右两个子区间。规律是区间 DP 枚举最后戳的 mid，而不是第一个戳的气球。把这个特例继续拆开看：先问暴力搜索里哪些不同路径到达同一个后续问题，再给这个后续问题命名为状态。状态必须包含未来决策需要的全部信息，转移只从更小且已经算清的状态来。复盘时把这个特例压成状态字段：状态下标、状态值语义、初始化、转移来源、遍历顺序；空间压缩时还要指出读到的是旧值还是新值；再用边界 补一的虚拟边界、区间长度遍历顺序、空区间收益为零、乘积可能较大 反查初始化、更新顺序和退出条件。

\textbf{状态回放。} 手算路线：手算时先写一个很小的 DP 表或记忆化搜索记录。每个格子旁边写它来自哪些更小状态。状态字段：状态下标、状态值语义、初始化、转移来源、遍历顺序；空间压缩时还要指出读到的是旧值还是新值。状态升级：把重复递归节点命名为状态；遍历顺序只是在保证依赖状态已经算完。

\textbf{证明和边界。} 证明从最后一步开始：任意最优解的最后一步都在转移枚举中，转移来源已经由更小状态正确计算。边界：补一的虚拟边界、区间长度遍历顺序、空区间收益为零、乘积可能较大。变体：很多区间 DP 都要换成最后一步思考，避免中间操作改变边界。

\noindent\textbf{LeetCode 329 矩阵中的最长递增路径。} 模型：每个格子的答案是从它出发能走出的最长严格递增路径。推导重点：把严格递增关系看成无环状态图，用记忆化 DFS 或按值排序的 DAG DP 复用后续格子结果。

\textbf{二刷读题。} 读题时先找重复子问题，而不是直接背公式。本题可以压成“每个格子的答案是从它出发能走出的最长严格递增路径”，把严格递增关系看成无环状态图，用记忆化 DFS 或按值排序的 DAG DP 复用后续格子结果就是状态之间的依赖关系。先遮住答案，只保留题面，复述候选对象、合法性条件和输出目标。

\textbf{暴力回放。} 慢版本先写递归搜索树：节点表示处理位置、剩余资源和当前选择。只要不同路径反复到达同一个节点，就说明需要把这个节点命名成 DP 状态。二刷时先写慢版本或口述慢版本，用它确认不会漏答案。

\textbf{特例回放。} 拿矩阵中 1->2->3 的小路径手算，从 1 出发会用到从 2 出发的答案；如果另一个格子也走到 2，后缀路径被重复计算。严格递增保证不会成环。规律是每个格子记忆化保存从这里出发的最长路径。把这个特例继续拆开看：先问暴力搜索里哪些不同路径到达同一个后续问题，再给这个后续问题命名为状态。状态必须包含未来决策需要的全部信息，转移只从更小且已经算清的状态来。复盘时把这个特例压成状态字段：状态下标、状态值语义、初始化、转移来源、遍历顺序；空间压缩时还要指出读到的是旧值还是新值；再用边界 平台相等不能走、四方向越界、递归深度、允许不下降时会产生环 反查初始化、更新顺序和退出条件。

\textbf{状态回放。} 手算路线：手算时先写一个很小的 DP 表或记忆化搜索记录。每个格子旁边写它来自哪些更小状态。状态字段：状态下标、状态值语义、初始化、转移来源、遍历顺序；空间压缩时还要指出读到的是旧值还是新值。状态升级：把重复递归节点命名为状态；遍历顺序只是在保证依赖状态已经算完。

\textbf{证明和边界。} 证明从最后一步开始：任意最优解的最后一步都在转移枚举中，转移来源已经由更小状态正确计算。边界：平台相等不能走、四方向越界、递归深度、允许不下降时会产生环。变体：若改成求最短路径，严格递增的 DP 模型不再直接适用。

\noindent\textbf{LeetCode 714 买卖股票含手续费。} 模型：手续费改变交易收益，仍可用持有和不持有状态。推导重点：卖出时扣费或买入时扣费都可以，但不要重复扣。

\textbf{二刷读题。} 读题时先找重复子问题，而不是直接背公式。本题可以压成“手续费改变交易收益，仍可用持有和不持有状态”，卖出时扣费或买入时扣费都可以，但不要重复扣就是状态之间的依赖关系。先遮住答案，只保留题面，复述候选对象、合法性条件和输出目标。

\textbf{暴力回放。} 慢版本先写递归搜索树：节点表示处理位置、剩余资源和当前选择。只要不同路径反复到达同一个节点，就说明需要把这个节点命名成 DP 状态。二刷时先写慢版本或口述慢版本，用它确认不会漏答案。

\textbf{特例回放。} 拿 [1,3,2,8,4,9]、fee=2 手算，8 卖出时净赚 5；手续费让频繁小波动不一定值得交易。规律是每天维护 hold 和 cash，卖出时扣手续费，买入时从 cash 转到 hold。把这个特例继续拆开看：先问暴力搜索里哪些不同路径到达同一个后续问题，再给这个后续问题命名为状态。状态必须包含未来决策需要的全部信息，转移只从更小且已经算清的状态来。复盘时把这个特例压成状态字段：状态下标、状态值语义、初始化、转移来源、遍历顺序；空间压缩时还要指出读到的是旧值还是新值；再用边界 手续费为零、价格震荡、小利润不足手续费 反查初始化、更新顺序和退出条件。

\textbf{状态回放。} 手算路线：手算时先写一个很小的 DP 表或记忆化搜索记录。每个格子旁边写它来自哪些更小状态。状态字段：状态下标、状态值语义、初始化、转移来源、遍历顺序；空间压缩时还要指出读到的是旧值还是新值。状态升级：把重复递归节点命名为状态；遍历顺序只是在保证依赖状态已经算完。

\textbf{证明和边界。} 证明从最后一步开始：任意最优解的最后一步都在转移枚举中，转移来源已经由更小状态正确计算。边界：手续费为零、价格震荡、小利润不足手续费。变体：这是状态机比局部贪心更稳的例子。

\noindent\textbf{LeetCode 1143 最长公共子序列。} 模型：状态表示两个前缀的最长公共子序列长度。推导重点：末尾相同取左上加一，不同取上方和左方最大。

\textbf{二刷读题。} 读题时先找重复子问题，而不是直接背公式。本题可以压成“状态表示两个前缀的最长公共子序列长度”，末尾相同取左上加一，不同取上方和左方最大就是状态之间的依赖关系。先遮住答案，只保留题面，复述候选对象、合法性条件和输出目标。

\textbf{暴力回放。} 慢版本先写递归搜索树：节点表示处理位置、剩余资源和当前选择。只要不同路径反复到达同一个节点，就说明需要把这个节点命名成 DP 状态。二刷时先写慢版本或口述慢版本，用它确认不会漏答案。

\textbf{特例回放。} 拿 abcde 和 ace 手算，末尾 e 相等时答案来自两个前缀 abcd 和 ac 加一；若末尾不同，就取删掉一侧末尾后的较大值。规律是 dp[i][j] 表示两个前缀的 LCS 长度，转移只看左、上、左上。把这个特例继续拆开看：先问暴力搜索里哪些不同路径到达同一个后续问题，再给这个后续问题命名为状态。状态必须包含未来决策需要的全部信息，转移只从更小且已经算清的状态来。复盘时把这个特例压成状态字段：状态下标、状态值语义、初始化、转移来源、遍历顺序；空间压缩时还要指出读到的是旧值还是新值；再用边界 空串、完全相同、无公共字符、子序列不是子串 反查初始化、更新顺序和退出条件。

\textbf{状态回放。} 手算路线：手算时先写一个很小的 DP 表或记忆化搜索记录。每个格子旁边写它来自哪些更小状态。状态字段：状态下标、状态值语义、初始化、转移来源、遍历顺序；空间压缩时还要指出读到的是旧值还是新值。状态升级：把重复递归节点命名为状态；遍历顺序只是在保证依赖状态已经算完。

\textbf{证明和边界。} 证明从最后一步开始：任意最优解的最后一步都在转移枚举中，转移来源已经由更小状态正确计算。边界：空串、完全相同、无公共字符、子序列不是子串。变体：若要求还原序列，需要从表右下角反向追踪选择。

\noindent\textbf{LeetCode 1235 规划兼职工作。} 模型：每份工作有开始、结束和收益，选择不冲突工作使收益最大。推导重点：按结束时间排序，dp[i] 在不选当前和选当前加上兼容前驱之间取最大，前驱用二分找。

\textbf{二刷读题。} 读题时先找重复子问题，而不是直接背公式。本题可以压成“每份工作有开始、结束和收益，选择不冲突工作使收益最大”，按结束时间排序，dp[i] 在不选当前和选当前加上兼容前驱之间取最大，前驱用二分找就是状态之间的依赖关系。先遮住答案，只保留题面，复述候选对象、合法性条件和输出目标。

\textbf{暴力回放。} 慢版本先写递归搜索树：节点表示处理位置、剩余资源和当前选择。只要不同路径反复到达同一个节点，就说明需要把这个节点命名成 DP 状态。二刷时先写慢版本或口述慢版本，用它确认不会漏答案。

\textbf{特例回放。} 拿 jobs (1,3,50),(2,4,10),(3,5,40),(3,6,70) 手算，选 (1,3,50) 后还能接 (3,6,70)，总 120。规律是按结束时间排序，dp[i] 在不选当前和选当前加上前一个不冲突工作之间取最大。把这个特例继续拆开看：先问暴力搜索里哪些不同路径到达同一个后续问题，再给这个后续问题命名为状态。状态必须包含未来决策需要的全部信息，转移只从更小且已经算清的状态来。复盘时把这个特例压成状态字段：状态下标、状态值语义、初始化、转移来源、遍历顺序；空间压缩时还要指出读到的是旧值还是新值；再用边界 结束等于开始是否兼容、收益相同、工作排序、空列表 反查初始化、更新顺序和退出条件。

\textbf{状态回放。} 手算路线：手算时先写一个很小的 DP 表或记忆化搜索记录。每个格子旁边写它来自哪些更小状态。状态字段：状态下标、状态值语义、初始化、转移来源、遍历顺序；空间压缩时还要指出读到的是旧值还是新值。状态升级：把重复递归节点命名为状态；遍历顺序只是在保证依赖状态已经算完。

\textbf{证明和边界。} 证明从最后一步开始：任意最优解的最后一步都在转移枚举中，转移来源已经由更小状态正确计算。边界：结束等于开始是否兼容、收益相同、工作排序、空列表。变体：这是排序、二分和 DP 的组合，不是按收益贪心。

\subsection{自测标准}

二刷时不要只确认自己见过这道题。请遮住答案，先讲题面模型，再讲暴力基线和瓶颈，然后说出优化结构保存了什么、不变量如何排除错误候选、边界实验如何打破错误实现。

\section{背包模型：二刷复盘卡}
这一大节只围绕一个题型复盘，不把题号直接铺成目录。阅读时先看复盘目标，再读核心题卡，最后用自测标准检查自己是否能独立重讲。

\subsection{复盘目标}

追问通常会问倒序和正序的区别、组合和排列的区别、为什么复杂度是伪多项式。请准备一个循环方向写反的最小反例。

\subsection{核心题卡}

\noindent\textbf{LeetCode 279 完全平方数。} 模型：完全平方数可以看成无限使用的物品，目标是凑出 n 的最少数量。推导重点：完全背包最少物品数，或把数字看成图节点做 BFS。

\textbf{二刷读题。} 读题时先找阶段、容量、物品使用次数和状态值类型。本题可以压成“完全平方数可以看成无限使用的物品，目标是凑出 n 的最少数量”，循环方向由这些条件决定。先遮住答案，只保留题面，复述候选对象、合法性条件和输出目标。

\textbf{暴力回放。} 慢版本枚举每个物品选不选、选几次，并记录阶段和剩余容量。相同阶段和容量反复出现时，就可以把指数搜索压成表。二刷时先写慢版本或口述慢版本，用它确认不会漏答案。

\textbf{特例回放。} 拿 n=12 手算，12 可以由 4+4+4 组成，不能只贪心选 9 后剩 3。规律是完全背包，dp[x] 表示和为 x 的最少平方数，枚举每个平方数时允许重复使用。把这个特例继续拆开看：阶段是物品，容量是资源，状态值是可达、最值还是计数。一个物品能否在同一轮再次使用，直接决定容量循环方向；组合和排列也由循环顺序表达。复盘时把这个特例压成状态字段：物品阶段、容量、状态值、循环方向、取模或不可达哨兵；组合和排列必须用不同循环顺序表达；再用边界 n 为零或一、完全平方数本身、较大 n 反查初始化、更新顺序和退出条件。

\textbf{状态回放。} 手算路线：手算时用一个容量很小的例子比较正序和倒序。只要循环方向改变后同一物品能否重复使用发生变化，就说明方向不是细节。状态字段：物品阶段、容量、状态值、循环方向、取模或不可达哨兵；组合和排列必须用不同循环顺序表达。状态升级：把选择树压成阶段和容量；物品是否可重复使用决定一维表的遍历方向。

\textbf{证明和边界。} 证明从当前物品使用次数开始：0-1、完全、计数和最值的循环语义必须和题意一致。边界：n 为零或一、完全平方数本身、较大 n。变体：数学四平方定理可给更强解法，但面试 DP 更通用。

\noindent\textbf{LeetCode 322 零钱兑换。} 模型：金额是容量，硬币可无限使用，目标是最少硬币数。推导重点：dp[value] 从 value 减 coin 的已知状态转移。

\textbf{二刷读题。} 读题时先找阶段、容量、物品使用次数和状态值类型。本题可以压成“金额是容量，硬币可无限使用，目标是最少硬币数”，循环方向由这些条件决定。先遮住答案，只保留题面，复述候选对象、合法性条件和输出目标。

\textbf{暴力回放。} 慢版本枚举每个物品选不选、选几次，并记录阶段和剩余容量。相同阶段和容量反复出现时，就可以把指数搜索压成表。二刷时先写慢版本或口述慢版本，用它确认不会漏答案。

\textbf{特例回放。} 拿 coins=[1,2,5]、amount=11 手算，最后一步可以用 1、2 或 5，若用 5，则来源是 dp[6]+1。规律是 dp[x] 表示凑成 x 的最少硬币数，完全背包允许同一硬币重复使用；不可达状态不能参与加一。把这个特例继续拆开看：阶段是物品，容量是资源，状态值是可达、最值还是计数。一个物品能否在同一轮再次使用，直接决定容量循环方向；组合和排列也由循环顺序表达。复盘时把这个特例压成状态字段：物品阶段、容量、状态值、循环方向、取模或不可达哨兵；组合和排列必须用不同循环顺序表达；再用边界 无法凑出、amount 为零、硬币面额大于目标、哨兵溢出 反查初始化、更新顺序和退出条件。

\textbf{状态回放。} 手算路线：手算时用一个容量很小的例子比较正序和倒序。只要循环方向改变后同一物品能否重复使用发生变化，就说明方向不是细节。状态字段：物品阶段、容量、状态值、循环方向、取模或不可达哨兵；组合和排列必须用不同循环顺序表达。状态升级：把选择树压成阶段和容量；物品是否可重复使用决定一维表的遍历方向。

\textbf{证明和边界。} 证明从当前物品使用次数开始：0-1、完全、计数和最值的循环语义必须和题意一致。边界：无法凑出、amount 为零、硬币面额大于目标、哨兵溢出。变体：零钱兑换 II 问组合数，循环语义不同。

\noindent\textbf{LeetCode 377 组合总和 IV。} 模型：题目统计排列数，选择顺序不同算不同方案。推导重点：外层遍历容量，内层遍历数字，让每个容量的方案按最后一个数累加。

\textbf{二刷读题。} 读题时先找阶段、容量、物品使用次数和状态值类型。本题可以压成“题目统计排列数，选择顺序不同算不同方案”，循环方向由这些条件决定。先遮住答案，只保留题面，复述候选对象、合法性条件和输出目标。

\textbf{暴力回放。} 慢版本枚举每个物品选不选、选几次，并记录阶段和剩余容量。相同阶段和容量反复出现时，就可以把指数搜索压成表。二刷时先写慢版本或口述慢版本，用它确认不会漏答案。

\textbf{特例回放。} 拿 nums=[1,2,3]、target=4 手算，1+3 和 3+1 是不同排列，所以 dp[4] 要按最后一步来自 1、2、3 分别累加。规律是组合总和 IV 计数排列，容量在外层、数字在内层，循环顺序不能写成组合题。把这个特例继续拆开看：阶段是物品，容量是资源，状态值是可达、最值还是计数。一个物品能否在同一轮再次使用，直接决定容量循环方向；组合和排列也由循环顺序表达。复盘时把这个特例压成状态字段：物品阶段、容量、状态值、循环方向、取模或不可达哨兵；组合和排列必须用不同循环顺序表达；再用边界 target 为零、数字大于目标、方案数溢出、循环顺序 反查初始化、更新顺序和退出条件。

\textbf{状态回放。} 手算路线：手算时用一个容量很小的例子比较正序和倒序。只要循环方向改变后同一物品能否重复使用发生变化，就说明方向不是细节。状态字段：物品阶段、容量、状态值、循环方向、取模或不可达哨兵；组合和排列必须用不同循环顺序表达。状态升级：把选择树压成阶段和容量；物品是否可重复使用决定一维表的遍历方向。

\textbf{证明和边界。} 证明从当前物品使用次数开始：0-1、完全、计数和最值的循环语义必须和题意一致。边界：target 为零、数字大于目标、方案数溢出、循环顺序。变体：若顺序不同不算新方案，就变成完全背包组合数，循环顺序要交换。

\noindent\textbf{LeetCode 416 分割等和子集。} 模型：总和一半是背包容量，每个数只能使用一次。推导重点：0-1 背包可达性，一维容量必须倒序。

\textbf{二刷读题。} 读题时先找阶段、容量、物品使用次数和状态值类型。本题可以压成“总和一半是背包容量，每个数只能使用一次”，循环方向由这些条件决定。先遮住答案，只保留题面，复述候选对象、合法性条件和输出目标。

\textbf{暴力回放。} 慢版本枚举每个物品选不选、选几次，并记录阶段和剩余容量。相同阶段和容量反复出现时，就可以把指数搜索压成表。二刷时先写慢版本或口述慢版本，用它确认不会漏答案。

\textbf{特例回放。} 拿 [1,5,11,5] 手算，总和 22，目标容量 11；可以用 11 单独凑成，也可以 5+5+1。规律是分割等和转成 0-1 背包是否能凑到 sum/2，容量必须倒序避免重复使用同一元素。把这个特例继续拆开看：阶段是物品，容量是资源，状态值是可达、最值还是计数。一个物品能否在同一轮再次使用，直接决定容量循环方向；组合和排列也由循环顺序表达。复盘时把这个特例压成状态字段：物品阶段、容量、状态值、循环方向、取模或不可达哨兵；组合和排列必须用不同循环顺序表达；再用边界 总和为奇数、目标为零、数值较大、复杂度依赖 target 反查初始化、更新顺序和退出条件。

\textbf{状态回放。} 手算路线：手算时用一个容量很小的例子比较正序和倒序。只要循环方向改变后同一物品能否重复使用发生变化，就说明方向不是细节。状态字段：物品阶段、容量、状态值、循环方向、取模或不可达哨兵；组合和排列必须用不同循环顺序表达。状态升级：把选择树压成阶段和容量；物品是否可重复使用决定一维表的遍历方向。

\textbf{证明和边界。} 证明从当前物品使用次数开始：0-1、完全、计数和最值的循环语义必须和题意一致。边界：总和为奇数、目标为零、数值较大、复杂度依赖 target。变体：负数会破坏普通容量模型，需要重新建模。

\noindent\textbf{LeetCode 474 一和零。} 模型：每个字符串消耗零和一两种容量，每个字符串最多选一次。推导重点：二维 0-1 背包，两个容量都要倒序遍历。

\textbf{二刷读题。} 读题时先找阶段、容量、物品使用次数和状态值类型。本题可以压成“每个字符串消耗零和一两种容量，每个字符串最多选一次”，循环方向由这些条件决定。先遮住答案，只保留题面，复述候选对象、合法性条件和输出目标。

\textbf{暴力回放。} 慢版本枚举每个物品选不选、选几次，并记录阶段和剩余容量。相同阶段和容量反复出现时，就可以把指数搜索压成表。二刷时先写慢版本或口述慢版本，用它确认不会漏答案。

\textbf{特例回放。} 拿 strs=[10,0,1]、m=1、n=1 手算，字符串 10 消耗一个 0 和一个 1，单独就得到长度 1；0 和 1 也能组合成长度 2。规律是二维 0-1 背包，容量是 0 的数量和 1 的数量，两维都要倒序。把这个特例继续拆开看：阶段是物品，容量是资源，状态值是可达、最值还是计数。一个物品能否在同一轮再次使用，直接决定容量循环方向；组合和排列也由循环顺序表达。复盘时把这个特例压成状态字段：物品阶段、容量、状态值、循环方向、取模或不可达哨兵；组合和排列必须用不同循环顺序表达；再用边界 字符串全零或全一、容量为零、倒序方向、计数预处理 反查初始化、更新顺序和退出条件。

\textbf{状态回放。} 手算路线：手算时用一个容量很小的例子比较正序和倒序。只要循环方向改变后同一物品能否重复使用发生变化，就说明方向不是细节。状态字段：物品阶段、容量、状态值、循环方向、取模或不可达哨兵；组合和排列必须用不同循环顺序表达。状态升级：把选择树压成阶段和容量；物品是否可重复使用决定一维表的遍历方向。

\textbf{证明和边界。} 证明从当前物品使用次数开始：0-1、完全、计数和最值的循环语义必须和题意一致。边界：字符串全零或全一、容量为零、倒序方向、计数预处理。变体：这是背包从一维容量扩展到二维容量的代表。

\noindent\textbf{LeetCode 494 目标和。} 模型：正负号选择可代数转换为选若干数凑出特定正和。推导重点：若 S 加 target 为奇数或目标绝对值过大，直接无解。

\textbf{二刷读题。} 读题时先找阶段、容量、物品使用次数和状态值类型。本题可以压成“正负号选择可代数转换为选若干数凑出特定正和”，循环方向由这些条件决定。先遮住答案，只保留题面，复述候选对象、合法性条件和输出目标。

\textbf{暴力回放。} 慢版本枚举每个物品选不选、选几次，并记录阶段和剩余容量。相同阶段和容量反复出现时，就可以把指数搜索压成表。二刷时先写慢版本或口述慢版本，用它确认不会漏答案。

\textbf{特例回放。} 拿 [1,1,1,1,1]、target=3 手算，给其中一个 1 加负号，其余为正可以得到 3，共有 5 种位置选择。也可转成选正数子集和为 (sum+target)/2。规律是符号分配变成 0-1 计数背包，奇偶和范围先判定。把这个特例继续拆开看：阶段是物品，容量是资源，状态值是可达、最值还是计数。一个物品能否在同一轮再次使用，直接决定容量循环方向；组合和排列也由循环顺序表达。复盘时把这个特例压成状态字段：物品阶段、容量、状态值、循环方向、取模或不可达哨兵；组合和排列必须用不同循环顺序表达；再用边界 零元素会让方案数翻倍、目标为负、容量为零 反查初始化、更新顺序和退出条件。

\textbf{状态回放。} 手算路线：手算时用一个容量很小的例子比较正序和倒序。只要循环方向改变后同一物品能否重复使用发生变化，就说明方向不是细节。状态字段：物品阶段、容量、状态值、循环方向、取模或不可达哨兵；组合和排列必须用不同循环顺序表达。状态升级：把选择树压成阶段和容量；物品是否可重复使用决定一维表的遍历方向。

\textbf{证明和边界。} 证明从当前物品使用次数开始：0-1、完全、计数和最值的循环语义必须和题意一致。边界：零元素会让方案数翻倍、目标为负、容量为零。变体：也可用哈希 DP 保存所有可能和，但背包转换更高效。

\noindent\textbf{LeetCode 518 零钱兑换 II。} 模型：硬币无限使用，问组合数而不是最少数量。推导重点：外层遍历硬币、内层金额正序，保证组合按固定硬币顺序统计。

\textbf{二刷读题。} 读题时先找阶段、容量、物品使用次数和状态值类型。本题可以压成“硬币无限使用，问组合数而不是最少数量”，循环方向由这些条件决定。先遮住答案，只保留题面，复述候选对象、合法性条件和输出目标。

\textbf{暴力回放。} 慢版本枚举每个物品选不选、选几次，并记录阶段和剩余容量。相同阶段和容量反复出现时，就可以把指数搜索压成表。二刷时先写慢版本或口述慢版本，用它确认不会漏答案。

\textbf{特例回放。} 拿 amount=5、coins=[1,2,5] 手算，组合 5、2+2+1、2+1+1+1、全 1 一共 4 种；1+2 和 2+1 不能重复算。规律是组合计数时外层枚举硬币，内层容量正序。把这个特例继续拆开看：阶段是物品，容量是资源，状态值是可达、最值还是计数。一个物品能否在同一轮再次使用，直接决定容量循环方向；组合和排列也由循环顺序表达。复盘时把这个特例压成状态字段：物品阶段、容量、状态值、循环方向、取模或不可达哨兵；组合和排列必须用不同循环顺序表达；再用边界 amount 为零、无硬币、组合数较大、循环顺序反例 反查初始化、更新顺序和退出条件。

\textbf{状态回放。} 手算路线：手算时用一个容量很小的例子比较正序和倒序。只要循环方向改变后同一物品能否重复使用发生变化，就说明方向不是细节。状态字段：物品阶段、容量、状态值、循环方向、取模或不可达哨兵；组合和排列必须用不同循环顺序表达。状态升级：把选择树压成阶段和容量；物品是否可重复使用决定一维表的遍历方向。

\textbf{证明和边界。} 证明从当前物品使用次数开始：0-1、完全、计数和最值的循环语义必须和题意一致。边界：amount 为零、无硬币、组合数较大、循环顺序反例。变体：若外层金额内层硬币，会统计排列数。

\noindent\textbf{LeetCode 879 盈利计划。} 模型：每个项目消耗人数并贡献利润，利润达到下限后可截断。推导重点：二维 DP 用人数和已达到利润表示方案数，利润维度超过目标就压到目标。

\textbf{二刷读题。} 读题时先找阶段、容量、物品使用次数和状态值类型。本题可以压成“每个项目消耗人数并贡献利润，利润达到下限后可截断”，循环方向由这些条件决定。先遮住答案，只保留题面，复述候选对象、合法性条件和输出目标。

\textbf{暴力回放。} 慢版本枚举每个物品选不选、选几次，并记录阶段和剩余容量。相同阶段和容量反复出现时，就可以把指数搜索压成表。二刷时先写慢版本或口述慢版本，用它确认不会漏答案。

\textbf{特例回放。} 拿 n=5、minProfit=3、group=[2,2]、profit=[2,3] 手算，第二个项目单独就达标，两个项目一起也达标且人数不超。规律是状态包含已用人数和被截断到 minProfit 的利润，项目只能选一次。把这个特例继续拆开看：阶段是物品，容量是资源，状态值是可达、最值还是计数。一个物品能否在同一轮再次使用，直接决定容量循环方向；组合和排列也由循环顺序表达。复盘时把这个特例压成状态字段：物品阶段、容量、状态值、循环方向、取模或不可达哨兵；组合和排列必须用不同循环顺序表达；再用边界 minProfit 为零、人数不足、方案数取模、项目只能选一次 反查初始化、更新顺序和退出条件。

\textbf{状态回放。} 手算路线：手算时用一个容量很小的例子比较正序和倒序。只要循环方向改变后同一物品能否重复使用发生变化，就说明方向不是细节。状态字段：物品阶段、容量、状态值、循环方向、取模或不可达哨兵；组合和排列必须用不同循环顺序表达。状态升级：把选择树压成阶段和容量；物品是否可重复使用决定一维表的遍历方向。

\textbf{证明和边界。} 证明从当前物品使用次数开始：0-1、完全、计数和最值的循环语义必须和题意一致。边界：minProfit 为零、人数不足、方案数取模、项目只能选一次。变体：这是 0-1 背包的计数扩展，目标维度需要饱和处理。

\noindent\textbf{LeetCode 956 最高的广告牌。} 模型：两组钢筋高度相等时得到广告牌，高度差是关键状态。推导重点：DP 保存某个高度差下较矮一侧的最大高度，新增钢筋可放高侧、低侧或不用。

\textbf{二刷读题。} 读题时先找阶段、容量、物品使用次数和状态值类型。本题可以压成“两组钢筋高度相等时得到广告牌，高度差是关键状态”，循环方向由这些条件决定。先遮住答案，只保留题面，复述候选对象、合法性条件和输出目标。

\textbf{暴力回放。} 慢版本枚举每个物品选不选、选几次，并记录阶段和剩余容量。相同阶段和容量反复出现时，就可以把指数搜索压成表。二刷时先写慢版本或口述慢版本，用它确认不会漏答案。

\textbf{特例回放。} 拿 rods=[1,2,3,6] 手算，左塔用 6，右塔用 1+2+3，也能达到高度 6。规律是状态保存两边高度差对应的较矮塔最大高度，加入一根杆可以放左、放右或不用。把这个特例继续拆开看：阶段是物品，容量是资源，状态值是可达、最值还是计数。一个物品能否在同一轮再次使用，直接决定容量循环方向；组合和排列也由循环顺序表达。复盘时把这个特例压成状态字段：物品阶段、容量、状态值、循环方向、取模或不可达哨兵；组合和排列必须用不同循环顺序表达；再用边界 空数组、无法平衡、重复长度、差值更新顺序 反查初始化、更新顺序和退出条件。

\textbf{状态回放。} 手算路线：手算时用一个容量很小的例子比较正序和倒序。只要循环方向改变后同一物品能否重复使用发生变化，就说明方向不是细节。状态字段：物品阶段、容量、状态值、循环方向、取模或不可达哨兵；组合和排列必须用不同循环顺序表达。状态升级：把选择树压成阶段和容量；物品是否可重复使用决定一维表的遍历方向。

\textbf{证明和边界。} 证明从当前物品使用次数开始：0-1、完全、计数和最值的循环语义必须和题意一致。边界：空数组、无法平衡、重复长度、差值更新顺序。变体：这是背包从容量转成差值状态的代表。

\noindent\textbf{LeetCode 1049 最后一块石头的重量 II。} 模型：把石头分成两堆，最终差值等于两堆重量差。推导重点：0-1 背包找不超过总和一半的最大可达重量。

\textbf{二刷读题。} 读题时先找阶段、容量、物品使用次数和状态值类型。本题可以压成“把石头分成两堆，最终差值等于两堆重量差”，循环方向由这些条件决定。先遮住答案，只保留题面，复述候选对象、合法性条件和输出目标。

\textbf{暴力回放。} 慢版本枚举每个物品选不选、选几次，并记录阶段和剩余容量。相同阶段和容量反复出现时，就可以把指数搜索压成表。二刷时先写慢版本或口述慢版本，用它确认不会漏答案。

\textbf{特例回放。} 拿 [2,7,4,1,8,1] 手算，把石头分成两堆，最后差值越小结果越小；总和 23，尽量凑接近 11 的一堆。规律是 0-1 背包求不超过 sum/2 的最大可达重量。把这个特例继续拆开看：阶段是物品，容量是资源，状态值是可达、最值还是计数。一个物品能否在同一轮再次使用，直接决定容量循环方向；组合和排列也由循环顺序表达。复盘时把这个特例压成状态字段：物品阶段、容量、状态值、循环方向、取模或不可达哨兵；组合和排列必须用不同循环顺序表达；再用边界 单块石头、总和为奇偶、多个相同重量、容量方向 反查初始化、更新顺序和退出条件。

\textbf{状态回放。} 手算路线：手算时用一个容量很小的例子比较正序和倒序。只要循环方向改变后同一物品能否重复使用发生变化，就说明方向不是细节。状态字段：物品阶段、容量、状态值、循环方向、取模或不可达哨兵；组合和排列必须用不同循环顺序表达。状态升级：把选择树压成阶段和容量；物品是否可重复使用决定一维表的遍历方向。

\textbf{证明和边界。} 证明从当前物品使用次数开始：0-1、完全、计数和最值的循环语义必须和题意一致。边界：单块石头、总和为奇偶、多个相同重量、容量方向。变体：它和分割等和子集同源，只是目标从是否可达变成最小差。

\noindent\textbf{LeetCode 1155 掷骰子等于目标和的方法数。} 模型：多个骰子每个只能取一到 faces 的点数，问目标和方案数。推导重点：按骰子数量分阶段，容量为当前点数和，枚举当前骰子的点数转移。

\textbf{二刷读题。} 读题时先找阶段、容量、物品使用次数和状态值类型。本题可以压成“多个骰子每个只能取一到 faces 的点数，问目标和方案数”，循环方向由这些条件决定。先遮住答案，只保留题面，复述候选对象、合法性条件和输出目标。

\textbf{暴力回放。} 慢版本枚举每个物品选不选、选几次，并记录阶段和剩余容量。相同阶段和容量反复出现时，就可以把指数搜索压成表。二刷时先写慢版本或口述慢版本，用它确认不会漏答案。

\textbf{特例回放。} 拿 2 个 6 面骰、target=7 手算，第一颗为 1 时第二颗必须 6，第一颗为 2 时第二颗必须 5。规律是阶段是骰子个数，容量是当前点数和，每个骰子枚举 1..k 的面值。把这个特例继续拆开看：阶段是物品，容量是资源，状态值是可达、最值还是计数。一个物品能否在同一轮再次使用，直接决定容量循环方向；组合和排列也由循环顺序表达。复盘时把这个特例压成状态字段：物品阶段、容量、状态值、循环方向、取模或不可达哨兵；组合和排列必须用不同循环顺序表达；再用边界 target 太小或太大、取模、骰子数为一、滚动数组方向 反查初始化、更新顺序和退出条件。

\textbf{状态回放。} 手算路线：手算时用一个容量很小的例子比较正序和倒序。只要循环方向改变后同一物品能否重复使用发生变化，就说明方向不是细节。状态字段：物品阶段、容量、状态值、循环方向、取模或不可达哨兵；组合和排列必须用不同循环顺序表达。状态升级：把选择树压成阶段和容量；物品是否可重复使用决定一维表的遍历方向。

\textbf{证明和边界。} 证明从当前物品使用次数开始：0-1、完全、计数和最值的循环语义必须和题意一致。边界：target 太小或太大、取模、骰子数为一、滚动数组方向。变体：这是有界选择计数，和完全背包不能混。

\noindent\textbf{LeetCode 1449 数位成本和为目标值的最大数字。} 模型：每个数字可重复选择，容量是成本，目标是组成字典序和长度最大的数字。推导重点：完全背包先最大化位数，再按数字从大到小贪心还原答案。

\textbf{二刷读题。} 读题时先找阶段、容量、物品使用次数和状态值类型。本题可以压成“每个数字可重复选择，容量是成本，目标是组成字典序和长度最大的数字”，循环方向由这些条件决定。先遮住答案，只保留题面，复述候选对象、合法性条件和输出目标。

\textbf{暴力回放。} 慢版本枚举每个物品选不选、选几次，并记录阶段和剩余容量。相同阶段和容量反复出现时，就可以把指数搜索压成表。二刷时先写慢版本或口述慢版本，用它确认不会漏答案。

\textbf{特例回放。} 拿 cost 中数字 7 和 2 成本较低、target=9 手算，先最大化能组成的位数，再从 9 到 1 尝试还原字典序最大的数字。规律是完全背包求最大位数，重建时优先放大数字且保证剩余目标可达。把这个特例继续拆开看：阶段是物品，容量是资源，状态值是可达、最值还是计数。一个物品能否在同一轮再次使用，直接决定容量循环方向；组合和排列也由循环顺序表达。复盘时把这个特例压成状态字段：物品阶段、容量、状态值、循环方向、取模或不可达哨兵；组合和排列必须用不同循环顺序表达；再用边界 无法组成目标、多个数字同成本、target 为零、字符串比较 反查初始化、更新顺序和退出条件。

\textbf{状态回放。} 手算路线：手算时用一个容量很小的例子比较正序和倒序。只要循环方向改变后同一物品能否重复使用发生变化，就说明方向不是细节。状态字段：物品阶段、容量、状态值、循环方向、取模或不可达哨兵；组合和排列必须用不同循环顺序表达。状态升级：把选择树压成阶段和容量；物品是否可重复使用决定一维表的遍历方向。

\textbf{证明和边界。} 证明从当前物品使用次数开始：0-1、完全、计数和最值的循环语义必须和题意一致。边界：无法组成目标、多个数字同成本、target 为零、字符串比较。变体：这题说明背包状态值不一定是数值收益，也可以是长度和字典序。

\subsection{自测标准}

二刷时不要只确认自己见过这道题。请遮住答案，先讲题面模型，再讲暴力基线和瓶颈，然后说出优化结构保存了什么、不变量如何排除错误候选、边界实验如何打破错误实现。

\section{贪心与交换证明：二刷复盘卡}
这一大节只围绕一个题型复盘，不把题号直接铺成目录。阅读时先看复盘目标，再读核心题卡，最后用自测标准检查自己是否能独立重讲。

\subsection{复盘目标}

追问通常会要求证明。请准备交换论证或领先性论证，并给出一个看似合理但错误的贪心规则反例。

\subsection{核心题卡}

\noindent\textbf{LeetCode 45 跳跃游戏 II。} 模型：每一次跳跃可以看成 BFS 的一层，当前层内所有位置共享同一次跳数。推导重点：扫描当前层时维护下一层最远边界，到达当前层末尾才增加跳数。

\textbf{二刷读题。} 读题时先找局部选择消耗什么资源，以及选择之后给未来留下什么边界。本题可以压成“每一次跳跃可以看成 BFS 的一层，当前层内所有位置共享同一次跳数”，扫描当前层时维护下一层最远边界，到达当前层末尾才增加跳数必须能被证明安全。先遮住答案，只保留题面，复述候选对象、合法性条件和输出目标。

\textbf{暴力回放。} 慢版本枚举选择顺序或用 DP 保留多个可能方案，先确认最优值存在。随后才检查局部选择是否能通过交换或领先性固定下来。二刷时先写慢版本或口述慢版本，用它确认不会漏答案。

\textbf{特例回放。} 拿 [2,3,1,1,4] 手算，起点这一层能到下标 1 和 2；扫描这一层时，下一层最远可到 4，走完当前层末尾才把步数加一。规律是把可达区间看成 BFS 层，当前层结束时切到下一层；不能在每个位置都加跳数。把这个特例继续拆开看：先构造一个非贪心选择，再比较两者留下的资源、右边界或剩余时间。只有能把非贪心第一步交换成贪心第一步且不变差，局部选择才是规律。复盘时把这个特例压成状态字段：排序后的候选、当前已选方案、剩余资源或最远边界、答案计数；每次局部选择后要能说明当前状态不落后；再用边界 数组长度一、刚好到达、存在多个可选跳点、不要在每个位置都加跳数 反查初始化、更新顺序和退出条件。

\textbf{状态回放。} 手算路线：手算时同时写一个非贪心选择，与贪心选择比较剩余资源。若无法说明贪心后剩余资源不差，就不能直接下结论。状态字段：排序后的候选、当前已选方案、剩余资源或最远边界、答案计数；每次局部选择后要能说明当前状态不落后。状态升级：把指数级选择压成一个可证明安全的局部选择；状态通常是当前最远边界、最早结束时间、最小代价或剩余资源。

\textbf{证明和边界。} 证明从交换或领先性开始：任意最优解都能替换成当前贪心选择而不变差，或贪心状态始终不落后。边界：数组长度一、刚好到达、存在多个可选跳点、不要在每个位置都加跳数。变体：若问能否到达，只需维护全局最远可达位置。

\noindent\textbf{LeetCode 55 跳跃游戏。} 模型：扫描过程中只关心当前能到达的最远位置。推导重点：如果下标超过最远可达，任何之前路径都不能到达它；否则用它扩展边界。

\textbf{二刷读题。} 读题时先找局部选择消耗什么资源，以及选择之后给未来留下什么边界。本题可以压成“扫描过程中只关心当前能到达的最远位置”，如果下标超过最远可达，任何之前路径都不能到达它；否则用它扩展边界必须能被证明安全。先遮住答案，只保留题面，复述候选对象、合法性条件和输出目标。

\textbf{暴力回放。} 慢版本枚举选择顺序或用 DP 保留多个可能方案，先确认最优值存在。随后才检查局部选择是否能通过交换或领先性固定下来。二刷时先写慢版本或口述慢版本，用它确认不会漏答案。

\textbf{特例回放。} 拿 [2,3,1,1,4] 手算，下标 0 能覆盖到 2；走到下标 1 时最远可达更新到 4，已经覆盖终点。再拿 [0,1]，扫描到下标 1 时已经超过最远可达，失败。规律是只维护当前能到达的最远边界，任何超过边界的位置都不可能由前面路径到达。把这个特例继续拆开看：先构造一个非贪心选择，再比较两者留下的资源、右边界或剩余时间。只有能把非贪心第一步交换成贪心第一步且不变差，局部选择才是规律。复盘时把这个特例压成状态字段：排序后的候选、当前已选方案、剩余资源或最远边界、答案计数；每次局部选择后要能说明当前状态不落后；再用边界 第一个元素为零、数组长度一、恰好到达最后、远超最后 反查初始化、更新顺序和退出条件。

\textbf{状态回放。} 手算路线：手算时同时写一个非贪心选择，与贪心选择比较剩余资源。若无法说明贪心后剩余资源不差，就不能直接下结论。状态字段：排序后的候选、当前已选方案、剩余资源或最远边界、答案计数；每次局部选择后要能说明当前状态不落后。状态升级：把指数级选择压成一个可证明安全的局部选择；状态通常是当前最远边界、最早结束时间、最小代价或剩余资源。

\textbf{证明和边界。} 证明从交换或领先性开始：任意最优解都能替换成当前贪心选择而不变差，或贪心状态始终不落后。边界：第一个元素为零、数组长度一、恰好到达最后、远超最后。变体：求最少跳数时把可达区间看成 BFS 层。

\noindent\textbf{LeetCode 122 买卖股票 II。} 模型：允许多次交易且不能同时持有，所有正收益差都可以收集。推导重点：把长上涨段拆成每天买卖收益相同，因此累加正差值最优。

\textbf{二刷读题。} 读题时先找局部选择消耗什么资源，以及选择之后给未来留下什么边界。本题可以压成“允许多次交易且不能同时持有，所有正收益差都可以收集”，把长上涨段拆成每天买卖收益相同，因此累加正差值最优必须能被证明安全。先遮住答案，只保留题面，复述候选对象、合法性条件和输出目标。

\textbf{暴力回放。} 慢版本枚举选择顺序或用 DP 保留多个可能方案，先确认最优值存在。随后才检查局部选择是否能通过交换或领先性固定下来。二刷时先写慢版本或口述慢版本，用它确认不会漏答案。

\textbf{特例回放。} 拿 [7,1,5,3,6,4] 手算，1 到 5 的利润 4，加上 3 到 6 的利润 3，总利润 7。把 1 到 6 的上涨段拆成每天正差值，收益相同。规律是无限次交易且无手续费时，所有正差值都可收集；下降段和平台不贡献收益。把这个特例继续拆开看：先构造一个非贪心选择，再比较两者留下的资源、右边界或剩余时间。只有能把非贪心第一步交换成贪心第一步且不变差，局部选择才是规律。复盘时把这个特例压成状态字段：排序后的候选、当前已选方案、剩余资源或最远边界、答案计数；每次局部选择后要能说明当前状态不落后；再用边界 下降段、平台价格、只有一天、手续费不存在 反查初始化、更新顺序和退出条件。

\textbf{状态回放。} 手算路线：手算时同时写一个非贪心选择，与贪心选择比较剩余资源。若无法说明贪心后剩余资源不差，就不能直接下结论。状态字段：排序后的候选、当前已选方案、剩余资源或最远边界、答案计数；每次局部选择后要能说明当前状态不落后。状态升级：把指数级选择压成一个可证明安全的局部选择；状态通常是当前最远边界、最早结束时间、最小代价或剩余资源。

\textbf{证明和边界。} 证明从交换或领先性开始：任意最优解都能替换成当前贪心选择而不变差，或贪心状态始终不落后。边界：下降段、平台价格、只有一天、手续费不存在。变体：有手续费时不能简单累加正差，需要状态机或调整贪心阈值。

\noindent\textbf{LeetCode 134 加油站。} 模型：绕行一圈是否可行取决于总油量差和局部剩余油量。推导重点：总和为负必不可行；扫描时若当前剩余为负，起点只能放到下一站。

\textbf{二刷读题。} 读题时先找局部选择消耗什么资源，以及选择之后给未来留下什么边界。本题可以压成“绕行一圈是否可行取决于总油量差和局部剩余油量”，总和为负必不可行；扫描时若当前剩余为负，起点只能放到下一站必须能被证明安全。先遮住答案，只保留题面，复述候选对象、合法性条件和输出目标。

\textbf{暴力回放。} 慢版本枚举选择顺序或用 DP 保留多个可能方案，先确认最优值存在。随后才检查局部选择是否能通过交换或领先性固定下来。二刷时先写慢版本或口述慢版本，用它确认不会漏答案。

\textbf{特例回放。} 拿 gas=[1,2,3,4,5]、cost=[3,4,5,1,2] 手算，从 0 出发中途油量为负，0 到失败点之间任何站都不能作为起点；从 3 出发可以绕完。规律是总油量差为负则无解，局部剩余为负时起点跳到下一站。把这个特例继续拆开看：先构造一个非贪心选择，再比较两者留下的资源、右边界或剩余时间。只有能把非贪心第一步交换成贪心第一步且不变差，局部选择才是规律。复盘时把这个特例压成状态字段：排序后的候选、当前已选方案、剩余资源或最远边界、答案计数；每次局部选择后要能说明当前状态不落后；再用边界 总和刚好为零、多个可行起点、某站油量为零、消耗大于补给 反查初始化、更新顺序和退出条件。

\textbf{状态回放。} 手算路线：手算时同时写一个非贪心选择，与贪心选择比较剩余资源。若无法说明贪心后剩余资源不差，就不能直接下结论。状态字段：排序后的候选、当前已选方案、剩余资源或最远边界、答案计数；每次局部选择后要能说明当前状态不落后。状态升级：把指数级选择压成一个可证明安全的局部选择；状态通常是当前最远边界、最早结束时间、最小代价或剩余资源。

\textbf{证明和边界。} 证明从交换或领先性开始：任意最优解都能替换成当前贪心选择而不变差，或贪心状态始终不落后。边界：总和刚好为零、多个可行起点、某站油量为零、消耗大于补给。变体：它的领先性证明是前一段任意站点都不能作为起点。

\noindent\textbf{LeetCode 135 分发糖果。} 模型：相邻评分约束需要同时满足左邻和右邻两个方向。推导重点：左到右满足递增关系，右到左满足反向递增关系，最后取两侧要求最大值。

\textbf{二刷读题。} 读题时先找局部选择消耗什么资源，以及选择之后给未来留下什么边界。本题可以压成“相邻评分约束需要同时满足左邻和右邻两个方向”，左到右满足递增关系，右到左满足反向递增关系，最后取两侧要求最大值必须能被证明安全。先遮住答案，只保留题面，复述候选对象、合法性条件和输出目标。

\textbf{暴力回放。} 慢版本枚举选择顺序或用 DP 保留多个可能方案，先确认最优值存在。随后才检查局部选择是否能通过交换或领先性固定下来。二刷时先写慢版本或口述慢版本，用它确认不会漏答案。

\textbf{特例回放。} 拿 ratings=[1,0,2] 手算，中间孩子评分最低，左右都要比它多糖，结果 [2,1,2]。只从左到右会漏掉右侧约束。规律是左到右满足左邻递增，右到左满足右邻递增，最终取两侧要求最大值。把这个特例继续拆开看：先构造一个非贪心选择，再比较两者留下的资源、右边界或剩余时间。只有能把非贪心第一步交换成贪心第一步且不变差，局部选择才是规律。复盘时把这个特例压成状态字段：排序后的候选、当前已选方案、剩余资源或最远边界、答案计数；每次局部选择后要能说明当前状态不落后；再用边界 评分相等、严格递增、严格递减、山峰和谷底 反查初始化、更新顺序和退出条件。

\textbf{状态回放。} 手算路线：手算时同时写一个非贪心选择，与贪心选择比较剩余资源。若无法说明贪心后剩余资源不差，就不能直接下结论。状态字段：排序后的候选、当前已选方案、剩余资源或最远边界、答案计数；每次局部选择后要能说明当前状态不落后。状态升级：把指数级选择压成一个可证明安全的局部选择；状态通常是当前最远边界、最早结束时间、最小代价或剩余资源。

\textbf{证明和边界。} 证明从交换或领先性开始：任意最优解都能替换成当前贪心选择而不变差，或贪心状态始终不落后。边界：评分相等、严格递增、严格递减、山峰和谷底。变体：这是局部约束两遍扫描，不能只用一个方向贪心。

\noindent\textbf{LeetCode 169 多数元素。} 模型：超过一半的元素可以和其他元素两两抵消后仍然剩下。推导重点：Boyer-Moore 投票维护候选和计数。

\textbf{二刷读题。} 读题时先找局部选择消耗什么资源，以及选择之后给未来留下什么边界。本题可以压成“超过一半的元素可以和其他元素两两抵消后仍然剩下”，Boyer-Moore 投票维护候选和计数必须能被证明安全。先遮住答案，只保留题面，复述候选对象、合法性条件和输出目标。

\textbf{暴力回放。} 慢版本枚举选择顺序或用 DP 保留多个可能方案，先确认最优值存在。随后才检查局部选择是否能通过交换或领先性固定下来。二刷时先写慢版本或口述慢版本，用它确认不会漏答案。

\textbf{特例回放。} 拿 [2,2,1,1,1,2,2] 手算，候选和不同元素互相抵消，最后剩下的候选是 2。规律是多数元素超过一半，和其他元素两两抵消后仍会剩下；题目不保证存在时需要二次计数验证。把这个特例继续拆开看：先构造一个非贪心选择，再比较两者留下的资源、右边界或剩余时间。只有能把非贪心第一步交换成贪心第一步且不变差，局部选择才是规律。复盘时把这个特例压成状态字段：排序后的候选、当前已选方案、剩余资源或最远边界、答案计数；每次局部选择后要能说明当前状态不落后；再用边界 题目是否保证存在多数、计数归零、全相同 反查初始化、更新顺序和退出条件。

\textbf{状态回放。} 手算路线：手算时同时写一个非贪心选择，与贪心选择比较剩余资源。若无法说明贪心后剩余资源不差，就不能直接下结论。状态字段：排序后的候选、当前已选方案、剩余资源或最远边界、答案计数；每次局部选择后要能说明当前状态不落后。状态升级：把指数级选择压成一个可证明安全的局部选择；状态通常是当前最远边界、最早结束时间、最小代价或剩余资源。

\textbf{证明和边界。} 证明从交换或领先性开始：任意最优解都能替换成当前贪心选择而不变差，或贪心状态始终不落后。边界：题目是否保证存在多数、计数归零、全相同。变体：若找超过三分之一的元素，需要维护两个候选。

\noindent\textbf{LeetCode 316 去除重复字母。} 模型：每个字符必须保留一次，同时结果字典序尽量小。推导重点：单调栈维护当前结果，若栈顶更大且后面还会出现，就可以弹出再等后面放回。

\textbf{二刷读题。} 读题时先找局部选择消耗什么资源，以及选择之后给未来留下什么边界。本题可以压成“每个字符必须保留一次，同时结果字典序尽量小”，单调栈维护当前结果，若栈顶更大且后面还会出现，就可以弹出再等后面放回必须能被证明安全。先遮住答案，只保留题面，复述候选对象、合法性条件和输出目标。

\textbf{暴力回放。} 慢版本枚举选择顺序或用 DP 保留多个可能方案，先确认最优值存在。随后才检查局部选择是否能通过交换或领先性固定下来。二刷时先写慢版本或口述慢版本，用它确认不会漏答案。

\textbf{特例回放。} 拿 cbacdcbc 手算，遇到 b 时，前面的 c 后面还会出现且 c>b，所以可以弹出 c；但遇到只剩最后一次的字符就不能弹。规律是单调栈维护结果字典序，弹出条件必须同时满足栈顶更大且后面还会出现。把这个特例继续拆开看：先构造一个非贪心选择，再比较两者留下的资源、右边界或剩余时间。只有能把非贪心第一步交换成贪心第一步且不变差，局部选择才是规律。复盘时把这个特例压成状态字段：排序后的候选、当前已选方案、剩余资源或最远边界、答案计数；每次局部选择后要能说明当前状态不落后；再用边界 字符只出现一次、重复字符、弹出后 visited 恢复、字典序比较 反查初始化、更新顺序和退出条件。

\textbf{状态回放。} 手算路线：手算时同时写一个非贪心选择，与贪心选择比较剩余资源。若无法说明贪心后剩余资源不差，就不能直接下结论。状态字段：排序后的候选、当前已选方案、剩余资源或最远边界、答案计数；每次局部选择后要能说明当前状态不落后。状态升级：把指数级选择压成一个可证明安全的局部选择；状态通常是当前最远边界、最早结束时间、最小代价或剩余资源。

\textbf{证明和边界。} 证明从交换或领先性开始：任意最优解都能替换成当前贪心选择而不变差，或贪心状态始终不落后。边界：字符只出现一次、重复字符、弹出后 visited 恢复、字典序比较。变体：402 移掉 K 位数字也是单调栈删除，但删除预算和必须保留集合不同。

\noindent\textbf{LeetCode 406 根据身高重建队列。} 模型：高个子的位置先确定后，不会被矮个子影响其前方高个数量。推导重点：按身高降序、k 升序排序，再按 k 插入当前结果位置。

\textbf{二刷读题。} 读题时先找局部选择消耗什么资源，以及选择之后给未来留下什么边界。本题可以压成“高个子的位置先确定后，不会被矮个子影响其前方高个数量”，按身高降序、k 升序排序，再按 k 插入当前结果位置必须能被证明安全。先遮住答案，只保留题面，复述候选对象、合法性条件和输出目标。

\textbf{暴力回放。} 慢版本枚举选择顺序或用 DP 保留多个可能方案，先确认最优值存在。随后才检查局部选择是否能通过交换或领先性固定下来。二刷时先写慢版本或口述慢版本，用它确认不会漏答案。

\textbf{特例回放。} 拿 [7,0]、[7,1]、[5,0] 手算，先放高个子，矮个子插入不会改变高个子前面高个数量；[5,0] 插到最前也不影响 [7,0] 的统计。规律是按身高降序、k 升序排序，再按 k 插入。把这个特例继续拆开看：先构造一个非贪心选择，再比较两者留下的资源、右边界或剩余时间。只有能把非贪心第一步交换成贪心第一步且不变差，局部选择才是规律。复盘时把这个特例压成状态字段：排序后的候选、当前已选方案、剩余资源或最远边界、答案计数；每次局部选择后要能说明当前状态不落后；再用边界 相同身高、k 为零、插入到末尾、结果稳定性 反查初始化、更新顺序和退出条件。

\textbf{状态回放。} 手算路线：手算时同时写一个非贪心选择，与贪心选择比较剩余资源。若无法说明贪心后剩余资源不差，就不能直接下结论。状态字段：排序后的候选、当前已选方案、剩余资源或最远边界、答案计数；每次局部选择后要能说明当前状态不落后。状态升级：把指数级选择压成一个可证明安全的局部选择；状态通常是当前最远边界、最早结束时间、最小代价或剩余资源。

\textbf{证明和边界。} 证明从交换或领先性开始：任意最优解都能替换成当前贪心选择而不变差，或贪心状态始终不落后。边界：相同身高、k 为零、插入到末尾、结果稳定性。变体：这是排序后插入的构造型贪心，证明依赖高个子先固定。

\noindent\textbf{LeetCode 435 无重叠区间。} 模型：保留最多不重叠区间等价于删除最少区间。推导重点：按终点升序选择，结束越早给后续留下空间越多。

\textbf{二刷读题。} 读题时先找局部选择消耗什么资源，以及选择之后给未来留下什么边界。本题可以压成“保留最多不重叠区间等价于删除最少区间”，按终点升序选择，结束越早给后续留下空间越多必须能被证明安全。先遮住答案，只保留题面，复述候选对象、合法性条件和输出目标。

\textbf{暴力回放。} 慢版本枚举选择顺序或用 DP 保留多个可能方案，先确认最优值存在。随后才检查局部选择是否能通过交换或领先性固定下来。二刷时先写慢版本或口述慢版本，用它确认不会漏答案。

\textbf{特例回放。} 拿 [[1,2],[2,3],[3,4],[1,3]] 手算，选 [1,2] 后还能接 [2,3] 和 [3,4]；若先选 [1,3]，会挡住更多区间。规律是按右端点升序选择，结束越早给未来留下的空间越多，删除数等于总数减保留数。把这个特例继续拆开看：先构造一个非贪心选择，再比较两者留下的资源、右边界或剩余时间。只有能把非贪心第一步交换成贪心第一步且不变差，局部选择才是规律。复盘时把这个特例压成状态字段：排序后的候选、当前已选方案、剩余资源或最远边界、答案计数；每次局部选择后要能说明当前状态不落后；再用边界 端点相接是否重叠、完全覆盖、相同终点、空列表 反查初始化、更新顺序和退出条件。

\textbf{状态回放。} 手算路线：手算时同时写一个非贪心选择，与贪心选择比较剩余资源。若无法说明贪心后剩余资源不差，就不能直接下结论。状态字段：排序后的候选、当前已选方案、剩余资源或最远边界、答案计数；每次局部选择后要能说明当前状态不落后。状态升级：把指数级选择压成一个可证明安全的局部选择；状态通常是当前最远边界、最早结束时间、最小代价或剩余资源。

\textbf{证明和边界。} 证明从交换或领先性开始：任意最优解都能替换成当前贪心选择而不变差，或贪心状态始终不落后。边界：端点相接是否重叠、完全覆盖、相同终点、空列表。变体：用交换论证说明最早结束的选择不劣于其他第一选择。

\noindent\textbf{LeetCode 452 用最少数量的箭引爆气球。} 模型：一支箭的位置可以覆盖所有包含该位置的区间。推导重点：按右端点排序，当前箭放在最早结束气球的右端。

\textbf{二刷读题。} 读题时先找局部选择消耗什么资源，以及选择之后给未来留下什么边界。本题可以压成“一支箭的位置可以覆盖所有包含该位置的区间”，按右端点排序，当前箭放在最早结束气球的右端必须能被证明安全。先遮住答案，只保留题面，复述候选对象、合法性条件和输出目标。

\textbf{暴力回放。} 慢版本枚举选择顺序或用 DP 保留多个可能方案，先确认最优值存在。随后才检查局部选择是否能通过交换或领先性固定下来。二刷时先写慢版本或口述慢版本，用它确认不会漏答案。

\textbf{特例回放。} 拿 [[10,16],[2,8],[1,6],[7,12]] 手算，按右端排序后第一箭射在 6，可以覆盖 [1,6] 和 [2,8]；下一箭射在 12。规律是每次把箭放在最早结束气球的右端，能最大化后续空间。把这个特例继续拆开看：先构造一个非贪心选择，再比较两者留下的资源、右边界或剩余时间。只有能把非贪心第一步交换成贪心第一步且不变差，局部选择才是规律。复盘时把这个特例压成状态字段：排序后的候选、当前已选方案、剩余资源或最远边界、答案计数；每次局部选择后要能说明当前状态不落后；再用边界 端点相同、负坐标、区间包含、闭区间边界 反查初始化、更新顺序和退出条件。

\textbf{状态回放。} 手算路线：手算时同时写一个非贪心选择，与贪心选择比较剩余资源。若无法说明贪心后剩余资源不差，就不能直接下结论。状态字段：排序后的候选、当前已选方案、剩余资源或最远边界、答案计数；每次局部选择后要能说明当前状态不落后。状态升级：把指数级选择压成一个可证明安全的局部选择；状态通常是当前最远边界、最早结束时间、最小代价或剩余资源。

\textbf{证明和边界。} 证明从交换或领先性开始：任意最优解都能替换成当前贪心选择而不变差，或贪心状态始终不落后。边界：端点相同、负坐标、区间包含、闭区间边界。变体：它和无重叠区间是同一类最早结束贪心。

\noindent\textbf{LeetCode 630 课程表 III。} 模型：每门课有持续时间和截止日，已选课程总时长不能超过当前截止日。推导重点：按截止日排序，先选当前课；若超时，就用大顶堆丢掉已选中耗时最长的课。

\textbf{二刷读题。} 读题时先找局部选择消耗什么资源，以及选择之后给未来留下什么边界。本题可以压成“每门课有持续时间和截止日，已选课程总时长不能超过当前截止日”，按截止日排序，先选当前课；若超时，就用大顶堆丢掉已选中耗时最长的课必须能被证明安全。先遮住答案，只保留题面，复述候选对象、合法性条件和输出目标。

\textbf{暴力回放。} 慢版本枚举选择顺序或用 DP 保留多个可能方案，先确认最优值存在。随后才检查局部选择是否能通过交换或领先性固定下来。二刷时先写慢版本或口述慢版本，用它确认不会漏答案。

\textbf{特例回放。} 拿课程 [100,200]、[200,1300]、[1000,1250] 手算，按截止日排序后若总时长超出当前截止日，就丢掉已选课程里耗时最长的。规律是大顶堆保存已选课程时长，替换最长课能给未来留下最多时间。把这个特例继续拆开看：先构造一个非贪心选择，再比较两者留下的资源、右边界或剩余时间。只有能把非贪心第一步交换成贪心第一步且不变差，局部选择才是规律。复盘时把这个特例压成状态字段：排序后的候选、当前已选方案、剩余资源或最远边界、答案计数；每次局部选择后要能说明当前状态不落后；再用边界 课程无法单独完成、截止日相同、替换已选课程、空列表 反查初始化、更新顺序和退出条件。

\textbf{状态回放。} 手算路线：手算时同时写一个非贪心选择，与贪心选择比较剩余资源。若无法说明贪心后剩余资源不差，就不能直接下结论。状态字段：排序后的候选、当前已选方案、剩余资源或最远边界、答案计数；每次局部选择后要能说明当前状态不落后。状态升级：把指数级选择压成一个可证明安全的局部选择；状态通常是当前最远边界、最早结束时间、最小代价或剩余资源。

\textbf{证明和边界。} 证明从交换或领先性开始：任意最优解都能替换成当前贪心选择而不变差，或贪心状态始终不落后。边界：课程无法单独完成、截止日相同、替换已选课程、空列表。变体：这是反悔贪心：接受后在约束被破坏时删最差候选。

\noindent\textbf{LeetCode 646 最长数对链。} 模型：每个数对像一个区间，选择链要求前一个右端小于后一个左端。推导重点：按右端升序选择，当前右端越小越容易接后续数对。

\textbf{二刷读题。} 读题时先找局部选择消耗什么资源，以及选择之后给未来留下什么边界。本题可以压成“每个数对像一个区间，选择链要求前一个右端小于后一个左端”，按右端升序选择，当前右端越小越容易接后续数对必须能被证明安全。先遮住答案，只保留题面，复述候选对象、合法性条件和输出目标。

\textbf{暴力回放。} 慢版本枚举选择顺序或用 DP 保留多个可能方案，先确认最优值存在。随后才检查局部选择是否能通过交换或领先性固定下来。二刷时先写慢版本或口述慢版本，用它确认不会漏答案。

\textbf{特例回放。} 拿 [[1,2],[2,3],[3,4]] 手算，[1,2] 不能接 [2,3]，因为要求前一个右端小于后一个左端；可以接 [3,4]。规律是按右端升序贪心选择，右端越小越容易接后续数对。把这个特例继续拆开看：先构造一个非贪心选择，再比较两者留下的资源、右边界或剩余时间。只有能把非贪心第一步交换成贪心第一步且不变差，局部选择才是规律。复盘时把这个特例压成状态字段：排序后的候选、当前已选方案、剩余资源或最远边界、答案计数；每次局部选择后要能说明当前状态不落后；再用边界 端点相等不允许连接、完全覆盖、负数端点、空数组 反查初始化、更新顺序和退出条件。

\textbf{状态回放。} 手算路线：手算时同时写一个非贪心选择，与贪心选择比较剩余资源。若无法说明贪心后剩余资源不差，就不能直接下结论。状态字段：排序后的候选、当前已选方案、剩余资源或最远边界、答案计数；每次局部选择后要能说明当前状态不落后。状态升级：把指数级选择压成一个可证明安全的局部选择；状态通常是当前最远边界、最早结束时间、最小代价或剩余资源。

\textbf{证明和边界。} 证明从交换或领先性开始：任意最优解都能替换成当前贪心选择而不变差，或贪心状态始终不落后。边界：端点相等不允许连接、完全覆盖、负数端点、空数组。变体：也可以 DP，但贪心用最早结束证明更简洁。

\noindent\textbf{LeetCode 678 有效的括号字符串。} 模型：星号可以作为左括号、右括号或空，使未匹配左括号数量变成一个可能区间。推导重点：贪心维护可能未匹配左括号数的下界和上界，遇到星号扩大区间并把下界截到零。

\textbf{二刷读题。} 读题时先找局部选择消耗什么资源，以及选择之后给未来留下什么边界。本题可以压成“星号可以作为左括号、右括号或空，使未匹配左括号数量变成一个可能区间”，贪心维护可能未匹配左括号数的下界和上界，遇到星号扩大区间并把下界截到零必须能被证明安全。先遮住答案，只保留题面，复述候选对象、合法性条件和输出目标。

\textbf{暴力回放。} 慢版本枚举选择顺序或用 DP 保留多个可能方案，先确认最优值存在。随后才检查局部选择是否能通过交换或领先性固定下来。二刷时先写慢版本或口述慢版本，用它确认不会漏答案。

\textbf{特例回放。} 拿 (*) 手算，星号可以当空或左括号；拿 (*))，星号必须当左括号才能匹配后面的右括号。与其枚举星号三种选择，不如维护可能未匹配左括号数的区间 [low, high]。规律是左括号让区间加一，右括号让区间减一，星号让下界减一上界加一，且下界不能低于零。把这个特例继续拆开看：先构造一个非贪心选择，再比较两者留下的资源、右边界或剩余时间。只有能把非贪心第一步交换成贪心第一步且不变差，局部选择才是规律。复盘时把这个特例压成状态字段：排序后的候选、当前已选方案、剩余资源或最远边界、答案计数；每次局部选择后要能说明当前状态不落后；再用边界 上界变负、最终下界不为零、连续星号、前缀右括号过多 反查初始化、更新顺序和退出条件。

\textbf{状态回放。} 手算路线：手算时同时写一个非贪心选择，与贪心选择比较剩余资源。若无法说明贪心后剩余资源不差，就不能直接下结论。状态字段：排序后的候选、当前已选方案、剩余资源或最远边界、答案计数；每次局部选择后要能说明当前状态不落后。状态升级：把指数级选择压成一个可证明安全的局部选择；状态通常是当前最远边界、最早结束时间、最小代价或剩余资源。

\textbf{证明和边界。} 证明从交换或领先性开始：任意最优解都能替换成当前贪心选择而不变差，或贪心状态始终不落后。边界：上界变负、最终下界不为零、连续星号、前缀右括号过多。变体：也可用双栈保存左括号和星号位置，但区间贪心更能说明状态压缩。

\noindent\textbf{LeetCode 763 划分字母区间。} 模型：每个片段必须覆盖其中所有字符的最后出现位置。推导重点：扫描维护当前片段最远右端，到达右端即可切分。

\textbf{二刷读题。} 读题时先找局部选择消耗什么资源，以及选择之后给未来留下什么边界。本题可以压成“每个片段必须覆盖其中所有字符的最后出现位置”，扫描维护当前片段最远右端，到达右端即可切分必须能被证明安全。先遮住答案，只保留题面，复述候选对象、合法性条件和输出目标。

\textbf{暴力回放。} 慢版本枚举选择顺序或用 DP 保留多个可能方案，先确认最优值存在。随后才检查局部选择是否能通过交换或领先性固定下来。二刷时先写慢版本或口述慢版本，用它确认不会漏答案。

\textbf{特例回放。} 拿 ababcbacadefegdehijhklij 手算，字符 a 的最后位置在 8，所以第一个片段至少到 8；扫描中 b、c 的最后位置也不超过 8，到 8 才能切。规律是当前片段右端是片段内所有字符最后出现位置的最大值。把这个特例继续拆开看：先构造一个非贪心选择，再比较两者留下的资源、右边界或剩余时间。只有能把非贪心第一步交换成贪心第一步且不变差，局部选择才是规律。复盘时把这个特例压成状态字段：排序后的候选、当前已选方案、剩余资源或最远边界、答案计数；每次局部选择后要能说明当前状态不落后；再用边界 单字符、所有字符相同、字符多次跨段 反查初始化、更新顺序和退出条件。

\textbf{状态回放。} 手算路线：手算时同时写一个非贪心选择，与贪心选择比较剩余资源。若无法说明贪心后剩余资源不差，就不能直接下结论。状态字段：排序后的候选、当前已选方案、剩余资源或最远边界、答案计数；每次局部选择后要能说明当前状态不落后。状态升级：把指数级选择压成一个可证明安全的局部选择；状态通常是当前最远边界、最早结束时间、最小代价或剩余资源。

\textbf{证明和边界。} 证明从交换或领先性开始：任意最优解都能替换成当前贪心选择而不变差，或贪心状态始终不落后。边界：单字符、所有字符相同、字符多次跨段。变体：这是最早合法切分贪心，能最大化片段数量。

\subsection{自测标准}

二刷时不要只确认自己见过这道题。请遮住答案，先讲题面模型，再讲暴力基线和瓶颈，然后说出优化结构保存了什么、不变量如何排除错误候选、边界实验如何打破错误实现。

\section{区间、排序与合并：二刷复盘卡}
这一大节只围绕一个题型复盘，不把题号直接铺成目录。阅读时先看复盘目标，再读核心题卡，最后用自测标准检查自己是否能独立重讲。

\subsection{复盘目标}

追问通常会问排序键为什么这样、端点相等算不算重叠、比较器是否满足严格弱序。请准备说明排序后哪些候选可以被安全丢弃。

\subsection{核心题卡}

\noindent\textbf{LeetCode 56 合并区间。} 模型：排序后可能与当前区间相交的只会是结果尾部。推导重点：按起点排序，扫描时维护结果最后区间的右端点。

\textbf{二刷读题。} 读题时先定义端点语义和冲突条件。本题可以压成“排序后可能与当前区间相交的只会是结果尾部”，排序或有序表只是在利用端点关系。先遮住答案，只保留题面，复述候选对象、合法性条件和输出目标。

\textbf{暴力回放。} 慢版本对新区间和所有旧区间两两比较，手工判断相交、覆盖或冲突。它能验证端点语义，但排序后很多远离当前边界的区间无需再看。二刷时先写慢版本或口述慢版本，用它确认不会漏答案。

\textbf{特例回放。} 拿 [[1,3],[2,6],[8,10]] 手算，排序后 [2,6] 只可能和结果尾部 [1,3] 相交，合并成 [1,6]；[8,10] 的起点大于尾部右端，另开新区间。规律是排序后当前区间只需要看结果尾部，端点相接是否合并要按闭区间语义决定。把这个特例继续拆开看：排序以后，真正还能影响当前区间的候选必须贴近当前端点、结果尾部或资源堆顶。某个区间一旦被覆盖、合并或弹出，就要说明它为什么不会和后续区间重新产生更优关系。复盘时把这个特例压成状态字段：排序后的区间、当前结果尾、扫描右端或资源堆、端点比较规则；动态版本还要保存前驱后继；再用边界 端点相接、完全覆盖、无重叠、空列表、输入未排序 反查初始化、更新顺序和退出条件。

\textbf{状态回放。} 手算路线：手算时先排序，再逐个区间写当前结果尾部、当前右端或资源堆。每个区间离开视野时都要说明它不会再影响后续。状态字段：排序后的区间、当前结果尾、扫描右端或资源堆、端点比较规则；动态版本还要保存前驱后继。状态升级：把两两比较压成排序后的相邻关系、当前边界或动态有序结构；远离当前边界的区间要被安全归档。

\textbf{证明和边界。} 证明从排序后的邻接关系开始：已经合并、选择或丢弃的区间不会被更后面的区间重新影响。边界：端点相接、完全覆盖、无重叠、空列表、输入未排序。变体：若区间是半开区间，端点相等不一定合并。

\noindent\textbf{LeetCode 57 插入区间。} 模型：新区间只会影响与它相交的一段连续区间。推导重点：先追加左侧完全不相交区间，再合并重叠段，最后追加右侧。

\textbf{二刷读题。} 读题时先定义端点语义和冲突条件。本题可以压成“新区间只会影响与它相交的一段连续区间”，排序或有序表只是在利用端点关系。先遮住答案，只保留题面，复述候选对象、合法性条件和输出目标。

\textbf{暴力回放。} 慢版本对新区间和所有旧区间两两比较，手工判断相交、覆盖或冲突。它能验证端点语义，但排序后很多远离当前边界的区间无需再看。二刷时先写慢版本或口述慢版本，用它确认不会漏答案。

\textbf{特例回放。} 拿 intervals=[[1,3],[6,9]]、new=[2,5] 手算，左侧完全小于新区间的先放入；[1,3] 与 [2,5] 相交，合成 [1,5]；[6,9] 在右侧直接追加。规律是新区间影响的只是一段连续重叠区间，先左、再合并、再右，原列表为空或新区间覆盖全部也服从同一流程。把这个特例继续拆开看：排序以后，真正还能影响当前区间的候选必须贴近当前端点、结果尾部或资源堆顶。某个区间一旦被覆盖、合并或弹出，就要说明它为什么不会和后续区间重新产生更优关系。复盘时把这个特例压成状态字段：排序后的区间、当前结果尾、扫描右端或资源堆、端点比较规则；动态版本还要保存前驱后继；再用边界 新区间在最前、最后、覆盖全部、刚好相邻、原列表为空 反查初始化、更新顺序和退出条件。

\textbf{状态回放。} 手算路线：手算时先排序，再逐个区间写当前结果尾部、当前右端或资源堆。每个区间离开视野时都要说明它不会再影响后续。状态字段：排序后的区间、当前结果尾、扫描右端或资源堆、端点比较规则；动态版本还要保存前驱后继。状态升级：把两两比较压成排序后的相邻关系、当前边界或动态有序结构；远离当前边界的区间要被安全归档。

\textbf{证明和边界。} 证明从排序后的邻接关系开始：已经合并、选择或丢弃的区间不会被更后面的区间重新影响。边界：新区间在最前、最后、覆盖全部、刚好相邻、原列表为空。变体：若输入不是有序且不重叠，必须先排序或换模型。

\noindent\textbf{LeetCode 252 会议室。} 模型：每个会议占用一个时间区间，任意重叠都表示不能全部参加。推导重点：按开始时间排序，只需检查相邻会议是否冲突。

\textbf{二刷读题。} 读题时先定义端点语义和冲突条件。本题可以压成“每个会议占用一个时间区间，任意重叠都表示不能全部参加”，排序或有序表只是在利用端点关系。先遮住答案，只保留题面，复述候选对象、合法性条件和输出目标。

\textbf{暴力回放。} 慢版本对新区间和所有旧区间两两比较，手工判断相交、覆盖或冲突。它能验证端点语义，但排序后很多远离当前边界的区间无需再看。二刷时先写慢版本或口述慢版本，用它确认不会漏答案。

\textbf{特例回放。} 拿 [[0,30],[5,10],[15,20]] 手算，排序后 [5,10] 的开始 5 小于上一场结束 30，会议冲突。规律是按开始时间排序后，只需要检查相邻会议是否重叠；端点相等是否冲突取决于半开区间语义。把这个特例继续拆开看：排序以后，真正还能影响当前区间的候选必须贴近当前端点、结果尾部或资源堆顶。某个区间一旦被覆盖、合并或弹出，就要说明它为什么不会和后续区间重新产生更优关系。复盘时把这个特例压成状态字段：排序后的区间、当前结果尾、扫描右端或资源堆、端点比较规则；动态版本还要保存前驱后继；再用边界 端点相等是否冲突、空列表、单会议、完全重叠 反查初始化、更新顺序和退出条件。

\textbf{状态回放。} 手算路线：手算时先排序，再逐个区间写当前结果尾部、当前右端或资源堆。每个区间离开视野时都要说明它不会再影响后续。状态字段：排序后的区间、当前结果尾、扫描右端或资源堆、端点比较规则；动态版本还要保存前驱后继。状态升级：把两两比较压成排序后的相邻关系、当前边界或动态有序结构；远离当前边界的区间要被安全归档。

\textbf{证明和边界。} 证明从排序后的邻接关系开始：已经合并、选择或丢弃的区间不会被更后面的区间重新影响。边界：端点相等是否冲突、空列表、单会议、完全重叠。变体：若问最少会议室，需要维护同时进行的会议数量。

\noindent\textbf{LeetCode 253 会议室 II。} 模型：同时进行的会议数量决定最少房间数。推导重点：按开始时间扫描，用小顶堆维护当前会议的最早结束时间。

\textbf{二刷读题。} 读题时先定义端点语义和冲突条件。本题可以压成“同时进行的会议数量决定最少房间数”，排序或有序表只是在利用端点关系。先遮住答案，只保留题面，复述候选对象、合法性条件和输出目标。

\textbf{暴力回放。} 慢版本对新区间和所有旧区间两两比较，手工判断相交、覆盖或冲突。它能验证端点语义，但排序后很多远离当前边界的区间无需再看。二刷时先写慢版本或口述慢版本，用它确认不会漏答案。

\textbf{特例回放。} 拿 [[0,30],[5,10],[15,20]] 手算，0-30 占一间，5-10 来时最早结束仍是 30，所以需要第二间；15-20 来时 10 已结束，可复用。规律是小顶堆保存当前占用会议的最早结束时间，堆大小就是房间数。把这个特例继续拆开看：排序以后，真正还能影响当前区间的候选必须贴近当前端点、结果尾部或资源堆顶。某个区间一旦被覆盖、合并或弹出，就要说明它为什么不会和后续区间重新产生更优关系。复盘时把这个特例压成状态字段：排序后的区间、当前结果尾、扫描右端或资源堆、端点比较规则；动态版本还要保存前驱后继；再用边界 会议端点相接、多个同一时间开始、空列表、堆顶释放条件 反查初始化、更新顺序和退出条件。

\textbf{状态回放。} 手算路线：手算时先排序，再逐个区间写当前结果尾部、当前右端或资源堆。每个区间离开视野时都要说明它不会再影响后续。状态字段：排序后的区间、当前结果尾、扫描右端或资源堆、端点比较规则；动态版本还要保存前驱后继。状态升级：把两两比较压成排序后的相邻关系、当前边界或动态有序结构；远离当前边界的区间要被安全归档。

\textbf{证明和边界。} 证明从排序后的邻接关系开始：已经合并、选择或丢弃的区间不会被更后面的区间重新影响。边界：会议端点相接、多个同一时间开始、空列表、堆顶释放条件。变体：也可以把开始和结束分别排序，用双指针统计并发数。

\noindent\textbf{LeetCode 352 将数据流变为多个不相交区间。} 模型：数字逐个到达，需要动态维护已覆盖的有序不相交区间。推导重点：用有序表找到新值附近的前驱和后继，判断是否连接、合并或新建区间。

\textbf{二刷读题。} 读题时先定义端点语义和冲突条件。本题可以压成“数字逐个到达，需要动态维护已覆盖的有序不相交区间”，排序或有序表只是在利用端点关系。先遮住答案，只保留题面，复述候选对象、合法性条件和输出目标。

\textbf{暴力回放。} 慢版本对新区间和所有旧区间两两比较，手工判断相交、覆盖或冲突。它能验证端点语义，但排序后很多远离当前边界的区间无需再看。二刷时先写慢版本或口述慢版本，用它确认不会漏答案。

\textbf{特例回放。} 拿数据流依次加入 1、3、7、2、6 手算，加入 2 时把 [1,1] 和 [3,3] 连成 [1,3]；加入 6 时和 [7,7] 连成 [6,7]。规律是有序表只需要看新值的前驱和后继，重复插入不能改变区间。把这个特例继续拆开看：排序以后，真正还能影响当前区间的候选必须贴近当前端点、结果尾部或资源堆顶。某个区间一旦被覆盖、合并或弹出，就要说明它为什么不会和后续区间重新产生更优关系。复盘时把这个特例压成状态字段：排序后的区间、当前结果尾、扫描右端或资源堆、端点比较规则；动态版本还要保存前驱后继；再用边界 重复插入、连接左右两个区间、插入最小最大值、相邻端点 反查初始化、更新顺序和退出条件。

\textbf{状态回放。} 手算路线：手算时先排序，再逐个区间写当前结果尾部、当前右端或资源堆。每个区间离开视野时都要说明它不会再影响后续。状态字段：排序后的区间、当前结果尾、扫描右端或资源堆、端点比较规则；动态版本还要保存前驱后继。状态升级：把两两比较压成排序后的相邻关系、当前边界或动态有序结构；远离当前边界的区间要被安全归档。

\textbf{证明和边界。} 证明从排序后的邻接关系开始：已经合并、选择或丢弃的区间不会被更后面的区间重新影响。边界：重复插入、连接左右两个区间、插入最小最大值、相邻端点。变体：这类动态区间不能用普通数组扫描长期维护。

\noindent\textbf{LeetCode 436 寻找右区间。} 模型：每个区间需要寻找起点不小于当前终点的最小起点区间。推导重点：把所有起点和原下标排序，对每个终点做 lower bound。

\textbf{二刷读题。} 读题时先定义端点语义和冲突条件。本题可以压成“每个区间需要寻找起点不小于当前终点的最小起点区间”，排序或有序表只是在利用端点关系。先遮住答案，只保留题面，复述候选对象、合法性条件和输出目标。

\textbf{暴力回放。} 慢版本对新区间和所有旧区间两两比较，手工判断相交、覆盖或冲突。它能验证端点语义，但排序后很多远离当前边界的区间无需再看。二刷时先写慢版本或口述慢版本，用它确认不会漏答案。

\textbf{特例回放。} 拿 [[1,2],[2,3],[3,4]] 手算，区间 [1,2] 的右区间是起点最小且不小于 2 的 [2,3]。规律是把所有起点排序，对每个终点做 lower\_bound；不存在时返回 -1。把这个特例继续拆开看：排序以后，真正还能影响当前区间的候选必须贴近当前端点、结果尾部或资源堆顶。某个区间一旦被覆盖、合并或弹出，就要说明它为什么不会和后续区间重新产生更优关系。复盘时把这个特例压成状态字段：排序后的区间、当前结果尾、扫描右端或资源堆、端点比较规则；动态版本还要保存前驱后继；再用边界 不存在右区间、起点相同、负端点、单区间 反查初始化、更新顺序和退出条件。

\textbf{状态回放。} 手算路线：手算时先排序，再逐个区间写当前结果尾部、当前右端或资源堆。每个区间离开视野时都要说明它不会再影响后续。状态字段：排序后的区间、当前结果尾、扫描右端或资源堆、端点比较规则；动态版本还要保存前驱后继。状态升级：把两两比较压成排序后的相邻关系、当前边界或动态有序结构；远离当前边界的区间要被安全归档。

\textbf{证明和边界。} 证明从排序后的邻接关系开始：已经合并、选择或丢弃的区间不会被更后面的区间重新影响。边界：不存在右区间、起点相同、负端点、单区间。变体：如果区间动态插入，则需要有序表而不是一次排序数组。

\noindent\textbf{LeetCode 729 我的日程安排表 I。} 模型：每次插入新区间时，需要判断它是否和已有区间重叠。推导重点：有序表按起点存区间，只检查新区间的前驱和后继。

\textbf{二刷读题。} 读题时先定义端点语义和冲突条件。本题可以压成“每次插入新区间时，需要判断它是否和已有区间重叠”，排序或有序表只是在利用端点关系。先遮住答案，只保留题面，复述候选对象、合法性条件和输出目标。

\textbf{暴力回放。} 慢版本对新区间和所有旧区间两两比较，手工判断相交、覆盖或冲突。它能验证端点语义，但排序后很多远离当前边界的区间无需再看。二刷时先写慢版本或口述慢版本，用它确认不会漏答案。

\textbf{特例回放。} 拿 book(10,20) 成功，再 book(15,25) 手算，它和已有区间重叠应失败；book(20,30) 与 [10,20) 端点相接可成功。规律是半开区间下只需检查前驱和后继是否重叠。把这个特例继续拆开看：排序以后，真正还能影响当前区间的候选必须贴近当前端点、结果尾部或资源堆顶。某个区间一旦被覆盖、合并或弹出，就要说明它为什么不会和后续区间重新产生更优关系。复盘时把这个特例压成状态字段：排序后的区间、当前结果尾、扫描右端或资源堆、端点比较规则；动态版本还要保存前驱后继；再用边界 端点相接、插入最前最后、完全覆盖已有区间、动态多次插入 反查初始化、更新顺序和退出条件。

\textbf{状态回放。} 手算路线：手算时先排序，再逐个区间写当前结果尾部、当前右端或资源堆。每个区间离开视野时都要说明它不会再影响后续。状态字段：排序后的区间、当前结果尾、扫描右端或资源堆、端点比较规则；动态版本还要保存前驱后继。状态升级：把两两比较压成排序后的相邻关系、当前边界或动态有序结构；远离当前边界的区间要被安全归档。

\textbf{证明和边界。} 证明从排序后的邻接关系开始：已经合并、选择或丢弃的区间不会被更后面的区间重新影响。边界：端点相接、插入最前最后、完全覆盖已有区间、动态多次插入。变体：哈希表无法回答前驱后继，动态区间要用有序结构。

\noindent\textbf{LeetCode 731 我的日程安排表 II。} 模型：允许双重预订但不允许三重预订，插入会改变区间重叠层数。推导重点：保存已有预订和已经双重预订的区间；新预订若碰到双重区间则失败，否则记录新产生的重叠。

\textbf{二刷读题。} 读题时先定义端点语义和冲突条件。本题可以压成“允许双重预订但不允许三重预订，插入会改变区间重叠层数”，排序或有序表只是在利用端点关系。先遮住答案，只保留题面，复述候选对象、合法性条件和输出目标。

\textbf{暴力回放。} 慢版本对新区间和所有旧区间两两比较，手工判断相交、覆盖或冲突。它能验证端点语义，但排序后很多远离当前边界的区间无需再看。二刷时先写慢版本或口述慢版本，用它确认不会漏答案。

\textbf{特例回放。} 拿 book(10,20)、book(15,25) 后，重叠段 [15,20) 已经是双重预订；再 book(17,22) 会碰到双重段，必须失败。规律是保存所有单次预订和已形成的双重预订，新区间不能与任何双重区间相交。把这个特例继续拆开看：排序以后，真正还能影响当前区间的候选必须贴近当前端点、结果尾部或资源堆顶。某个区间一旦被覆盖、合并或弹出，就要说明它为什么不会和后续区间重新产生更优关系。复盘时把这个特例压成状态字段：排序后的区间、当前结果尾、扫描右端或资源堆、端点比较规则；动态版本还要保存前驱后继；再用边界 端点相接、多个局部重叠、完全覆盖、回滚失败插入 反查初始化、更新顺序和退出条件。

\textbf{状态回放。} 手算路线：手算时先排序，再逐个区间写当前结果尾部、当前右端或资源堆。每个区间离开视野时都要说明它不会再影响后续。状态字段：排序后的区间、当前结果尾、扫描右端或资源堆、端点比较规则；动态版本还要保存前驱后继。状态升级：把两两比较压成排序后的相邻关系、当前边界或动态有序结构；远离当前边界的区间要被安全归档。

\textbf{证明和边界。} 证明从排序后的邻接关系开始：已经合并、选择或丢弃的区间不会被更后面的区间重新影响。边界：端点相接、多个局部重叠、完全覆盖、回滚失败插入。变体：扫描线差分也能做，但要处理失败时恢复状态。

\noindent\textbf{LeetCode 986 区间列表的交集。} 模型：两个有序且不重叠的区间列表，交集只可能来自当前两个区间。推导重点：每次比较两个当前区间的重叠部分，然后推进右端较小的区间。

\textbf{二刷读题。} 读题时先定义端点语义和冲突条件。本题可以压成“两个有序且不重叠的区间列表，交集只可能来自当前两个区间”，排序或有序表只是在利用端点关系。先遮住答案，只保留题面，复述候选对象、合法性条件和输出目标。

\textbf{暴力回放。} 慢版本对新区间和所有旧区间两两比较，手工判断相交、覆盖或冲突。它能验证端点语义，但排序后很多远离当前边界的区间无需再看。二刷时先写慢版本或口述慢版本，用它确认不会漏答案。

\textbf{特例回放。} 拿 A=[0,2],[5,10] 和 B=[1,5],[8,12] 手算，当前两个区间交集是 [1,2]，然后推进右端较小的 [0,2]。规律是两个有序列表只需要比较当前区间，谁先结束谁不可能再和对方后续产生交集。把这个特例继续拆开看：排序以后，真正还能影响当前区间的候选必须贴近当前端点、结果尾部或资源堆顶。某个区间一旦被覆盖、合并或弹出，就要说明它为什么不会和后续区间重新产生更优关系。复盘时把这个特例压成状态字段：排序后的区间、当前结果尾、扫描右端或资源堆、端点比较规则；动态版本还要保存前驱后继；再用边界 空列表、端点相接、一个区间覆盖多个区间、无交集 反查初始化、更新顺序和退出条件。

\textbf{状态回放。} 手算路线：手算时先排序，再逐个区间写当前结果尾部、当前右端或资源堆。每个区间离开视野时都要说明它不会再影响后续。状态字段：排序后的区间、当前结果尾、扫描右端或资源堆、端点比较规则；动态版本还要保存前驱后继。状态升级：把两两比较压成排序后的相邻关系、当前边界或动态有序结构；远离当前边界的区间要被安全归档。

\textbf{证明和边界。} 证明从排序后的邻接关系开始：已经合并、选择或丢弃的区间不会被更后面的区间重新影响。边界：空列表、端点相接、一个区间覆盖多个区间、无交集。变体：这是区间双指针，不需要把两个列表合并排序。

\noindent\textbf{LeetCode 1094 拼车。} 模型：乘客上下车事件构成沿路线位置变化的载客量。推导重点：用差分数组或扫描线记录每个位置人数变化，累计人数不能超过容量。

\textbf{二刷读题。} 读题时先定义端点语义和冲突条件。本题可以压成“乘客上下车事件构成沿路线位置变化的载客量”，排序或有序表只是在利用端点关系。先遮住答案，只保留题面，复述候选对象、合法性条件和输出目标。

\textbf{暴力回放。} 慢版本对新区间和所有旧区间两两比较，手工判断相交、覆盖或冲突。它能验证端点语义，但排序后很多远离当前边界的区间无需再看。二刷时先写慢版本或口述慢版本，用它确认不会漏答案。

\textbf{特例回放。} 拿 trips=[[2,1,5],[3,3,7]]、capacity=4 手算，位置 3 到 5 之间车上人数是 5，超过容量。规律是上车点加人数、下车点减人数，按位置前缀累计检查容量。把这个特例继续拆开看：排序以后，真正还能影响当前区间的候选必须贴近当前端点、结果尾部或资源堆顶。某个区间一旦被覆盖、合并或弹出，就要说明它为什么不会和后续区间重新产生更优关系。复盘时把这个特例压成状态字段：排序后的区间、当前结果尾、扫描右端或资源堆、端点比较规则；动态版本还要保存前驱后继；再用边界 同站上下车、容量刚好、站点范围、输入未排序 反查初始化、更新顺序和退出条件。

\textbf{状态回放。} 手算路线：手算时先排序，再逐个区间写当前结果尾部、当前右端或资源堆。每个区间离开视野时都要说明它不会再影响后续。状态字段：排序后的区间、当前结果尾、扫描右端或资源堆、端点比较规则；动态版本还要保存前驱后继。状态升级：把两两比较压成排序后的相邻关系、当前边界或动态有序结构；远离当前边界的区间要被安全归档。

\textbf{证明和边界。} 证明从排序后的邻接关系开始：已经合并、选择或丢弃的区间不会被更后面的区间重新影响。边界：同站上下车、容量刚好、站点范围、输入未排序。变体：航班预订统计也是差分思想，只是只需要最终每段累加结果。

\noindent\textbf{LeetCode 1109 航班预订统计。} 模型：每条预订给一个连续航班区间增加座位数。推导重点：差分数组在区间起点加人数，终点后一位减人数，最后前缀还原每航班总数。

\textbf{二刷读题。} 读题时先定义端点语义和冲突条件。本题可以压成“每条预订给一个连续航班区间增加座位数”，排序或有序表只是在利用端点关系。先遮住答案，只保留题面，复述候选对象、合法性条件和输出目标。

\textbf{暴力回放。} 慢版本对新区间和所有旧区间两两比较，手工判断相交、覆盖或冲突。它能验证端点语义，但排序后很多远离当前边界的区间无需再看。二刷时先写慢版本或口述慢版本，用它确认不会漏答案。

\textbf{特例回放。} 拿 bookings=[[1,2,10],[2,3,20]]、n=3 手算，第 1 班加 10，第 2 班加 30，第 3 班加 20。规律是区间加法用差分：起点加，终点后一位减，最后前缀还原。把这个特例继续拆开看：排序以后，真正还能影响当前区间的候选必须贴近当前端点、结果尾部或资源堆顶。某个区间一旦被覆盖、合并或弹出，就要说明它为什么不会和后续区间重新产生更优关系。复盘时把这个特例压成状态字段：排序后的区间、当前结果尾、扫描右端或资源堆、端点比较规则；动态版本还要保存前驱后继；再用边界 区间到最后一个航班、单点区间、多条重叠预订、下标从一开始 反查初始化、更新顺序和退出条件。

\textbf{状态回放。} 手算路线：手算时先排序，再逐个区间写当前结果尾部、当前右端或资源堆。每个区间离开视野时都要说明它不会再影响后续。状态字段：排序后的区间、当前结果尾、扫描右端或资源堆、端点比较规则；动态版本还要保存前驱后继。状态升级：把两两比较压成排序后的相邻关系、当前边界或动态有序结构；远离当前边界的区间要被安全归档。

\textbf{证明和边界。} 证明从排序后的邻接关系开始：已经合并、选择或丢弃的区间不会被更后面的区间重新影响。边界：区间到最后一个航班、单点区间、多条重叠预订、下标从一开始。变体：这是静态区间加法，不需要线段树。

\noindent\textbf{LeetCode 1229 安排会议日程。} 模型：两个人的可用时间段都有序，目标是找最早长度足够的交集。推导重点：双指针比较当前两个区间交集，若长度不足就推进结束更早的一侧。

\textbf{二刷读题。} 读题时先定义端点语义和冲突条件。本题可以压成“两个人的可用时间段都有序，目标是找最早长度足够的交集”，排序或有序表只是在利用端点关系。先遮住答案，只保留题面，复述候选对象、合法性条件和输出目标。

\textbf{暴力回放。} 慢版本对新区间和所有旧区间两两比较，手工判断相交、覆盖或冲突。它能验证端点语义，但排序后很多远离当前边界的区间无需再看。二刷时先写慢版本或口述慢版本，用它确认不会漏答案。

\textbf{特例回放。} 拿 slots1=[10,50],[60,120]，slots2=[0,15],[60,70]，duration=8 手算，第一组交集 [10,15] 不够，第二组交集 [60,70] 足够，返回 [60,68]。规律是双指针比较当前交集，不够时推进结束更早的一侧。把这个特例继续拆开看：排序以后，真正还能影响当前区间的候选必须贴近当前端点、结果尾部或资源堆顶。某个区间一旦被覆盖、合并或弹出，就要说明它为什么不会和后续区间重新产生更优关系。复盘时把这个特例压成状态字段：排序后的区间、当前结果尾、扫描右端或资源堆、端点比较规则；动态版本还要保存前驱后继；再用边界 无可行时间、刚好等于 duration、区间端点相接、输入需要先排序 反查初始化、更新顺序和退出条件。

\textbf{状态回放。} 手算路线：手算时先排序，再逐个区间写当前结果尾部、当前右端或资源堆。每个区间离开视野时都要说明它不会再影响后续。状态字段：排序后的区间、当前结果尾、扫描右端或资源堆、端点比较规则；动态版本还要保存前驱后继。状态升级：把两两比较压成排序后的相邻关系、当前边界或动态有序结构；远离当前边界的区间要被安全归档。

\textbf{证明和边界。} 证明从排序后的邻接关系开始：已经合并、选择或丢弃的区间不会被更后面的区间重新影响。边界：无可行时间、刚好等于 duration、区间端点相接、输入需要先排序。变体：它和区间列表交集相同，但额外带最早可行和长度约束。

\noindent\textbf{LeetCode 1288 删除被覆盖区间。} 模型：被覆盖要求另一区间起点不晚且终点不早。推导重点：按起点升序、终点降序排序，扫描维护最大右端。

\textbf{二刷读题。} 读题时先定义端点语义和冲突条件。本题可以压成“被覆盖要求另一区间起点不晚且终点不早”，排序或有序表只是在利用端点关系。先遮住答案，只保留题面，复述候选对象、合法性条件和输出目标。

\textbf{暴力回放。} 慢版本对新区间和所有旧区间两两比较，手工判断相交、覆盖或冲突。它能验证端点语义，但排序后很多远离当前边界的区间无需再看。二刷时先写慢版本或口述慢版本，用它确认不会漏答案。

\textbf{特例回放。} 拿 [[1,4],[3,6],[2,8]] 手算，[3,6] 被 [2,8] 覆盖，应删除；同起点时要先看更长区间，否则短区间会干扰判断。规律是起点升序、终点降序排序，扫描维护最大右端。把这个特例继续拆开看：排序以后，真正还能影响当前区间的候选必须贴近当前端点、结果尾部或资源堆顶。某个区间一旦被覆盖、合并或弹出，就要说明它为什么不会和后续区间重新产生更优关系。复盘时把这个特例压成状态字段：排序后的区间、当前结果尾、扫描右端或资源堆、端点比较规则；动态版本还要保存前驱后继；再用边界 同起点区间、完全覆盖、无覆盖、端点相等 反查初始化、更新顺序和退出条件。

\textbf{状态回放。} 手算路线：手算时先排序，再逐个区间写当前结果尾部、当前右端或资源堆。每个区间离开视野时都要说明它不会再影响后续。状态字段：排序后的区间、当前结果尾、扫描右端或资源堆、端点比较规则；动态版本还要保存前驱后继。状态升级：把两两比较压成排序后的相邻关系、当前边界或动态有序结构；远离当前边界的区间要被安全归档。

\textbf{证明和边界。} 证明从排序后的邻接关系开始：已经合并、选择或丢弃的区间不会被更后面的区间重新影响。边界：同起点区间、完全覆盖、无覆盖、端点相等。变体：终点降序是关键，否则同起点小区间会先出现造成误判。

\subsection{自测标准}

二刷时不要只确认自己见过这道题。请遮住答案，先讲题面模型，再讲暴力基线和瓶颈，然后说出优化结构保存了什么、不变量如何排除错误候选、边界实验如何打破错误实现。

\section{位运算与状态压缩：二刷复盘卡}
这一大节只围绕一个题型复盘，不把题号直接铺成目录。阅读时先看复盘目标，再读核心题卡，最后用自测标准检查自己是否能独立重讲。

\subsection{复盘目标}

追问通常会问每一位代表什么、负数和移位是否安全、状态相同但资源不同能否合并。请准备说明位运算只是编码，真正关键仍是状态语义。

\subsection{核心题卡}

\noindent\textbf{LeetCode 29 两数相除。} 模型：除法可以看成不断减去除数的倍增值。推导重点：把除数按二进制倍增，优先减去不超过被除数的最大倍数并累加商。

\textbf{二刷读题。} 读题时先用普通状态解释每一项布尔信息。本题可以压成“除法可以看成不断减去除数的倍增值”，位运算只是更紧凑的编码。先遮住答案，只保留题面，复述候选对象、合法性条件和输出目标。

\textbf{暴力回放。} 慢版本先用集合、数组或布尔表保存状态，逐步执行题目动作。确认每个布尔字段的含义后，才能把它压缩成位，而不是直接写位运算公式。二刷时先写慢版本或口述慢版本，用它确认不会漏答案。

\textbf{特例回放。} 拿 43 除以 5 手算，5 的倍增序列是 5、10、20、40，先减 40 商加 8，余 3 停止。规律是用二进制倍增找不超过被除数的最大除数倍数，符号和 int 最小值溢出要单独处理。把这个特例继续拆开看：每一位或每个掩码位都要对应题目里的一个布尔事实。位运算只是压缩表示；如果两个状态在位置相同但掩码不同，未来能力不同，就不能合并。复盘时把这个特例压成状态字段：掩码含义、当前位、目标位集合、状态转移和答案读取；负数或高位参与时要说明整数宽度；再用边界 除数为一、被除数为 int 最小值、结果溢出、符号处理 反查初始化、更新顺序和退出条件。

\textbf{状态回放。} 手算路线：手算时把整数写成二进制，并在每一位旁边标注含义。状态转移只允许改变题目动作允许改变的那些位。状态字段：掩码含义、当前位、目标位集合、状态转移和答案读取；负数或高位参与时要说明整数宽度。状态升级：把集合状态压成掩码；包含、加入、删除、枚举子集都变成位操作，但每一位的题目含义不能丢。

\textbf{证明和边界。} 证明从逐位独立或子集归纳开始：被压缩到位里的信息足以决定未来转移。边界：除数为一、被除数为 int 最小值、结果溢出、符号处理。变体：它和快速幂类似，都利用二进制拆分减少重复加减。

\noindent\textbf{LeetCode 136 只出现一次的数字。} 模型：成对元素可以用异或抵消，剩下单个元素。推导重点：异或满足交换结合，扫描顺序不影响答案。

\textbf{二刷读题。} 读题时先用普通状态解释每一项布尔信息。本题可以压成“成对元素可以用异或抵消，剩下单个元素”，位运算只是更紧凑的编码。先遮住答案，只保留题面，复述候选对象、合法性条件和输出目标。

\textbf{暴力回放。} 慢版本先用集合、数组或布尔表保存状态，逐步执行题目动作。确认每个布尔字段的含义后，才能把它压缩成位，而不是直接写位运算公式。二刷时先写慢版本或口述慢版本，用它确认不会漏答案。

\textbf{特例回放。} 拿 [4,1,2,1,2] 手算，1 和 1 异或抵消，2 和 2 抵消，最后剩 4。规律是异或满足交换结合，且相同数异或后变成零，扫描顺序不影响结果；零和负数也只是按位参与同样运算。把这个特例继续拆开看：每一位或每个掩码位都要对应题目里的一个布尔事实。位运算只是压缩表示；如果两个状态在位置相同但掩码不同，未来能力不同，就不能合并。复盘时把这个特例压成状态字段：掩码含义、当前位、目标位集合、状态转移和答案读取；负数或高位参与时要说明整数宽度；再用边界 零、负数、唯一值在开头或结尾、所有其他值恰好两次 反查初始化、更新顺序和退出条件。

\textbf{状态回放。} 手算路线：手算时把整数写成二进制，并在每一位旁边标注含义。状态转移只允许改变题目动作允许改变的那些位。状态字段：掩码含义、当前位、目标位集合、状态转移和答案读取；负数或高位参与时要说明整数宽度。状态升级：把集合状态压成掩码；包含、加入、删除、枚举子集都变成位操作，但每一位的题目含义不能丢。

\textbf{证明和边界。} 证明从逐位独立或子集归纳开始：被压缩到位里的信息足以决定未来转移。边界：零、负数、唯一值在开头或结尾、所有其他值恰好两次。变体：若其他值出现三次，要改为按位计数取模。

\noindent\textbf{LeetCode 137 只出现一次的数字 II。} 模型：除一个数字外，其余数字都出现三次，单纯异或不能抵消。推导重点：逐位统计一的数量，对三取模后还原目标数字。

\textbf{二刷读题。} 读题时先用普通状态解释每一项布尔信息。本题可以压成“除一个数字外，其余数字都出现三次，单纯异或不能抵消”，位运算只是更紧凑的编码。先遮住答案，只保留题面，复述候选对象、合法性条件和输出目标。

\textbf{暴力回放。} 慢版本先用集合、数组或布尔表保存状态，逐步执行题目动作。确认每个布尔字段的含义后，才能把它压缩成位，而不是直接写位运算公式。二刷时先写慢版本或口述慢版本，用它确认不会漏答案。

\textbf{特例回放。} 拿 [2,2,3,2] 手算，每一位上 2 的贡献出现三次，对 3 取模后消失，只剩 3 的位。规律是逐位计数并对 3 取模，符号位也必须按固定整数宽度处理。把这个特例继续拆开看：每一位或每个掩码位都要对应题目里的一个布尔事实。位运算只是压缩表示；如果两个状态在位置相同但掩码不同，未来能力不同，就不能合并。复盘时把这个特例压成状态字段：掩码含义、当前位、目标位集合、状态转移和答案读取；负数或高位参与时要说明整数宽度；再用边界 负数、符号位、目标为零、所有其他数恰好三次 反查初始化、更新顺序和退出条件。

\textbf{状态回放。} 手算路线：手算时把整数写成二进制，并在每一位旁边标注含义。状态转移只允许改变题目动作允许改变的那些位。状态字段：掩码含义、当前位、目标位集合、状态转移和答案读取；负数或高位参与时要说明整数宽度。状态升级：把集合状态压成掩码；包含、加入、删除、枚举子集都变成位操作，但每一位的题目含义不能丢。

\textbf{证明和边界。} 证明从逐位独立或子集归纳开始：被压缩到位里的信息足以决定未来转移。边界：负数、符号位、目标为零、所有其他数恰好三次。变体：若其他数出现 k 次，思路仍是按位计数取模。

\noindent\textbf{LeetCode 190 颠倒二进制位。} 模型：每一位的新位置由原位置镜像决定。推导重点：循环三十二次，把结果左移并接上当前最低位，再右移输入。

\textbf{二刷读题。} 读题时先用普通状态解释每一项布尔信息。本题可以压成“每一位的新位置由原位置镜像决定”，位运算只是更紧凑的编码。先遮住答案，只保留题面，复述候选对象、合法性条件和输出目标。

\textbf{暴力回放。} 慢版本先用集合、数组或布尔表保存状态，逐步执行题目动作。确认每个布尔字段的含义后，才能把它压缩成位，而不是直接写位运算公式。二刷时先写慢版本或口述慢版本，用它确认不会漏答案。

\textbf{特例回放。} 拿低位是 101 的数手算，颠倒时每次把结果左移一位，再接上当前最低位；原数右移继续取下一位。规律是循环 32 次处理固定宽度，不能因为高位暂时为 0 就提前结束。把这个特例继续拆开看：每一位或每个掩码位都要对应题目里的一个布尔事实。位运算只是压缩表示；如果两个状态在位置相同但掩码不同，未来能力不同，就不能合并。复盘时把这个特例压成状态字段：掩码含义、当前位、目标位集合、状态转移和答案读取；负数或高位参与时要说明整数宽度；再用边界 最高位、最低位、全零、全一、无符号语义 反查初始化、更新顺序和退出条件。

\textbf{状态回放。} 手算路线：手算时把整数写成二进制，并在每一位旁边标注含义。状态转移只允许改变题目动作允许改变的那些位。状态字段：掩码含义、当前位、目标位集合、状态转移和答案读取；负数或高位参与时要说明整数宽度。状态升级：把集合状态压成掩码；包含、加入、删除、枚举子集都变成位操作，但每一位的题目含义不能丢。

\textbf{证明和边界。} 证明从逐位独立或子集归纳开始：被压缩到位里的信息足以决定未来转移。边界：最高位、最低位、全零、全一、无符号语义。变体：若多次调用，可以按字节预处理反转表。

\noindent\textbf{LeetCode 191 位 1 的个数。} 模型：每次清掉最低位的一可以让循环次数等于一的个数。推导重点：使用 n \& (n - 1) 删除最低位一，比逐位检查更直接。

\textbf{二刷读题。} 读题时先用普通状态解释每一项布尔信息。本题可以压成“每次清掉最低位的一可以让循环次数等于一的个数”，位运算只是更紧凑的编码。先遮住答案，只保留题面，复述候选对象、合法性条件和输出目标。

\textbf{暴力回放。} 慢版本先用集合、数组或布尔表保存状态，逐步执行题目动作。确认每个布尔字段的含义后，才能把它压缩成位，而不是直接写位运算公式。二刷时先写慢版本或口述慢版本，用它确认不会漏答案。

\textbf{特例回放。} 拿 11 的二进制 1011 手算，n\&(n-1) 第一次删掉最低位 1 得 1010，第二次得 1000，第三次得 0。规律是每次操作恰好删除一个 1，所以循环次数就是 1 的个数。把这个特例继续拆开看：每一位或每个掩码位都要对应题目里的一个布尔事实。位运算只是压缩表示；如果两个状态在位置相同但掩码不同，未来能力不同，就不能合并。复盘时把这个特例压成状态字段：掩码含义、当前位、目标位集合、状态转移和答案读取；负数或高位参与时要说明整数宽度；再用边界 零、只有最高位、一的个数很多、无符号输入 反查初始化、更新顺序和退出条件。

\textbf{状态回放。} 手算路线：手算时把整数写成二进制，并在每一位旁边标注含义。状态转移只允许改变题目动作允许改变的那些位。状态字段：掩码含义、当前位、目标位集合、状态转移和答案读取；负数或高位参与时要说明整数宽度。状态升级：把集合状态压成掩码；包含、加入、删除、枚举子集都变成位操作，但每一位的题目含义不能丢。

\textbf{证明和边界。} 证明从逐位独立或子集归纳开始：被压缩到位里的信息足以决定未来转移。边界：零、只有最高位、一的个数很多、无符号输入。变体：比特位计数可以把这个过程 DP 化到一整段数。

\noindent\textbf{LeetCode 201 数字范围按位与。} 模型：区间内所有数按位与后，只保留左右端公共二进制前缀。推导重点：不断右移 left 和 right，直到二者相等，再把公共前缀移回原位。

\textbf{二刷读题。} 读题时先用普通状态解释每一项布尔信息。本题可以压成“区间内所有数按位与后，只保留左右端公共二进制前缀”，位运算只是更紧凑的编码。先遮住答案，只保留题面，复述候选对象、合法性条件和输出目标。

\textbf{暴力回放。} 慢版本先用集合、数组或布尔表保存状态，逐步执行题目动作。确认每个布尔字段的含义后，才能把它压缩成位，而不是直接写位运算公式。二刷时先写慢版本或口述慢版本，用它确认不会漏答案。

\textbf{特例回放。} 拿区间 [5,7]，二进制是 101、110、111，只有最高公共前缀 1 能留下，低位在区间内会出现 0 和 1，被按位与清零。规律是不断右移 left 和 right 找公共前缀，再移回原位。把这个特例继续拆开看：每一位或每个掩码位都要对应题目里的一个布尔事实。位运算只是压缩表示；如果两个状态在位置相同但掩码不同，未来能力不同，就不能合并。复盘时把这个特例压成状态字段：掩码含义、当前位、目标位集合、状态转移和答案读取；负数或高位参与时要说明整数宽度；再用边界 left 等于 right、区间跨过二的幂边界、结果为零 反查初始化、更新顺序和退出条件。

\textbf{状态回放。} 手算路线：手算时把整数写成二进制，并在每一位旁边标注含义。状态转移只允许改变题目动作允许改变的那些位。状态字段：掩码含义、当前位、目标位集合、状态转移和答案读取；负数或高位参与时要说明整数宽度。状态升级：把集合状态压成掩码；包含、加入、删除、枚举子集都变成位操作，但每一位的题目含义不能丢。

\textbf{证明和边界。} 证明从逐位独立或子集归纳开始：被压缩到位里的信息足以决定未来转移。边界：left 等于 right、区间跨过二的幂边界、结果为零。变体：也可以不断清除 right 的最低位一，直到 right 不大于 left。

\noindent\textbf{LeetCode 231 2 的幂。} 模型：二的幂在二进制中只有一个一。推导重点：正数 n 满足 n 与 n 减一按位与为零。

\textbf{二刷读题。} 读题时先用普通状态解释每一项布尔信息。本题可以压成“二的幂在二进制中只有一个一”，位运算只是更紧凑的编码。先遮住答案，只保留题面，复述候选对象、合法性条件和输出目标。

\textbf{暴力回放。} 慢版本先用集合、数组或布尔表保存状态，逐步执行题目动作。确认每个布尔字段的含义后，才能把它压缩成位，而不是直接写位运算公式。二刷时先写慢版本或口述慢版本，用它确认不会漏答案。

\textbf{特例回放。} 拿 8 的二进制 1000 手算，8\&(7)=1000\&0111=0；拿 6 的 110，6\&5=100 不为 0。规律是 2 的幂只有一个 1，且 n 必须为正；0 和负数不能只靠按位与判断。把这个特例继续拆开看：每一位或每个掩码位都要对应题目里的一个布尔事实。位运算只是压缩表示；如果两个状态在位置相同但掩码不同，未来能力不同，就不能合并。复盘时把这个特例压成状态字段：掩码含义、当前位、目标位集合、状态转移和答案读取；负数或高位参与时要说明整数宽度；再用边界 n 为零、负数、int 最小值、最高位 反查初始化、更新顺序和退出条件。

\textbf{状态回放。} 手算路线：手算时把整数写成二进制，并在每一位旁边标注含义。状态转移只允许改变题目动作允许改变的那些位。状态字段：掩码含义、当前位、目标位集合、状态转移和答案读取；负数或高位参与时要说明整数宽度。状态升级：把集合状态压成掩码；包含、加入、删除、枚举子集都变成位操作，但每一位的题目含义不能丢。

\textbf{证明和边界。} 证明从逐位独立或子集归纳开始：被压缩到位里的信息足以决定未来转移。边界：n 为零、负数、int 最小值、最高位。变体：判断四的幂还要额外限制一所在的位置为偶数位。

\noindent\textbf{LeetCode 260 只出现一次的数字 III。} 模型：两个只出现一次的数异或后至少有一位不同。推导重点：用最低不同位把数组分成两组，每组内部再异或抵消。

\textbf{二刷读题。} 读题时先用普通状态解释每一项布尔信息。本题可以压成“两个只出现一次的数异或后至少有一位不同”，位运算只是更紧凑的编码。先遮住答案，只保留题面，复述候选对象、合法性条件和输出目标。

\textbf{暴力回放。} 慢版本先用集合、数组或布尔表保存状态，逐步执行题目动作。确认每个布尔字段的含义后，才能把它压缩成位，而不是直接写位运算公式。二刷时先写慢版本或口述慢版本，用它确认不会漏答案。

\textbf{特例回放。} 拿 [1,2,1,3,2,5] 手算，全部异或后只剩 3 和 5 的异或结果，取最低不同位可把 3 和 5 分到两组；每组内部重复数再异或抵消。规律是两个答案至少有一位不同，用这位分组后退化成两个只出现一次的问题。把这个特例继续拆开看：每一位或每个掩码位都要对应题目里的一个布尔事实。位运算只是压缩表示；如果两个状态在位置相同但掩码不同，未来能力不同，就不能合并。复盘时把这个特例压成状态字段：掩码含义、当前位、目标位集合、状态转移和答案读取；负数或高位参与时要说明整数宽度；再用边界 两个答案一正一负、最低位分组、其他数恰好两次 反查初始化、更新顺序和退出条件。

\textbf{状态回放。} 手算路线：手算时把整数写成二进制，并在每一位旁边标注含义。状态转移只允许改变题目动作允许改变的那些位。状态字段：掩码含义、当前位、目标位集合、状态转移和答案读取；负数或高位参与时要说明整数宽度。状态升级：把集合状态压成掩码；包含、加入、删除、枚举子集都变成位操作，但每一位的题目含义不能丢。

\textbf{证明和边界。} 证明从逐位独立或子集归纳开始：被压缩到位里的信息足以决定未来转移。边界：两个答案一正一负、最低位分组、其他数恰好两次。变体：它是在异或抵消基础上增加分组位。

\noindent\textbf{LeetCode 268 丢失的数字。} 模型：下标零到 n 与数组值相比，缺失值会在异或抵消后剩下。推导重点：把所有下标和所有值一起异或，重复出现的数抵消。

\textbf{二刷读题。} 读题时先用普通状态解释每一项布尔信息。本题可以压成“下标零到 n 与数组值相比，缺失值会在异或抵消后剩下”，位运算只是更紧凑的编码。先遮住答案，只保留题面，复述候选对象、合法性条件和输出目标。

\textbf{暴力回放。} 慢版本先用集合、数组或布尔表保存状态，逐步执行题目动作。确认每个布尔字段的含义后，才能把它压缩成位，而不是直接写位运算公式。二刷时先写慢版本或口述慢版本，用它确认不会漏答案。

\textbf{特例回放。} 拿 [3,0,1] 手算，0、1、3 与完整下标 0、1、2、3 一起异或，重复的 0、1、3 抵消，剩 2。规律是缺失数可以由异或抵消得到，也可用求和但要注意溢出。把这个特例继续拆开看：每一位或每个掩码位都要对应题目里的一个布尔事实。位运算只是压缩表示；如果两个状态在位置相同但掩码不同，未来能力不同，就不能合并。复盘时把这个特例压成状态字段：掩码含义、当前位、目标位集合、状态转移和答案读取；负数或高位参与时要说明整数宽度；再用边界 缺失零、缺失 n、数组长度一、也可用求和但要防溢出 反查初始化、更新顺序和退出条件。

\textbf{状态回放。} 手算路线：手算时把整数写成二进制，并在每一位旁边标注含义。状态转移只允许改变题目动作允许改变的那些位。状态字段：掩码含义、当前位、目标位集合、状态转移和答案读取；负数或高位参与时要说明整数宽度。状态升级：把集合状态压成掩码；包含、加入、删除、枚举子集都变成位操作，但每一位的题目含义不能丢。

\textbf{证明和边界。} 证明从逐位独立或子集归纳开始：被压缩到位里的信息足以决定未来转移。边界：缺失零、缺失 n、数组长度一、也可用求和但要防溢出。变体：异或法适合展示值域完整且只有一个缺口的结构。

\noindent\textbf{LeetCode 287 寻找重复数。} 模型：数组值域使下标到值形成函数图，重复数是环入口。推导重点：快慢指针不修改数组且常数空间；值域二分用抽屉原理。

\textbf{二刷读题。} 读题时先用普通状态解释每一项布尔信息。本题可以压成“数组值域使下标到值形成函数图，重复数是环入口”，位运算只是更紧凑的编码。先遮住答案，只保留题面，复述候选对象、合法性条件和输出目标。

\textbf{暴力回放。} 慢版本先用集合、数组或布尔表保存状态，逐步执行题目动作。确认每个布尔字段的含义后，才能把它压缩成位，而不是直接写位运算公式。二刷时先写慢版本或口述慢版本，用它确认不会漏答案。

\textbf{特例回放。} 拿 [1,3,4,2,2] 手算，把值当成 next 指针会形成 1->3->2->4->2 的环，入环点 2 就是重复数。规律是长度 n+1、值域 1..n 强制形成环；不能修改数组时可用快慢指针或值域二分。把这个特例继续拆开看：每一位或每个掩码位都要对应题目里的一个布尔事实。位运算只是压缩表示；如果两个状态在位置相同但掩码不同，未来能力不同，就不能合并。复盘时把这个特例压成状态字段：掩码含义、当前位、目标位集合、状态转移和答案读取；负数或高位参与时要说明整数宽度；再用边界 重复数出现多次、不能排序、不能改数组、长度为 n 加一 反查初始化、更新顺序和退出条件。

\textbf{状态回放。} 手算路线：手算时把整数写成二进制，并在每一位旁边标注含义。状态转移只允许改变题目动作允许改变的那些位。状态字段：掩码含义、当前位、目标位集合、状态转移和答案读取；负数或高位参与时要说明整数宽度。状态升级：把集合状态压成掩码；包含、加入、删除、枚举子集都变成位操作，但每一位的题目含义不能丢。

\textbf{证明和边界。} 证明从逐位独立或子集归纳开始：被压缩到位里的信息足以决定未来转移。边界：重复数出现多次、不能排序、不能改数组、长度为 n 加一。变体：快慢指针要求值都能作为合法下标。

\noindent\textbf{LeetCode 318 最大单词长度乘积。} 模型：两个单词没有公共字母时才可相乘，字母集合可压成位掩码。推导重点：每个单词生成 26 位掩码，两个掩码按位与为零表示无交集。

\textbf{二刷读题。} 读题时先用普通状态解释每一项布尔信息。本题可以压成“两个单词没有公共字母时才可相乘，字母集合可压成位掩码”，位运算只是更紧凑的编码。先遮住答案，只保留题面，复述候选对象、合法性条件和输出目标。

\textbf{暴力回放。} 慢版本先用集合、数组或布尔表保存状态，逐步执行题目动作。确认每个布尔字段的含义后，才能把它压缩成位，而不是直接写位运算公式。二刷时先写慢版本或口述慢版本，用它确认不会漏答案。

\textbf{特例回放。} 拿 abcw 和 xtfn 手算，两个单词的 26 位掩码按位与为 0，说明没有公共字符，乘积 4*4 可作为候选。规律是每个单词压成字符集合掩码，重复字母不增加位，掩码相交才排除。把这个特例继续拆开看：每一位或每个掩码位都要对应题目里的一个布尔事实。位运算只是压缩表示；如果两个状态在位置相同但掩码不同，未来能力不同，就不能合并。复盘时把这个特例压成状态字段：掩码含义、当前位、目标位集合、状态转移和答案读取；负数或高位参与时要说明整数宽度；再用边界 重复字母、空字符串、多个单词同掩码、乘积范围 反查初始化、更新顺序和退出条件。

\textbf{状态回放。} 手算路线：手算时把整数写成二进制，并在每一位旁边标注含义。状态转移只允许改变题目动作允许改变的那些位。状态字段：掩码含义、当前位、目标位集合、状态转移和答案读取；负数或高位参与时要说明整数宽度。状态升级：把集合状态压成掩码；包含、加入、删除、枚举子集都变成位操作，但每一位的题目含义不能丢。

\textbf{证明和边界。} 证明从逐位独立或子集归纳开始：被压缩到位里的信息足以决定未来转移。边界：重复字母、空字符串、多个单词同掩码、乘积范围。变体：若字符集大于整数位数，需要 bitset 或哈希集合。

\noindent\textbf{LeetCode 338 比特位计数。} 模型：每个数的一的个数可由去掉最低位一后的较小数转移。推导重点：dp[x] = dp[x \& (x - 1)] + 1，或 dp[x] = dp[x >> 1] + (x \& 1)。

\textbf{二刷读题。} 读题时先用普通状态解释每一项布尔信息。本题可以压成“每个数的一的个数可由去掉最低位一后的较小数转移”，位运算只是更紧凑的编码。先遮住答案，只保留题面，复述候选对象、合法性条件和输出目标。

\textbf{暴力回放。} 慢版本先用集合、数组或布尔表保存状态，逐步执行题目动作。确认每个布尔字段的含义后，才能把它压缩成位，而不是直接写位运算公式。二刷时先写慢版本或口述慢版本，用它确认不会漏答案。

\textbf{特例回放。} 拿 n=5 手算，5 的二进制 101，比 4 多一个最低位 1，所以 bits[5]=bits[4]+1；也可用 5\&(4)=4 得到 bits[5]=bits[4]+1。规律是每个数的答案来自去掉最低位 1 或右移一位后的更小数。把这个特例继续拆开看：每一位或每个掩码位都要对应题目里的一个布尔事实。位运算只是压缩表示；如果两个状态在位置相同但掩码不同，未来能力不同，就不能合并。复盘时把这个特例压成状态字段：掩码含义、当前位、目标位集合、状态转移和答案读取；负数或高位参与时要说明整数宽度；再用边界 n 为零、二的幂、连续区间、表达式优先级 反查初始化、更新顺序和退出条件。

\textbf{状态回放。} 手算路线：手算时把整数写成二进制，并在每一位旁边标注含义。状态转移只允许改变题目动作允许改变的那些位。状态字段：掩码含义、当前位、目标位集合、状态转移和答案读取；负数或高位参与时要说明整数宽度。状态升级：把集合状态压成掩码；包含、加入、删除、枚举子集都变成位操作，但每一位的题目含义不能丢。

\textbf{证明和边界。} 证明从逐位独立或子集归纳开始：被压缩到位里的信息足以决定未来转移。边界：n 为零、二的幂、连续区间、表达式优先级。变体：它展示位运算和 DP 的结合，而不是逐个数循环数位。

\noindent\textbf{LeetCode 371 两整数之和。} 模型：不用加号时，加法可以拆成无进位和与进位。推导重点：异或得到无进位和，与运算后左移得到进位，迭代直到进位为零。

\textbf{二刷读题。} 读题时先用普通状态解释每一项布尔信息。本题可以压成“不用加号时，加法可以拆成无进位和与进位”，位运算只是更紧凑的编码。先遮住答案，只保留题面，复述候选对象、合法性条件和输出目标。

\textbf{暴力回放。} 慢版本先用集合、数组或布尔表保存状态，逐步执行题目动作。确认每个布尔字段的含义后，才能把它压缩成位，而不是直接写位运算公式。二刷时先写慢版本或口述慢版本，用它确认不会漏答案。

\textbf{特例回放。} 拿 5+7 手算，异或得到无进位和 2，按位与左移得到进位 10；继续迭代直到进位为 0。规律是加法拆成无进位和与进位两部分，负数参与时要用足够明确的整数宽度处理。把这个特例继续拆开看：每一位或每个掩码位都要对应题目里的一个布尔事实。位运算只是压缩表示；如果两个状态在位置相同但掩码不同，未来能力不同，就不能合并。复盘时把这个特例压成状态字段：掩码含义、当前位、目标位集合、状态转移和答案读取；负数或高位参与时要说明整数宽度；再用边界 正负数、进位传播、有符号溢出、需要使用无符号中间值 反查初始化、更新顺序和退出条件。

\textbf{状态回放。} 手算路线：手算时把整数写成二进制，并在每一位旁边标注含义。状态转移只允许改变题目动作允许改变的那些位。状态字段：掩码含义、当前位、目标位集合、状态转移和答案读取；负数或高位参与时要说明整数宽度。状态升级：把集合状态压成掩码；包含、加入、删除、枚举子集都变成位操作，但每一位的题目含义不能丢。

\textbf{证明和边界。} 证明从逐位独立或子集归纳开始：被压缩到位里的信息足以决定未来转移。边界：正负数、进位传播、有符号溢出、需要使用无符号中间值。变体：这是位级算术模拟，和异或抵消类题目的语义不同。

\noindent\textbf{LeetCode 393 UTF-8 编码验证。} 模型：每个字节的高位模式决定一个字符需要的后续字节数量。推导重点：用位掩码判断首字节类型，并维护还需要多少个 continuation 字节。

\textbf{二刷读题。} 读题时先用普通状态解释每一项布尔信息。本题可以压成“每个字节的高位模式决定一个字符需要的后续字节数量”，位运算只是更紧凑的编码。先遮住答案，只保留题面，复述候选对象、合法性条件和输出目标。

\textbf{暴力回放。} 慢版本先用集合、数组或布尔表保存状态，逐步执行题目动作。确认每个布尔字段的含义后，才能把它压缩成位，而不是直接写位运算公式。二刷时先写慢版本或口述慢版本，用它确认不会漏答案。

\textbf{特例回放。} 拿 [197,130,1] 手算，197 的高位模式表示还需要 1 个 continuation 字节，130 满足 10xxxxxx，之后 1 是单字节。规律是状态保存剩余后续字节数量，非法首字节或后续字节不足都失败。把这个特例继续拆开看：每一位或每个掩码位都要对应题目里的一个布尔事实。位运算只是压缩表示；如果两个状态在位置相同但掩码不同，未来能力不同，就不能合并。复盘时把这个特例压成状态字段：掩码含义、当前位、目标位集合、状态转移和答案读取；负数或高位参与时要说明整数宽度；再用边界 非法首字节、后续字节不足、后续字节格式错误、单字节字符 反查初始化、更新顺序和退出条件。

\textbf{状态回放。} 手算路线：手算时把整数写成二进制，并在每一位旁边标注含义。状态转移只允许改变题目动作允许改变的那些位。状态字段：掩码含义、当前位、目标位集合、状态转移和答案读取；负数或高位参与时要说明整数宽度。状态升级：把集合状态压成掩码；包含、加入、删除、枚举子集都变成位操作，但每一位的题目含义不能丢。

\textbf{证明和边界。} 证明从逐位独立或子集归纳开始：被压缩到位里的信息足以决定未来转移。边界：非法首字节、后续字节不足、后续字节格式错误、单字节字符。变体：这是位模式解析题，重点是协议规则和状态机。

\noindent\textbf{LeetCode 421 数组中两个数的最大异或值。} 模型：最大异或值由高位到低位尽量取一决定。推导重点：逐位尝试把当前前缀答案设为一，用哈希集合验证是否存在两个前缀配对。

\textbf{二刷读题。} 读题时先用普通状态解释每一项布尔信息。本题可以压成“最大异或值由高位到低位尽量取一决定”，位运算只是更紧凑的编码。先遮住答案，只保留题面，复述候选对象、合法性条件和输出目标。

\textbf{暴力回放。} 慢版本先用集合、数组或布尔表保存状态，逐步执行题目动作。确认每个布尔字段的含义后，才能把它压缩成位，而不是直接写位运算公式。二刷时先写慢版本或口述慢版本，用它确认不会漏答案。

\textbf{特例回放。} 拿 [3,10,5,25,2,8] 手算，25 和 5 的异或为 28。逐位构造答案时，先假设当前位能为 1，再用哈希集合检查是否存在两个前缀配对。规律是从高位到低位贪心验证前缀，不需要枚举所有数对。把这个特例继续拆开看：每一位或每个掩码位都要对应题目里的一个布尔事实。位运算只是压缩表示；如果两个状态在位置相同但掩码不同，未来能力不同，就不能合并。复盘时把这个特例压成状态字段：掩码含义、当前位、目标位集合、状态转移和答案读取；负数或高位参与时要说明整数宽度；再用边界 零、重复数、最高位、负数语义 反查初始化、更新顺序和退出条件。

\textbf{状态回放。} 手算路线：手算时把整数写成二进制，并在每一位旁边标注含义。状态转移只允许改变题目动作允许改变的那些位。状态字段：掩码含义、当前位、目标位集合、状态转移和答案读取；负数或高位参与时要说明整数宽度。状态升级：把集合状态压成掩码；包含、加入、删除、枚举子集都变成位操作，但每一位的题目含义不能丢。

\textbf{证明和边界。} 证明从逐位独立或子集归纳开始：被压缩到位里的信息足以决定未来转移。边界：零、重复数、最高位、负数语义。变体：也可以建二进制 Trie，每个数尽量走相反位。

\subsection{自测标准}

二刷时不要只确认自己见过这道题。请遮住答案，先讲题面模型，再讲暴力基线和瓶颈，然后说出优化结构保存了什么、不变量如何排除错误候选、边界实验如何打破错误实现。

\begin{principlebox}{二刷标准}
二刷不是重新看一遍答案，而是把题目变成可讲、可证、可调试、可迁移的知识块。每道题至少回答：暴力瓶颈、优化依据、不变量、复杂度、边界、C++ 成本和一个变体。
\end{principlebox}
